\documentclass[11pt]{article}
\usepackage[a4paper,left=24mm,right=30mm,top=22mm,bottom=22mm,headheight=14pt]{geometry}
\usepackage{fontspec,multirow,array,graphicx}
\usepackage[unicode,hidelinks]{hyperref}
\usepackage{polyglossia}
\setmainlanguage{english}
\setotherlanguages{arabic}
\setmainfont{Linux Libertine O}[Ligatures=TeX]
\newfontfamily\arabicfont{Amiri}[Script=Arabic]
\newfontfamily\ArBody{Amiri}[Script=Arabic,Scale=1.08]
\newfontfamily\ArN{Amiri}[Script=Arabic,Scale=0.95]
\newfontfamily\ArT{Amiri}[Script=Arabic,Scale=0.9]
\newfontfamily\SgFont{Yu Gothic}
\newfontfamily\LatinF{Linux Libertine O}
\newcommand{\Sg}[1]{{\SgFont #1}}
\hypersetup{pdftitle={al-Battani, Commentary on Ptolemy's Tetrabiblos (Escorial, RBMSL, ar. 969, ff. 20r-80v). After the PAL transcription by Nadine Loehr. S14},pdfauthor={al-Battani; transcription: Nadine Loehr, Ptolemaeus Arabus et Latinus (CC BY 4.0)}}
\setlength{\parindent}{0pt}\setlength{\parskip}{0pt}
\tracinglostchars=2
\makeatletter
\def\ps@ms{\def\@oddhead{\hbox to\textwidth{\footnotesize AL-BATTĀNĪ — COMMENTARY ON THE TETRABIBLOS (ESCORIAL ÁR. 969)\hfil\leftmark}}%
\let\@evenhead\@oddhead\def\@oddfoot{\hbox to\textwidth{\footnotesize\rightmark\hfil\thepage}}\let\@evenfoot\@oddfoot}
\makeatother
\pagestyle{ms}
\newcommand{\Folio}[2]{\par\bigskip\penalty-200\markboth{f. #2}{AB03-F#1}\pdfbookmark[1]{f. #2}{f-#1}%
 \hypertarget{AB03-F#1}{}{\large\bfseries f.\,#2}\par\smallskip}
\newcommand{\Chap}[2]{\pdfbookmark[2]{#2}{ch-#1}}
\newsavebox{\LBox}
% one line of the manuscript, right to left, its number in the right margin; a line longer than the measure (a
% marginal addition set in the line) runs on, right-aligned, to the next line of the page
\newcommand{\MsL}[3]{\par\noindent\hypertarget{#2}{}\hbox to\linewidth{\parbox[t]{\linewidth}{\setRTL\ArBody\raggedleft #3\par}\rlap{\hspace{3mm}\footnotesize #1}}\par}
\newcommand{\MsLc}[3]{\par\noindent\hypertarget{#2}{}\hbox to\linewidth{\parbox[t]{\linewidth}{\setRTL\ArBody\centering #3\par}\rlap{\hspace{3mm}\footnotesize #1}}\par}
\newcommand{\Hd}[1]{{\bfseries #1}}
\newcommand{\PN}[1]{\LR{\textsuperscript{\fontsize{7}{7}\selectfont #1}}}
\newcommand{\BF}[1]{\LR{\fontsize{8}{8}\selectfont\,/#1/\,}}
\newenvironment{PNotes}{\par\smallskip\noindent\rule{40mm}{0.4pt}\par\footnotesize}{\par}
\newcommand{\PNote}[3]{\par\hangindent=6mm\hypertarget{#2}{}\textbf{#1}\enspace #3}
\begin{document}

\thispagestyle{empty}
\begin{center}\Large AL-BATTĀNĪ\par\vspace{4mm}\LARGE Commentary on Ptolemy’s \textit{Tetrabiblos}\par\vspace{3mm}\large\RL{\ArBody شرح أبي عبد الله محمّد بن جابر بن سنان الحرّانيّ المعروف بالبتّانيّ على كتاب الأربع مقالات لبطلميوس}\par\vspace{5mm}\large Escorial, Real Biblioteca de San Lorenzo, árabe 969, ff.\,20r–80v\par\vspace{6mm}\normalsize After the transcription by Nadine Löhr, \textit{Ptolemaeus Arabus et Latinus. Texts}\\(https://ptolemaeus.badw.de/text/M1016), CC BY 4.0, commit \texttt{762614d992ac}\par\vspace{4mm}\large S14 --- edition of the source, line by line\end{center}\vspace{6mm}
\noindent Every line of this edition is one line of the manuscript as Nadine Löhr’s transcription divides it (her «||»), numbered in the right margin and anchored \texttt{AB03-F020r-L01} and so on. The text, the editorial signs, the notes, the folio correspondences of the Berlin manuscript (\texttt{/B1r/}) and the chapter headings are hers, unchanged; the notes stand after each page of the manuscript, with her numbers. This conversion was made by an AI (Claude) without the images of the manuscript; no person has reviewed it. The transcription is licensed CC BY 4.0, and so is this edition of it. Her introduction follows; the fields that PAL’s website fills from its catalogue (the siglum, author and title of the work and the name of the witness), empty in the files, are filled from PAL’s page of the work, C.2.2 (in italics).\par\medskip\sloppy
\par\smallskip\noindent Work 	extit{C.2.2}
\par\smallskip\noindent transcribed by Nadine Löhr
\par\smallskip\noindent Nadine Löhr, ‘Transcription of 	extit{al-Battānī}, 	extit{Sharḥ Kitāb al-Arbaʿ maqālāt fī aḥkām ʿilm al-nujūm}’ (update 25.05.2021), Ptolemaeus Arabus et Latinus. Texts, URL = https://ptolemaeus.badw.de/text/M1016
\par\smallskip\noindent This is a transcription of al-Battānī’s commentary on Ptolemy’s Tetrabiblos based on manuscript Escorial árabe 969. It comments on some chapters of Books I and II and on all chapters of Books III and IV.
\par\smallskip\noindent I indicate the folio correspondences between the Escorial manuscript (e.g., /1r/) and the Berlin Sprenger 1840 manuscript (including siglum B; e.g., /B1r/) which is acephalous and missing the table at the end of Book VI.
\par\smallskip\noindent I have consulted Yamamoto’s preliminary edition of the Arabic Tetrabiblos for some difficult or ambiguous passages. However, I only indicate textual variants to this edition when single words in the Ptolemaic quotes are either significantly different or apparently without the intended meaning in both witnesses (Berlin and Escorial).
\par\smallskip\noindent Note: There are numerous scribal mistakes concerning geographical information in both manuscripts, particularly in Chapter 10. For practical reasons I decided against an inclusion of the vast amount of related footnotes in this chapter and only rendered what is written in the base manuscript. (For identification of places see Robbins:1940, Yamamoto/Mimura or contact the editor.)
\par\smallskip\noindent I am very grateful to Walid Sousa and Pouyan Rezvani for the correction of some mistakes, and to Benno van Dalen for helping me with the table on f. 80r.
\par\smallskip\noindent I have normalized in many places, without particular indications according to our transcription principles:
\par\hangindent=4mm\noindent\textbullet\ the writing of final radicals wāw and yāʾ
\par\hangindent=4mm\noindent\textbullet\ as well as hamzas and maddas.
\par\hangindent=4mm\noindent\textbullet\ I added shaddas wherever they were missing, except on solar letters and relative pronouns.
\par\hangindent=4mm\noindent\textbullet\ I added diacritics in case they were vacant.
\par\smallskip\noindent Symbols
\par\hangindent=4mm\noindent\textbullet\ || indicates a new line in the manuscript.
\par\hangindent=4mm\noindent\textbullet\ // indicates the start of a new page, e.g.: “/2v/”.
\par\hangindent=4mm\noindent\textbullet\ \Sg{〈}\Sg{〉} around a passage added by the editor. At the beginning of a heading, the \Sg{〈}\Sg{〉}-brackets surrounding a number indicate chapter numbers which were not provided by the scribes.
\par\hangindent=4mm\noindent\textbullet\ [] around a passage found in the text but deemed as to be deleted, additions by the scribe, eulogies etc.
\par\hangindent=4mm\noindent\textbullet\ \{\} around a passage whose reading is uncertain or conjectural.
\par\hangindent=4mm\noindent\textbullet\ \{…\} dots inside the curly brackets indicate an illegible lacuna in the manuscript.
\par\hangindent=4mm\noindent\textbullet\ [[…]] is used to indicate words and letters that were intentionally and visibly deleted by the scribe or in a different hand mostly by means of crossing out.
\par\hangindent=4mm\noindent\textbullet\ \Sg{〈}\Sg{〈}…\Sg{〉}\Sg{〉} is used to indicate marginal and interlinear additions by the scribe or in a different hand
\par\smallskip\noindent Sigla
\par\hangindent=4mm\noindent\textbullet\ B: Berlin, Staatsbibliothek Preußischer Kulturbesitz, Sprenger 1840
\par\hangindent=4mm\noindent\textbullet\ T: Preliminary Edition of the Tetrabiblos (Ibrāhīm b. al-Ṣalt/Ḥunayn b. Isḥāq) by Prof. Keiji Yamamoto
\par\clearpage
\Folio{020r}{20r}
\MsLc{1}{AB03-F020r-L01}{كتاب}
\MsLc{2}{AB03-F020r-L02}{الأربع مقالات في أحكام على النجوم}
\MsLc{3}{AB03-F020r-L03}{سیّد الحکماء بطلميوس}
\MsLc{4}{AB03-F020r-L04}{شرح أبي عبد الله}
\MsLc{5}{AB03-F020r-L05}{محمّد بن جابر}
\MsLc{6}{AB03-F020r-L06}{أبو\PN{1} سنان}
\MsLc{7}{AB03-F020r-L07}{الحرّانيّ}
\MsLc{8}{AB03-F020r-L08}{المعروف بالبتّانيّ}
\begin{PNotes}
\PNote{1}{AB03-n001}{sic}
\end{PNotes}
\Folio{020v}{20v}
\MsL{1}{AB03-F020v-L01}{[بسم الله الرحمن الرحيم وهو حسبي]}
\Chap{I.1}{<1> معرفة قوّة الشمس والقمر والكواكب المتحیّرة والثابتة}
\MsL{2}{AB03-F020v-L02}{\Hd{\LR{\LatinF \Sg{〈}1\Sg{〉}}}\Hd{ معرفة قوّة الشمس والقمر والكواكب المتحیّرة والثابتة}}
\MsL{3}{AB03-F020v-L03}{قال أبو عبد الله محمّد بن جابر الحرّانيّ المعروف بالبتّاني: فعل جوهر الشمس أن يجفّف}
\MsL{4}{AB03-F020v-L04}{ويسخّن تجفيفًا يسيرًا. وأمّا القمر فأكثر فعل قوّته الترطيب وهو يغيّر الأجسام تغيّرًا بيّنًا لإنضاجه}
\MsL{5}{AB03-F020v-L05}{وتعفّنه أكثر ما فيها وهو أيضًا يشارك الشمس في الحرارة مشاركة يسيّره لقبوله الضوء}
\MsL{6}{AB03-F020v-L06}{منها. وأمّا كوكب زحل فأكثر فعله في التبريد وهو يجفّف يسيّرًا بسبب برده ويبسه\PN{2}. وأمّا}
\MsL{7}{AB03-F020v-L07}{كوكب المشتري فمعتدل المزاج ليسخّن ویرطّب معًا وهو فاعل للرياح المولدة. وأمّا كوكب}
\MsL{8}{AB03-F020v-L08}{المرّيخ فإنّه يجفّف خاصّة وأكثر فعل قوّته الإحراق وذلك أنّ طبيعته محرقة لمشاكلة لونه}
\MsL{9}{AB03-F020v-L09}{لون العار\PN{3}. وأمّا كوكب الزهرة فمعتدل المزاج غير أنّه يرطّب ترطيبًا كثيرًا ويسخّن إسخانًا يسيّرًا.}
\MsL{10}{AB03-F020v-L10}{وأمّا كوكب عطارد فالأمر في ترطيبه وتجفيفه متساوٍ وذلك أنّه في بعض الأوقات يجفّف}
\MsL{11}{AB03-F020v-L11}{وینشّف وفي بعضها یرطّب ويغيّره إلى کلتي الحالين تغيّرًا سريعًا وهو مولّد للرياح. وأمّا}
\MsL{12}{AB03-F020v-L12}{قوى الكواكب الثابتة من مزاجاتها إذا قيّست إلى قوى الكواكب المتحیّرة وطبائعها التي}
\MsL{13}{AB03-F020v-L13}{وجدت لها لمّا رصدت فقد أثبتنا ما في جدوال مواضع الكواكب الثابتة من غير هذا}
\MsL{14}{AB03-F020v-L14}{الكتاب ورسمنا إذا كلّ كوكب منها رسمًا تعرف فيه الكواكب المتحیّرة التي تشبهها طبائعها}
\MsL{15}{AB03-F020v-L15}{فيها ما يشبه مزاجه مزاج كوكب واحد من المتحیّرة ومنها ما یشبه مزاجه مزاج كوكبين}
\MsL{16}{AB03-F020v-L16}{بالسويّة ومنها ما یشبه أحد الكوكبين شبهًا كثيرًا. ویشبه الثاني شبهًا قليلًا فعلى حسب}
\MsL{17}{AB03-F020v-L17}{الرسوم الموجودة هناك تكون مشابهة طبائعها ومن قبل ذلك استغنينا عن \BF{B1r}ذكر}
\Chap{I.2}{<2> في قوى الكواكب حسب الأشكال}
\MsL{18}{AB03-F020v-L18}{قواها في هذا الموضع. \Hd{\LR{\LatinF \Sg{〈}2\Sg{〉}}}\Hd{ في قوى الكواكب حسب الأشكال} التي تكون لها بقياس من}
\MsL{19}{AB03-F020v-L19}{الشمس وقد يقبل القمر والكواكب الثلاثة العلويّة الزيادة والنقصان في قواها}
\MsL{20}{AB03-F020v-L20}{الخاصّيّة\PN{4} بحسب الأشكال التي تكون لها عند الشمس وذلك أنّ القمر منذ الوقت الذي}
\MsL{21}{AB03-F020v-L21}{يهلّ فيه إلى وقت انتصافه. الأوّل في الضوء يكون فاعلًا للرطوبة خاصّة ومنذ وقت}
\begin{PNotes}
\PNote{2}{AB03-n002}{T: \RL{\ArN يشبّ}}
\PNote{3}{AB03-n003}{T: \RL{\ArN النار}}
\PNote{4}{AB03-n004}{B: \RL{\ArN الخاصّة}}
\end{PNotes}
\Folio{021r}{21r}
\MsL{1}{AB03-F021r-L01}{انتصافه الأوّل إلى وقت امتلائه مقابلته للشمس يكون فاعلًا للحرارة ومنذ وقت امتلائه}
\MsL{2}{AB03-F021r-L02}{إلى وقت انتصافه الثاني في الضوء يكون فاعلًا اليبوسة ومنذ وقت انتصافه الثاني إلى}
\MsL{3}{AB03-F021r-L03}{وقت اختفائه يكون فاعلًا للبرودة. وأمّا زحل والمشتري والمرّيخ فهي وقت طلوعها}
\MsL{4}{AB03-F021r-L04}{بالغدوات وظهورها من تحت شعاع الشمس إلى وقت وقوفها الأوّل ترطّب أكثر. ومنذ}
\MsL{5}{AB03-F021r-L05}{وقت وقوفها الأوّل إلى وقت طلوعها في أوّل الليل وذلك حين تقابل الشمس تسخّن أكثر.}
\MsL{6}{AB03-F021r-L06}{ومنذ وقت طلوعها في أوّل الليل إلى وقوفها الثاني تجفّف أكثر. ومنذ وقت وقوفها الثاني إلى وقت}
\MsL{7}{AB03-F021r-L07}{غروبها واختفائها تبرّد أكثر. وإذا مازج بعضها بعضًا تحدث في الهواء أصنافًا كثيرة مختلفة الأحوال}
\Chap{I.3}{<3> في قوى الأربعة الأزمنة}
\MsL{8}{AB03-F021r-L08}{جدًّا وتسمّى هذه الأحوال أسباب الكواكب. \Hd{\LR{\LatinF \Sg{〈}3\Sg{〉}}}\Hd{ في قوى الأربعة الأزمنة} من السنة وأسنان}
\MsL{9}{AB03-F021r-L09}{الحيوان [[الرابع]] الغالب على وقت الربيع الرطوبة والغالب على وقت الصيف الحرارة والغالب على}
\MsL{10}{AB03-F021r-L10}{وقت الخريف اليبوسة والغالب على وقت الشتاء البرودة. وكذلك أوّل السنّ في الحيوان}
\MsL{11}{AB03-F021r-L11}{الغالب عليه الجوهر \BF{B1v}الرطب والغالب عليه في السنّ الثانية إلى وقت المنتهى الحرارة وفي السنّ}
\MsL{12}{AB03-F021r-L12}{الثالثة إذا قارب الأخطاط وبدأ في النقصان اليبس وفي آخر الأسنان إذا قارب الانحلال.}
\Chap{I.4}{<4> في قوى الأربعة زوايا والأربعة}
\MsL{13}{AB03-F021r-L13}{البرد بمنزلة الربيع والصيف والخريف والشتاء. \Hd{\LR{\LatinF \Sg{〈}4\Sg{〉}}}\Hd{ في قوى الأربعة زوايا والأربعة}}
\MsL{14}{AB03-F021r-L14}{مواضع من الأفق وكذلك أيضًا حال الأربعة مواضع من الأفق والزوايا التي منها هبوب الرياح}
\MsL{15}{AB03-F021r-L15}{إلى جميع النواحيّ فإنّ الغالب على الناحية التي تلي المشرق اليبس والريح التي تهبّ منها وهي}
\MsL{16}{AB03-F021r-L16}{الشرقيّة التي تسمّى ابليوطس\PN{5} عديمة الرطوبة مجفّفة. وأمّا الرياح التي تهبّ من جانب الجنوب}
\MsL{17}{AB03-F021r-L17}{التي تسمّى رياحًا جنوبيّة فحارّة محرقة. وأمّا ناحيها\PN{6} المغرب فرطبة. والرياح التي تهبّ منها التي}
\MsL{18}{AB03-F021r-L18}{تسمّى الغربيّة تليّن ترطب. وأمّا ناحية الشمال فباردة جدًّا أو الرياح التي تهبّ من هذه}
\MsL{19}{AB03-F021r-L19}{الناحية وهي التي تسمّى الشماليّة باردة مكثّفة. ومعرفة هذه الأمور أيضًا نافعة ومعيّنة على}
\MsL{20}{AB03-F021r-L20}{تمييز المزاجات فإنّه من البيّن أنّ القوّة التي في الكواكب تتغيّر أفعالها عن أحوال أوقات}
\MsL{21}{AB03-F021r-L21}{السنة أو عن الأسنان أو الزوايا وإنّها إذا كانت الأحوال مشاكلة لها كانت كيفياتها}
\begin{PNotes}
\PNote{5}{AB03-n005}{T: \RL{\ArN اقيلوطس}}
\PNote{6}{AB03-n006}{B: \RL{\ArN ناحية}}
\end{PNotes}
\Folio{021v}{21v}
\MsL{1}{AB03-F021v-L01}{خالصة محضة لا يخالطها شيء ويكون فعلها أقوى مثل ما تكون الكواكب المسخنة بالطبع}
\MsL{2}{AB03-F021v-L02}{في الأحوال الحارّة أقوى. والكواكب الرطبة في الأحوال الرطبة أقوى وتكون قوّتها في}
\MsL{3}{AB03-F021v-L03}{أضداد هذه الأحوال ليست محضة خالصة لأنّ ممزوجة بغيّرها ضعيفة كما تكون قوّة الكواكب}
\MsL{4}{AB03-F021v-L04}{المسخنة في الأحوال الباردة أضعف وقوّة الكواكب الرطبة في الأحوال اليابسة أضعف}
\MsL{5}{AB03-F021v-L05}{وتزيد في ذلك وتنقص وتغيّر تغيّرًا ليس باليسيّر \BF{B2r}طبائع الكواكب الثابتة التي تمازج}
\MsL{6}{AB03-F021v-L06}{تلك الكواكب أو تشترك معها في أحد الأشكال التسعة التي تكون لها لا سيّما عند الشمس}
\MsL{7}{AB03-F021v-L07}{وهي المذكورة فيما تقدّم من هذا الكتاب. وكذلك يكون التغيّر في سائر الأحوال يقاس}
\Chap{I.5}{<5> في البروج الثابتة والمنقلبة وذوات الجسدين}
\MsL{8}{AB03-F021v-L08}{امتزاج الكيفية\LR{\LatinF . }\Hd{\LR{\LatinF \Sg{〈}5\Sg{〉}}}\Hd{ في البروج الثابتة والمنقلبة وذوات الجسدين} من البروج برجان.}
\MsL{9}{AB03-F021v-L09}{منقلبان وبرجان معتدلان فالمنقلبان هما برجًا السرطان والجدي والمعتدلان}
\MsL{10}{AB03-F021v-L10}{هما برجًا الحمل والميزان وإنّما سمّيت هذه البروج بهذه الأسماء لأنّ الشمس إذا صارت}
\MsL{11}{AB03-F021v-L11}{إلى نقطة المنقلب الصیفيّ أحدث الصيف وإذا صارت إلى نقطة المنقلب}
\MsL{12}{AB03-F021v-L12}{الشتويّ أحدثت الشتاء وإذا صارت إلى نقطة الاعتدال الربيعيّ أحدث الربيع. وإذا صارت}
\MsL{13}{AB03-F021v-L13}{إلى نقطة الاعتدال الخريفيّ أحدث الخريف واعتدل الليل والنهار في هذين النقطتين}
\MsL{14}{AB03-F021v-L14}{في جميع الأرض ومن الثمانيّة البروج الباقيّة التي ما بين نقطة الاعتدالين والانقلابين}
\MsL{15}{AB03-F021v-L15}{أربعة ثابتة. وإنّما سمّيت ثابتة لأنّ الزمان يثبت فيها على حالة واحدة وهي التي تتلوا برجي}
\MsL{16}{AB03-F021v-L16}{الاعتدالين وبرجي الانقلابين وتسمّى الثور والعقرب والأسد والدلو. والبروج}
\MsL{17}{AB03-F021v-L17}{ذات الجسدين هي التي تتلوا البروج الثابتة وإنّما سمّيت بهذا الاسم لأنّها بين البروج}
\MsL{18}{AB03-F021v-L18}{الثابتة وبروج الاعتدال والانقلاب وكأنّها تشارك بأوائلها البروج الثابتة}
\MsL{19}{AB03-F021v-L19}{وبأواخرها البروج المتغيّرة وهي برج الجوزاء السنبلة والقوس والسمكة.}
\Chap{I.6}{<6> في اشتراك البروج في الأشكال}
\MsL{20}{AB03-F021v-L20}{\Hd{\LR{\LatinF \Sg{〈}6\Sg{〉}}}\Hd{ في اشتراك البروج في الأشكال} \BF{B2v}أمّا أمور الاشتراكات في الأشكال التي}
\MsL{21}{AB03-F021v-L21}{هي المقارنة والمقابلة والتربيع والتسديس والتثليث والتي ينظر بعضها إلى بعض}
\Folio{022r}{22r}
\MsL{1}{AB03-F022r-L01}{ويستوي في القوّة والأمره المطيعة من البروج التي لا ترتبط بعضها ببعض}
\MsL{2}{AB03-F022r-L02}{ولا تشترك في الأشكال فقد بيّناها فيما تقدّم من هذا الكتاب إذ كان السبب الذي}
\MsL{3}{AB03-F022r-L03}{يدلّ على استعمال هذه الأبعاد فقط دون غيرها من الأشياء الهندسيّة وكذلك غيرها}
\MsL{4}{AB03-F022r-L04}{من البروج المتضادّة والمتقارنة هي التي تقابل بعضها بعضًا والبروج المتباغضة}
\MsL{5}{AB03-F022r-L05}{هي التي تكون بعضها على تربيع بعض. والبروج المتحابّة هي التي يكون بعضها على تثليث}
\MsL{6}{AB03-F022r-L06}{بعض والتي تكون أيضًا على تسديس إذ أنّ هذين الشكلين يقعان من بروج}
\MsL{7}{AB03-F022r-L07}{متّفقة في الجنس إمّا مذكّرة وإمّا مؤنّثه والشكلان الأوّلان على خلاف ذلك.}
\MsL{8}{AB03-F022r-L08}{وأيضًا فأن بيوت الكوكبين الجنسين على مقابلة بيوت النيّرين وعلى تربيعها وبيوت}
\Chap{I.7}{<7> في البروج المذكّرة والمؤنّثة}
\MsL{9}{AB03-F022r-L09}{الكوكبين السعدين على تثليث وتسديس بيوت النيّرين. \Hd{\LR{\LatinF \Sg{〈}7\Sg{〉}}}\Hd{ في البروج المذكّرة والمؤنّثة}}
\MsL{10}{AB03-F022r-L10}{وغير ذلك من حالتها وما سمّيت به. أمّا البروج المؤنّثة والمذكّرة بالطبع فقد بيّناها}
\MsL{11}{AB03-F022r-L11}{عند ذكر المثلّثات وقد استعمل بعض الناس في ترتيب البروج المذكّرة والمؤنّثة ضربًا}
\MsL{12}{AB03-F022r-L12}{آخر من برج الطالع فعدّوا واحدًا مذكّرًا وواحدًا مؤنّثًا وصیّروا الابتداء من}
\MsL{13}{AB03-F022r-L13}{المذكّر كما أنّ بعض الناس يأخذ البروج المنقلبة ويبتدئ بها من برج القمر لأنّ}
\MsL{14}{AB03-F022r-L14}{القمر يعرض له الانقلاب لسرعة حركته أكثر من سائر الكواكب ومنهم من قسم البروج على}
\MsL{15}{AB03-F022r-L15}{أربعة أجزاء ورأى أنّ البروج المنسوبة إلى الغدوات المذكّرة التي هي \BF{B3r}من الطالع}
\MsL{16}{AB03-F022r-L16}{إلى وسط السماء والتي يقابلها من المغرب إلى وتد الأرض ورأوا أنّ الربعين الباقيين}
\MsL{17}{AB03-F022r-L17}{منسوبان إلى العشیّات مؤنّثين فجعلوا البروج التي تقع في هذين الربعين مؤنّثة}
\MsL{18}{AB03-F022r-L18}{منسوبة إلى العشیّات. وقد لقّبوا البروج بالقاب كثيرة مأخوذة من الصور التي}
\MsL{19}{AB03-F022r-L19}{تنسب إليها وذلك انّهم سمّوا بعضها ذوات الأربع قوائم وبعضها بريئة وبعضها}
\MsL{20}{AB03-F022r-L20}{مدبّرة وأهليّة وكثيرة الزرع وقليلة الزرع وما أشبه ذلك فقالوا إنّ البروج الأهليّة}
\MsL{21}{AB03-F022r-L21}{هي الحمل والثور والبريئة هي الأسد والقوس والناطقة التي تشاكل الناس}
\Folio{022v}{22v}
\MsL{1}{AB03-F022v-L01}{هي الجوزاء والدلو والميزان والسنبلة والقوس النصف الأوّل منه. وذوات}
\MsL{2}{AB03-F022v-L02}{الأجنحة السنبلة والقوس والجوزاء والحوت والمعقومة الحمل والثور والعقرب}
\MsL{3}{AB03-F022v-L03}{والجدي والسرطان والقليلة النسل والزرع والحمل والثور والميزان والقوس}
\MsL{4}{AB03-F022v-L04}{والجدي والدلو. والكثيرة النسل السرطان والعقرب والحوت والعقيمة الجوزاء}
\MsL{5}{AB03-F022v-L05}{والأسد والسنبلة. والمقطوعة الأعضاء الحمل والثور والأسد والعقرب. والكثيرة}
\MsL{6}{AB03-F022v-L06}{الغضب الحمل والأسد. والدالّة على الجمال والصباحة الجوزاء والميزان والسنبلة}
\MsL{7}{AB03-F022v-L07}{والعقرب والقوس والحوت والتي تجتمع وتمتلئ الحمل ومثلّثة والتي تعطي اليسار}
\MsL{8}{AB03-F022v-L08}{الكثير الجدي ومثلّثة والتي تصب وتستفرغ الميزان ومثلّثته والتي نقبض}
\MsL{9}{AB03-F022v-L09}{وتأخذ السرطان ومثلّثته والبروج الشديدة الصوت الجوزاء والسنبلة والميزان.}
\MsL{10}{AB03-F022v-L10}{والمعتدلة التي لها نصف صوت الحمل والثور والأسد والقوس \BF{B3v}والضعيفة الصوت}
\MsL{11}{AB03-F022v-L11}{الجدي والدلو. والتي لا صوت لها السرطان والحوت والعقرب. والدالّة على شدّة}
\MsL{12}{AB03-F022v-L12}{الهمّ الأسد والعقرب والجدي. والقليلة الهمّ السنبلة والميزان لظلمتها وذوات}
\MsL{13}{AB03-F022v-L13}{المكر والخداع الأسد والقوس والجدي والحوت. وقالوا إنّ الحمل والثور والجوزاء}
\MsL{14}{AB03-F022v-L14}{متى وافقت أربابها فيها \LR{\Sg{〉}}فهي\LR{\Sg{〈}} متحرّكة والسرطان والأسد والسنبلة متى وافقت أربابها فهي}
\MsL{15}{AB03-F022v-L15}{ساكنة وكذلك الميزان والعقرب والقوس أيضًا. إذا وافقت أربابها فيها فهي متحرّكات}
\MsL{16}{AB03-F022v-L16}{والجدي والدلو والحوت كذلك أيضًا إذا وافقت أربابها فيها فهي ساكنات.}
\Chap{I.8}{<8> في تذكير الكواكب وتأنيثها بقياسها}
\MsL{17}{AB03-F022v-L17}{\Hd{\LR{\LatinF \Sg{〈}8\Sg{〉}}}\Hd{ في تذكير الكواكب وتأنيثها بقياسها} إلى الشمس وإلى الأفق أيضًا. ويزعمون}
\MsL{18}{AB03-F022v-L18}{أنّ الكواكب تذکّر وتؤنّث بحسب الأشكال التي تكون لها بقياسها إلى الشمس}
\MsL{19}{AB03-F022v-L19}{وذلك لأنّها إذا كانت مشرقة متقدّمة للشمس فهي مذكّرة. وإذا كانت مغربة}
\MsL{20}{AB03-F022v-L20}{تابعة لها فهي مؤنّثة. وقد يكون ذلك أيضًا بقياسها إلى الأفق وذلك أنّها إذا}
\MsL{21}{AB03-F022v-L21}{كانت في الأشكال التي من المشرق إلى وسط السماء أو من المغرب إلى وتد}
\Folio{023r}{23r}
\MsL{1}{AB03-F023r-L01}{الأرض \LR{\Sg{〉}}\LR{\Sg{〉}} فهي مذكّرة لأنّ هاذين الربعين مذكّرين والأوّل منها يسمّى الشربع والثاني وسط السرعة. وإذا كانت فيما بين وسط السماء والمغرب أو فيما بين وتد الأرض \LR{\Sg{〈}}\LR{\Sg{〈}}\PN{7} والطالع فهي مؤنّثة لأنّ هذين الربعين مؤنّثين والأوّل منها يسمّى وسط}
\MsL{2}{AB03-F023r-L02}{\{الإبطاء\} والثاني الذي بين الطالع ووتد الأرض يسمّى البطن الثقيل.}
\Chap{I.9}{<9> في المواجهة والإنارة وما أشبههما}
\MsL{3}{AB03-F023r-L03}{\BF{B4r}\Hd{\LR{\LatinF \Sg{〈}9\Sg{〉}}}\Hd{ في المواجهة والإنارة وما أشبههما} وقد يقال إنّ للكواكب مواجة خاصّيّة إذا}
\MsL{4}{AB03-F023r-L04}{كان لكلّ واحد منها مشاركة للشمس في الشكل والقمر\PN{8} على هذه السبيل وهو أن يكون بين}
\MsL{5}{AB03-F023r-L05}{كلّ واحد منها وبين الشمس أو القمر من البعد مثل الذي بين بيته وبين بيت الشمس}
\MsL{6}{AB03-F023r-L06}{أو القمر غير أنّه يكون مغربًا عن الشمس أعني في البروج التالية مشرقًا عن القمر في}
\MsL{7}{AB03-F023r-L07}{البروج المتقدّمة مثال ذلك أن تكون الزهرة على تسد يس أحد النيّرين إلّا أنّها تكون}
\MsL{8}{AB03-F023r-L08}{في البرج الثالث من برج الشمس [[أو في الثاني عشر من برج القمر وباقي الكواكب]]\PN{9} \LR{\Sg{〉}}\LR{\Sg{〉}} على}
\MsL{9}{AB03-F023r-L09}{توالي البروج أو في البرج الحادي عشر من البرج الذي فيه القمر. ومثل أن يكون}
\MsL{10}{AB03-F023r-L10}{عطارد في البرج الثاني من برج الشمس أو في الثاني عشر \LR{\Sg{〈}}\LR{\Sg{〈}}\PN{10} من برج القمر وباقي}
\MsL{11}{AB03-F023r-L11}{الكواكب فتحسب بيوتها على هذه الحال. فإذا كانت هذه حال الكواكب عند\PN{11} النيّرين\PN{12} مواجهة}
\MsL{12}{AB03-F023r-L12}{لها. وأمّا الإنارة فإنّه يقال إنّ الكواكب إذا كانت في حيّزها وما أشبهه من المواضع}
\MsL{13}{AB03-F023r-L13}{الموافقة لها. وإذا تهیّأ لأحدها حظان من حظوظها أو أكثر مثل أن يكون في بيته وشرفه}
\MsL{14}{AB03-F023r-L14}{ومثلّثة وحده وما يشاكل أنّه في موضع فرحة وكرامته وحينئذ تزيد قوبتها في أفعالها}
\MsL{15}{AB03-F023r-L15}{المتشابهة للبروج التي هي لها ومشاكلتها في الطبع وقد يقال أيضًا إنّها تفرح إذا}
\MsL{16}{AB03-F023r-L16}{كانت في حظوظ الكواكب التي من حيّزها وهذا أضعف من الأوّل لأنّه دونها في القوّة}
\MsL{17}{AB03-F023r-L17}{ولكنّها على حال تشاركها على هذه الجهة مشاركة كما أنّها إذا وجهت في مواضع غريبة}
\MsL{18}{AB03-F023r-L18}{أو في حيّز مخالف لحيّزها تضعف قوّتها الخاصّيّة ونضير مخالفة مزاجها لمزاج البروج}
\MsL{19}{AB03-F023r-L19}{التي هي فيها سببًا لأن يحدث لها فعل مرکّب من الطبيعتين \BF{B4v}المختلفتين.}
\Chap{I.10}{<10> في مشاكلة المثلّثات للبلدان والكواكب}
\MsL{20}{AB03-F023r-L20}{\Hd{\LR{\LatinF \Sg{〈}10\Sg{〉}}}\Hd{ في مشاكلة المثلّثات للبلدان والكواكب} قد قسم بطلميوس جميع الأرض المسكونة}
\MsL{21}{AB03-F023r-L21}{التي في ناحية الشمال من معدّل النهار بخطّين يقطعانها في الطول والعرض أرباعًا}
\begin{PNotes}
\PNote{7}{AB03-n007}{left margin}
\PNote{8}{AB03-n008}{B: \RL{\ArN للقمر}}
\PNote{9}{AB03-n009}{B: same section erased}
\PNote{10}{AB03-n010}{in the main text and again in the left margin [erased]}
\PNote{11}{AB03-n011}{B: add. \RL{\ArN أحد}}
\PNote{12}{AB03-n012}{B: add. \RL{\ArN فهي}}
\end{PNotes}
\Folio{023v}{23v}
\MsL{1}{AB03-F023v-L01}{فجعل كلّ ربع منها مشاكلًا للمثلّث الذي ينسب إليه من مثلّثات تلك البروج الأربعة}
\MsL{2}{AB03-F023v-L02}{وجعل ما يلي وسط العمر أنّ من جميع الأرض المسكونة من كلّ ربع يقبل المشاكلة من}
\MsL{3}{AB03-F023v-L03}{المثلّث [[الذي]] المدبّر للزاويّة المقابلة لتلك الزاويّة وأشرك كوكب عطارد مع الكواكب}
\MsL{4}{AB03-F023v-L04}{المدبّرة لأنّه من حيّز مشترك وجعل مشاكلة البلدان لتلك الكواكب التي هي أصحاب البيوت}
\MsL{5}{AB03-F023v-L05}{من تلك المثلّثات مع بروج المثلّثات أيضًا وفصّل كلّ ربع من هذه الأرباع بثلاثة}
\MsL{6}{AB03-F023v-L06}{فصول حذي بها تلك البلدان التي فيه فصار كلّ جزء من هذه الثلاثة الأجزاء يشاكل}
\MsL{7}{AB03-F023v-L07}{برجًا من بروج المثلّثة مع الكوكب الذي هو صاحب ذلك البروج. وبين طبائع الأمّم}
\MsL{8}{AB03-F023v-L08}{التي تسكن في بلد من هذه البلدان الجزئيّة وأخلاقها وحالاتها وعمّ كلّ واحد من الأرباع}
\MsL{9}{AB03-F023v-L09}{بما يخصّه\PN{13} المثلّثات والمشاكلة له والكواكب التي تدبّر تلك المثلّثات وبحسب طبائع}
\MsL{10}{AB03-F023v-L10}{البروج وأصحابها من الكواكب واحتاج في ذلك ولخصّه بما نحن مستغنون عن ذكره في}
\MsL{11}{AB03-F023v-L11}{هذا الموضع. ثمّ وصف مشاكلة كلّ واحد من البلدان التي في كلّ ربع منها وما يلي وسط}
\MsL{12}{AB03-F023v-L12}{الأرض المحمودة منه وصنّفه تصنيفًا مفردًا للبروج والكواكب التي هي أرباب تلك البروج}
\MsL{13}{AB03-F023v-L13}{ليسهل النظر فيه على ما تقدّم من قوله وقد اقتصرنا على ذلك و\{اتقناه\} على هذه الجهة}
\MsL{14}{AB03-F023v-L14}{وبدأنا بما يشاكل كلّ برج وصاحبه من البلدان المنسوبة إلى ذلك الثلث وأتبعناه}
\MsL{15}{AB03-F023v-L15}{بالبلدان \BF{B5r}التي تلي وسط العمران وفصّلنا بيّننا بما يعرف به كلّ واحد منها ثمّ أثبتنا}
\MsL{16}{AB03-F023v-L16}{مشاكلة الدرج للبروج والكواكب على ما يقال إنّها كانت طالعة في وقت بنائها أو}
\MsL{17}{AB03-F023v-L17}{مواضع الشمس والقمر أو وسط سماء الملوك الذي كانوا في ذلك الزمان الذي}
\MsL{18}{AB03-F023v-L18}{بنيت فيه. فصار الذي يشاكل الحمل من البلدان والكور بلاد قرطانيا وفرسطانيا}
\MsL{19}{AB03-F023v-L19}{وجرمانيا وعلاطيا [[أ]]و بلاد اليهود، سرررنا العميقة. فصارت بمشاركة هذه البلدان}
\MsL{20}{AB03-F023v-L20}{مع الحمل للمرّيخ وذلك يشاكل الثور بلاد فارتيا ومعديا وفارس وفعلادس وسواحل}
\MsL{21}{AB03-F023v-L21}{آسيا الصغرى وقبرس وصارت بمشاكلتها مع الثور للزهرة والتي تشاكل الجوزاء}
\begin{PNotes}
\PNote{13}{AB03-n013}{B: add. \RL{\ArN على سبّ}}
\end{PNotes}
\Folio{024r}{24r}
\MsL{1}{AB03-F024r-L01}{بلاد جربان وطبرستان وأرمينية والقيروان ومرماريقي التي أسفل بلاد مصر وفورينا.}
\MsL{2}{AB03-F024r-L02}{وصارت بمشاكلتها الجوزاء لعطارد والذي يشاكل السرطان يوميذىا وفارصرونيا وإفريقية}
\MsL{3}{AB03-F024r-L03}{وسوسا وفرحيا. وصارت بمشاكلتها السرطان للقمر والذي يشاكل الأسد إيظاليا}
\MsL{4}{AB03-F024r-L04}{\{وعاليلواىولىا\} وصقلّية وىوتيقا \{وحلدىىا\}. وصارت بمشاكلتها مع الأمر للشمس أيضًا والذي}
\MsL{5}{AB03-F024r-L05}{يشاكل السنبلة بابل والجزيرة وبلاد اىور أو الأمن \{واحاىىا\} واقريطش ويشاكل}
\MsL{6}{AB03-F024r-L06}{عطارد أيضًا والذي يشاكل الميزان ينظر بالي وفاسعيرى وسريقي وساوس وواسيلس}
\MsL{7}{AB03-F024r-L07}{\{وطوعلودوطىىعى\} ويشاكل الزهرة أيضًا مع الميزان الذي يشاكل العقرب \{باطاعونيطس\}}
\MsL{8}{AB03-F024r-L08}{\{وماوردىطاىىا\} \{وعاطولىا\} وسوريا وقوماس \{وفافا\} \{حرٯما\} والحجاز ويشاكل المرّيخ أيضًا}
\MsL{9}{AB03-F024r-L09}{والذي يشاكل القوس \{طرسوسا\} \{وىلطىٯى\} \{وسفاىىا\} وبلاد الأعراب العامرة أيضًا. وصارت}
\MsL{10}{AB03-F024r-L10}{مشاكلتها للمشتري والذي يشاكل الجدي بلاد الهند وأريناي، \{ادىا\} الصين \{وحدروسي\}}
\MsL{11}{AB03-F024r-L11}{وبراقى ماقادونيا \{واىللودس\} ويشاكل زحل والذي يشاكل الدلو \{سرماطىٯى\} \{وطلساي\}}
\MsL{12}{AB03-F024r-L12}{\{والسعد\} \{وارانياىا\} \{ىىوىا\} الوسطى يعني بلاد الحبش الوسطى وتشاكل زحل أيضًا}
\MsL{13}{AB03-F024r-L13}{والتي تشاكل الحوت \{ىاراىىا\} \{وٯعلاعرىىا\} \BF{B5v}\{وىاسيسامرىىطس\} وعارى ما طبقى \{ولوريا\}}
\MsL{14}{AB03-F024r-L14}{\{وفلىعىا\} \{وقالعولىا\} ويشاكل المشتري مع الحوت فجميع هذه ما يميّزه\PN{14} وتسعون قطرًا ممّا}
\MsL{15}{AB03-F024r-L15}{في الأربع المذكورة. وممّا يلي وسط العمران منها ويحبّ أن تضيف إلى ذلك أيضًا إنّ كلّ}
\MsL{16}{AB03-F024r-L16}{واحد من الكواكب الثابتة يشاكل البلدان التي يشاكلها أجزاء من تلك البروج وتؤثر}
\MsL{17}{AB03-F024r-L17}{جميعها فيها وهي الأجزاء التي تحاذيها الكواكب وتشامتها على دائرة واحدة من الدوائر}
\MsL{18}{AB03-F024r-L18}{التي تمرّ بقطبي تلك البروج وبيّن ذلك من عروض الكواكب إذا كان مجازها فوق}
\MsL{19}{AB03-F024r-L19}{سمت الرؤس وهذا ذكر ما وصفنا في المدن مدائن الحجاز للعقرب مكّة للميزان وزحل}
\MsL{20}{AB03-F024r-L20}{مدينة الرسول عليه السلام\PN{15} للعقرب والزهرة الربذة للثور والقمر التغلبيّ للقوس}
\MsL{21}{AB03-F024r-L21}{وعطارد ومدائن اليمن للقوس صنعاء للجوزاء\PN{16} أو عطارد الجرين للحوت والمرّيخ عمان}
\begin{PNotes}
\PNote{14}{AB03-n014}{B: \RL{\ArN ما يئتان}}
\PNote{15}{AB03-n015}{B: \RL{\ArN سلّم}}
\PNote{16}{AB03-n016}{B: \RL{\ArN للعور}}
\end{PNotes}
\Folio{024v}{24v}
\MsL{1}{AB03-F024v-L01}{للثور والزهرة مدائن العراق للأسد البصرة للحوت وعطارد الكوفة للعذراء والمشتري}
\MsL{2}{AB03-F024v-L02}{بغداد للثور وعطارد وسط\PN{17} للحوت والمشتري دسكرة الملك للأسد والمشتري}
\MsL{3}{AB03-F024v-L03}{المدائن للأسد والمشتري السوس للرأس والقمر مدائن \{ماه\} للثور حلوان للثور}
\MsL{4}{AB03-F024v-L04}{والقمر الدينور للرأس والشمس. نهاوند للعقرب والمرّيخ همذان للثور والزهرة}
\MsL{5}{AB03-F024v-L05}{الرئيسيّ للثور وزحل. نهاوند للرأس والقمر الديلم للعقرب والقمر طبرستان للجدي}
\MsL{6}{AB03-F024v-L06}{والقمر ساريّة طبرستان للحمل والزهرة مدينة جرجان للرأس والمشتري قومس للجدي}
\MsL{7}{AB03-F024v-L07}{والقمر أصفهان للرأس والشمس نيسابور للحمل والمشتري فسا للرأس والشمس أمّا ورد}
\MsL{8}{AB03-F024v-L08}{للميزان وعطارد طوس للحمل وعطارد \{مرخس\} للعذراء والقمر مرّ والرود للجوزاء وزحل}
\MsL{9}{AB03-F024v-L09}{الطالقان للميزان والمرّيخ مرد وللحوت وزحل بلخ للميزان والمرّيخ أحل خراسان}
\MsL{10}{AB03-F024v-L10}{للثور والقمر \BF{B6r}خوارزم للتوأمين وزحل. سمرقند للحمل وعطارد \{الدسقوته\} للميزان والمرّيخ}
\MsL{11}{AB03-F024v-L11}{الشاش للعذراء والشمس فرغانة للثور والمشتري قرميان للتوأمين والزهرة}
\MsL{12}{AB03-F024v-L12}{\{يهله\} للسرطان وزحل \{تاوعيس\} للتوأمين وعطارد سجستان للسرطان والمشتري}
\MsL{13}{AB03-F024v-L13}{كابل للحمل والمشتري داورد للتوأمين والعقرب زابل ستان للحمل والقمر مدينة ثبت}
\MsL{14}{AB03-F024v-L14}{للجدي والشمس مكران للجوزاء وزحل مدائن الهند للجدي شامل للميزان والجدي}
\MsL{15}{AB03-F024v-L15}{\{قودين\} للقوس والقمر مالود للجدي والمرّيخ. الهند مار للتوأمين والقمر قارورة للحوت}
\MsL{16}{AB03-F024v-L16}{والشمس طبريّة للحمل والمرّيخ طعمة للجدي والمشتري مدائن الصين اسعير للأسد}
\MsL{17}{AB03-F024v-L17}{والقمر كسعير للثور والقمر مدينة الصين للأسد والزهرة مدائن أثور وهو بلد}
\MsL{18}{AB03-F024v-L18}{الموصل للسنبلة نينوى للجدي والمشتري أربيل للعقرب والمشتري الحومة}
\MsL{19}{AB03-F024v-L19}{للعقرب والقمر \{برسيار\} للحوت والزهرة أليس للرأس وزحل تكريت للعذراء}
\MsL{20}{AB03-F024v-L20}{وزحل الكرخ للميزان والمرّيخ مدائن أذربيجان أردبيل للثور والزهرة}
\MsL{21}{AB03-F024v-L21}{ىٯدى سابور للتوأمين وعطارد فسا للميزان والزهرة مدينة كرمان للميزان والمرّيخ}
\begin{PNotes}
\PNote{17}{AB03-n017}{B: \RL{\ArN واسط}}
\end{PNotes}
\Folio{025r}{25r}
\MsL{1}{AB03-F025r-L01}{كرمان التي سرّ كرمان للدلو والمرّيخ مدائن الجزيرة للسنبلة والعقرب وعطارد}
\MsL{2}{AB03-F025r-L02}{نصيبين للأسد والمشاري سنجار للرأس والمشتري قزدين للسرطان والمشتري كاري}
\MsL{3}{AB03-F025r-L03}{للميزان والمشتري رأس العين للدلو الحوت تلّ موزن للعذراء وزحل آمد للدلو والقمر}
\MsL{4}{AB03-F025r-L04}{ميافارقين للسرطان والزهرة إرزن للتوأمين وعطارد الرقة للقوس والقمر قرقيسيا}
\MsL{5}{AB03-F025r-L05}{للعقرب والزهرة غانة للميزان والمشتري جيب \BF{B6v}للجدي والمشتري حران للتوأمين}
\MsL{6}{AB03-F025r-L06}{وعطارد الرها للأسد والشمس مدائن الباب مدينة باب الخزر للرأس وزحل}
\MsL{7}{AB03-F025r-L07}{مدينة صول للرأس وزحل ردعه للدلو وزحل آمد للدلو والقمر قبله للثور والمرّيخ}
\MsL{8}{AB03-F025r-L08}{كسل للعقرب والشمس \{حزران\} للأسد والشمس باب اللآن للميزان راسان للدلو}
\MsL{9}{AB03-F025r-L09}{وزحل بارقة للميزان اسيل للعذراء والزهرة الآن للحوت والمشتري \{هيل العين\} للحوت}
\MsL{10}{AB03-F025r-L10}{والمشتري رسل للعذراء وزحل أسري السرطان والمرّيخ ناليغلًا للجدي وعطارد ميلان}
\MsL{11}{AB03-F025r-L11}{للتوأمين والمرّيخ \{ماطيه\} للدلو والقمر قيساريّة للحوت وزحل سميساه للحوت وعطارد}
\MsL{12}{AB03-F025r-L12}{\{طراوىنه\} وبلوانها للحوت وزحل \{ٯدروٯىىه\} للجدي والقمر ساطلا للحوت وزحل طرسوس}
\MsL{13}{AB03-F025r-L13}{للميزان والمرّيخ الصيصة الأسد والقمر ادية للقوس والمشتري جرجرة للثور أنقرة للرأس}
\MsL{14}{AB03-F025r-L14}{والقمر سلوقيّة للسرطان والشمس \{حررىىى\} للميزان وعطارد ومويّا للدلو والقمر}
\MsL{15}{AB03-F025r-L15}{أفيوش للعقرب والقمر سابوس للدلو وزحل كاوس الجزيرة للسرطان والشمس}
\MsL{16}{AB03-F025r-L16}{دوحس الجزيرة للحوت والمرّيخ \{حلقدوىسه\} للدلو وعطارد \{ياقدروميه\} للثور وزحل}
\MsL{17}{AB03-F025r-L17}{أو قمر للميزان والقمر \{سعىه\} للحوت والزهرة \{عبوريه\} للحوت والزهرة سميساط للميزان}
\MsL{18}{AB03-F025r-L18}{والمرّيخ مدائن \{صوس\} للثور ىاىىوس للعذراء والمرّيخ امنوس للحمل والقمر قسطنطينيّة}
\MsL{19}{AB03-F025r-L19}{قبرس للثور وعطارد مدائن سوريّة للعقرب اللاذقيّة للعذراء والشمس صور للتوأمين}
\MsL{20}{AB03-F025r-L20}{والمشتري اطرابلس للعقرب وعطارد أنطاكيّة للحمل والمشتري حرقة للحوت وعطارد منبج \BF{B7r}}
\MsL{21}{AB03-F025r-L21}{للثور والشمس حلب للعقرب وعطارد دمشق للعقرب والشمس بيت جدي للسرطان}
\Folio{025v}{25v}
\MsL{1}{AB03-F025v-L01}{والقمر عسقلان للأسد والشمس صيد اللرأس والشمس معرة مصرين للعذراء وعطارد}
\MsL{2}{AB03-F025v-L02}{تدمر للعقرب والمرّيخ طبرية للحوت والزهرة نابلس للثور والزهرة بيت المقدس للسرطان والمرّيخ}
\MsL{3}{AB03-F025v-L03}{سحره للثور والمشتري قسمة لوبيرومي أرض مصر السفلى فراس للسرطان والمرّيخ الإسكندريّة}
\MsL{4}{AB03-F025v-L04}{للأسد والمرّيخ \{مند\} للتوأمين والشمس عين شمس للتوأمين وعطارد أسوان هذه أدلة للميزان}
\MsL{5}{AB03-F025v-L05}{وزحل و \{ىري\} للثور والقمر مرسله للعقرب وعطارد أرماط للتوأمين والمشتري}
\MsL{6}{AB03-F025v-L06}{مدائن البريدة للسرطان قرطبة للثور والمشتري \{ىىنٯا\} للتوأمين والقمر مدائن إنطابلس}
\MsL{7}{AB03-F025v-L07}{للتوأمين إفريقية للحمل والقمر أرسياني للدلو وزحل \{بطلىاىس\} للثور والمشتري ابلونيا}
\MsL{8}{AB03-F025v-L08}{للحوت والشمس فروس للحوت والقمر مراقية للسرطان والشمس مدائن إفريقية للسرطان}
\MsL{9}{AB03-F025v-L09}{سرّ للثور وعطارد \{ىىرطوبة\} للسرطان والشمس افرطعينه للرأس والمرّيخ لانطس}
\MsL{10}{AB03-F025v-L10}{الكبرى للحمل والمشتري سفر للأسد والقمر القيروان للعقرب والشمس مدائن طنجة}
\MsL{11}{AB03-F025v-L11}{للعقرب مدينة طنجة للدلو وزحل قيساريّة طنجة للحوت والزهرة وإليه للحمل والقمر}
\MsL{12}{AB03-F025v-L12}{لو يحين للسرطان والمشتري سوس طنجة للعذراء وللقمر روحا للسرطان والمشتري}
\MsL{13}{AB03-F025v-L13}{قسمه نواحيّ قسطنطينيّة مدينة الملك للسرطان والمرّيخ روميّة العظمى للأسد والشمس}
\MsL{14}{AB03-F025v-L14}{جزيرة يول للحمل وعطارد جزيرة \{ابيرسا\} للعقرب وزحل \{باي\} للرأس والمرّيخ}
\MsL{15}{AB03-F025v-L15}{مدائن الأندلس للقوس باريا قيساريا \BF{B7v}للعذراء أو المرّيخ إيحاء وسطًا للحوت والشمس}
\MsL{16}{AB03-F025v-L16}{البليس لليومان وعطارد قرطوبة للسرطان والمشتري أسطوريا \{ارغسطا\} للأسد}
\MsL{17}{AB03-F025v-L17}{والمشتري طرقونة للعقرب والقمر قلوبيّا للأسد والقمر ما رينيون للعذراء أو زحل}
\MsL{18}{AB03-F025v-L18}{كريبًا الجزيرة للسرطان والزهرة مسأليّة للميزان وعطارد \{ىاىوىىا\} للميزان والمقر}
\MsL{19}{AB03-F025v-L19}{مدائن جرمانيّة للحمل لوقة للتوأمين والمرّيخ \{دىاورىا\} للعذراء والمرّيخ سهند [ر]}
\MsL{20}{AB03-F025v-L20}{للجدي والمشتري \{بعطليون\} للعقرب وزحل \{إزالاطى\} للحمل والشمس \{اىاىا\} للحوت}
\MsL{21}{AB03-F025v-L21}{والمرّيخ سوبر طيون للحمل والمرّيخ \{لابارزطا\} للحمل وزحل سمريون للجدي والمرّيخ}
\Folio{026r}{26r}
\MsL{1}{AB03-F026r-L01}{سندون للعذراء والمرّيخ مدائن أنطالية أوّلها رومة وقد ذكرناها للعقرب وزحل}
\MsL{2}{AB03-F026r-L02}{طراقه للحمل والمرّيخ \{ماقلس\} للثور والزهرة رأىسا للرأس والشمس بليدس للرأس وزحل}
\MsL{3}{AB03-F026r-L03}{مدائن جزيرة ورأس الأدنى للتوأمين والقمر \{هاريني\} للعذراء والقمر مدائن}
\MsL{4}{AB03-F026r-L04}{جزيرة سردانيّة سالقو للسرطان وزحل واداس للثور والزهرة طورواس للحوت.}
\MsL{5}{AB03-F026r-L05}{والمرّيخ مدائن جزيرة صقلّيّة ىسىىدي الأسد والقمر \{اعا\} للسرطان والمشتري}
\MsL{6}{AB03-F026r-L06}{سوس للدلو والمشتري \{طرىىقي\} للدلو وعطارد سور اقريسا للحوت وعطارد}
\MsL{7}{AB03-F026r-L07}{قطاي للعقرب والمشتري باجور للجدي وزحل أسوس للثور وزحل لاوس للأسد}
\MsL{8}{AB03-F026r-L08}{والمرّيخ ليوفيتا للجدي وعطارد واربيوس للعقرب وعطارد \{يرسما\} ينسون}
\MsL{9}{AB03-F026r-L09}{للجدي وعطارد ديراقولسوس للسرطان والقمر رينون للأسد والزهرة أسالونيقي}
\MsL{10}{AB03-F026r-L10}{الأسد والشمس \{بالاهي\} مدينة الإسكندريّة للثور وعطارد ولارينا للعذراء وعطارد}
\MsL{11}{AB03-F026r-L11}{قساير وللعقرب والقمر طباس للثور وعطارد \BF{B8r}فارقور للعذراء وعطارد والمشتري}
\MsL{12}{AB03-F026r-L12}{\{ىفافالاما\} للحوت وزحل. وجدنا في هذه النسخة التي للمدائن مشاكلة الكور البروج}
\MsL{13}{AB03-F026r-L13}{على هذا الرسم الذي أذكر وترى أنّ هذه النسخة ممّا عني بتفسيرها وكأنّه يخالف}
\MsL{14}{AB03-F026r-L14}{بطلميوس في بعض الكور ويميل إلى بعض رأي المصريين في قسمة البلدان التي تشاكل}
\MsL{15}{AB03-F026r-L15}{الحمل من الكور فارس وأذربيجان وبريطانية وطعنه وجرمانية وفلسطين وبعض}
\MsL{16}{AB03-F026r-L16}{البلقا وتشاكل مع الحمل والمرّيخ والذي يشاكل الثور جميع بلاد ماه وأصبهان وما صغر}
\MsL{17}{AB03-F026r-L17}{من جزيرة بحر الروم وقبرس وأثار تشاكل مع الثور والزهرة وتشاكل الجوزاء \{حلان\}}
\MsL{18}{AB03-F026r-L18}{والديلم وجرجان وطبرستان وأرمينية الكبرى ومرم وأنطابلس ومراقيا ومصر}
\MsL{19}{AB03-F026r-L19}{وتشاكل مع الجوزاء عطارد ويشاكل السرطان أرض البربر وإفريقية وينبه إلى}
\MsL{20}{AB03-F026r-L20}{أرض الروم وفروقية ولاوقة ولاودية يشاكل مع السرطان القمر ويشاكل الأسد}
\MsL{21}{AB03-F026r-L21}{أرض الترك سمر ونطاليّا وقلقلية و\{انقبيه\} وأرض حمص ودمشق وسواد الكوفة}
\Folio{026v}{26v}
\MsL{1}{AB03-F026v-L01}{ويشاكل مع الأسد الشمس ويشاكل السنبلة أرض بابل ومواقيّة وأرض الموصل}
\MsL{2}{AB03-F026v-L02}{والجزيرة وأرض الرياس وإقريطية ويشاكل مع السنبلة عطارد ويشاكل الميزان}
\MsL{3}{AB03-F026v-L03}{أرض بخارا\LR{\Sg{〉}}ن\LR{\Sg{〈}} وحارستان وحمير وتبت وجزيرة سرداوس وبعض أرض الحبش الذي تدعي}
\MsL{4}{AB03-F026v-L04}{طراموية وصنعاء وأرض نجران وأرض كرمان ويشاكل مع الميزان الزهرة والذي يشاكل}
\MsL{5}{AB03-F026v-L05}{العقرب بلد الحجاز و\{امد\} وطنجة بأطالية وارزايّا والسماويّة وسريرا و\{قندو قية\}}
\MsL{6}{AB03-F026v-L06}{و\{ٯىلىىىعىه\} ويشاكل مع العقرب المرّيخ ويشاكل القوس ملطيقي والأندلس}
\MsL{7}{AB03-F026v-L07}{واليمن وجزيرة صقلّيّة التي في بحر الشام. ويشاكل مع القوس المشتري و\LR{\LatinF ]]}يشاكل مع}
\MsL{8}{AB03-F026v-L08}{القوس المشتري\LR{\LatinF [[} ويشاكل مع الجدي الهند والسوس ويكران وهراه وبراقيّة}
\MsL{9}{AB03-F026v-L09}{وسقالية وأنيلوس وتشاكل مع الجدي زحل وتشاكل الدلو سردا ونطيقي}
\MsL{10}{AB03-F026v-L10}{وهو بلخ والصغير وفرغانه و\{اننيما\} من البلقى ووسط الجمش ويشاكل \BF{B8v}مع الدلو}
\MsL{11}{AB03-F026v-L11}{زحل والذي يشاكل الحوت قاتومه وىٯلٯىوىىه وبسهوطس وحرفه و\{نبد لوقية\}}
\Chap{I.11}{<11> في معرفة الأحداث الكائنة في البلدان بالكسوفات}
\MsL{12}{AB03-F026v-L12}{ويشاكل مع الحوت المشتري. \Hd{\LR{\LatinF \Sg{〈}11\Sg{〉}}}\Hd{ في معرفة}\PN{18}\Hd{ الأحداث الكائنة في البلدان بالكسوفات}}
\MsL{13}{AB03-F026v-L13}{نبدأ بالفحص عن موضع الكسوف من فلك البروج ومن إلى المثلّثات هو والي أيّ النواحيّ}
\MsL{14}{AB03-F026v-L14}{نسب ذلك المثلّث ثمّ نعمد إلى البلدان والمدن التي شاكل كلّ الكسوف خاصّة}
\MsL{15}{AB03-F026v-L15}{من تلك المثلّثة على ما رسمنا وننظر هل يرى الكسوف في جميع تلك الأماكن أو بعضها.}
\MsL{16}{AB03-F026v-L16}{فإنّ الذي يظهر فيه الكسوف منها لا يحدث فيه شيء ونضع طالع وسط الكسوف وأوتاده}
\MsL{17}{AB03-F026v-L17}{في المواضع التي نعلم أنّه يرى بها من تلك البلدان أن نحسب الساعات وارتفاع القطب فيها}
\MsL{18}{AB03-F026v-L18}{ثمّ نستدلّ على زمان الحادث ومقدار مدّته من بعد ذلك بأن ننظر كم من ساعة معتدلة}
\MsL{19}{AB03-F026v-L19}{منذ طلّه الكسوف من ابتدائه إلى آخر الانجلاء في كلّ واحد من تلك المساكن. فإن كان}
\MsL{20}{AB03-F026v-L20}{الكسوف شمسيًّا دام الحادث بعدد الساعات سنينًا وإن كان قمريًّا دام بعددها شهورًا}
\MsL{21}{AB03-F026v-L21}{وإن كان موضع الكسوف في ناحية المشرق من الأفق علمنا أن ابتداء الحادث يكون في}
\begin{PNotes}
\PNote{18}{AB03-n018}{ink blot}
\end{PNotes}
\Folio{027r}{27r}
\MsL{1}{AB03-F027r-L01}{الثلث الأوّل من تلك المدّة وإنّ شدّة الأمر في الحادث يكون في ذلك الثلث الأوّل.}
\MsL{2}{AB03-F027r-L02}{وإذا كان في ناحية وسط السماء استدللنا أنّه يكون في الثلث الثاني من تلك المدّة وإذا}
\MsL{3}{AB03-F027r-L03}{كان في ناحية المغرب من الأفق كان في الثلث الآخر منها ونميّز موضع الكسوف من الأفق}
\MsL{4}{AB03-F027r-L04}{على ما أصف لك إن كان موضع الكسوف \BF{B9r}من الطالع إلى ناحية وسط السماء أقلّ من أربع}
\MsL{5}{AB03-F027r-L05}{ساعات زمانيّة فهي في ناحية المشرق. وإذا تجاوز الأربع ساعات إلي أن ينتهي إلى وتد}
\MsL{6}{AB03-F027r-L06}{السماء ويتجاوزه بساعتين زمانيتين فهو في ناحية وسط السماء وإذا تجاوز ذلك كان في}
\MsL{7}{AB03-F027r-L07}{ناحية المغرب وتعمل في ذلك على أنّه إذا كان في نفس الجزء الطالع كان في أوّل المدّة وإذا}
\MsL{8}{AB03-F027r-L08}{كان في جزء وسط السماء كان في وسط المدّة وإذا كان في نفس الجزء الغارب كان}
\MsL{9}{AB03-F027r-L09}{الحادث في آخر المدّة. وأمّا الزيادات والنقصانات الحدّ ربّه في الأحداث فإنّها تعرف من}
\MsL{10}{AB03-F027r-L10}{الاجتماعات والمقابلات التي تكون فيما بين ذلك إذا عرضت في المواضع التي هي سبب}
\MsL{11}{AB03-F027r-L11}{الحدثان أو في المواضع التي تشاركها في الشكل أعني مواضع الكسوف ومواضع الأوتاد}
\MsL{12}{AB03-F027r-L12}{والمواضع التي تشترك معها في الشكل ومن مواضع الكواكب المتحيّرة تفسيرها التي تلي التدبير}
\MsL{13}{AB03-F027r-L13}{إذا ظهرت مشرقة أو مغربة أو واقفة أو طالعة في أوّل الليل وكانت مشاركة في الشكل}
\MsL{14}{AB03-F027r-L14}{البروج الكسوف أو بروج الأوتاد التي هي سبب الحدثان وذلك إنّها كانت مشرقة أو واقفة}
\MsL{15}{AB03-F027r-L15}{كان عنها زيادة الأحداث. وإذا كانت معربة أو تحت الشعاع وكانت طالعة أوّل الليل}
\MsL{16}{AB03-F027r-L16}{راجعة كان عنها نقصان الأحداث وتنظر في نقصان الكواكب العلويّة وتباتها في كلّ}
\MsL{17}{AB03-F027r-L17}{واحد من الأزمان أعني وقوف زحل والمشتري والمرّيخ إذا وقفت وكانت هي أصحاب}
\MsL{18}{AB03-F027r-L18}{الدلالة باستيلائها على مواضع الكسوفات والأوتاد على الجهة التي سنذكرها. وأمّا الجنس}
\MsL{19}{AB03-F027r-L19}{الذي يحدث فيه الحادث فيعلم من الأحوال والصور التي تخصّ البروج التي تعرض فيها}
\MsL{20}{AB03-F027r-L20}{الكسوفات والتي تكون فيها الكواكب المتحيّرة وغير المتحيّرة التي تدبّر البرج}
\MsL{21}{AB03-F027r-L21}{الذي يكون فيه الكسوف ويعلم المدبّر من الكواكب وبرج الوتد الذي قبل الكسوف وهذا الكوكب الذي}
\Folio{027v}{27v}
\MsL{1}{AB03-F027v-L01}{يجتاز به النيّر \BF{B9v}المنكسف قبل وقت الكسوف أمّا أن يكون هو الوتد الطالع أو وتد}
\MsL{2}{AB03-F027v-L02}{وسط السماء أيهما اتّفق أن يمر به قبل الكسوف. ويعلم المدبّر من الكواكب المتحيّرة لهذين}
\MsL{3}{AB03-F027v-L03}{الموضعين على ما أصف لك وهو أن تنظر إلى الكوكب الذي حصص كبيرة في هذين البرجين}
\MsL{4}{AB03-F027v-L04}{اللذين ذكرنا أمّا أن يكون مقارنًا لدرجة الكسوف أو درجة الوتد أو يصيّر إليها أو قديًّا}
\MsL{5}{AB03-F027v-L05}{\{ورها\} قريبًا وإمّا أن تشاركها في الشكل من أحد الاتّصالات أعني التثليث والتسديس}
\MsL{6}{AB03-F027v-L06}{والتربيع والمقابلة وما شاكل ذلك بحسب ما له من الحظوظ في البروج من أن يكون}
\MsL{7}{AB03-F027v-L07}{صاحب البيت والشرف والمثلّثة والحدود فإن اجتمع ذلك بكوكب واحد فهو صاحب}
\MsL{8}{AB03-F027v-L08}{التدبير وحده فإن لم يوجد كوكب واحد بعينه هو صاحب الكسوف وصاحب الوتد.}
\MsL{9}{AB03-F027v-L09}{فينبغي أن يؤخذ الكوكبان جميعًا الذين لهما مشاكلات أكثر إلى كلّ واحد من}
\MsL{10}{AB03-F027v-L10}{الموضعين على ما تقدّم من القول ونختار أوّلًا برج الكسوف ونقدّمه فإن كانت}
\MsL{11}{AB03-F027v-L11}{هذه الكواكب التي لها حصص في هذين الموضعين أكثر من كوكبين وكان ما لكلّ واحد}
\MsL{12}{AB03-F027v-L12}{منهما في هذين الموضعين مساويًّا الآخر فإنّا نختار في التدبير في نقدّم ما كان منها أوّل}
\MsL{13}{AB03-F027v-L13}{بأن يقال إنّه في الوتد أو كان أقوى فعلًا وكان ذلك الحيّز أقوى. وأمّا المدبّرة من}
\MsL{14}{AB03-F027v-L14}{الكواكب الثابتة فإنّا نعلمها مع ذلك أيضًا وبندأ أوّلًا بالكوكب الذي كان مع موضع}
\MsL{15}{AB03-F027v-L15}{الكسوف في الوتد الذي كان قبل الكسوف ونعني الكوكب الذي طلع مع جزء الكسوف}
\MsL{16}{AB03-F027v-L16}{أو توسط السماء مع جزء الكسوف إلى هذين الوتدين كان اجتياز أحد النيّرين}
\MsL{17}{AB03-F027v-L17}{به قبل كسوفه من الكواكب النيّرة نريد الكواكب العظام \BF{B10r}على الناس وعلى الملوك}
\MsL{18}{AB03-F027v-L18}{وكذلك أيضًا إن كانت على المواضع التي وصفها في المشرق خاصّة في وقت الكسوف}
\MsL{19}{AB03-F027v-L19}{دلّت على أنّ الحادث يكون الثمار وفي سنّ الشباب وفي الأساسات وما كان}
\MsL{20}{AB03-F027v-L20}{حدث السنّ وإن كانت في وسط السماء من فوق الأرض دلّت على أنّ الحادث}
\MsL{21}{AB03-F027v-L21}{يكون في الهياكل والملوك ومن كان سنة السنّ الأوسط وإن كانت في المغرب دلّت}
\Folio{028r}{28r}
\MsL{1}{AB03-F028r-L01}{على أنّ الحادث يكون في أمر السنين وتغييراتها وفي من كانت سنة من المشايخ و\{الموتى\}.}
\MsL{2}{AB03-F028r-L02}{وأمّا معرفة مقدار الجزء الذي يحدث فيه الحادث من الجنس الموضوع فيعلم من قبل مقدار}
\MsL{3}{AB03-F028r-L03}{ظله الكسوف ومن قياس الكواكب المتواليّة للتدبير إلى موضع الكسوفات وذلك أنّها}
\MsL{4}{AB03-F028r-L04}{إذا كانت في أشكال منسوبة إلى العشيّات عند كسوفات الشمس أو كانت في أشكال}
\MsL{5}{AB03-F028r-L05}{منسوبة إلى الغدوات عند كسوفات القمر فإنّ الأحداث يكون في أقلّ ذلك}
\MsL{6}{AB03-F028r-L06}{الجنس وإن كانت مقابلة على القطر كأنّ الحادث على النصف من ذلك الجنس وإن}
\MsL{7}{AB03-F028r-L07}{كانت في أشكال منسوبة إلى الغدوات عند كسوفات الشمس أو في الأشكال المنسوبة}
\MsL{8}{AB03-F028r-L08}{إلى العشيّات عند كسوفات القمر كان ذلك في أكثر ذلك الجنس ويعلم ما يحدث هل هو خير}
\MsL{9}{AB03-F028r-L09}{أم ضدّ ذلك من طبيعة الكواكب المدبّرة الرئيسة التي هي الفاعلة للحادث ومن امتزاج}
\MsL{10}{AB03-F028r-L10}{بعضها ببعض ومن المواضع التي تكون فيها وذلك أنّ الشمس والقمر أمران رئيسان}
\MsL{11}{AB03-F028r-L11}{لسائر الكواكب فيما يصيّر أنّ سبب جميع الفعل وسبب تدبير الكواكب وسبب قوّة}
\MsL{12}{AB03-F028r-L12}{الكواكب المدبّرة وضعفها. وأمّا النظر في امتزاج الكواكب المتولّيّة للتدبير بعضها مع}
\MsL{13}{AB03-F028r-L13}{بعض فهو الذي يدلّ على حال الأحداث التي تحدث فزحل إذا كان صاحب التدبير}
\MsL{14}{AB03-F028r-L14}{وحده فهو بالجملة يكون سبب فساد يكون عن برد. وأمّا فعله الجزئيّ الخاصّ به فإنّه}
\MsL{15}{AB03-F028r-L15}{إذا كان الحادث في الناس فهو أحداث الأمراض الطويلة \BF{B10v}والسلّ والذوبان وتأذ}
\MsL{16}{AB03-F028r-L16}{برطوبات وانصباب الفضول وحمايات الربع والهرب والحيرة والحبس والحزن}
\MsL{17}{AB03-F028r-L17}{والخوف والموت سيّما لمن قد طعن في السنّ. وإذا كان الحادث في البهائم التي يستعملها}
\MsL{18}{AB03-F028r-L18}{الناس كثيرًا كانت الآفة فيها قلّتها وأنّه يعرض لمّا كان منها موجودًا فساد أجسامها وأمراض}
\MsL{19}{AB03-F028r-L19}{تعرض فيها فيعرض بسبب ذلك لمن استعمل أكل هذه البهائم من الناس شبيه بهذه العلل.}
\MsL{20}{AB03-F028r-L20}{وإن كان الحادث في الهواء عرض برد شديد هع جمد وضباب وحر يحدث وباء وردأة}
\MsL{21}{AB03-F028r-L21}{الهواء وكثافة الغيوم وظلمة ودمق مفسد ويتولد من ذلك هوام كثيرة مؤذيّة لطبيعة}
\Folio{028v}{28v}
\MsL{1}{AB03-F028v-L01}{الناس. وأمّا إذا كانت الدلالة في الأنهار والبحار عرض فيما يشتركان فيه هواء}
\MsL{2}{AB03-F028v-L02}{شات وغرق السفن وصعوبة في ركوب الماء وقلّة السمك وفساده. وأمّا في الجار خاصّة}
\MsL{3}{AB03-F028v-L03}{فيعرض الجزر والمدّ وأمّا في الأنهار خاصّة فزيادة مفرطة وردأة المياه وإذا كانت}
\MsL{4}{AB03-F028v-L04}{الأحداث في الثمار الذي تنبت في الأرض فيعرض لها عوز وقلّة وفساد لا سيّما في التي}
\MsL{5}{AB03-F028v-L05}{يحتاج الناس إليها ضرورة والذي يعرض لها من الفساد يكون من دود البقل أو}
\MsL{6}{AB03-F028v-L06}{من الجراد وأمّا من الغرق وأمّا من امطار كثيرة وأمّا من برد وما أشبه ذلك حتّى}
\MsL{7}{AB03-F028v-L07}{يعرض من ذلك احيانًا الجدب فيهلك بسببه الناس\LR{\LatinF .} وأمّا المشتري فإنّه إذا كان}
\MsL{8}{AB03-F028v-L08}{صاحب التدبير وحده وهو بالجملة ينمى الأشياء وأفعاله الخاصّيّة فيما يحدث أمّا في الناس}
\MsL{9}{AB03-F028v-L09}{فرفعة المراتب وحسب الحال والخصب وحال السلامة والأمن ونمو الأشياء التي}
\MsL{10}{AB03-F028v-L10}{تحتاج إليها وحسن حال الأبدان والأنفس وجوائز وعطايًّا من الملوك وتزيد}
\MsL{11}{AB03-F028v-L11}{أمور الملوك وعظم شأنهم وكبر أنفسهم وهو بالجملة سبب السعادات. وأمّا في}
\MsL{12}{AB03-F028v-L12}{الدوابّ إن كان الحادث فيها فهو أن ما كان منها يستعمله الناس كثيرًا فإنّه \BF{B11r}يكثير}
\MsL{13}{AB03-F028v-L13}{ويزيد حرًّا وما كان منها على ضدّ ذلك يفسد ويهلك. وأمّا دلالته في الهواء فيكون}
\MsL{14}{AB03-F028v-L14}{معتدلًا فاعلًا للصحة كثير الرياح رطبًا وينمي الثمار التي تنبت في الأرض ويكون}
\MsL{15}{AB03-F028v-L15}{مسير السفن في البحر سليمًا ومدود الأنهار معتدلة وجميع ما أشبه هذه الأشياء وكلّما}
\MsL{16}{AB03-F028v-L16}{كانت الدلالة في الحادث في أحد الأنواع المسماة أصلحها وكان ما يحدث منه حسنًا}
\MsL{17}{AB03-F028v-L17}{معتدلًا موافقًا محمودًا عائدًا بالصلاح\LR{\LatinF . }وأمّا المرّيخ إذا كان صاحب التدبير وحده}
\MsL{18}{AB03-F028v-L18}{فهو بالجملة سبب الفساد الذي يكون من اليبس وأمّا فعالة الجزئيّة إن كان}
\MsL{19}{AB03-F028v-L19}{الحادث في الناس فتحدث فيهم حروب وقتن تقع فيما بين القوم وبين أهل}
\MsL{20}{AB03-F028v-L20}{جنسهم وسبي وعبوديّة وتوثب جملة الناس وغضب الرؤساء ويعرض لبعض الناس من}
\MsL{21}{AB03-F028v-L21}{هذه الأشياء الموت الذي تكون فجأة وأمراض تكون مع حمايات غبّ وقذف دم}
\Folio{029r}{29r}
\MsL{1}{AB03-F029r-L01}{وموت وحي عنيف وبخاصّة لمن كان في سنّ المنتهى وكذلك أيضًا يصاب بعض بالغضب}
\MsL{2}{AB03-F029r-L02}{والهوان والتعدّي والخروج عن الشريعة وحريق وقتل الرجال وسلب وقطع الطريق.}
\MsL{3}{AB03-F029r-L03}{وأمّا إن كان الحادث في الهواء فيعرض حرّ شديد ورياح حارّة ونابيّه مؤذيّة وسقوط}
\MsL{4}{AB03-F029r-L04}{الصواعق والصاعقة التي يقال لها فريطس وقلّة الأمطار وإذا كان في البحر}
\MsL{5}{AB03-F029r-L05}{عرض غرق السفن بغتة عن رياح مختلفة مضربة وعن صواعق وما أشبه ذلك}
\MsL{6}{AB03-F029r-L06}{من عدد وحرق وما يشاكله. وأمّا في الأنهار فيعرض نقصان الماء وجفاف العيون}
\MsL{7}{AB03-F029r-L07}{وفساد الماء المشروب وأمّا في الأشياء التي تصلح لاستعمال الناس من الحيوان}
\MsL{8}{AB03-F029r-L08}{والنبات فيعرض فيها قلّة وفساد الثمار التي يتولد عن إحراق \{احرًّا ىاها\}}
\MsL{9}{AB03-F029r-L09}{\BF{B11v}وأمّا عن الوباء من الجراد وأمّا عن انتشار يعرض لها من الرياح وأمّا عن احتراقها}
\MsL{10}{AB03-F029r-L10}{في الخزائن التي يخزن فيها. وأمّا الزهرة فإنّها بالجملة تحدث شبيهًا بما يحدث عن}
\MsL{11}{AB03-F029r-L11}{المشتري ويكون فيما يحدث عنها مع ذلك تحريك أمر من أمور الجماع واللهو وأمور}
\MsL{12}{AB03-F029r-L12}{النساء وفعلها الخاصّ الجزئيّ أمّا في الناس فرفع المراتب وسرور وحسن الحال وصلاح}
\MsL{13}{AB03-F029r-L13}{أمر التزويج وكثرة الأولاد وحمد جميع أمور الازدواج وتزيد الملك والتدبير الظاهر}
\MsL{14}{AB03-F029r-L14}{التقيّ وحسن المذهب في أمور الدين وخصب الأبدان والاتّصال بالرؤساء والمصاهرة}
\MsL{15}{AB03-F029r-L15}{إلى الولاة\LR{\LatinF . }وأمّا في حال الهواء فتعرض رياح معتدلة رطبة كثيرة الغذاء ويكون الهواء}
\MsL{16}{AB03-F029r-L16}{جيّدًا صافيًّا وأمّا المياه المولّدة والأمطار فتجعلها غزيرة وأمّا السفن فتسلّم}
\MsL{17}{AB03-F029r-L17}{في البحر ويحدث عنها الإصابات في التجارات والأرياح ومدود\{ا\} لأنهار وامتلاؤها.}
\MsL{18}{AB03-F029r-L18}{وأمّا الحيوان الذي يستعمله الناس وثمار الأرض خاصّة فإنّه يكثر ويخصب يتنفع}
\MsL{19}{AB03-F029r-L19}{به وكلما وقع الحادث فيه من نوع من هذه الأنواع كانت حاله حالًا محمودة نافعة}
\MsL{20}{AB03-F029r-L20}{للناس وأمّا عطارد فإنّه إذ كان صاحب التدبير وحده وهو بالجملة إذا مازج}
\MsL{21}{AB03-F029r-L21}{كلّ واحد من الكواكب الباقيّة \{شاكل طبائعها\}\PN{19} وأمّا فعله الخاصّيّ فهو التحريك للأشياء}
\begin{PNotes}
\PNote{19}{AB03-n019}{ink blot}
\end{PNotes}
\Folio{029v}{29v}
\MsL{1}{AB03-F029v-L01}{أكثر من سائر الكواكب وأمّا فعله في الأحداث التي تحدث عن الناس فهو حادّ}
\MsL{2}{AB03-F029v-L02}{قويّ لطيف الحيلة في الشيء الذي يقصد له صاحب قطع طريق وسرقات وتلصص}
\MsL{3}{AB03-F029v-L03}{وغارات ويحدث أيضًا صعوبة في ركوب البحر عند مشاركته في الشكل للكواكب}
\MsL{4}{AB03-F029v-L04}{المنحسة وهو سبب الأمراض اليابسة والحمّيات\PN{20} التي تنوب\PN{21} في كلّ يوم والسعال}
\MsL{5}{AB03-F029v-L05}{والقذف والسلّ ولأمور الكهانة ولعبادة الله تبارك وتعالى ولمّا يفرع في غلّات}
\MsL{6}{AB03-F029v-L06}{الهياكل \BF{B12r}وتغيير السنن والعادات في بعض الأوقات على قدم\PN{22} ما يلازم ممازحته}
\MsL{7}{AB03-F029v-L07}{لكلّ واحد من الكواكب ولمّا كان مزاجه يابسًا سريع الحركة لقربه من الشمس ولسرعة}
\MsL{8}{AB03-F029v-L08}{رجوعه في أدواره صار يحدث في الهواء خاصّة رياحًا مضطربة سريعة سهاه}
\MsL{9}{AB03-F029v-L09}{التغيير والرعد والصاعقة التي تقال لها فريطس والخسف والزلازل}
\MsL{10}{AB03-F029v-L10}{والبرق وربّما حدث بسبب هذه الأشياء التي ذكرت فسادًا في الحيوان}
\MsL{11}{AB03-F029v-L11}{والنبات الذي يستعمله الناس كثيرًا ويحدث عند أوقات غيبوبته نقصان في}
\MsL{12}{AB03-F029v-L12}{المياه والأنهار. وأمّا عند أوقات طلوعه فإنّه يزيد فيها وكل واحد من هذه الكواكب}
\MsL{13}{AB03-F029v-L13}{إذا كان على حاله الطبيعيّة الخاصّيّة له فهو يحدث أمثال هذه الأشياء. فإذا مزجت}
\MsL{14}{AB03-F029v-L14}{بعضها ببعض بحسب الاشتراكات في الأشكال وبحسب اختلاف البروج وبحسب}
\MsL{15}{AB03-F029v-L15}{ظهورها الكائن بقياسها إلى الشمس وبحسب افتراقها واجتماعها ثمّ من قبل الامتزاج}
\MsL{16}{AB03-F029v-L16}{أيضًا الذي يكون في الأفعال صار فعلها الخاصّيّ مختلف الأنواع على حسب ذلك}
\MsL{17}{AB03-F029v-L17}{وصار ممتزجًا من الطبائع شاركت طبيعته. ولمّا كان وصف الفعل الخاصّيّ الذي}
\MsL{18}{AB03-F029v-L18}{يكون من كلّ واحد من أنواع الامتزاجات\PN{23} لا نهاية له وكان غير ممكن وكان بالجملة}
\MsL{19}{AB03-F029v-L19}{ذكر جميع ذلك في أيّ الأشكال اتّفق أن يكون قد يفهم على إيحاء كبيرة فقد يحبّ أن}
\MsL{20}{AB03-F029v-L20}{نترك هذا النوع للناظر في ذلك ليكون هو الذي يفحص عن الأشياء الجزئيّة}
\MsL{21}{AB03-F029v-L21}{التي فيه وينبغي أن تنظر\PN{24} في مشاكلة البلدان والمدن لتلك الكواكب التي عنها}
\begin{PNotes}
\PNote{20}{AB03-n020}{B: \RL{\ArN الحمايات}}
\PNote{21}{AB03-n021}{B: \RL{\ArN تبوب}}
\PNote{22}{AB03-n022}{B: \RL{\ArN قدر}}
\PNote{23}{AB03-n023}{B: \RL{\ArN الامتجازات}}
\PNote{24}{AB03-n024}{B: \RL{\ArN ننظر}}
\end{PNotes}
\Folio{030r}{30r}
\MsL{1}{AB03-F030r-L01}{يكون\PN{25} الحوادث التي تعرض فيها وذلك أنّه إذا كانت تلك الكواكب مفيدة وشاكلة\PN{26}}
\MsL{2}{AB03-F030r-L02}{المواضع التي تحدث الأحداث فيها ولم يستعل عليها الكواكب والمنحسة المضادة لها}
\MsL{3}{AB03-F030r-L03}{صار فعلها الخاصّ بطبائعها\PN{27} أكثر نفعًا كما إنّها إذا لم تشاكل تلك البلدان واستعلت}
\MsL{4}{AB03-F030r-L04}{عليها الكواكب والمنحسة كان نفعها\PN{28} \BF{B12v}أقلّ وإذا كان مزاجها منحسًا ضارًّا وكانت هي}
\MsL{5}{AB03-F030r-L05}{أصحاب الفعل الحادث فإنّها إن كانت مشاكلة لتلك البلدان التي تحدث فيها الأحداث}
\MsL{6}{AB03-F030r-L06}{واستعلت عليها الكواكب المسعدة كان ضررها يسيّرًا وإن لم تكن مشاكلة للبلدان ولم}
\MsL{7}{AB03-F030r-L07}{تستعل عليها الكواكب المسعدة كان تأثير مزاجها المفسد فيها أكثر. وقوله تستعلي}
\MsL{8}{AB03-F030r-L08}{عليها الكواكب يُجوز على معان وجهات كثيرة منها أن يكون الكوكب مقارنًا للكوكب}
\MsL{9}{AB03-F030r-L09}{فيكون أكثرها عرضًا في الشمال هو المستعلي على صاحبه فإن كانا في الجنوب}
\MsL{10}{AB03-F030r-L10}{فأقلهما عرضًا هو المستعلي ويكون أحدهما في الشمال والآخر في الجنوب فالشمال}
\MsL{11}{AB03-F030r-L11}{هو المستعلي على الذي في الجنوب ومنها أن يشاكل الكوكب للكوكب من شكل التربيع}
\MsL{12}{AB03-F030r-L12}{فالذي في البروج المتقدّمة هو المستعلي وذلك إنّه يكون في البرج العاشر من}
\MsL{13}{AB03-F030r-L13}{البرج الذي هو فيه الآخر وهو وسط سماء ذلك البرج ومنها أن يكون الاستعلاء}
\MsL{14}{AB03-F030r-L14}{على الجهة التي ذكرناها في صدر الكتاب فيكون أحدهما في البروج العالية والآخر}
\MsL{15}{AB03-F030r-L15}{في البروج المنخفضة وبعدهما من بعد نقطتي الاعتدالين بعدًا واحدًا ومنها أن يكون}
\MsL{16}{AB03-F030r-L16}{أحدهما في وتد والآخر زائل عن الوتد أو يكون أحدهما فوق الأرض والآخر تحت الأرض}
\MsL{17}{AB03-F030r-L17}{وبالجملة إذا كانت السعود تنظر إلى النحوس من أيّ الأشكال كانت تنقص من}
\MsL{18}{AB03-F030r-L18}{شرها وإذا نظرت النحوس إلى السعود نقضت من نفعها لا سيّما إذا كانت من}
\MsL{19}{AB03-F030r-L19}{التربيع والمقابلة وأكثر ما يعني بقوله الاستعلاء أن يكون الكواكب تابعًا للكوكب}
\MsL{20}{AB03-F030r-L20}{يرتفع إلى مكانه ذلك وذلك إذا كان يطلع بعده فيزول الأوّل ويبقى الآخر}
\MsL{21}{AB03-F030r-L21}{مكانه فيستعلي عليه وفي أكثر الأمر إنّما تعرض من الأعراض العامّيّة للناس إذا}
\begin{PNotes}
\PNote{25}{AB03-n025}{B: \RL{\ArN تكون}}
\PNote{26}{AB03-n026}{B: \RL{\ArN وشاكلت}}
\PNote{27}{AB03-n027}{B: \RL{\ArN بطيائعها}}
\PNote{28}{AB03-n028}{B: \RL{\ArN [[نفعها]]}}
\end{PNotes}
\Folio{030v}{30v}
\MsL{1}{AB03-F030v-L01}{في المواضع الاضطراريّة في مواليدهم أعني\PN{29} المواضع التي فيها النيّران أو مواضع للأوتاد\PN{30}}
\MsL{2}{AB03-F030v-L02}{في \BF{B13r}المواليد وذلك إنّها هي مواضع الأشياء التي هي أسباب الأعراض العامّيّة التي أسبابها}
\MsL{3}{AB03-F030v-L03}{الكسوفات الكائنة في تلك المواضع أو في المواضع التي تقابلها وأصعب ما يكون ذلك وأبعده}
\MsL{4}{AB03-F030v-L04}{من أن يتحرّز منه إذا كانت أجزاء موضع الكسوفات الكائنة في تلك المواضع أو في المواضع}
\MsL{5}{AB03-F030v-L05}{التي تقابلها هي أجزاء مواضع النيّرين في المولد وبعدها الأجزاء المقابلة لها وأصعبها كسوف}
\MsL{6}{AB03-F030v-L06}{الشمس لأنّه إنّما يقع في تحويل السنة ضرورة وقد ينكسف القمر في موضعه في الأصل وقت}
\Chap{I.12}{<12> في ألوان الكسوفات والكواكب ذوات الأذناب}
\MsL{7}{AB03-F030v-L07}{التحويل. \Hd{\LR{\LatinF \Sg{〈}12\Sg{〉}}}\Hd{ في ألوان الكسوفات والكواكب ذوات الأذناب} وينبغي أن ينظر أيضًا}
\MsL{8}{AB03-F030v-L08}{في أمر الألوان التي تظهر في أوقات الكسوفات أعني ألوان النيّرين أنفسهما وألوان}
\MsL{9}{AB03-F030v-L09}{الأشياء التي حولها إذا حدثت حولهما مثل التي تسمّى القضبان والعصا والهالة وما}
\MsL{10}{AB03-F030v-L10}{أشبهها وذلك أنّها إذا كانت سوداءً غلاظًا ومائلًا إلى الخضرة دلّت على الأشياء التي}
\MsL{11}{AB03-F030v-L11}{ذكرنا أنّها من طبيعة زحل وإذا كانت بيضاء دلّت على الأشياء التي هي من طبيعة المشتري}
\MsL{12}{AB03-F030v-L12}{وإذا كانت مائلة إلى الحمرة دلّت على الأشياء التي هي من طبيعة المرّيخ وإذا كانت صهباء}
\MsL{13}{AB03-F030v-L13}{ذهبيّة دلّت على الأشياء التي من طبيعة الزهرة وإذا كانت مختلفة الألوان دلّت}
\MsL{14}{AB03-F030v-L14}{على الأشياء التي من طبيعة عطارد. فإن كان اللون في جميع جرم النيّرين أو في جميع}
\MsL{15}{AB03-F030v-L15}{المواضع التي حولهما فإنّ الحادث يكون في أكثر الأجزاء البلدان وإن كان في جزء}
\MsL{16}{AB03-F030v-L16}{من أجزاء النيّرين أو ممّا حولهما أيّ جزء وكان فإنّ الحادث يعرض \BF{B13v}في تلك الناحية}
\MsL{17}{AB03-F030v-L17}{التي يحاذيها\PN{31} ذلك الجزء الذي فيه اللون فقط. وينبغي أن تنظر أيّ في الكواكب}
\MsL{18}{AB03-F030v-L18}{ذوات الأذناب أيضًا التي تكون في أوقات الكسوفات أوفي غيرها من الأوقات}
\MsL{19}{AB03-F030v-L19}{وإلى سائر الأحداث الكلّيّة التي تحدث من\PN{32} التي تسمّى العصًا والبوقات والحرّاي\PN{33}}
\MsL{20}{AB03-F030v-L20}{وما أشبهها وتقيم طبائعها فيما تحدثه مقام ما قلنا إنّ طبيعة المرّيخ وطبيعة عطارد}
\MsL{21}{AB03-F030v-L21}{تحدثه من الحروب والإحراق والحرّ الشديد والزلازل والأعراض التي تحدث عنها}
\begin{PNotes}
\PNote{29}{AB03-n029}{B: \RL{\ArN \{أذعني\}}}
\PNote{30}{AB03-n030}{B: \RL{\ArN بالأوتاد}}
\PNote{31}{AB03-n031}{T: \RL{\ArN يحازيها}}
\PNote{32}{AB03-n032}{BT: \RL{\ArN مثل}}
\PNote{33}{AB03-n033}{T: \RL{\ArN الخوابئ}}
\end{PNotes}
\Folio{031r}{31r}
\MsL{1}{AB03-F031r-L01}{ويستدلّ على المواضع التي تحدث فيها الحوادث بمواضعها من فلك البروج أو بالمواضع}
\MsL{2}{AB03-F031r-L02}{التي يحاذيها سمت ذنب الكواكب ويستدلّ من صورة الحادث الذي يظهر على نوع الحادث}
\MsL{3}{AB03-F031r-L03}{الذي يحدث وعلى الجنس الذي يحدث فيه الحادث بمواضع ذلك الشيء يظهر من فلك}
\MsL{4}{AB03-F031r-L04}{البروج ومن زمان مكث الحادث الذي يظهر على مدّة الأحداث ومن قياسها إلى الشمس}
\Chap{I.13}{<13> في معرفة أحوال السنة}
\MsL{5}{AB03-F031r-L05}{على وقت ابتداء الحادث لأنّها في أكثر الأمر إذا كان ظهورها وهي منسوبة \LR{\Sg{〉}}\LR{\Sg{〉}} إلى الغدوات كان الحادث عنها سريعًا. وإذا كان ظهورها وهي منسوبة \LR{\Sg{〈}}\LR{\Sg{〈}}\PN{34} إلى العشيّات كان الحادث عنها بطيئًا \Hd{\LR{\LatinF \Sg{〈}13\Sg{〉}}}\Hd{ في معرفة أحوال السنة} وينبغي لنا أن نستعمل في}
\MsL{6}{AB03-F031r-L06}{تمييز السنة دخول الشمس أوائل البروج المنقلبة والمعتدلة أعني نقطة الاعتدالين}
\MsL{7}{AB03-F031r-L07}{والانقلابين التي هي مبادئ فصول أزمنة السنة وتنعقد\PN{35} الاجتماعات والمقابلات\PN{36}}
\MsL{8}{AB03-F031r-L08}{القريبة من هذه المبادئ الكائنة للشمس والقمر قبل هذه الأوقات وبخاصّة ما يحدث}
\MsL{9}{AB03-F031r-L09}{فيه منها كسوف حتّى يعلم من الابتداء الذي يكون لدخول الشمس الحمل كيف حال}
\MsL{10}{AB03-F031r-L10}{الربيع ومن الابتداء الذي يوجد\PN{37} من دخولها\PN{38} الميزان كيف حال الخريف ومن الابتداء}
\MsL{11}{AB03-F031r-L11}{الذي يوجد من دخولها السرطان كيف \BF{B14r}حال الصيف ومن دخولها الجدي كيف حال}
\MsL{12}{AB03-F031r-L12}{الشتاء إذا كانت هذه الأوقات الكلّيّة وحالاتها إنّما تكون عن الشمس. وينبغي}
\MsL{13}{AB03-F031r-L13}{أن نستعمل في ذلك خواصّ البروج ويستدلّ بها على الرياح التي تهب من الجهات}
\MsL{14}{AB03-F031r-L14}{وعلى جملة طبائعها و\{أمّا\} التغيّرات التي تكون إلى الزيادة والنقصان في بعض الأوقات}
\MsL{15}{AB03-F031r-L15}{فقد يستدلّ عليها أيضًا بالجملة بالاجتماعات والمقابلات الكائنة من بعد دخول}
\MsL{16}{AB03-F031r-L16}{الشمس أوائل هذه البروج وهي الاجتماعات والمقابلات التي تلي هذه النقط}
\MsL{17}{AB03-F031r-L17}{بمشاركة الكواكب المتحيّرة لها في الأشكال ويستدلّ على الأمور الجزئيّة التي}
\MsL{18}{AB03-F031r-L18}{تكون في أحد الأرباع بالاجتماعات التي تكون في كلّ واحد من البروج وبمواضع}
\MsL{19}{AB03-F031r-L19}{الكواكب وذلك يمكن أن يسمّى الفحص عن الشهور ويحتاج في هذا الباب مع ما قلنا}
\MsL{20}{AB03-F031r-L20}{إلى معرفة الخواصّ الطبيعيّة الجزئيّة التي تكون لكلّ واحد من البروج على ما يليق}
\begin{PNotes}
\PNote{34}{AB03-n034}{left margin}
\PNote{35}{AB03-n035}{B: \RL{\ArN تتفقّد}}
\PNote{36}{AB03-n036}{T: \RL{\ArN والاستقبالات} (wherever it occurs in this chapter)}
\PNote{37}{AB03-n037}{T: \RL{\ArN يؤخذ} (wherever it occurs in this chapter)}
\PNote{38}{AB03-n038}{T: \RL{\ArN كينونتها} (wherever it occurs in this chapter)}
\end{PNotes}
\Folio{031v}{31v}
\MsL{1}{AB03-F031v-L01}{بحالات السنة فأمّا مزاجات الكواكب المتحيّرة والثابتة التي مزاجاها\PN{39} شبيهة بمزاج}
\MsL{2}{AB03-F031v-L02}{تلك المتحيّرة مشاكلتها للهواء والرياح فقد ذكرناها فيما تقدّم وكذلك مشاكلة جميع}
\MsL{3}{AB03-F031v-L03}{البروج للرياح وإنّما بقي الأخبار بطبائعها الجزئيّة ومشاكلتها للأوقات أعني}
\MsL{4}{AB03-F031v-L04}{مشاكلة كلّ برج فالحمل على دائرة معدّل النهار وطبع جملة برج الحمل الاعتدال الذي فيه}
\MsL{5}{AB03-F031v-L05}{مولّد الرعد والبرق وما كان منه شماليًّا فهو حارّ مفسد وما كان منه جنوبيًّا فهو مولّد}
\MsL{6}{AB03-F031v-L06}{للجمد بارد والثور إلى الدائرة الصيفيّة أقرب وجملته\PN{40} برج الثور دالّة على المزاجين}
\MsL{7}{AB03-F031v-L07}{جميعًا مائل إلى الحرّ وما كان منه شماليًّا فهو معتدل المزاج وما كان منه جنوبيًّا فهو مولّد}
\MsL{8}{AB03-F031v-L08}{للاضطراب والزلازل. وأمّا بروج التوأمين فهو إلى الدائرة الصيفيّة أقرب وطبع}
\MsL{9}{AB03-F031v-L09}{جملته فاعلة لاعتدال المزاج وما كان منه شماليًّا فهو مولّد للرياح فاعل \BF{B14v}للزلازل وما}
\MsL{10}{AB03-F031v-L10}{كان منه جنوبيًّا فهو ناري محترق. وأمّا برج السرطان فهو على دائرة الصيف وجملته}
\MsL{11}{AB03-F031v-L11}{[وجملته] فاعله للصحر حارّة وما كان منه شماليًّا أو جنوبيًّا فهو ناري محرق. وأمّا}
\MsL{12}{AB03-F031v-L12}{برج الأسد فهو إلى الدائرة الصيفيّة أقرب وجملته محرق للهواء الحارّ}
\MsL{13}{AB03-F031v-L13}{الراكد والوباء وما كان منه شماليًّا فهو محرّك ناري يريد محرّك للرياح الحارّة وما}
\MsL{14}{AB03-F031v-L14}{كان منه جنوبيًّا فهو مرطّب. وأمّا برج العذراء فهو إلى دائرة الاعتدال أقرب}
\MsL{15}{AB03-F031v-L15}{وجملته مرطّبة محدثة للرعد وما كان منه شماليًّا فهو محدث للرياح وما كان منه}
\MsL{16}{AB03-F031v-L16}{جنوبيًّا فهو معتدل المزاج. وأمّا برج الميزان فهو على دائرة معدّل النهار}
\MsL{17}{AB03-F031v-L17}{وجملته مغيّرة منقلبة ما كان منه شماليًّا فهو مولّد للرمل وما كان منه جنوبيًّا}
\MsL{18}{AB03-F031v-L18}{فهو مرطّب محدث للوباء. وأمّا برج العقرب فإنّه إلى الدائرة الشتويّة}
\MsL{19}{AB03-F031v-L19}{أقرب وطبع جملته محدثة للرعد ناريّة وما كان منه شماليًّا فهو ناري محرق وما كان}
\MsL{20}{AB03-F031v-L20}{منه جنوبيًّا فهو مرطّب. وأمّا برج القوس فهو إلى الدائرة الشتويّة أقرب}
\MsL{21}{AB03-F031v-L21}{وجملته محدثة للرياح وما كان منه شماليًّا فهو محدث للرياح وما كان منه جنوبيًّا}
\begin{PNotes}
\PNote{39}{AB03-n039}{i.e. \RL{\ArN مزاجاتها}}
\PNote{40}{AB03-n040}{B: add. \RL{\ArN فاعلة}}
\end{PNotes}
\Folio{032r}{32r}
\MsL{1}{AB03-F032r-L01}{فهو مرطّب جدًّا مغيّر. وأمّا برج الجدي فإنّه على الدائرة الشتويّة وجملته مرطّبة}
\MsL{2}{AB03-F032r-L02}{جدًّا وما كان من شماليًّا أو جنوبيًّا فهو مرطّب مفسد جدًّا. وأمّا برج الدلو فهو إلى}
\MsL{3}{AB03-F032r-L03}{الدائرة الشتويّة أقرب وجملته باردة مائيّة وما كان منه شماليًّا فهو محرق وما كان منه جنوبيًّا}
\MsL{4}{AB03-F032r-L04}{يحدث الحريق. وأمّا برج السمكين فإنّه إلى دائرة معدّل النهار أقرب وجملته}
\MsL{5}{AB03-F032r-L05}{باردة مولّدة للرياح وما كان منه شماليًّا فهو مولّد للرياح وما كان منه جنوبيًّا فهو مائي}
\MsL{6}{AB03-F032r-L06}{وأمّا طبع أجزاء البروج فيعرف من مزاج الكواكب الثابتة التي في كلّ برج}
\MsL{7}{AB03-F032r-L07}{من الكواكب الثابتة التي في كلّ برج من الكواكب التي في صور فلك البروج والخارجة}
\MsL{8}{AB03-F032r-L08}{عنها بحسب \BF{B15r}أقدارها في العظم ومشاكلة طبائعها لطبائع الكواكب المتحيّرة وذلك}
\MsL{9}{AB03-F032r-L09}{أن ما كان من الكواكب السماويّة التي في فلك البروج مثل الثريّا التي مزاجها الأوّل من}
\MsL{10}{AB03-F032r-L10}{القمر ثمّ المرّيخ الثاني فإنّها تحدث الزلازل والرياح والضباب والسحابيّ الذي على}
\MsL{11}{AB03-F032r-L11}{صورة السرطان التي يسمّى العلف ومزاجه الأوّل مزاج المرّيخ ثمّ القمر الثاني فاعل}
\MsL{12}{AB03-F032r-L12}{للهواء الحارّ الراكد والزلازل والظلمة. وأمّا السحابيّ الذي يتلوا شولة العقرب}
\MsL{13}{AB03-F032r-L13}{ومزاجه مزاج الشمس ثمّ المرّيخ فاعل الزلازل وكذلك المرّيخ السحابيّ الذي على عين}
\MsL{14}{AB03-F032r-L14}{الرأس وما كان مزاجه مزاج المرّيخ ثمّ الدبران الذي هو في العظم الأكبر فإنّه ناري}
\MsL{15}{AB03-F032r-L15}{مولّد للبرق الصواعق. ومثل قلب الأسد الذي مزاجه أوّلًا مزاج المرّيخ ثمّ المشتري}
\MsL{16}{AB03-F032r-L16}{فيحدث الهواء الراكد الحارّ والوباء وما كان منها على مزاج الزهرة وحدّها فهو يدلّ}
\MsL{17}{AB03-F032r-L17}{على كثرة الرطوبة والأمطار وما كان منها على مزاج المشتري فهو يدلّ على الحرارة والرطوبة}
\MsL{18}{AB03-F032r-L18}{المعتدلة وما كان من مزاج زحل مفردًا فهو يدلّ على البرد وما كان على مزاج عطارد}
\MsL{19}{AB03-F032r-L19}{مفردًا فتارة يجفف ومرة يرطب ويولد الرياح على حسب ما قد ذكرنا من قوى}
\MsL{20}{AB03-F032r-L20}{الكواكب المتحيّرة وبمشاكلته للرياح ويمزج بعد ذلك طبائعها بطبائع البروج التي}
\MsL{21}{AB03-F032r-L21}{هي فيها وبطبائع أزمان السنة وذلك أنّ الكواكب الرطبة في ناحية الغرب وكذلك}
\Folio{032v}{32v}
\MsL{1}{AB03-F032v-L01}{الكواكب الحارّة في الصيف تكثر حرارتها لمشاكلتها للزمان وإذا كانت في الناحية}
\MsL{2}{AB03-F032v-L02}{الجنوبيّة كان أكثر حرها. وإذا كانت البرج التي هي فيها ملائمة لطبائعها كان ذلك}
\MsL{3}{AB03-F032v-L03}{أزيد في فعلها وكذلك إذا مالت في العرض إلى الشمال أو الجنوب كان ذلك ممّا يزيد أو}
\MsL{4}{AB03-F032v-L04}{ينقص من فعلها بحسب ملائمة طبائعها لتلك الجهة وإذا كان مزاج الكو[[ا]]كب}
\MsL{5}{AB03-F032v-L05}{منسوبًا إلى الكوكبين من المتحيّرة كان الأوّل منهما أقوى فعلًا أو يمزج الثاني معه}
\MsL{6}{AB03-F032v-L06}{من بعد على حسب الرسم وذلك أنّ منها ما يشبه مزاجه للكوكب الثاني شبهًا}
\MsL{7}{AB03-F032v-L07}{قليلًا وبحسب ما يتّصل من الكواكب المتحيّرة أحدهما بصاحبه إذا وليت التدبير}
\MsL{8}{AB03-F032v-L08}{في \BF{B15v}مشاكلة طبائعها بعضها ببعض ومضادتها حتّى تقف من جملة ما يجتمع من}
\MsL{9}{AB03-F032v-L09}{هذه الامتزاجات على الحال التي يكون وتغلب في القوّة ما كان أظهر فعلًا}
\MsL{10}{AB03-F032v-L10}{وأقوى مكانًا. وإذ قدّمنا وضع هذه الأشياء على هذه الجهة فإنّ مذاهب الدلائل}
\MsL{11}{AB03-F032v-L11}{الجزئيّة تكون على ما أصف أمّا أحد هذه المذاهب فهو النظر في جملة الأمر وهو الذي}
\MsL{12}{AB03-F032v-L12}{يعلم من النظر من أمر ربع من الأرباع وهو الذي يحتاج أن ينظر فيه ممّا قلنا}
\MsL{13}{AB03-F032v-L13}{ويتفقّد أمر أقرب الاجتماعات الشمس والقمر واستقبالهما الذي يكون قبل}
\MsL{14}{AB03-F032v-L14}{نقطة الانقلابين والاعتدالين وأن تضع الأوتاد وقت الاجتماع والاستقبال}
\MsL{15}{AB03-F032v-L15}{في كلّ واحد من الأقاليم والمدن التي يقصد للفحص عنها ثمّ تأخذ بعد ذلك الكواكب}
\MsL{16}{AB03-F032v-L16}{المدبّرة للمواضع التي يكون فيها الاجتماع أو التي يكون فيها الاستقبال المدبّرة للوتد}
\MsL{17}{AB03-F032v-L17}{التي يتلوا هذه المواضع وهو الوتد التي يجوز فيه أحد النيّرين قبل وقت الاجتماع}
\MsL{18}{AB03-F032v-L18}{والامتلاء قريبًا على ما أوضحنا وبيناه في باب الكسوفات حتّى يعلم المدبّر من الكواكب}
\MsL{19}{AB03-F032v-L19}{المتحيّرة وغير المتحيّرة على تلك الجهة حتّى نعلم جملة االأمر من خواصّ الأرباع ونعلم الزيادة}
\MsL{20}{AB03-F032v-L20}{والنقصان في شدّة تلك الجملة ولينها من طبيعة الكواكب المدبّرة بأن تنظر في كيفيتها}
\MsL{21}{AB03-F032v-L21}{من أيّ الكيفيات هي وأيّ الأحوال تتحرّك عنها والمذهب الثاني هو النظر في أموار}
\Folio{033r}{33r}
\MsL{1}{AB03-F033r-L01}{الشهور وهو الذي يحتاج أن ينظر فيه في الاجتماعات والاستقبالات التي تكون}
\MsL{2}{AB03-F033r-L02}{في كلّ واحد من البروج على الجهة المذكورة بعينها بعدان تكون حافظين لمّا إذا ذاكره}
\MsL{3}{AB03-F033r-L03}{وهو أنّه إذا وقع الاجتماع بالقرب من نقطة الانقلاب التي جاز والاعتدال\PN{41} استعملنا}
\MsL{4}{AB03-F033r-L04}{الاجتماعات إلى أن ينتهي إلى الربع التي يتلوا ذلك الربع وكذلك أيضًا ينبغي أن}
\MsL{5}{AB03-F033r-L05}{ينظر في الأوتاد وفي الكواكب المدبّرة للموضعين جميعًا الذين ذكرنا}
\MsL{6}{AB03-F033r-L06}{وأن نفصح\PN{42} خاصّة عن ظهور الكواكب القريبة \BF{B16r}وعن اتّصالات الكواكب المتحيّرة}
\MsL{7}{AB03-F033r-L07}{وانضر افاتها وعن خواصّها وحالاتها وأيّ الرياح تحرّك هذه الكواكب وأجزاء البروج}
\MsL{8}{AB03-F033r-L08}{التي يتّفق أن يكون فيها. وأيضًا قلنا ننظر إلى مهب أي ريح يميل عرض القمر إذا مال}
\MsL{9}{AB03-F033r-L09}{عن منطقة تلك البرج حتّى تقدّم فيعلم من جميع هذه الأشياء ما الغالب في أكبر الأمر}
\MsL{10}{AB03-F033r-L10}{من أحوال الشهور ورياحها. والمذهب الثالث وهو أن تتفقّد ألطف الدلائل التي تدلّ}
\MsL{11}{AB03-F033r-L11}{على فتور الأمر وصعوبته وذلك شيء قد يعلم من اشتراك الشمس والقمر في الأشكال}
\MsL{12}{AB03-F033r-L12}{الجزئيّة ليس في أشكال الاجتماع والاستقبال فقط لكن في الذي تكون عند اتصاف}
\MsL{13}{AB03-F033r-L13}{القمر في الضوء أيضًا ويكون ابتداء التفسير الذي قد استدلّ عليه في أكثر الأمر قبل ثلاثة}
\MsL{14}{AB03-F033r-L14}{أيّام من مسير القمر إلى الشمس وربّما كان ذلك بعد ثلاثة أيّام وذلك على بعد أجزاء}
\MsL{15}{AB03-F033r-L15}{متساويّة. ويعلم ذلك أيضًا من مشاركتها للشكل في الكواكب المتحيّرة في كلّ واحد}
\MsL{16}{AB03-F033r-L16}{من هذه الأحوال ومن غيرها مثل حال الثليث والتسديس وذلك أنّ هذه}
\MsL{17}{AB03-F033r-L17}{الأشياء تابعة للمجرى الطبيعيّ. وأمّا حال التغيير فإنّا ندرك معرفتها من اتباعنا}
\MsL{18}{AB03-F033r-L18}{للمشاكلات الطبيعيّة التي تكون الكواكب التي توجد والبروج للهواء المحيط وللرياح}
\MsL{19}{AB03-F033r-L19}{وقوّة هذه الحالات الجزئيّة التي تكون في يوم يكون خاصّة إذا كان للكواكب التي}
\MsL{20}{AB03-F033r-L20}{هي أعظم مقدارًا وأقوى فعلًا من الكواكب الثابتة ظهور بالغدوات أو بالعشيّات}
\MsL{21}{AB03-F033r-L21}{في المشرق أو في المغرب بقياس الشمس وذلك أنّها تغيّر وتقلب الأحوال الجزئيّة}
\begin{PNotes}
\PNote{41}{AB03-n041}{B: \RL{\ArN أو الاعتدال}}
\PNote{42}{AB03-n042}{B: \RL{\ArN نفحص}}
\end{PNotes}
\Folio{033v}{33v}
\MsL{1}{AB03-F033v-L01}{في أكثر الأمر إلى طبائعها وقد يكون ذلك أيضًا يسير بدون ما قلنا عند ممرّ النيّرين}
\MsL{2}{AB03-F033v-L02}{ببعض الأوتاد وذلك إنّ الفتور والشدّة التي تكون للحالات في وقت يكون}
\MsL{3}{AB03-F033v-L03}{عند مصيّر النيّرين إلى هذه المواضع كما أنّ الجزر والمدّ يكون عند مصيّر القمر إلى}
\MsL{4}{AB03-F033v-L04}{هذه وتغيّر\PN{43} الرياح خاصّة إنّما يكون عند مسير النيّرين إلى الأوتاد ويكون}
\MsL{5}{AB03-F033v-L05}{تغيّرها إلى الريح التي يميل إليها عرض القمر وقد \BF{B16v}ينبغي أن يعمل في كلّ موضع على أنّ}
\MsL{6}{AB03-F033v-L06}{السبب الأوّل الكلّيّ هوالمتقدّم وعلى أنّ التابع هو سبب الأعراض الجزئيّة وعلى}
\MsL{7}{AB03-F033v-L07}{أنّ الفعل يصح خاصّة ويقوى إذا وجدت الكواكب التي كانت المدبّرة للطبائع}
\Chap{I.14}{<14> في الدلائل التي تؤخذ من الآثار العلويّة}
\MsL{8}{AB03-F033v-L08}{الكلّيّة مشاركة للطبائع الجزئيّة في الشكل. \Hd{\LR{\LatinF \Sg{〈}14\Sg{〉}}}\Hd{ في الدلائل التي تؤخذ من الآثار العلويّة}}
\MsL{9}{AB03-F033v-L09}{قد ينتفع في تقدمة [[المعرفة]] الدلائل الجزئيّة بتقفّد العلامات التي تكون حول الشمس}
\MsL{10}{AB03-F033v-L10}{والقمر والكواكب وقتها وينبغي أن تتفقّد ما يعرض من ذلك الشمس في طلب الحالات}
\MsL{11}{AB03-F033v-L11}{التي تكون بالنهار إذا طلعت وفي طلب معرفة الدلالات\PN{44} التي تكون باليل إذا غابب}
\MsL{12}{AB03-F033v-L12}{وبتفقّد في معرفة الأحوال التي تطول مدّتها أوقات مشاركة الشمس لا إذا طلعت}
\MsL{13}{AB03-F033v-L13}{أو غابت وهي صافيّة من غير أن يكون شيء يسترها إلى القمر في الشكل إذا كان كلّ}
\MsL{14}{AB03-F033v-L14}{شكل من الأشكال يدلّ في أكثر الأمر على الحال التي تكون في ذلك الوقت إلى الوقت}
\MsL{15}{AB03-F033v-L15}{الذي ينتهي فيه القمر وإلى الشكل الذي يتلوه. وذلك أنّ الشمس إذا طلعت أو غابت}
\MsL{16}{AB03-F033v-L16}{وهي صافيّة من غير أن يكون شيء يسترها وهي نقيّة [[من]] بلا غيم دلّت على الصحو. فإن}
\MsL{17}{AB03-F033v-L17}{كانت دائرتها مختلفة الألوان أو كانت مائلة إلى الحمرة الناريّة أو كان الشعاع}
\MsL{18}{AB03-F033v-L18}{الذي ينبعث منها والذي حولها أحمر وكانت السحابات التي تسمّى شموسًا من جهة}
\MsL{19}{AB03-F033v-L19}{واحدة منها وكانت شبيهة بلون السحاب المائل إلى الحمرة وكان شعاع الشمس}
\MsL{20}{AB03-F033v-L20}{ممتدّا طويلًا دلّت على رياح شديدة وإن هذه الرياح تهبّ إلى الزوايا التي}
\MsL{21}{AB03-F033v-L21}{تكون فيها \BF{B17r}العلامات التي تقدّم ذكرها وقد يجوّز بدّل قوله السحابات التي\PN{45} شموسًا}
\begin{PNotes}
\PNote{43}{AB03-n043}{B: add. \RL{\ArN حالات}}
\PNote{44}{AB03-n044}{T: \RL{\ArN الحالات}}
\PNote{45}{AB03-n045}{B: add. \RL{\ArN تسمّى}}
\end{PNotes}
\Folio{034r}{34r}
\MsL{1}{AB03-F034r-L01}{أن يقال ما يجاوز الشمس. فإن طلعت سوداء أو مائلة إلى الخضرة أو غابت مع سحاب}
\MsL{2}{AB03-F034r-L02}{أو كانت حولها حالة\PN{46} واحدة أو حالتان أو كانت عن جنبها السحابات التي تسمّى}
\MsL{3}{AB03-F034r-L03}{شموسًا وكان شعاعه مائلًا إلى الخضرة أو أسود دلّت على هواء شتائيّ ومطر. وأمّا}
\MsL{4}{AB03-F034r-L04}{القمر فينبغي أن يتفقّد موضعه في اليوم الذي قبل الاجتماع بثلاثة أيّام والذي}
\MsL{5}{AB03-F034r-L05}{قبل الاستقبال بثلاثة أيّام والذي قبل الانتصاف في الضوء بثلاثة أيّام بعد هذه}
\MsL{6}{AB03-F034r-L06}{الأوقات بثلاثة أيّام وذلك أنّه إن ظهر فيها\PN{47} صفاء ولم يكن حوله شيء دلّ على صحو. وإن}
\MsL{7}{AB03-F034r-L07}{كان دقيقًا أحمر وكان جميع ما لم يكن فيه ضوء من دائرة كأنّه يتحرّك دلّ على هبوب}
\MsL{8}{AB03-F034r-L08}{الرياح من النواحيّ التي تميل إليها فإن كان أسود أو أخضر غليظًا دلّ على هواء شتائيّ}
\MsL{9}{AB03-F034r-L09}{وعلى أمطار. وقد ينبغي أن يتفقّد أيضًا الحالة التي تكون حوله وذلك أنّها إن}
\MsL{10}{AB03-F034r-L10}{كانت واحدة ثمّ كانت صافيّة تضمحلّ قليلًا دلّت على هواء صاف وإن كانت إثنتين}
\MsL{11}{AB03-F034r-L11}{أو ثلاثًا دلّت على هواء شتائيّ. فإن كانت مائلة إلى الحمرة الناصعة وكانت كالمنقطعة دلّت}
\MsL{12}{AB03-F034r-L12}{على هواء شتائيّ يكون بسبب رياح شديدة. فإن كانت كدرة غليظة دلّت على هواء شتائيّ}
\MsL{13}{AB03-F034r-L13}{تكون بسبب الزمن\PN{48} وإن كانت مائلة إلى الخضرة سوداء منقطعة دلّت على هواء شتائيّ}
\MsL{14}{AB03-F034r-L14}{تكون في الأمرين جميعًا وكلّما كانت الحالات أكثر فعلى حسب ذلك تكثر الأشياء}
\MsL{15}{AB03-F034r-L15}{التي ذكرنا. وأيضًا فإنّ الحالة التي تكون حول الكواكب المتحيّرة وحول النيّرة التي}
\MsL{16}{AB03-F034r-L16}{ليست بمتحيّرة تدلّ على الأشياء الألوان تلك الحالات ولطبائع ما يحيط بها. وينبغي}
\MsL{17}{AB03-F034r-L17}{أن تتفقّد أيضًا الكواكب التي ليست بمتحيّرة فينظر إلى الجماعات منها التي فيها}
\MsL{18}{AB03-F034r-L18}{كواكب كثير قريبة بعضها لبعض وتتفقّد ألوانها ومقاديرها \BF{B17v}وذلك}
\MsL{19}{AB03-F034r-L19}{أنّها إذا ظهرت وهي أكثر ضياء وأعظم ممّا جرت به العادة فيها فإنّه يدلّ على هبوب ريح\PN{49}}
\MsL{20}{AB03-F034r-L20}{من تلك الناحية التي فيها الكواكب التي تظهر فيها أيّ ناحية كانت من الفلك. وكذلك}
\MsL{21}{AB03-F034r-L21}{أيضًا \LR{\LatinF ]}وكذلك أيضًا\LR{\LatinF [} إذا كانت الكواكب التي تشبه اللطخ خاصّة مثل الكوكب}
\begin{PNotes}
\PNote{46}{AB03-n046}{T: \RL{\ArN هالة} (wherever it occurs in this chapter)}
\PNote{47}{AB03-n047}{T: \RL{\ArN دقيقًا}}
\PNote{48}{AB03-n048}{T: \RL{\ArN الدمق}}
\PNote{49}{AB03-n049}{BT: \RL{\ArN رياح}}
\end{PNotes}
\Folio{034v}{34v}
\MsL{1}{AB03-F034v-L01}{الذي يسمّى القلب\PN{50} أعني الذي على صورة السرطان وما أشبهه في وقت الصحو خفيّة}
\MsL{2}{AB03-F034v-L02}{لا تظهر أو الأشياء لتي تكون في بعض الأزمنة من الآثار العلويّة فإنّ الكواكب المذنبة}
\MsL{3}{AB03-F034v-L03}{اللطخيّة تدّل دائمًا على يبس الهواء على رياح قوّة ذلك بحسب كثرة الأشياء}
\MsL{4}{AB03-F034v-L04}{وبحسب طول مكثها. وأمّا انفضاض\PN{51} الكواكب وجريها فإنّه إن كان من زويّة واحدة}
\MsL{5}{AB03-F034v-L05}{فإنّه يدلّ على هبوب ريح من تلك الناحية وإن كان انفضاضها\PN{52} من زوايا مختلفة دلّت}
\MsL{6}{AB03-F034v-L06}{على رياح تجري على غير نظام وترتيب فإن كان ذلك من الأربع زوايا دلّت على هواء شتائيّ}
\MsL{7}{AB03-F034v-L07}{كثيرة الاختلاف حتّى ينتهي ذلك إلى حدوث الرعد والبرق وما أشبهه. وكذلك أيضًا}
\MsL{8}{AB03-F034v-L08}{حال الغيوم إذا كانت مشبهة بحرّار صوف ربّما دلّت على هواء شتائيّ وقوس قزح الذي}
\MsL{9}{AB03-F034v-L09}{يكون في بعض الأوقات فإنّه إذا ظهر في وقت صحو دلّ على هواء شتائيّ [[وقوس قزح الذي}
\MsL{10}{AB03-F034v-L10}{يكون]] وإذا ظهر في وقت هواء شتائيّ دلّ على هواء صاح\PN{53} وإنّا أقول قولًا مجملًا أو في أكثر}
\MsL{11}{AB03-F034v-L11}{الأمر الألوان الخاصّيّة التي تظهر في الهواء تدلّ على أحوال ملائة للأعراض التي دلّت عليها}
\Chap{I.15}{<15> معرفة أمور الموليد وما یستدلّ به فيها}
\MsL{12}{AB03-F034v-L12}{الألوان التي تقدّم ذكرها. \Hd{\LR{\LatinF \Sg{〈}15\Sg{〉}}}\Hd{ معرفة أمور الموليد وما یستدلّ به فيها}}
\MsL{13}{AB03-F034v-L13}{أوّل ما ينبغي أن تنظر\PN{54} فيه من أمور المواليد لمّا كان للناس ابتداء زمانيّ في كونهم وكان}
\MsL{14}{AB03-F034v-L14}{هذا الابتداء \BF{B18r}من هذا النوع من النظر بعینه لمن أراد أن يدقّق النظر ويغيص في الفحص}
\Chap{I.16}{<16> معرفة درج الطالع للمولود}
\MsL{15}{AB03-F034v-L15}{عن الخواصّ التي تكون في نفس التركيب وحدها. \Hd{\LR{\LatinF \Sg{〈}16\Sg{〉}}}\Hd{ معرفة درج الطالع للمولود}}
\MsL{16}{AB03-F034v-L16}{وإذ قد يقع الشدّ كثيرًا في الأمر لأوّل الذي يعتمد عليه وهو درجة الطالع في وقت خروج}
\MsL{17}{AB03-F034v-L17}{الجنين وذلك أنّما يقف على معرفة الساعات\PN{55} في وقت خروج الجنين باستقصاء في}
\MsL{18}{AB03-F034v-L18}{أكثر الأمر الذين يرصدون ذلك بذات الصفائح\PN{56} وحدها من بين الجميع آلات}
\MsL{19}{AB03-F034v-L19}{الساعات. فأمّا السائر آلات الساعات التي تقيس بها كثير ممّن يعني بذلك فقد}
\MsL{20}{AB03-F034v-L20}{يمكن أن يكذب كثيرًا ویجید عن الحقّ. فأمّا ما كان منها ممّا يقاس به الساعات}
\MsL{21}{AB03-F034v-L21}{الشمس فبسبب ما يعرض من ميلها في الموضع ومن اعوجاج مقايسها وأمّا آلات}
\begin{PNotes}
\PNote{50}{AB03-n050}{T: \RL{\ArN المعلف }}
\PNote{51}{AB03-n051}{T: \RL{\ArN انقضاض}}
\PNote{52}{AB03-n052}{T: \RL{\ArN انقضاضها}}
\PNote{53}{AB03-n053}{BT: \RL{\ArN صحو}}
\PNote{54}{AB03-n054}{B: \RL{\ArN ننظر}}
\PNote{55}{AB03-n055}{B: \RL{\ArN الساعا}}
\PNote{56}{AB03-n056}{cf. T: \RL{\ArN بالأسطرولاب}}
\end{PNotes}
\Folio{035r}{35r}
\MsL{1}{AB03-F035r-L01}{ساعات الماء فبسبب\PN{57} ما يعرض للماء من الاحتباس\PN{58} ومن اختلاف السيل الذي يكون}
\MsL{2}{AB03-F035r-L02}{عن أسباب مختلفة وعن أدنى سبب يعرض له. وجب ضرورة أن يتقدّم فنخبر بالجهة}
\MsL{3}{AB03-F035r-L03}{التي تمكن الإنسان بها أن يعلم الدرجة من فلك البروج الذي يجب بحسب القياس}
\MsL{4}{AB03-F035r-L04}{الطبيعيّ اللازم أن تكون هي الطالعة من بعد أن تكون قد علمنا من الساعة التي قد أعطيناها}
\MsL{5}{AB03-F035r-L05}{بالتقريب درجة الطالع من [[رب الفلك]] باب المطالع فنقول إنّه ينبغي أن نأخذ من}
\MsL{6}{AB03-F035r-L06}{الاجتماع أو الاستقبال الذين كانا قبل الولادة أقربهما منه أيّهما كان الأقرب ويستقصي}
\MsL{7}{AB03-F035r-L07}{معرفة درجته. فإن كان الذي أخذنا اجتماعًا أخذنا درجة النيّرين وإن كان}
\MsL{8}{AB03-F035r-L08}{المأخوذ استقبالًا أخذنا درجة أيّ النيّرين اتّفق أن يكون فوق الأرض.}
\MsL{9}{AB03-F035r-L09}{ثمّ تنظر أيّ الكواكب يستحقّ أن يكون له تدبير تلك الدرجة في وقت المولّد والوجة}
\MsL{10}{AB03-F035r-L10}{\BF{B18v}بالجملة الذي يستحقّ الكوكب التدبير يتمّ به الخمسة أشياء أعني التثليث والبيت}
\MsL{11}{AB03-F035r-L11}{والشرف والجدّ والظهور القريب من الشعاع أو مشاركة الشمس في الشكل فيكون}
\MsL{12}{AB03-F035r-L12}{تدبير الموضع للكوكب الذي له في ذلك الموضع أكثر من خصلة واحدة من هذه الخصال}
\MsL{13}{AB03-F035r-L13}{أو جميعها. فإن وجدنا لكوكب واحد جميع هذه الأشياء أو أكثر نظرنا إلى درجته}
\MsL{14}{AB03-F035r-L14}{المأخوذة بالاستقصاء من البرج الذي هو فيه وقت المولّد فحكمنا بأنّ الدرجة المساويّة}
\MsL{15}{AB03-F035r-L15}{لها في القدر من البرج الذي كنّا أخذناه بالتقريب من الطالع هي درجة أحد أوتاد}
\MsL{16}{AB03-F035r-L16}{الطالع. وإن وجدنا أنّه قد اشترك في التدبير كوكبان أو أكثر من كوكبين نظرنا}
\MsL{17}{AB03-F035r-L17}{إلى عدد ما ساره كلّ كوكب واحد منها في وقت المولد فأیّها كان عدده أقرب إلى عدد}
\MsL{18}{AB03-F035r-L18}{درجة الطالع التي كانت خرجت لنا بالطالع فهو الذي يطلع في وقت المولد.}
\MsL{19}{AB03-F035r-L19}{وإن وجدنا عدد درج كوكبين أو أكثر من ذلك قريبًا من عدد درجة الطالع}
\MsL{20}{AB03-F035r-L20}{التي خرجت لنا بالطالع اتّبعنا في ذلك الكوكب الذي له حصص أكثر في الأوتاد}
\MsL{21}{AB03-F035r-L21}{وفي أنّ الحيّز حيّزه. فإن كان بعد درجة الكو[ا]كب المدبّر من درج الطالع الذي}
\begin{PNotes}
\PNote{57}{AB03-n057}{B: \RL{\ArN بسبب}}
\PNote{58}{AB03-n058}{T: \RL{\ArN أجناس}}
\end{PNotes}
\Folio{035v}{35v}
\MsL{1}{AB03-F035v-L01}{أخذ على الأمر المجمل أكثر من بعدها من درج وسط السماء. فأقمنا بها سائر الأوتاد}
\MsL{2}{AB03-F035v-L02}{وينبغي أن تنظر\PN{59} إلى عرض الكوكب إذ كان لا يثبت على نطاق البروج. فيعلم}
\MsL{3}{AB03-F035v-L03}{الدرجة التي تطلع معها الكوكب من فلك البروج والدرجة التي معها یتوسّط}
\MsL{4}{AB03-F035v-L04}{السماء من قبل العرض على ما شرحنا في صدر الكتاب. ثمّ نعرف درجة\PN{60} الطالع}
\MsL{5}{AB03-F035v-L05}{ودرجات وسط السماء التي عرفناها من المطالع بالتقريب. فإن كان درجات}
\MsL{6}{AB03-F035v-L06}{الكوكب أقرب إلى درجات السماء منها في العدد منها إلى درجات الطالع صیّرنا}
\MsL{7}{AB03-F035v-L07}{عدد درجات وسط السماء بقدر\PN{61} درجات الكوكب الذي معها یتوسّط السماء}
\MsL{8}{AB03-F035v-L08}{وقد يجوّز أن يقع فيما يؤخذ بذات \BF{B19r}الصفائح الخلاف بجهات شتّى ولذلك}
\Chap{I.17}{<17> في قسمة القول في المواليد}
\MsL{9}{AB03-F035v-L09}{صار هذا الباب عيارًا على الجميع. \Hd{\LR{\LatinF \Sg{〈}17\Sg{〉}}}\Hd{ في قسمة القول في المواليد}}
\MsL{10}{AB03-F035v-L10}{وإذ قدّمنا وضع هذه الأشياء فإنّا نقول إنّه إذ أراد مريد أن يقسم جملة النظر في}
\MsL{11}{AB03-F035v-L11}{المواليد قسمة على ترتيب ونظام وجد أنّ بعض الأشياء التي يمكن إدراكها على}
\MsL{12}{AB03-F035v-L12}{المجرى الطبيعيّ هي الأعراض التي تعرض قبل الولادة مثل القول في الوالدين}
\MsL{13}{AB03-F035v-L13}{وبعضها هي الأعراض التي تعرض قبل الولادة وبعدها مثل القول في الإخوة.}
\MsL{14}{AB03-F035v-L14}{وأمّا التي تكون في الولاد نفسه فليس أمرها أمرًا وحدًا بسيطًا كالذي وصفنا.}
\MsL{15}{AB03-F035v-L15}{وآخر ما يمكن إدراكه من الأعراض هو ما يحدث مثله بعد الولادة وهذا النوع}
\MsL{16}{AB03-F035v-L16}{كثير الفنون والأجزاء. والذي تراد معرفة ممّا في وقت الولادة نفسها هو القول}
\MsL{17}{AB03-F035v-L17}{في الذكورة والأناث والقول فيما يولد من التوأمين وأكثر من التوأمين \LR{\Sg{〉}}\LR{\Sg{〉}} والقول في المشهور هي الخلق \LR{\Sg{〈}}\LR{\Sg{〈}}\PN{62} والقول}
\MsL{18}{AB03-F035v-L18}{في الذين لا یغتدون والذي يرا[د] معرفته ممّا بعد الولادة هو القول في العمر}
\MsL{19}{AB03-F035v-L19}{إذ كان [[ممّا]] ليس ممّا يتّصل بالقول في الذين لا یغتدون ومن\PN{63} ذلك القول}
\MsL{20}{AB03-F035v-L20}{في صورة البدن والعلل والآفات\PN{64} التي تعرض للبدن. ثمّ من بعد ذلك القول}
\MsL{21}{AB03-F035v-L21}{في أحوال النفس وفي الأعراض النفسانیّة. ثمّ من\PN{65} بعد ذلك القول في بخت المولود}
\begin{PNotes}
\PNote{59}{AB03-n059}{B: \RL{\ArN ننظر}}
\PNote{60}{AB03-n060}{B: add. \RL{\ArN الكتاب}}
\PNote{61}{AB03-n061}{B: add. \Sg{〈}\Sg{〈} \RL{\ArN عدد} \Sg{〉}\Sg{〉}}
\PNote{62}{AB03-n062}{right margin; B: \RL{\ArN والقول في مسوّى الحلق}, corrected to \RL{\ArN والقول في مشوهي الحلق}}
\PNote{63}{AB03-n063}{B: add. \RL{\ArN بعد}}
\PNote{64}{AB03-n064}{B: \RL{\ArN الأوقات}}
\PNote{65}{AB03-n065}{B: om.}
\end{PNotes}
\Folio{036r}{36r}
\MsL{1}{AB03-F036r-L01}{في أمر الملك وفي أمر المرتبة والرفعة ثمّ من\PN{66} بعد ذلك القول في حاله فيما يعالج}
\MsL{2}{AB03-F036r-L02}{من الأعمال. ثمّ من بعد ذلك القول في تزويجه ومقامه مع زوجته وفيما يوالد له}
\MsL{3}{AB03-F036r-L03}{من الأولاد وفي مخالطته الناس ومشاركته لهم وفي أصدقائه. ثمّ من\PN{67} بعد ذلك القول}
\MsL{4}{AB03-F036r-L04}{في الغربة. ثمّ بآخرته القول في حال منيته\PN{68} التي هي بالقوّة مشاكله للقول \BF{B19v}في العمر}
\MsL{5}{AB03-F036r-L05}{ومرتبتها یجب أن تكون بعد هذا كلّه. ويتّصف\PN{69} أمر كلّ واحد من الأشياء على جهة}
\MsL{6}{AB03-F036r-L06}{الجمل. فإذ قد شرحنا أمر درجة الطالع وتصحيحها على الجهة فنضیف إلى أمر كلّ}
\MsL{7}{AB03-F036r-L07}{ما يحدث من أحوال المواليد ونذکر ذلك مع ذكرنا للقوى المجرّدة الفاعلة كما}
\MsL{8}{AB03-F036r-L08}{قلنا آنفًا ونطرح ما وصفه كثير من الناس ممّا لا حجّة فيه ولا قول مقنع يليق بالأسباب}
\MsL{9}{AB03-F036r-L09}{التي تكون عن الطبيعة الأولى ونفحص عمّا يمكن إدراكه منها لا بالفرع والأعداد}
\MsL{10}{AB03-F036r-L10}{التي ليس لها سبب يوجبها لکن بالنظر في أشكال الكواكب وملائمتها لمواضعها}
\MsL{11}{AB03-F036r-L11}{ولئلًّا يحتاج إلى إعادة القول وتكريره في كلّ واحد من الأبواب التي ذكرنا}
\MsL{12}{AB03-F036r-L12}{ينبغي أن تقول في جميعها قولًا عامًّا. إنّا ننظر أوّلًا في المواضع من فلك البرج}
\MsL{13}{AB03-F036r-L13}{الذي يلائم نوع ذلك الباب المطلوب من أبواب النظر في المولد كملائمة}
\MsL{14}{AB03-F036r-L14}{موضع وسط السماء مثلًا لأمن\PN{70} الأعمال التي تعالج أعني \{الصنائع\} أوكملائمة موضع}
\MsL{15}{AB03-F036r-L15}{الشمس لأمن\PN{71} الأب. ثمّ تنظر\PN{72} بعد ذلك الكواكب التي تستحقّ أن يكون لها تدبير}
\MsL{16}{AB03-F036r-L16}{ذلك الموضع بالخمسة الأنواع التي ذكرناها فيما تقدّم من قولنا فإن كان ربّ}
\MsL{17}{AB03-F036r-L17}{الموضع كوكبًا واحدًا من جميع الجهات جعلنا\PN{73} تدبير ذلك الأمر له وإن كان كوكبين}
\MsL{18}{AB03-F036r-L18}{أو ثلاثة كواكب جعلنا التدبير لمّا كان له منها حصص أكثر عددًا. ثمّ نجعل نظرنا بعد}
\MsL{19}{AB03-F036r-L19}{ذلك في حال الحادث وفحصنا عنه من طبائع الكواكب أنفسها التي قد استحقّت}
\MsL{20}{AB03-F036r-L20}{التدبير ومن طبائع البروج التي تكون فيها الكواكب المواضع المشاكلة. ثمّ نجعل فحصنا}
\MsL{21}{AB03-F036r-L21}{عن الحادث ومقداره من قوّتنا فننظر هل هي قویّة في حالنا في العالم أو في حالها}
\begin{PNotes}
\PNote{66}{AB03-n066}{B: om.}
\PNote{67}{AB03-n067}{B: om.}
\PNote{68}{AB03-n068}{B: \RL{\ArN متنية}}
\PNote{69}{AB03-n069}{B: \RL{\ArN سنصف}}
\PNote{70}{AB03-n070}{T: \RL{\ArN لأمر}}
\PNote{71}{AB03-n071}{T: \RL{\ArN لأمر}}
\PNote{72}{AB03-n072}{B: \RL{\ArN ننظر}}
\PNote{73}{AB03-n073}{B: add. \RL{\ArN جميع}}
\end{PNotes}
\Folio{036v}{36v}
\MsL{1}{AB03-F036v-L01}{في المولد أو الأمر على ضدّ ذلك وأقوى ما يكون في حالنا في العالم إذا كانت في}
\MsL{2}{AB03-F036v-L02}{المواضع التي تخصّنا أو تلاءمها. وإذا كانت مشرقة وكانت متزیّدة في السير}
\MsL{3}{AB03-F036v-L03}{وأقوى ما تكون في حالها في المولد إذا كانت في الأوتاد \BF{B20r}أو في المواضع الذي}
\MsL{4}{AB03-F036v-L04}{ترتفع إلى الأوتاد وبخاصّة الطالع ووسط السماء وأضعف ما تكون في حالها في}
\MsL{5}{AB03-F036v-L05}{العالم إذا كانت في مواضع غربيّة أو كانت مغربة أو راجعة أو متناقصة في المسير.}
\MsL{6}{AB03-F036v-L06}{وأضعف ما تكون في حالها في المواليد إذا زالت عن الأوتاد وتجعل نظرك في مدّة}
\MsL{7}{AB03-F036v-L07}{زمان الحادث بأن تفحص عنها هل هي منسوبة إلى الغدوات أم منسوبة إلى العشیّات}
\MsL{8}{AB03-F036v-L08}{بقياسها لها إلى الشمس أو إلى الطالع وذلك أنّ الربعين المتقدّمين لكلّ واحد}
\MsL{9}{AB03-F036v-L09}{منها والذين يقابلانهما على القطر منسوبة إلى الغدوات والأرباع الباقيّة}
\MsL{10}{AB03-F036v-L10}{التي تتلوها تكون منسوبة إلى العشیّات وتنظر هل هي في الأوتاد\PN{74} صاعدة}
\MsL{11}{AB03-F036v-L11}{إلى الأوتاد فإن كانت منسوبة إلى الغدوات وكانت على الأوتاد. فهي أحد فعلًا}
\MsL{12}{AB03-F036v-L12}{وإن كانت منسوبة إلى العشيّات وكانت صاعدة إلى الأوتاد فهي إبطاء فعلًا}
\MsL{13}{AB03-F036v-L13}{وينبغي أن تحفظ\PN{75} أيضًا في كلّ الأمور التي تحدث من أمور المواليد أنّه تجب ضرورة}
\MsL{14}{AB03-F036v-L14}{في الكواكب التي تستحدث شيئًا ما في الموضع المطلوب أن يكون مشاكلًا له منذ}
\MsL{15}{AB03-F036v-L15}{الابتداء. فإنّه أن\PN{76} يكن مشاكلًا له لم يمكنه أن يحدث فيه شيئًا منه إذ كان لم يشارك}
\MsL{16}{AB03-F036v-L16}{الابتداء الأوّل يعني أنّ الكوكب إذا لم يرتبط البرج الذي هو فيه بالبرج المطلوب}
\MsL{17}{AB03-F036v-L17}{الذي يقع الكوكب فيه لم یتهیّأ أن يحدث في ذلك البرج شيئًا من طبيعته}
\MsL{18}{AB03-F036v-L18}{إذ كان غير مرتبط به في المولد. فأمّا زمان كون الشيء الحادث فليس شبّه التدبير}
\MsL{19}{AB03-F036v-L19}{الأوّل لكن شبيه حال الكوكب المستوليّ على التدبير عند الشمس وعند زوايا}
\Chap{I.18}{<18> في الوالدين}
\MsL{20}{AB03-F036v-L20}{العالم. \Hd{\LR{\LatinF \Sg{〈}18\Sg{〉}}}\Hd{ في الوالدين} \BF{B20v}إنّ الشمس وزحل مشاكلان للآباء بالطبع}
\MsL{21}{AB03-F036v-L21}{والقمر والزهرة مشاكلان فيحسب ما يوجد عليه حال هذه الكواكب بعضها عند}
\begin{PNotes}
\PNote{74}{AB03-n074}{B: add. \RL{\ArN أو هي}}
\PNote{75}{AB03-n075}{B: \RL{\ArN نحفظ}}
\PNote{76}{AB03-n076}{B: add. \RL{\ArN لم} (!)}
\end{PNotes}
\Folio{037r}{37r}
\MsL{1}{AB03-F037r-L01}{بعض وعند غيرها كذلك ينبغي أن تتوهّم\PN{77} أنّ حال الوالدين \LR{\Sg{〉}}\LR{\Sg{〉}} يكون \LR{\Sg{〈}}\LR{\Sg{〈}}\PN{78} وذلك أنّه ينبغي أن}
\MsL{2}{AB03-F037r-L02}{تنظر في \{نجتهم\} وملكهم من احتفاف الكواكب النیّرین وإنّه متى أحاطت بهما الكواكب}
\MsL{3}{AB03-F037r-L03}{المسعودة التي هي من حيّزها أو كانت تلك الكواكب إمّا في برجي النيّرين أنفسهما}
\MsL{4}{AB03-F037r-L04}{وإمّا في البروج التي تليهما دلّت على أنّ أمر الوالدين أمر بهي نبيل لا سيّما متى}
\MsL{5}{AB03-F037r-L05}{كانت الكواكب \{المحتفّة\} بالشمس منسوبة إلى الغدوات والمحتفّة بالقمر}
\MsL{6}{AB03-F037r-L06}{منسوبة إلى العشيّات وذلك أن تطلع قبل الشمس وبعد القمر. فيكون وضعها}
\MsL{7}{AB03-F037r-L07}{لذلك وضعًا محمودًا. فإن كان أيضًا زحل والزهرة مشرقين أو كانا على المواجهة}
\MsL{8}{AB03-F037r-L08}{التي تخصّهما أو كانا على الأوتاد فينبغي لنا أن نتوهّم\PN{79} من ذلك سعادة الوالدين}
\MsL{9}{AB03-F037r-L09}{ظاهرة على قدر ما يلازم\PN{80} كلّ واحد منهما وإن كان الأمر على خلاف ذلك وإن كان}
\MsL{10}{AB03-F037r-L10}{النيّرين خالي السیر ولم يحفّها شيء من الكواكب دلّا على ضعة الوالدين}
\MsL{11}{AB03-F037r-L11}{وخساستهما ولا سيّما متى كان وضع الزهرة وزحل وضعًا غير محمود. فإن احتفّتهما}
\MsL{12}{AB03-F037r-L12}{الكواكب المنحسة التي ليست من حيّزهما كما يعرض إذا طلع المرّيخ بعد الشمس}
\MsL{13}{AB03-F037r-L13}{وطلع زحل بعد القمر ولم يكن موضع المسعدة من النيّرين موضعًا محمودًا ولا كانت}
\MsL{14}{AB03-F037r-L14}{من حيّزهما. فينبغي أن يستدلّ من ذلك على خساسة عيش الوالدين واضطراب}
\MsL{15}{AB03-F037r-L15}{حالهما. فإن كان السهم الدالّ على البخت وهو سهم السعادة الذي سنذكره}
\MsL{16}{AB03-F037r-L16}{موافقًا للمو[[ا]]ل[[ي]]د\PN{81} للكواكب المحتفّة بالشمس أو القمر موافقة محمودة سلم ملك}
\MsL{17}{AB03-F037r-L17}{الوالدين. فإن لم يكن السهم على ما قلنا وكان الأمر على خلاف ذلك. ثمّ لم يحفّ}
\MsL{18}{AB03-F037r-L18}{النيّرين شيء من الكواكب المسعدة أو كان الذي يحفّهما المنحسة لم ينتفع}
\MsL{19}{AB03-F037r-L19}{بملك الوالدين وصار ضارًّا لهما. \BF{B21r}فأمّا طول عمر الوالدين وقصره فينظر فيه}
\MsL{20}{AB03-F037r-L20}{من سائر اشتراكات الكواكب في الأشكال. فأمّا عمر الأب فإنّ المشتري والزهرة}
\MsL{21}{AB03-F037r-L21}{إذا\PN{82} شاركًا الشمس أو الزحل في الشكل أيّ جهة\PN{83} كان ذلك أو كان زحل}
\begin{PNotes}
\PNote{77}{AB03-n077}{T: \RL{\ArN نتقدّم}}
\PNote{78}{AB03-n078}{interlinear}
\PNote{79}{AB03-n079}{B: \RL{\ArN يتوهّم}}
\PNote{80}{AB03-n080}{T: \RL{\ArN يلائم}}
\PNote{81}{AB03-n081}{B: \RL{\ArN للمولد}}
\PNote{82}{AB03-n082}{B: \RL{\ArN إن}}
\PNote{83}{AB03-n083}{T: \RL{\ArN وجه}}
\end{PNotes}
\Folio{037v}{37v}
\MsL{1}{AB03-F037v-L01}{أيضًا مشاركًا للشمس وكانت مشاركته للشمس مشاركة اتّفاق إمّا أن يكون معها}
\MsL{2}{AB03-F037v-L02}{أو في تثليثها أو في تسديسها وكانا قویّین دلّ \LR{\LatinF ]]}ذلك\LR{\LatinF [[} على طول عمر الأب. فإن لم}
\MsL{3}{AB03-F037v-L03}{يكونا قويّين لم يكون الأمر كذلك إلّا أنّهما لا يدلّان على قصر العمر لا محالة فإن لم يكن}
\MsL{4}{AB03-F037v-L04}{الأمر كذلك وكان المرّيخ عاليًا على الشمس أو على زحل أو كان أيضًا تابعًا لهما}
\MsL{5}{AB03-F037v-L05}{يعني يطلع بعدهما ولم يكن زحل أيضًا موافقًا للشمس لكنّ على تربيعها أو على مقابلتها}
\MsL{6}{AB03-F037v-L06}{على القطر أو كانا جميعًا زائلين عن الأوتاد حدث عنهما ضعف الآباء فإن كان}
\MsL{7}{AB03-F037v-L07}{على الأوتاد أو حيث يرتفعان إلى الأوتاد صيّر الآباء قصيري الأعمار أو أحدثًا بهما}
\MsL{8}{AB03-F037v-L08}{آفة. أمّا قصر العمر فيحدثانه إن كانا على الطالع أو على وسط السماء. وأمّا الآفة}
\MsL{9}{AB03-F037v-L09}{فيحدثانها إن كانا على الوتدين الباقيين أعني الغارب ووتد الأرض أو بحيث}
\MsL{10}{AB03-F037v-L10}{يرتفعان [[إليهما]]. وذلك أنّ المرّيخ إذا نظر إلى الشمس على الجهة التي ذكرتها}
\MsL{11}{AB03-F037v-L11}{في صدر الكتاب وهو الأجزاء التي ذكرنا \{أنّهما\} تنظر بعضها إلى بعض ويستوفي}
\MsL{12}{AB03-F037v-L12}{في القوّة مات الأب فجأة أو حدثت عليه آفة في عينيه. فإن كان نظر المرّيخ}
\MsL{13}{AB03-F037v-L13}{هو إلى زحل حدث على الأب إمّا موت وإمّا الحمّى التي تكون مع نافض وإمّا آفة}
\MsL{14}{AB03-F037v-L14}{تناله بسبب أقطع أو إحراق. وزحل أيضًا إذا كان شكله عند الشمس شكلًا}
\MsL{15}{AB03-F037v-L15}{رديئًا حدث عنه موت الآباء بالأمراض وبعلل تعرض لهم من رطوبات مؤذیّة وأمّا}
\MsL{16}{AB03-F037v-L16}{عمر الأمّ فإنّ المشتري إنّ شارك القمر أو الزهرة في الشكل على أيّ جهة}
\MsL{17}{AB03-F037v-L17}{كان ذلك أو وافقت الزهرة أيضًا القمر إمّا بأن يكون معه أو تسديسها أو}
\MsL{18}{AB03-F037v-L18}{تثليثها. ثمّ كانا قویّین دلًّا \BF{B21v}على طول عمر الأمّ. وأمّا المرّيخ فإنّه أن ینظر إلى القمر أو}
\MsL{19}{AB03-F037v-L19}{كان تابعًا للزهرة أو على التربيع منها أو المقابلة أو كانت حال زحل عند}
\MsL{20}{AB03-F037v-L20}{القمر هذه الحال فإنّهما إن كانا ناقصيّ السير راجعين أو زائلين عن}
\MsL{21}{AB03-F037v-L21}{الأوتاد \{اجتلبًا\} على الأمّهات وقوعًا في العلل والأسقام فقط وإن كانا}
\Folio{038r}{38r}
\MsL{1}{AB03-F038r-L01}{متزيّدي السير أو كانا على الأوتاد قصر أعمار الأمّهات وأحدثًا بهنّ. أمّا العمر}
\MsL{2}{AB03-F038r-L02}{فيقصر أنّه إذا كانا على الوتدين الشرقيّين أو بحيث يرتفعان إليهما وأمّا}
\MsL{3}{AB03-F038r-L03}{الآفة. فيحدثانها إذا كانا على الوتدين الغربيّین أو يحدث يصيّران إليهما. وذلك أنّ}
\MsL{4}{AB03-F038r-L04}{المرّيخ إذا نظر إلى القمر على الجهة التي ذكرنا. ثمّ كان القمر مشرقًا فأحدث عن ذلك}
\MsL{5}{AB03-F038r-L05}{موت الأمّهات فجأة أو آفة تعرض لهنّ في الأعين. فأمّا إن كان القمر غربيًّا فإنّه}
\MsL{6}{AB03-F038r-L06}{تحدث عنه موتهنّ من إسقاط جنين أو ما أشبه أو آفة تعرض لهنّ من القطع أو}
\MsL{7}{AB03-F038r-L07}{الإحراق ويجوّز بدّل قوله غربيًّا أن يقال زائلًا. فإن نظر المرّيخ إلى الزهرة حدث}
\MsL{8}{AB03-F038r-L08}{عن ذلك موتهنّ بالحمّيات وعلل تنالهنّ من ظلمة البصر ومن موادّ تنصب بغتة.}
\MsL{9}{AB03-F038r-L09}{وإذا نظر زحل إلى القمر حدث عن ذلك موت وعلل تعرض لهنّ أمّا أن كان القمر}
\MsL{10}{AB03-F038r-L10}{مشرقيًّا\PN{84}. فإنّ موتهنّ يكون عن حمیّات مع ناقص وأمّا أن كان غربيًّا فإنّ ذلك}
\MsL{11}{AB03-F038r-L11}{يكون عن أورام ساعيّة وتأكّل يكونان في الرحم. وقد ينبغي لنا أن نستعمل في}
\MsL{12}{AB03-F038r-L12}{معرفة أنواع الآفات الجزئيّة والعلل أو الموت العارض للوالدين مع ما ذكرنا}
\MsL{13}{AB03-F038r-L13}{أحوال البروج التي توجد فيها الكواكب التي هي على\PN{85} الفاعلة للحال الخاصّيّة وسنذكر}
\MsL{14}{AB03-F038r-L14}{أحوال البروج ونشرح أمرها عند ذكرنا الأمر المولود وإن تتفقّد بالنهار خاصّة}
\MsL{15}{AB03-F038r-L15}{الشمس والزهرة وبالليل خاصّة زحل والقمر والذي يليق بنا في الفحص عن أمور}
\MsL{16}{AB03-F038r-L16}{الوالدين الجزئيّة الباقيّة ويلزمنا هو أن نقيم موضع الأب أو الأمّ الذي للخير}
\MsL{17}{AB03-F038r-L17}{مقام الطالع. وتنظر في سائر أمورها كما يفعل في مولدي الوالدين \BF{B22r}أنفسهما}
\MsL{18}{AB03-F038r-L18}{على ما سنبيّن من بعد من المذاهب العرضیّة والفاعلة التي توجد في الأنواع}
\MsL{19}{AB03-F038r-L19}{العامّیّة. أمّا جهات تراكب الأشياء\PN{86} بعضها مع بعض. فينبغي أن يكون ذاكرين لها}
\MsL{20}{AB03-F038r-L20}{في هذا الموضع وفي جميع المواضع التي تأتي من بعد في أمور المواليد بأن تنظر\PN{87}}
\MsL{21}{AB03-F038r-L21}{إذا لم تكن الكواكب التي هي أرباع\PN{88} المواضع المطلوبة من نوع واحد لکنّ من أنواع}
\begin{PNotes}
\PNote{84}{AB03-n084}{BT: \RL{\ArN شرقيًّا}}
\PNote{85}{AB03-n085}{BT: om.}
\PNote{86}{AB03-n086}{T: \RL{\ArN الأسباب}}
\PNote{87}{AB03-n087}{BT: \RL{\ArN ننظر}}
\PNote{88}{AB03-n088}{T: \RL{\ArN أرباب}}
\end{PNotes}
\Folio{038v}{38v}
\MsL{1}{AB03-F038v-L01}{مختلفة. فاعلة للأضداد أيّ كوكب من الكواكب التي أشدّ قوّة واستعلاء\PN{89} وغلبة\PN{90} نجده أكثر}
\MsL{2}{AB03-F038v-L02}{حصصًا حتّى يكون هو أولى بما يحدث فيلزم في الفحص عنها دائمًا طبائع هذه الأشياء}
\MsL{3}{AB03-F038v-L03}{وإن كانت الحصص مستاويّة الوزن. فإنّ الكواكب العالية إن كانت مع ما}
\MsL{4}{AB03-F038v-L04}{وصفنا\PN{91} إلى معرفة الشيء المجتمع من امتزاج هذه الطبائع المختلفة بسهولة وإن}
\MsL{5}{AB03-F038v-L05}{كانت متفرّقة فأنّا نعطي كلّ واحد منها من الأعراض ما يلائمه على حسب الأوقات}
\MsL{6}{AB03-F038v-L06}{الجزئيّة ونعطي أوّلًا ما كان منها منسوبًا إلى العشیّات وقال غير بطلميوس\PN{92}}
\MsL{7}{AB03-F038v-L07}{فيما\PN{93} دلائل الآباء مثل ذلك غير أنّهمم أضافوا إلى ذلك بالليل والنهار سهم الأب}
\MsL{8}{AB03-F038v-L08}{والأمّ وبرج الآباء ونظروا في قوّة جميع الأولاد وضعفها ومواضعها من الأوتاد وأرباب}
\MsL{9}{AB03-F038v-L09}{هذه المواضع. فإنّ وجدوها قويّة استدلّوا بذلك على فضل الآباء وحسن حالهم وإن}
\MsL{10}{AB03-F038v-L10}{وجدوها ساقطة ضعيفة استدلّوا بها على حساستهم واستعملوا في ذلك أيضًا أرباب مثلّثات}
\MsL{11}{AB03-F038v-L11}{الأدلّاء. فقالوا إن كانت مع فساد الأدلّاء منحوسة فاسدة معها دلّت على زمانة}
\MsL{12}{AB03-F038v-L12}{الآباء. وإن وجد سعد في بيت الآباء وأرباب البيوت والحدود وحدود النيّرين}
\MsL{13}{AB03-F038v-L13}{ناظرة إليها حكموا بشرف الآباء وإن نظرت النحوس إلى بيت الآباء من تربيع}
\MsL{14}{AB03-F038v-L14}{أو مقابلة أو كانت فيه ولم تنظر السعود حكموا بأنّهم عبيد أو شبيه بالعبيد. \BF{B22v}فإن}
\MsL{15}{AB03-F038v-L15}{كان برج الآباء بيت الشمس كانت الدلالة له الأب وإن كان في بيت القمر كانت}
\MsL{16}{AB03-F038v-L16}{الدلالة للأمّ وإن وجدت ربّ بيت\PN{94} الثاني عشر مع القمر في الثالث عشر من}
\MsL{17}{AB03-F038v-L17}{الوتد كانت أمّه أمّة وإن كانت هذه الحال للشمس كان أبوه\PN{95} عبدًا ومن كانت}
\MsL{18}{AB03-F038v-L18}{أرباب مثلّثات دليلة منها في حظوظ أنفسهما في مواضع محمود\PN{96} في الفلك بريئة}
\MsL{19}{AB03-F038v-L19}{من النحوس تنظر إلى السعود دلّت على الفائدة والمنزلة فيما يستأنف ولا سيّما}
\MsL{20}{AB03-F038v-L20}{متى كان سهم السعادة وربّه منحوس\PN{97} سعودين مع سهم الآباء. فإن ذلك}
\begin{PNotes}
\PNote{89}{AB03-n089}{T: \RL{\ArN استيلاء}}
\PNote{90}{AB03-n090}{T: \RL{\ArN عليه}}
\PNote{91}{AB03-n091}{T: \RL{\ArN وصلنا}}
\PNote{92}{AB03-n092}{B: \RL{\ArN بطلميوس} (wherever it occurs)}
\PNote{93}{AB03-n093}{B: \RL{\ArN في}}
\PNote{94}{AB03-n094}{B: \RL{\ArN للبيت}}
\PNote{95}{AB03-n095}{B: \RL{\ArN أبوه}}
\PNote{96}{AB03-n096}{B: \RL{\ArN محمودة}}
\PNote{97}{AB03-n097}{B: om.}
\end{PNotes}
\Folio{039r}{39r}
\MsL{1}{AB03-F039r-L01}{يدلّ على دوام الأمر وتوريثه للأبناء فإن كان الأمر على ضدّ ذلك لم ينتفع الوالدين}
\MsL{2}{AB03-F039r-L02}{\{بلكهما\} ودلّا على سقوط منزلتهما ونظروا في موافقة الآباء للمولود ومحبِّتهم إیّاه. ينظر}
\MsL{3}{AB03-F039r-L03}{إلى الأغلب على الطالع والأغلب على أمر الأب بكثرة الحصص فإن اتّفقا وكانا متلاؤمين}
\MsL{4}{AB03-F039r-L04}{كانوا متحابّین وإن كان الأمر على ضدّ ذلك كانوا متباغضين وكذلك إن كانت الشمس}
\MsL{5}{AB03-F039r-L05}{في حدّ المرّيخ بالنهار وجد زحل بالليل لا سيّما إذا نظرت النحوس إلى الشمس ولم تنظر}
\MsL{6}{AB03-F039r-L06}{السعود دلّ على التباغض أو يكون المولود محدّ بالمنزلة وبإله وإن كان القمر على ما وصفت}
\MsL{7}{AB03-F039r-L07}{من حال الشمس لقيت الأمّ من ولدها بالًا ومن وجد سهمه منها في الثاني عشر}
\MsL{8}{AB03-F039r-L08}{فإن المولود يكون عدوًّا له ومتی كانت\PN{98} الشمس في وتد ونظر إليها المرّيخ من المقابلة}
\MsL{9}{AB03-F039r-L09}{أبغض الأب ولده وكذلك إذا كان القمر على هذه الحال أبغضت الأمّ الولد وإن}
\MsL{10}{AB03-F039r-L10}{تقابل المشتري وزحل ويشرف بعضهما على بعض دلّ على التباغض وإن كانا}
\MsL{11}{AB03-F039r-L11}{من التربيع فكذلك. وأمّا من التسديس والتثليث فإنّه يدلّ على المحبّة فإن لم ينظر}
\MsL{12}{AB03-F039r-L12}{صاحب بيت الآباء إلى الآباء ولا صاحب الشمس إليهما ولا صاحب سهم الآباء إلى}
\MsL{13}{AB03-F039r-L13}{السهم \BF{B23r}كان المولود \{سعود\}\PN{99} وشدّة فإن نحست الشمس بالنهار من المقابلة والتربيع}
\MsL{14}{AB03-F039r-L14}{وكانت في مكان \{رديّ\} أبغض الأب ولده ونجاة عنه يريد انتقامه وقالوا في عمر}
\MsL{15}{AB03-F039r-L15}{الآباء وإن كانت الشمس تنظر إلى الطالع بالنهار سيّرت لعمر الأب وإن نظر زحل إلى}
\MsL{16}{AB03-F039r-L16}{الطالع كانت التسيير منه أن سقطت الشمس وأن سقط زحل. فدرجة وسط السماء بالنهار}
\MsL{17}{AB03-F039r-L17}{وبالليل درجة وتد الليل الأرض وبالليل والنهار سهم الآباء وقس ما يخرج من}
\MsL{18}{AB03-F039r-L18}{التسيير إلى سني الكوكب الأكبر خصص في الموضع الذي تسييره للأب وإن وافقها}
\MsL{19}{AB03-F039r-L19}{أو كانت قريبة من عددها دلّ على هلال الأب واستعمل\PN{100} أدلّاء الأمّ مثل ذلك}
\MsL{20}{AB03-F039r-L20}{ومنهم من كان ينظر بالنهار ويبدأ بالشمس للأب. فإن كانت مرتبطة بالطالع}
\MsL{21}{AB03-F039r-L21}{ينظر إلیها أحد أصحاب الحصص التي تلي موضعها استعملها في عمر الأب فإن لم}
\begin{PNotes}
\PNote{98}{AB03-n098}{B: add. \RL{\ArN نظر}}
\PNote{99}{AB03-n099}{B: \RL{\ArN لغمر؟}}
\PNote{100}{AB03-n100}{B: add. \RL{\ArN في}}
\end{PNotes}
\Folio{039v}{39v}
\MsL{1}{AB03-F039v-L01}{نجدها كذلك نظر إلى زحل فإن لم نجده بهذه الحال نظر إلى وسط وتد السماء}
\MsL{2}{AB03-F039v-L02}{بعد أن لا يكون لسهم الأب لهذه الحال وتبتدئ بالليل، ثمّ بزحل، ثمّ بالشمس، ثمّ بسهم}
\MsL{3}{AB03-F039v-L03}{الأب، ثمّ بوتد الأرض وتجعل الكواكب الذي له حصّة ونظر معيّنًا لدليل القمر بحسب}
\Chap{I.19}{<19> في الإخوة}
\MsL{4}{AB03-F039v-L04}{سنيه وتستعمل في أدلّاء الأمّهات مثل ذلك. \Hd{\LR{\LatinF \Sg{〈}19\Sg{〉}}}\Hd{ في الإخوة} وأمّا أمر الإخوة}
\MsL{5}{AB03-F039v-L05}{فإنّه أن فحص الإنسان عنه فحصًا مجملًا ولم يتجاوز فيه المقدار الممكن لطلب عدد}
\MsL{6}{AB03-F039v-L06}{الإخوة على الاستقصاء واحدًا واحدًا أمكنه أن يعلم أمر الإخوة الذين هم لأمّ واحدة}
\MsL{7}{AB03-F039v-L07}{على المذهب الطبيعيّ من البرج الذي في وسط السماء من موضع الأمر\PN{101} أعني \BF{B23v}من موضع}
\MsL{8}{AB03-F039v-L08}{الزهرة بالنهار وموضع القمر بالليل من البروج. ولمّا كانت هذه البروج والبرج}
\MsL{9}{AB03-F039v-L09}{الذي يرتفع إلى مكانه ممّا\PN{102} الموضع الذي تعرف منه ولدًا لأمّ وحب أن يكون}
\MsL{10}{AB03-F039v-L10}{ذلك موضعًا اللإخوة. وإنّ شاركت السعود هذه المواضع في الشكل حكمنا من ذلك}
\MsL{11}{AB03-F039v-L11}{على كثرة الإخوة وجعلنا حديثنا من ذلك من عدد هذه الكواكب ونظرنا هل هي}
\MsL{12}{AB03-F039v-L12}{في بروج ذوات صورة واحدة أم في بروج ذوات جسدين. فإن استعلت على}
\MsL{13}{AB03-F039v-L13}{الموضع الكواكب المنحسة أو قابلت على القطر وكانت في البيت المقابل له دلّت}
\MsL{14}{AB03-F039v-L14}{على قلّة الإخوة لا سيّما متى دخلت معها الشمس وإن كانت الشمس على الأوتاد}
\MsL{15}{AB03-F039v-L15}{وبخاصّة على الطالع. ثمّ كان الكوكب الذي في موضع من هذه المواضع زحل كان}
\MsL{16}{AB03-F039v-L16}{دليل ذلك في الأبكار وإن كان المرّيخ قلل عدد الإخوة بالموت. فأمّا الكواكب}
\MsL{17}{AB03-F039v-L17}{المعطيّة الإخوة فإنّ حالها في وضعها من العالم إن كانت حالًا محمودة. فينبغي أن}
\MsL{18}{AB03-F039v-L18}{\{نظنّ\} بالإخوة أنّها\PN{103} يكونون ذوي أقدار أشرافًا. وإن كان الأمر على ضدّ ذلك}
\MsL{19}{AB03-F039v-L19}{كانوا وضعاء \{خاملي\} الذكر وإن استعلت الكواكب المنحسة على المعطيّة أو كانت}
\MsL{20}{AB03-F039v-L20}{تابعة لها كانوا أقصار الأعمار أيضًا والذي يعطي منها الذكور هي الكواكب المذكّرة}
\MsL{21}{AB03-F039v-L21}{في حالها في العالم والتي تعطي منها للأناث هي المؤنّثة في حالها في العالم وأيضًا}
\begin{PNotes}
\PNote{101}{AB03-n101}{T: \RL{\ArN الأمّ}}
\PNote{102}{AB03-n102}{B: \RL{\ArN هما}}
\PNote{103}{AB03-n103}{BT: \RL{\ArN أنّهم}}
\end{PNotes}
\Folio{040r}{40r}
\MsL{1}{AB03-F040r-L01}{فإنّ الشرقيّة تعطي الإخوة الأوّل والغربيّة تعطي الإخوة الأواخر ومع ذلك أيضًا}
\MsL{2}{AB03-F040r-L02}{فإنّ الكواكب المعطيّة إن كانت مشاركة في الشكل للبرج المتولّيّ لأمر الإخوة مشاركة}
\MsL{3}{AB03-F040r-L03}{متّفقة دلّت على تحابّ الإخوة. وإن كانت موافقة لسهم السعادة دلّت على إشراك}
\MsL{4}{AB03-F040r-L04}{الإخوة في أمورهم وإن كانت في البروج التي لارتباط بينها وكانت\PN{104} على الحال المضادّة}
\MsL{5}{AB03-F040r-L05}{للحال التي ذكرنا دلّت على تعادي الإخوة تحاسدهم وأنّهم أصحاب مكر في أكثر الأمر. فإن}
\MsL{6}{AB03-F040r-L06}{أراد الإنسان أن يفحص عن كلّ واحد من أمور الإخوة بأكثر من هذا المقدار أمكنه}
\MsL{7}{AB03-F040r-L07}{أن يحدّث\PN{105} على ذلك أيضًا بأن يقيم الكوكب \BF{B24r}المعطي مقام الطالع وينظر في سائر أمورهم}
\MsL{8}{AB03-F040r-L08}{كما ينظر في المولد. وأمّا غير بطليموس من القدماء فإنّهم كانوا ينظرون في أمور الإخوة}
\MsL{9}{AB03-F040r-L09}{من البيت الثالث من الطالع والمرّيخ وسهم الإخوة وأربابها. فإن كان حلّ هذه الأدلّاء}
\MsL{10}{AB03-F040r-L10}{لا سيّما الكواكب الكبيرة الحصص في هذه المواضع في البروج الكثيرة النسل وبرج بيت}
\MsL{11}{AB03-F040r-L11}{الإخوة كذلك فإنّها تدلّ على كثرة الإخوة. وإن كانت على خلاف ذلك دلّت على قلّتهم}
\MsL{12}{AB03-F040r-L12}{وإن كان جلّها لاسيّما الكثيرة الحصص منها في البروج الذكرة كان أكثرهم ذكورًا}
\MsL{13}{AB03-F040r-L13}{وإن كان في إناث كان أكثرهم إناثًا وقل في سلامتهم وموتهم وإصلاح حالهم}
\MsL{14}{AB03-F040r-L14}{وفسادها من صلاح حال الأولاد وبيت الإخوة وفسادهم وانظر إلى قوى}
\MsL{15}{AB03-F040r-L15}{\{زلي\} مثلّثة الطالع. فإن كان في الطالع فإنّه بكر وإن كان في وسط السماء فإنّه}
\MsL{16}{AB03-F040r-L16}{بكر أورابع وإن كان في السابع بكر أو سابع وإن كان في غير وتد فوق الأرض}
\MsL{17}{AB03-F040r-L17}{فقدمته إلى الطالع فإن كان بينهما نحس دلّ على أنّه سقط أو مات بعض}
\MsL{18}{AB03-F040r-L18}{الأولاد وإن كانوا سعودًا أو كانوا في بعض حصص أنفسهم وكانوا أقوياء}
\MsL{19}{AB03-F040r-L19}{وبروج موافقة في أماكن محمودة دلّوا على عدد ما ولد قبله. وإن كان ربّ المثلّثة}
\MsL{20}{AB03-F040r-L20}{تحت الأرض فيفهن الطالع إليه على ما ذكرنا أوّلًا ونظروا أيضًا لعدد ما ولد}
\MsL{21}{AB03-F040r-L21}{قبله من عدد الكواكب التي توجد فيما بين وسط السماء والطالع فإن لم يوجد}
\begin{PNotes}
\PNote{104}{AB03-n104}{B: \RL{\ArN أو كانت}}
\PNote{105}{AB03-n105}{B: \RL{\ArN يحدس}}
\end{PNotes}
\Folio{040v}{40v}
\MsL{1}{AB03-F040v-L01}{هناك كوكب فإنّه بكر وإن كان ولد قبله مات وفي عدد ما يولد بعده من عدد}
\MsL{2}{AB03-F040v-L02}{الكواكب التي توجد فيما بين الطالع إلى الرابع وقالوا إنّ القمر وطالع المولود}
\MsL{3}{AB03-F040v-L03}{إذا كانا في بروج كثيرة النسل كانت الأمّ كثيرة الأولاد. وإن وجد القمر في}
\MsL{4}{AB03-F040v-L04}{العقرب في ثلاثة أجزاء منها لم يكن لها ولد وإن كانوا لم يبقوا وماتوا وإن}
\MsL{5}{AB03-F040v-L05}{وجدنا \{ربّي\} مثلّثة الطالع كلّيهما في الثالث من الأوتاد لم يكن له إخوة وإن}
\MsL{6}{AB03-F040v-L06}{كانوا تفرقوا ولم يجتمعوا في الأرض. وإن وجدت ربّ بيت الإخوة في الطالع}
\MsL{7}{AB03-F040v-L07}{فإنّه يكون فردًا لا أخ له. وإن \BF{B24v}وجد في وسط السماء فإن له أخًا أكبر منه وإن}
\MsL{8}{AB03-F040v-L08}{وجد في الآباء كان له أخ أصغر منه. وإن وجد بيت الإخوة\PN{106} فأنّه يدلّ على أنّه}
\MsL{9}{AB03-F040v-L09}{[[يكون]] برجًا ذا حسيدين وصاحبه أيضًا في برج ذي حسيدين لاسيّما إذا}
\MsL{10}{AB03-F040v-L10}{كان البروج الذي فيه سهم الإخوة كذلك فإنّه يدلّ على أنّه يكون له إخوة من غير}
\MsL{11}{AB03-F040v-L11}{أبيه أو من غير أمّه فلينطر في تحابهم ومعاداتهم إلى سهم الإخوة وربّه فإن نظر}
\MsL{12}{AB03-F040v-L12}{بعضها إلى بعض من تثليث أو تسديس دلّ على التحاب. وإن كان من تربيع}
\MsL{13}{AB03-F040v-L13}{دلّ على صداقة وسطه إن كان من مقابلة\PN{107} دلّ على المعادات وإن لم يتناظرًا}
\MsL{14}{AB03-F040v-L14}{دلًّا\PN{108} على الاغتراب وقلّة الاجتماع وكذلك يعلم من نظر الأغلب على الطالع}
\MsL{15}{AB03-F040v-L15}{بحصصه إلى الأغلب على بيت الإخوة وسهم الإخوة والمرّيخ وإن كان}
\MsL{16}{AB03-F040v-L16}{صاحب سهم السعادة مع سهم الإخوة وصاحب سهم الإخوة مع صاحب سهم السعادة}
\MsL{17}{AB03-F040v-L17}{دلّ على نفع بعضهم لبعض وربّما أصابوا من إخوانهم مالًا كثيرًا. وينظر لموت}
\MsL{18}{AB03-F040v-L18}{الإخوة من نحو سهم صاحب سهم الإخوة بمقارنة وتربيع ومقابلة النحوس لاسيّما}
\MsL{19}{AB03-F040v-L19}{إذا كان النحس في مقامه الأوّل أو راجعًا وخاصّة إذا لم ينظر المشتري}
\MsL{20}{AB03-F040v-L20}{إلى سهم فإن نظر أحد النحسين إلى السهم إذا كان في حصص عطارد}
\MsL{21}{AB03-F040v-L21}{دلّ على قتل إخوته وإن كان أحد النحسين في البروج الذي فيه عطارد}
\begin{PNotes}
\PNote{106}{AB03-n106}{B: add. \RL{\ArN كذلك}}
\PNote{107}{AB03-n107}{B: \RL{\ArN المقابلة}}
\PNote{108}{AB03-n108}{B: \RL{\ArN دلّ}}
\end{PNotes}
\Folio{041r}{41r}
\MsL{1}{AB03-F041r-L01}{ونظر إلى سهم الإخوة من المقابلة والتربيع فكذلك نفعل ويستعمل في}
\Chap{I.20}{<20> في الذكور والإناث}
\MsL{2}{AB03-F041r-L02}{الإخوات الزهرة مكان عطارد. \Hd{\LR{\LatinF \Sg{〈}20\Sg{〉}}}\Hd{ في الذكور والإناث} وأمّا معرفة}
\MsL{3}{AB03-F041r-L03}{الذكورة والإناث من المواليد فإنّما تكون من النيّرين جمعيًّا ومن الطالع}
\MsL{4}{AB03-F041r-L04}{ومن الكواكب التي لها في ذلك حصص ويعلم ذلك خاصّة من الحال في}
\MsL{5}{AB03-F041r-L05}{وقت سقوط النطفة. فأمّا على الجملة فيعلم من \BF{B25r}الحال في وقت الولادة أيضًا.}
\MsL{6}{AB03-F041r-L06}{وينبغي أن نتفقّد هذه المواضع التي ذكّرنا والكواكب التي تدبّرها فإن كانت}
\MsL{7}{AB03-F041r-L07}{جميعًا أو أكثرها مذكّرة كان المولود ذكرًا وإن كانت مؤنّثة فأنثى وينبغي}
\MsL{8}{AB03-F041r-L08}{أن يفعل من المواضع والكواكب المذكّرة والمؤنّثة على الجملة التي ذكرناها فيما}
\MsL{9}{AB03-F041r-L09}{تقدّم من القول في صدر هذا الكتاب أعني من طبيعة البروج التي هي فيها ومن}
\MsL{10}{AB03-F041r-L10}{الطبيعة الكواكب أنفسها ومن حالها بقياسها إلى العالم وذلك أنّها تذكر إذا}
\MsL{11}{AB03-F041r-L11}{كانت شرقيّة من الأفق وتأنث إذا كانت غربيّة ومن حالها أيضًا بقياسها}
\MsL{12}{AB03-F041r-L12}{إلى الشمس لأنّها تذكر إذا كانت منسوبة إلى الغدوات وتؤنث إذا كانت منسوبة}
\MsL{13}{AB03-F041r-L13}{إلى العشيّات فيحدث من جميع ذلك أيّ جنس يغلب على المولود في الأمر}
\Chap{I.21}{<21> في ولادة التوأمين}
\MsL{14}{AB03-F041r-L14}{الأكثر. \Hd{\LR{\LatinF \Sg{〈}21\Sg{〉}}}\Hd{ في ولادة التوأمين} وكذلك ينبغي أن تتفقّد في ولادة}
\MsL{15}{AB03-F041r-L15}{التوأمين أو أكثر من التوأمين أمر هذه الثلاثة مواضع. أعني النيّرين والطالع}
\MsL{16}{AB03-F041r-L16}{ومن شأن\PN{109} هذا لأمر أن يعرض عند التركيبات فإذا كان موضعان من هذه}
\MsL{17}{AB03-F041r-L17}{المواضع أو ثلاثتها بروجًا ذوات جسدين لا سيّما إذا عرضت للكواكب المدبّرة لها}
\MsL{18}{AB03-F041r-L18}{أيضًا هذا العارض \{بعينه\} أو كان بعضها في بروج ذوات جسدين أو كان}
\MsL{19}{AB03-F041r-L19}{بعضها إثنين إثنين أو أكثر من إثنين في التدبير فإن كانت المواضع الرئيسيّة}
\MsL{20}{AB03-F041r-L20}{بروجًا ذوات جسدين وكانت مشاركة في الشكل وكانت مع ذلك الكثير}
\MsL{21}{AB03-F041r-L21}{من الكواكب. فإن المولودين يكونون أكثر من إثنين وقد يزيد في عددهم الكوكب}
\begin{PNotes}
\PNote{109}{AB03-n109}{T: \RL{\ArN شيئين}}
\end{PNotes}
\Folio{041v}{41v}
\MsL{1}{AB03-F041v-L01}{الذي له هذه الخاصّيّة ويزيد في معرفة التذكير والتأنيث منهم جنس}
\MsL{2}{AB03-F041v-L02}{الكواكب التي شاركت الشمس والقمر والطالع في الشكل الذي يزيد في}
\MsL{3}{AB03-F041v-L03}{توليد الذكور والإناث التي تقدّم ذكرها إن لم تكن هذه الحال شاملة}
\MsL{4}{AB03-F041v-L04}{لبرج الطالع \BF{B25v}مع النيّرين لكن يوجد بدله على هذه الحال برج وسط السماء}
\MsL{5}{AB03-F041v-L05}{كان ما بلده\PN{110} الأمّهات في أكثر الأمر توأمًا أو أكثر من ذلك وممّا يخصّ في}
\MsL{6}{AB03-F041v-L06}{ولود ثلاثة ذكورة تامّة سويّة الخلق أن يكون الأمر على مثله ويكون}
\MsL{7}{AB03-F041v-L07}{لها نطروس\PN{111} وهم الملوك بأن يكون زحل والمشتري والمرّيخ في}
\MsL{8}{AB03-F041v-L08}{المواضع التي ذكرنا في بروج ذوات جسدين ويخصّ في ولود ثلاث}
\MsL{9}{AB03-F041v-L09}{إناث أن يكون الأمر على مثله في كون الحارطين\PN{112} أن تكون الزهرة}
\MsL{10}{AB03-F041v-L10}{والقمر مع عطارد وهو مؤنّث في المواضع التي ذكرنا ويخصّ في مولود}
\MsL{11}{AB03-F041v-L11}{ذكرين وأنثى أن يكون الأمر على مثله في كون الذي يقال لهم \{دياسفودا\}\PN{113}}
\MsL{12}{AB03-F041v-L12}{بأن يكون زحل والمشتري والزهرة في هذه المواضع ولدين\PN{114} أنثيين}
\MsL{13}{AB03-F041v-L13}{وذكرًا واحدًا إذ كان الأمر على مثله في كون الذي يقال له لمنظر\PN{115} وفودس\PN{116}}
\MsL{14}{AB03-F041v-L14}{بأن تكون الزهرة والقمر والمرّيخ في هذه المواضع ويعرض في أكثر}
\MsL{15}{AB03-F041v-L15}{الأمر في هذه الحالات أن لا يكون المولود تامًّا وأن تكون}
\MsL{16}{AB03-F041v-L16}{ولادة مع بعض الآفات العارضة في البدن وأن يحث في}
\MsL{17}{AB03-F041v-L17}{بعض أعضائه أشياء من غير سبب ظاهرًا ويكون ممّا لا يتوقّع}
\Chap{I.22}{<22> في المشوّهي الخلق}
\MsL{18}{AB03-F041v-L18}{أو يكون به بعض هذه الأعراض على جهة الشيء الذي ينظهر\PN{117}. \Hd{\LR{\LatinF \Sg{〈}22\Sg{〉}}}\Hd{ في المشوّهي الخلق} وممّا ليس بخارج عن هذا العرض القول}
\MsL{19}{AB03-F041v-L19}{في مشوّهي الخلق فنقول أوّلًا إنّ النيّرين يوجدان في أكثر}
\MsL{20}{AB03-F041v-L20}{الأمر في مثل هؤلاء زائلين أو غير مربوطين بالطالع ويحتوي}
\begin{PNotes}
\PNote{110}{AB03-n110}{T: \RL{\ArN تلد}}
\PNote{111}{AB03-n111}{T: \RL{\ArN أناقطرس}}
\PNote{112}{AB03-n112}{T: \RL{\ArN الحارطس}}
\PNote{113}{AB03-n113}{T: \RL{\ArN ديوسقور}}
\PNote{114}{AB03-n114}{T: \RL{\ArN فيلد}}
\PNote{115}{AB03-n115}{T: \RL{\ArN ديميطر}}
\PNote{116}{AB03-n116}{T: \RL{\ArN وقوريس}}
\PNote{117}{AB03-n117}{T: \RL{\ArN يعرض}}
\end{PNotes}
\Folio{042r}{42r}
\MsL{1}{AB03-F042r-L01}{على الأوتاد الكواكب المنحسة. فإذا وجدنا الأمر على هذه الحال فقد}
\MsL{2}{AB03-F042r-L02}{ينبغي إذ كان ذلك سببًا\PN{118} يكون كثيرًا في المواليد الوضيعة الشقيّة}
\MsL{3}{AB03-F042r-L03}{البخت. فإن لم تكن تلك المواليد لمشوّهي الخلق إن تنظر في الاجتماع}
\MsL{4}{AB03-F042r-L04}{والاستقبال التقدّم فيعلم الكوكب الذي\PN{119} لموضع أحدهما}
\MsL{5}{AB03-F042r-L05}{ومدي\PN{120} في موضع النيّرين في وقت المولد فإن كان مواضعها}
\MsL{6}{AB03-F042r-L06}{في \BF{B26r}المولد وموضع القمر أو موضع الطالع إمّا جميعها وإمّا أكثرها}
\MsL{7}{AB03-F042r-L07}{غير مرتبطة بموضع الاجتماع أو الاستقبال فإنّ المولود}
\MsL{8}{AB03-F042r-L08}{يكون مبدّل الخلقة. فإن كانت هذه الأشياء كذلك ووجدنا}
\MsL{9}{AB03-F042r-L09}{النيّرين مع ذلك في البروج ذوات الأربع قوائم أو في بروج}
\MsL{10}{AB03-F042r-L10}{سبعيّة وكان ما على الأوتاد من الكواكب كواكب منحسة فإنّ}
\MsL{11}{AB03-F042r-L11}{المولود ليس يكون من إنسان فإن لم يشهد للنيّرين شيء}
\MsL{12}{AB03-F042r-L12}{من المسعدة وشهدت لها النحوس كانت المولود غير مستأنس}
\MsL{13}{AB03-F042r-L13}{وكانت طبيعته من الطبائع الوحشيّة الضارة. فإن شهد}
\MsL{14}{AB03-F042r-L14}{النيّرين إمّا المشتري وإمّا الزهرة كان المولود من الحيوان}
\MsL{15}{AB03-F042r-L15}{الذي يألف الناس مثل الكلاب والنموس وما أشبهها فإن شهد}
\MsL{16}{AB03-F042r-L16}{عطارد وكان المولود من الحيوان الذي يحتاج إليه اناس}
\MsL{17}{AB03-F042r-L17}{وليستعملونها مثل الخنازير والدجاج والبقر والمعز وما}
\MsL{18}{AB03-F042r-L18}{أشبهها. فإن كان النيّران في بروج صورها على صور}
\MsL{19}{AB03-F042r-L19}{الناس وكانت سائر الأشياء على مثل ذلك كان المولود}
\MsL{20}{AB03-F042r-L20}{من \LR{\Sg{〉}}\LR{\LatinF …}\LR{\Sg{〈}}\PN{121} أو ممّا يجاور الناس إلّا أنّه يكون مشوّه الخلقة}
\MsL{21}{AB03-F042r-L21}{مبدّلًا ومعرفة الحال في ذلك وخاصّة في هذا الموضع أيضًا}
\begin{PNotes}
\PNote{118}{AB03-n118}{T: \RL{\ArN شيئًا}}
\PNote{119}{AB03-n119}{BT: \RL{\ArN المدبّر}}
\PNote{120}{AB03-n120}{T: \RL{\ArN مدبّري}}
\PNote{121}{AB03-n121}{space reserved; B: om.; T: \RL{\ArN الناس}}
\end{PNotes}
\Folio{042v}{42v}
\MsL{1}{AB03-F042v-L01}{تكون من صور البروج التي يكون فيها النحوس المحتويّة\PN{122} على النيّرين}
\MsL{2}{AB03-F042v-L02}{أو على الأوتاد وإن لم يشهد في هذا الموضع أيضًا شيء من}
\MsL{3}{AB03-F042v-L03}{المسعدة لشيء من المواضع التي ذكرنا كان المولود غير ناطق}
\MsL{4}{AB03-F042v-L04}{وكانت خلقته مبدّلة بالحقيقة على التمام فإن شهد المشتري}
\MsL{5}{AB03-F042v-L05}{أو الزهرة كانت خاصّة المولود الذي به العاهة أن يكون مكرمًا\PN{123}}
\MsL{6}{AB03-F042v-L06}{موقّرًا وذلك أنّه يعرض من ذلك الصورة التي يقال لها ارماطروديطس\PN{124}}
\MsL{7}{AB03-F042v-L07}{ومعناها الصورة المنسوبة إلى عطارد والزهرة والصورة التي تسمّى}
\MsL{8}{AB03-F042v-L08}{ادفوفراطيافوا\PN{125} وما أشبهها فإن شهد عطارد أيضًا كان}
\MsL{9}{AB03-F042v-L09}{المولود مع هذه الأشياء ميسر\PN{126} الأحكام\PN{127} وكان \LR{\Sg{〉}}\LR{\Sg{〉}} كسبه \LR{\Sg{〈}}\LR{\Sg{〈}}\PN{128} واختياله}
\MsL{10}{AB03-F042v-L10}{من هذه الأشياء إلّا أنّه يكون أدرد أصمّ وأمّا \BF{B26v}في سائر الأشياء}
\Chap{I.23}{<23> في الذين لا يغذون}
\MsL{11}{AB03-F042v-L11}{فإنّه يكون جيّد الطبع صاحب أمر\PN{129} وفكر\PN{130}. \Hd{\LR{\LatinF \Sg{〈}23\Sg{〉}}}\Hd{ في الذين لا يغذون}}
\MsL{12}{AB03-F042v-L12}{\Hd{والذين لا يتربون} وإذ قد بقي فيما يراد معرفته ممّا في وقت المولود}
\MsL{13}{AB03-F042v-L13}{القول فيمن لا يغتذي فقد ينبغي أن يعلم أنّ هذا القول في}
\MsL{14}{AB03-F042v-L14}{بعض الأحايين\PN{131} مقرون بالقول في العمر لأنّ ذلك ليس}
\MsL{15}{AB03-F042v-L15}{ببعيد من كلّ واحد في هذين الصنفين وفي بعض الأحاين}
\MsL{16}{AB03-F042v-L16}{يكون ممّا يناله \PN{132}على أنّ قوّة هذه الفحص تختلف بعض الاختلاف}
\MsL{17}{AB03-F042v-L17}{وذلك أنّ القول في العمر إنّما يكون في جميع من كان له عمر محسوس}
\MsL{18}{AB03-F042v-L18}{وهم الذين لا يقصر عمرهم عن دورة\PN{133} واحدة من أدوار الشمس.}
\MsL{19}{AB03-F042v-L19}{يعني الحول والعمر الذي بالغوة\PN{134} هو ما كان أقلّ من هذا}
\MsL{20}{AB03-F042v-L20}{أيضًا وهي الشهور والأيّام والساعات. وأمّا القول في الذين}
\MsL{21}{AB03-F042v-L21}{لا يغتذون فهو القول في الذين لايبلغون شيئًا من الأزمان التي}
\begin{PNotes}
\PNote{122}{AB03-n122}{T: \RL{\ArN المحفوفة}}
\PNote{123}{AB03-n123}{T: \RL{\ArN مكرها}}
\PNote{124}{AB03-n124}{T: \RL{\ArN ارمامرويطس}}
\PNote{125}{AB03-n125}{T: \RL{\ArN ارقوقراطاقوى}}
\PNote{126}{AB03-n126}{T: \RL{\ArN مفسّر}}
\PNote{127}{AB03-n127}{BT: \RL{\ArN الأحلام}}
\PNote{128}{AB03-n128}{right margin}
\PNote{129}{AB03-n129}{T: \RL{\ArN أدب}}
\PNote{130}{AB03-n130}{T: \RL{\ArN مكر}}
\PNote{131}{AB03-n131}{T: \RL{\ArN الأجناس} (wherever it occurs in this chapter)}
\PNote{132}{AB03-n132}{T: \RL{\ArN متباينا له}}
\PNote{133}{AB03-n133}{T: \RL{\ArN دون}}
\PNote{134}{AB03-n134}{T: \RL{\ArN بالقوّة}}
\end{PNotes}
\Folio{043r}{43r}
\MsL{1}{AB03-F043r-L01}{ذكرناها لإفراط الشرّ والمضرّة عليهم وبهذا صار الفحص عن العمر}
\MsL{2}{AB03-F043r-L02}{كثير الأجزاء وصار الفحص عمّن لا يغتذي مجملًا بسيطًا. فبالجملة}
\MsL{3}{AB03-F043r-L03}{متى كان أحد النيّرين على وتد من الأوتاد وكان أيضًا النحسين}
\MsL{4}{AB03-F043r-L04}{مشاركًا له في الطول إمّا بأن يكون معه درجة بدرجة وإمّا}
\MsL{5}{AB03-F043r-L05}{مقابلة وإمّا بأنّه على شكل متساويّ الساقين ولم يشاركه}
\MsL{6}{AB03-F043r-L06}{في الشكل شيء من المسعدة. وكان مدبّر موضع النيّرين}
\MsL{7}{AB03-F043r-L07}{موجودًا مع الكواكب المنحسة فإنّ المولود لا يغتذي ويموت}
\MsL{8}{AB03-F043r-L08}{من ساعتة. فإن لم يكن ذلك على شكل متساويّ الساقين}
\MsL{9}{AB03-F043r-L09}{لكن كان مطرح شعاع الكواكب المنحسة قريبًا من موضع النيّرين}
\MsL{10}{AB03-F043r-L10}{صاعدًا إليهما يريد أن يكون النيّران يقارباهن في الأجزاء}
\MsL{11}{AB03-F043r-L11}{التي بعد أجزاء النيّرين من البروج الذي هي فيه أوكانت}
\MsL{12}{AB03-F043r-L12}{المنحسة إثنين وكانت تضرّ بأحد النيّرين أو بهما جميعًا}
\MsL{13}{AB03-F043r-L13}{إمّا بأن يكونا صاعدين إلى موضعهما أو موضع أحدهما}
\MsL{14}{AB03-F043r-L14}{على ما ذكرنا بالمقابلة وإمّا \BF{B27r}بأن يقابلانه أو يقابلانهما جميعًا}
\MsL{15}{AB03-F043r-L15}{على القطر وإمّا بأن يكون كلّ وحد من المنحسين يضرّ}
\MsL{16}{AB03-F043r-L16}{بواحد من النيّرين أو يكون أحدهما مقابلًا لأحد النيّرين}
\MsL{17}{AB03-F043r-L17}{والآخر صاعدًا إلى موضع الآخر على جهة المقارنة وهو عند}
\MsL{18}{AB03-F043r-L18}{ذلك على ما وصفنًا في الأجزاء التي تتلوا أجزاء النيّرين لم}
\MsL{19}{AB03-F043r-L19}{يكن للمولود عمر وذلك أنّ كثرة ما في هذا الأمر يفسد}
\MsL{20}{AB03-F043r-L20}{ومن المضرّة يفسد ويخفي المنفعة والصلاح الذي يأتي من}
\MsL{21}{AB03-F043r-L21}{قبل بعد الموضع الذي يصعد إلى \LR{\Sg{〉}}\LR{\Sg{〉}} موضع النيّرين في باقاءالعمر والذي يضيّربالشمس خاصّة ممّا يكون في الموضع الذي يصعد إلى موضعها \LR{\Sg{〈}}\LR{\Sg{〈}} \PN{135}وبالقمر وزحل}
\begin{PNotes}
\PNote{135}{AB03-n135}{left margin}
\end{PNotes}
\Folio{043v}{43v}
\MsL{1}{AB03-F043v-L01}{وأمّا فيما يقابلهما على القطر أو فيما يعلوا عنهما فإنّ الأمر على عكس}
\MsL{2}{AB03-F043v-L02}{ذلك وهو أن يضرّ بالشمس زحل وبالقمر المرّيخ وخاصّة إن}
\MsL{3}{AB03-F043v-L03}{استولت على مكان النيّرين والطالع بأن تكون مدبّرة لها}
\MsL{4}{AB03-F043v-L04}{فإن كان لها مقابلتان على القطر وكان النيّران والنحسان}
\MsL{5}{AB03-F043v-L05}{على الأوتاد وعلى شكل متساويّ الساقين فإنّ الأجنّة}
\MsL{6}{AB03-F043v-L06}{حينئذ يولدون إمّا وهم موتى وإمّا وهم في حال بين الحيوة}
\MsL{7}{AB03-F043v-L07}{والموت. وإذا كانت هذه الأشياء كذلك وكان النيّران مضرّين}
\MsL{8}{AB03-F043v-L08}{عن أحد الكواكب المسعدة أو كانا يشاركانه في الشكل وكان}
\MsL{9}{AB03-F043v-L09}{الشعاع يصل إلى الأجزاء [[المتساويّة]] المتقدّمة فإنّ المولود}
\MsL{10}{AB03-F043v-L10}{يعيش بدّل\PN{136} عدد البروج التي فيما بين الهيلاج وبين}
\MsL{11}{AB03-F043v-L11}{شعاع الشمس\PN{137} الذي هو أقرب مثل ذلك من عدد الشهور}
\MsL{12}{AB03-F043v-L12}{أو الأيّام أو الساعات بحسب مقدار رداءة الحال وقوّة}
\MsL{13}{AB03-F043v-L13}{الأشياء التي هي أسباب ذلك. فإن كان شاع المنحسة}
\MsL{14}{AB03-F043v-L14}{يصل إلى الأجزاء المتقدّمة فإنّ المولود يعيش بدّل عدد}
\MsL{15}{AB03-F043v-L15}{الدرج التي فيما بين الهيلاج وبين شعاع النحس الذي هو}
\MsL{16}{AB03-F043v-L16}{أقرب مثل ذلك من عدد الشهور أو الأيّام أو الساعات}
\MsL{17}{AB03-F043v-L17}{بحسب مقدار رداءة الحال وقوّة الأشياء \BF{B27v}التي هي أسباب ذلك}
\MsL{18}{AB03-F043v-L18}{وإن كان شعاع المنحسة يصل إلى الأشياء\PN{138} المتقدّمة للنيّرين}
\MsL{19}{AB03-F043v-L19}{وكان شعاع السعود يصل إلى الأجزاء التي تتلوها فإنّ}
\MsL{20}{AB03-F043v-L20}{المولود يغتذي ويعيش وأيضًا فإن المنحسة من الكواكب}
\MsL{21}{AB03-F043v-L21}{إن كانت عالية\PN{139} على المسعدة المشاركة في الشكل صار [المولود]}
\begin{PNotes}
\PNote{136}{AB03-n136}{T: \RL{\ArN مدّة}}
\PNote{137}{AB03-n137}{T: \RL{\ArN المنحس}}
\PNote{138}{AB03-n138}{B: \RL{\ArN الأجزاء}}
\PNote{139}{AB03-n139}{T: \RL{\ArN غالية}}
\end{PNotes}
\Folio{044r}{44r}
\MsL{1}{AB03-F044r-L01}{سيئ\PN{140} الحال قليلًا\PN{141}. فإن استعلت المسعدة صار المولود إلى غير والديه بأن طلع}
\MsL{2}{AB03-F044r-L02}{أحد المسعدة مع القمر أو اتّصل به وكان أحد المنحسة في المغيب فإنّ المولود}
\MsL{3}{AB03-F044r-L03}{يربّيه والده وعلى هذه الجهة يكون النظر فيما\PN{142} يكثر من الأولاد. فإنّه إن كان}
\MsL{4}{AB03-F044r-L04}{أحد الكواكب التي لها اشتراك في الأشكال على جهتين أو أكثر من ذلك في المغرب}
\MsL{5}{AB03-F044r-L05}{فإنّ المولود يولّد وهو بين الميت والحيّ أو يولّد بمنزلة بضعة من لحم أو ناقص غير}
\MsL{6}{AB03-F044r-L06}{تامّ الخلقة. فإن استعلى عليه أحد المنحسة لم يترب المولود ولم يكن له عمر بحسب}
\MsL{7}{AB03-F044r-L07}{ما يوجبه ذلك. وقالت الأوائل في التربية أيضًا أنظر إلى مثلّثات الطالع وإلى}
\MsL{8}{AB03-F044r-L08}{مثلّثات الشمس بالنهار وإلى أرباب مثلّثات القمر ومثلّثات سهم السعادة بالليل}
\MsL{9}{AB03-F044r-L09}{فإذا وجدت أكبرها أو الكبير الولايّة منها في مواضع جيّدة من الفلك والبروج}
\MsL{10}{AB03-F044r-L10}{\LR{\Sg{〉}}\LR{\Sg{〉}} سيّما \LR{\Sg{〈}}\LR{\Sg{〈}} \PN{143}وسط السماء والحادي عشر والخامس بقيًا من المناحس والعيوب فإنّه يدلّ على التربية}
\MsL{11}{AB03-F044r-L11}{والعمر. فإن كانت الزهرة بقيت من العيوب وكانت في الخمس عشرة درجة التي}
\MsL{12}{AB03-F044r-L12}{أعلى الطالع وينظر إليها المشتري وهو نقيّ فإنّه يدلّ على التربية وإن لم يكن القمر}
\MsL{13}{AB03-F044r-L13}{منحوسًا ولم يكن في الأوتاد شيء من النحوس وكان في الطالع أو في وسط السماء}
\MsL{14}{AB03-F044r-L14}{أو الحادي عشر كوكب من كواكب النهاريّة بالنهار والليليّة بالليل فإنّم يدلّ على}
\MsL{15}{AB03-F044r-L15}{التربية. وإذا كانت أرباب المثلّثات في وسط السماء والحادي عشر والخامس \BF{B28r}}
\MsL{16}{AB03-F044r-L16}{فإنّها تدلّ على اليسر في التربية. وإذا كان في السابع \{والفنّ\} على مراده التربية لا سيّما}
\MsL{17}{AB03-F044r-L17}{إذا كانت نحوسًا وإذا كانت الأوتاد قائمة والكواكب بريئة من المناحس نقيّة من العيب}
\MsL{18}{AB03-F044r-L18}{في صحة من مزاجاتها وأرباب البيوت في مواضع جياد في حدود السعود وشعاعها}
\MsL{19}{AB03-F044r-L19}{النيّران يوافق بعضها بعضًا والطالع فإنّه يدلّ على طول العمر. وانظر إلى أكبر الكواكب}
\MsL{20}{AB03-F044r-L20}{ولايّة من أرباب مثلّثات الطالع ومثلّثات الشمس بالنهار والقمر وسهم السعادة بالليل}
\MsL{21}{AB03-F044r-L21}{واحدًا كان أو إثنين فإن اتّفق واحدًا وكان ساقطًا مقارنًا لدرجة الحسن وكانت}
\begin{PNotes}
\PNote{140}{AB03-n140}{T: \RL{\ArN شيء}}
\PNote{141}{AB03-n141}{T: \RL{\ArN دليلًا}}
\PNote{142}{AB03-n142}{B: \RL{\ArN وفيما}}
\PNote{143}{AB03-n143}{left margin}
\end{PNotes}
\Folio{044v}{44v}
\MsL{1}{AB03-F044v-L01}{درجة الطالع مقارنة للقمر أيضًا فإنّه دليل على فلة البقاء وكذلك إذا قارن النحس}
\MsL{2}{AB03-F044v-L02}{درجة الطالع أو درجة النيّر الذي له الحيّز وكذلك أرباب مثلّثات النيّر وكذلك}
\MsL{3}{AB03-F044v-L03}{إذا زالت أرباب مثلّثات الطالع كلّها وزال النيّران وأربابها وزال الاجتماع}
\MsL{4}{AB03-F044v-L04}{والامتلاء وربّه فإنّه الإبقاء للمولود وإن كان المستوليّ ساقطًا عن الأوتاد منحوسًا إلّا}
\MsL{5}{AB03-F044v-L05}{أنّ بينه وبين النحس بعدة الدرجات فإن المولود يبقى بقدر ما بين المستوليّ وبين}
\MsL{6}{AB03-F044v-L06}{النحس من الدرج أيّامًا أو شهورًا أو سنينًا. وكذلك إذا كان النحس في الوتد}
\MsL{7}{AB03-F044v-L07}{والمستوليّ ساقط عن الوتد أو درجة الطالع أو درجة القمر تنظر\PN{144} بنحس فإنّه}
\MsL{8}{AB03-F044v-L08}{يبقى بقدر تلك الدرجة\PN{145} أيّامًا أو شهورًا أو سنينًا. وإذا كان القمر منحوسًا ونحن}
\MsL{9}{AB03-F044v-L09}{في وتد لا سيّما الطالع ووسط السماء ولم ينظر إليه سعد دلّ على قلّة البقاء لا سيّما إذا كان}
\MsL{10}{AB03-F044v-L10}{النحس مع درجة القمر أو درجة الطالع أو تربيعهما أو مقابلتهما في الدرجة والحدّ}
\MsL{11}{AB03-F044v-L11}{لا سيّما إذا كان مع درجة الطالع والقمر جميعًا فإنّه إذا كان كذلك ثمّ نظرت إليه}
\MsL{12}{AB03-F044v-L12}{السعود لم يعش. فإن كان القمر في السابع منحوسًا من محاسدة أو تربيع أو}
\MsL{13}{AB03-F044v-L13}{مقابلة درجة نحس ولم ينظر \BF{B28v}إليه سعد فبت عليه أنّه لا يعيس. وإن كان القمر في الوتد}
\MsL{14}{AB03-F044v-L14}{الرابع وكان النحس قابلته قابلته ولا سيّما إذا\PN{146} كان نحس آخر مجامعًا للقمر لم ينج المولود}
\MsL{15}{AB03-F044v-L15}{وأمّه جميعًا في اليوم الذي يولد فيه من الموت إلّا أن يكون القمر في بيته أو شرفه}
\MsL{16}{AB03-F044v-L16}{أو مثلّثته وينظر سعد بقوّة إلى ذلك النحس من التثليث فإنّه يدلّ على نجاتهما}
\MsL{17}{AB03-F044v-L17}{لأن السعد في موضع خطوطه يزداد قوّة وخيرًا والشمس في موضع الخطوط الذي}
\MsL{18}{AB03-F044v-L18}{يزداد نقصانًا من الشرّ ولا سيّما إذا كان المرّيخ في موج أنثى وزحل في برج}
\MsL{19}{AB03-F044v-L19}{ذكر وإن كان القمر كما وصفت وهو محصور بين نحسين ولم ينظر إليه سعد وكان}
\MsL{20}{AB03-F044v-L20}{مع ذلك ناقص الضوء فإنّه مع ما يدلّ على موت الولد\PN{147} يدلّ على ذهاب بصره قبل ذلك.}
\MsL{21}{AB03-F044v-L21}{وإن كان نحس في الطالع ونحس في السابع والقمر في وتد الطالع أو وسط السماء}
\begin{PNotes}
\PNote{144}{AB03-n144}{B: \RL{\ArN تنظر بل؟}}
\PNote{145}{AB03-n145}{B: \RL{\ArN الدرج}}
\PNote{146}{AB03-n146}{B: \RL{\ArN إن}}
\PNote{147}{AB03-n147}{B: \RL{\ArN المولود}}
\end{PNotes}
\Folio{045r}{45r}
\MsL{1}{AB03-F045r-L01}{ولم تنظر إليه السعود فإنّه يدلّ مع الموت على إخراجه من بيت الآباء قبل موته}
\MsL{2}{AB03-F045r-L02}{لا سيّما إذا كان القمر في حدّ النحوس وربّ ذلك الحدّ في وتد فإنّه أردى وأضرّ}
\MsL{3}{AB03-F045r-L03}{وإن سقطت أكثر المثلّثات لا سيّما في السادس والثاني عشر وبخاصّة الكبير الولايّة}
\MsL{4}{AB03-F045r-L04}{منها فإنّها تدلّ على قلّة البقاء والوقت في ذلك أن تنظر إلى السعد إذا كان في}
\MsL{5}{AB03-F045r-L05}{وتد أو مكان جيّد قويّ أو إلى الكوكب المتّصل بدرجة الطالع والمضيين وسهم}
\MsL{6}{AB03-F045r-L06}{السعادة فإنّه يدلّ على قلّة مقدار دوره سنين أو شهور أو أيّامًا. وإذا كان النحوس}
\MsL{7}{AB03-F045r-L07}{في الطالع قريبة منه بلدرجات فإنّه وإن كانت هناك سعود أقرب بالدرجات}
\MsL{8}{AB03-F045r-L08}{إلى درجة الطالع فإن كانت سائر الأدلّاء قويّة على التربية والعمر فإنّه يدلّ منقوص العقل}
\MsL{9}{AB03-F045r-L09}{والجوارح كثير الآفات إلى آخر عمره وعلى حال لا يكون إلّا قليل العمر وإذا وجدت}
\MsL{10}{AB03-F045r-L10}{البرج الذي فيه الاجتماع \BF{B29r}أو الامتلاء والناظر من أصحاب الحدود الذي يسمّى المعنق}
\MsL{11}{AB03-F045r-L11}{وربّ حدّ النيّرين والطالع وأرباب المثلّثات التي ذكرنا لا سيّما مثلّثة الشمس}
\MsL{12}{AB03-F045r-L12}{بالنهار والقمر بالليل وسهم السعادة كلّها منحوسة فاقض عليه أنّه يخرج من بطن أمّه}
\Chap{I.24}{<24> في العمر}
\MsL{13}{AB03-F045r-L13}{ميتًا أو يخرج ميتًا من بطن أمّه. \Hd{\LR{\LatinF \Sg{〈}24\Sg{〉}}}\Hd{ في العمر} وأمّا القول في الأعراض التي تكون}
\MsL{14}{AB03-F045r-L14}{بعد الولادة فإنّ المقدّم منها هو القول في العمر والنظر في ذلك ليس بمطلق ولا}
\MsL{15}{AB03-F045r-L15}{مرسل بل يؤخذ ويعلم من استيلاء المواضع الرئيسيّة وغلبتها وذلك شيء يستعمل فيه}
\MsL{16}{AB03-F045r-L16}{جهات كثيرة. وأمّا الجهة التي توافق رأينا وتكون تابعة للمجرى الطبيعيّ فهي ما إنّا}
\MsL{17}{AB03-F045r-L17}{ذاكره وقائل في ذلك أنّ النظر في ذلك يكون من مواضع الهيلاج ومن المواضع المستوليّة}
\MsL{18}{AB03-F045r-L18}{على مواضع الهيلاج ومن المواضع أو الكواكب المقابلة أيضًا ويميّز كلّ واحد من هذه الأمور}
\MsL{19}{AB03-F045r-L19}{على ما أصف لك أوّل ما ينبغي أن نعلم مواضع الهيلاج هي التي يجب أن تكون فيها لا محالة}
\MsL{20}{AB03-F045r-L20}{الكواكب الذي تصيّر له ولايّة الهيلاج وهي برج الطالع الذي من خمس درج يتقدّم}
\MsL{21}{AB03-F045r-L21}{فتطلع من الأفق قبله أو\PN{148} إلى الخامس والعشرين الدرجة الباقيّة التي تطلع من}
\begin{PNotes}
\PNote{148}{AB03-n148}{B: om.}
\end{PNotes}
\Folio{045v}{45v}
\MsL{1}{AB03-F045v-L01}{بعد وتصعد والدرجة التي عن يمينه هذه الثلاثين درجة على تسديسها وهي}
\MsL{2}{AB03-F045v-L02}{لبيت السعادة والتي على تربيعها وهي لوسط السماء والتي على تثليثها وهي}
\MsL{3}{AB03-F045v-L03}{للموضع الذي يسمّى باوس\PN{149} وهو بيت \{الدين\} والمسعد\PN{150} والتي على مقابلتها وهي للغارب}
\MsL{4}{AB03-F045v-L04}{والذي ينبغي أن يقدّم من هذه أيضًا لقوّته \BF{B29v}في الاستيلاء ما كان في وسط السماء ثمّ ما كان}
\MsL{5}{AB03-F045v-L05}{في الطالع ثمّ ما كان في بيت السعادة وهو الذي يصعد إلى وسط السماء ثمّ ما كان}
\MsL{6}{AB03-F045v-L06}{في الغارب ثمّ ما كان متقدّمًا \LR{\Sg{〉}}\LR{\Sg{〉}} لوسط السماء \LR{\Sg{〈}}\LR{\Sg{〈}}\PN{151} وهو التاسع الذي قلنا إنّه بيت. وأمّا جميع ما تحت}
\MsL{7}{AB03-F045v-L07}{الأرض فينبغي أن نترك استعماله في مثل هذا الأمر الشريف ما خلا الخمس}
\MsL{8}{AB03-F045v-L08}{والعشرين درجة التي هي دائمًا تطلع وتظهر فوق الأرض. وأمّا ممّا فوق الأرض}
\MsL{9}{AB03-F045v-L09}{فينبغي أن يستعمل البرج الذي ليس له ارتباط بالطالع ولا البرج}
\MsL{10}{AB03-F045v-L10}{الذي قد طلع قبل الطالع وهو الذي يسمّى بيت الشقاوة وذلك أنّه مع زواله}
\MsL{11}{AB03-F045v-L11}{عن الوتد فإنّ القوّة التي تبعث منه الأرض تكدرها البخار الغليظ المظلم}
\MsL{12}{AB03-F045v-L12}{المتصاعد إليه من رطوبات الأرض قبل طلوع الشمس مع أنّ ذلك يكون ممّا ترى}
\MsL{13}{AB03-F045v-L13}{من ألوان ما في هذا البيت من الكواكب ومقادير أجرامها على غير ما هي عليه بالطبع. ثمّ}
\MsL{14}{AB03-F045v-L14}{ينبغي أن نجعل الهيالجات الأربعة الأشياء الرئيسيّة أعني الشمس والقمر والطالع}
\MsL{15}{AB03-F045v-L15}{و\{بيتهم\} البخت الذي هو سهم السعادة ومدبّرات مواضع هذه الأشياء. فأمّا سهم البخت}
\MsL{16}{AB03-F045v-L16}{فيعلم بأن يؤخذ أبدًا العدد الذي يحصل من موضع الشمس إلى موضع القمر بالنهار}
\MsL{17}{AB03-F045v-L17}{وبالليل جميعًا وتبتدئ من الطالع فنعدّ منه مثل ذلك العدد آخذين إلى الجهة}
\MsL{18}{AB03-F045v-L18}{التي تتلوا من البروج فنكون قد عرفنا موضع هذا السهم وإنّما فعلنا ذلك}
\MsL{19}{AB03-F045v-L19}{ليكون حال الشمس في قياسها إلى الطالع كحال القمر في قياسه إلى سهم البخت}
\MsL{20}{AB03-F045v-L20}{ويصيّر السهم بمنزلة طالع القمر وخليق أن يكون إلى هذا المعنى بعينه ذهب الذي}
\MsL{21}{AB03-F045v-L21}{قال إنّه ينبغي أن يعدّ في المواليد التي تكون في الليل من القمر \BF{B30r}إلى الشمس جهات\PN{152}}
\begin{PNotes}
\PNote{149}{AB03-n149}{i.e. \RL{\ArN قوس}; T: \RL{\ArN ثاوس}}
\PNote{150}{AB03-n150}{B: \RL{\ArN بيت الذين والسعر}; T: \RL{\ArN بيت الأسد}}
\PNote{151}{AB03-n151}{right margin}
\PNote{152}{AB03-n152}{B: add. \RL{\ArN ما يفعل}}
\end{PNotes}
\Folio{046r}{46r}
\MsL{1}{AB03-F046r-L01}{بالنهار ثمّ يبتدئ من الطالع فيعدّ منه مثل ذلك العدد الذي حفاق}
\MsL{2}{AB03-F046r-L02}{الجهة التي عددناها بالنهار أعني إلى جهة المتقدّمة من البروج. فإنّا إذا فعلنا ذلك}
\MsL{3}{AB03-F046r-L03}{خرج لنا موضع سهم البخت على هذه الجهة الموضع نفسه الذي كان لنا وصارت الحال}
\MsL{4}{AB03-F046r-L04}{في مشاركة في الشكل تلك الحال. وينبغي أن يختار من هذه الأشياء أيضًا أمّا}
\MsL{5}{AB03-F046r-L05}{بالنهار وبالشمس أوّلًا إن كانت في مواضع الهيلاج وإن لم يكن كذلك اختير}
\MsL{6}{AB03-F046r-L06}{القمر وإلّا فالكوكب الذي له حصص أكثر في تدبير موضع الشمس وموضع}
\MsL{7}{AB03-F046r-L07}{الاجتماع المتقدّم قبل ذلك والطالع أعني إذا كانت له من الحصص الخمس}
\MsL{8}{AB03-F046r-L08}{التي يكون التدبير عليها ثلاث حصص في أحد المواضع التي ذكرنا}
\MsL{9}{AB03-F046r-L09}{أو أكثر من ثلاث أو كان ذلك في جميعها له فإن لم يتهيّأ له ذلك}
\MsL{10}{AB03-F046r-L10}{أخذنا درجة الطالع. وأمّا في الليل فينبغي أن يختار أوّلًا القمر}
\MsL{11}{AB03-F046r-L11}{ثمّ الشمس إن كانت في البروج التي تطلع وتصعد فوق الأرض}
\MsL{12}{AB03-F046r-L12}{من بيت الطالع ثمّ الكواكب\PN{153} الذي له حصص أكثر في موضع القمر}
\MsL{13}{AB03-F046r-L13}{وموضع الاستقبال الذي كان قبل ذلك وموضع سهم البخت.}
\MsL{14}{AB03-F046r-L14}{فإن لم يتهيّأ ذلك أخذنا في \LR{\Sg{〉}}\LR{\Sg{〉}} آخر \LR{\Sg{〈}}\LR{\Sg{〈}}\PN{154} الأمر إن كان الذي تقدّم المولد}
\MsL{15}{AB03-F046r-L15}{اجتماعًا فالطالع وإن كان الذي تقدّم استقبالًا بسهم البخت.}
\MsL{16}{AB03-F046r-L16}{وإن كان النيّران جميعًا والمدبّر الذي بحسب الحيّز الملائم في}
\MsL{17}{AB03-F046r-L17}{بعض مواضع الهيلاج فينبغي أن نأخذ ما كان من النيّرين}
\MsL{18}{AB03-F046r-L18}{في موضع أكبر\PN{155} رئاسة وأقوى وذلك بيّن فيما رتبنا في مواضع}
\MsL{19}{AB03-F046r-L19}{الهيلاج. وأمّا المدبّر فإنّما ينبغي أن تختاره على النيّر إذا كان}
\MsL{20}{AB03-F046r-L20}{في موضع أكبر\PN{156} رئاسة وأقوى من مواضع التميّز\PN{157} على ذلك الترتيب}
\MsL{21}{AB03-F046r-L21}{وكان له بحسب الخيّرين جميعًا حصص في التدبير فبهذا}
\begin{PNotes}
\PNote{153}{AB03-n153}{B: corrected to \RL{\ArN الكوكب}}
\PNote{154}{AB03-n154}{\RL{\ArN سائر} corrected interlinear}
\PNote{155}{AB03-n155}{B: \RL{\ArN أكثر}}
\PNote{156}{AB03-n156}{B: \RL{\ArN أكثر}}
\PNote{157}{AB03-n157}{B: \RL{\ArN النيّر}}
\end{PNotes}
\Folio{046v}{46v}
\MsL{1}{AB03-F046v-L01}{العمل يعلم الهيلاج الدالّ على العمر. وأمّا غير بطليموس من الأوائل}
\MsL{2}{AB03-F046v-L02}{فإنّهم أجمعوا على أنّه لا يستحقّ أحد هذه التي ذكرنا أن يكون}
\MsL{3}{AB03-F046v-L03}{\BF{B30v}هيلاجًا حتّى تشهد له أصحاب الحصص في موضعه والشهادة}
\MsL{4}{AB03-F046v-L04}{هي أنّ تشاركه في الشكل من أحد الاشتراكات أيّها كان إمّا}
\MsL{5}{AB03-F046v-L05}{من التربيع أو المقابلة أو التثليث والتسديس أو النظر أو}
\MsL{6}{AB03-F046v-L06}{الاستواء أو القوّة أو من الإمرة والمطيعة التي تسمّى العالية}
\MsL{7}{AB03-F046v-L07}{والمتحفّظة فالذي يكون منها هذه الشهادة هو الهيلاج}
\MsL{8}{AB03-F046v-L08}{والذي يشترك معهنّ أصحاب الحصص في موضعه هو}
\MsL{9}{AB03-F046v-L09}{العين الذي يعطي السنين وقد رأينا بطليموس ذكر في هذه}
\MsL{10}{AB03-F046v-L10}{المواضع أصحاب الحصص وهي المتولّيّة لتدبير هذه المواضع}
\MsL{11}{AB03-F046v-L11}{ولم يستشهد بها الهيلاج فيما بعد وقد يجوّز أن يكون قد}
\MsL{12}{AB03-F046v-L12}{أومأ إلى ذلك إنّما لم يخرج في الترجمة وما أشبه ذلك والذين}
\MsL{13}{AB03-F046v-L13}{قالوا بهذا القول جعلوا درجة الاجتماع أو الامئلًا الذي}
\MsL{14}{AB03-F046v-L14}{يقع قبل المولد في الهيلاجيّة مكان الكوكب الذي تكون له}
\MsL{15}{AB03-F046v-L15}{الحصص في تلك المواضع ولم يفرقوا فيما بين فوق}
\MsL{16}{AB03-F046v-L16}{الأرض وتحتها في مواضع الهيلاج وسنشرح رأيهم بعد}
\MsL{17}{AB03-F046v-L17}{فراغنا من رأي بطليموس في هذا الباب. قال بطليموس}
\MsL{18}{AB03-F046v-L18}{وإذ كنّا قد عرفنا من الهيلاج فينبغي أن نستعمل نوعي \{التكيسير\}\PN{158}}
\MsL{19}{AB03-F046v-L19}{وهما النوع الذي على توالي البروج والذي على تواليهما}
\MsL{20}{AB03-F046v-L20}{وخلاف التوالي. أمّا التوالي في البروج فقط فينبغي أن يستعمل}
\MsL{21}{AB03-F046v-L21}{عند الحال التي تسمّى مطرح الشعاع وذلك إذا كان الهيلاج في}
\begin{PNotes}
\PNote{158}{AB03-n158}{B: \RL{\ArN التسير}; T: \RL{\ArN الهيلاج}}
\end{PNotes}
\Folio{047r}{47r}
\MsL{1}{AB03-F047r-L01}{المواضع الشرقيّة أعني التي من وسط السماء إلى الطالع وأمّا التي تستعمل}
\MsL{2}{AB03-F047r-L02}{على توالي البروج وعلى الجهة المتقدّمة من البروج جميعًا معًا}
\MsL{3}{AB03-F047r-L03}{ففي الحال التي يقال لها الوقتيّة وذلك يكون إذا كان الهيلاج}
\MsL{4}{AB03-F047r-L04}{في المو[[ا]]ضع التاسع الذي يسمّى ناوس الزائل عن وسط السماء.}
\MsL{5}{AB03-F047r-L05}{وإذا كان ذلك كذلك فإن البرج الذي يدلّ على القتل يكون}
\MsL{6}{AB03-F047r-L06}{إمّا بحسب الهيلاج الذي يكون عن النوع الذي إلى الجهة}
\MsL{7}{AB03-F047r-L07}{المتقدّمة من البروج فدرجة الغارب فقط لأنّها تخفي صاحب}
\MsL{8}{AB03-F047r-L08}{الحياة\PN{159}. وأمّا درج الكواكب التي تلقى الهيلاج \BF{B31r}أو تشهد له}
\MsL{9}{AB03-F047r-L09}{فإنّما ينقص ويزيد فقط في الأعداد التي تجتمع إلى وقت}
\MsL{10}{AB03-F047r-L10}{غيبوبة الهيلاج من غير أن يقتل لأنّها ليست هي التي تصيّر}
\MsL{11}{AB03-F047r-L11}{إلى موضع الهيلاج بل هو الذي يصيّر إلى مواضعها والتي تزيد}
\MsL{12}{AB03-F047r-L12}{هي الكواكب المسعدة والتي تنقص هي الكواكب المنحسة.}
\MsL{13}{AB03-F047r-L13}{وأمّا عطارد أيضًا فيصيّر مع الكواكب التي تشاركها في}
\MsL{14}{AB03-F047r-L14}{الشكل من تلك الكواكب فإن كانت تزيد زاد وإن كانت}
\MsL{15}{AB03-F047r-L15}{تنقص نقص وعدد الزيادة والنقصان يعلم من وضع درجة}
\MsL{16}{AB03-F047r-L16}{كلّ واحد منها أو شعاعه وذلك أنّه بحسب أزمنة الساعات}
\MsL{17}{AB03-F047r-L17}{التي لكلّ درجة منها. أمّا إذا كان الوقت نهاريًّا بأزمنة الساعات}
\MsL{18}{AB03-F047r-L18}{النهاريّة. وإن كان الوقت ليلًا بأزمنة\PN{160} الليليّة يكون عدد السنين}
\MsL{19}{AB03-F047r-L19}{التامّة وذلك شيء ينبغي أن يفهم أنّه قيل على أنّها في الطالع}
\MsL{20}{AB03-F047r-L20}{ثمّ تنقص منه على حسب تباعده عن الطالع حتّى ينتهي}
\MsL{21}{AB03-F047r-L21}{إلى الغارب ولا يبقى من ذلك شيء ومعرفة ذلك بأن}
\begin{PNotes}
\PNote{159}{AB03-n159}{B: \RL{\ArN الحيوة}}
\PNote{160}{AB03-n160}{B: add. \RL{\ArN الساعات}}
\end{PNotes}
\Folio{047v}{47v}
\MsL{1}{AB03-F047v-L01}{تأخذ بعد الدرجة من الغارب على ما بيّنا في الباب الأوّل من هذا}
\MsL{2}{AB03-F047v-L02}{الكتاب في ساعات بعد الدرج من الأوتاد فما كان من ساعة}
\MsL{3}{AB03-F047v-L03}{زمانيّة ما في نصف سدس أزمان ساعات تلك الدرجة}
\MsL{4}{AB03-F047v-L04}{الليليّة أو النهاريّة بحسب الوقت فما بلغ ذلك فهو مقدار}
\MsL{5}{AB03-F047v-L05}{الزيادة النقصان. وأمّا بحسب النوع الذي يكون إلى الجهة}
\MsL{6}{AB03-F047v-L06}{التالية من موضع الهيلاج فإن الذي يقتل هو مواضع الكواكب}
\MsL{7}{AB03-F047v-L07}{المنحسة أعني موضع زحل والمرّيخ. إمّا إذا لقيت الهيلاج}
\MsL{8}{AB03-F047v-L08}{بأجسامها وأمّا إذا لقيت شعاعها من أيّ موضع من}
\MsL{9}{AB03-F047v-L09}{هذين الموضعين كأن أعني موضع التربيع وموضع المقابلة}
\MsL{10}{AB03-F047v-L10}{وربّما كان ذلك من التسديسات التي تسمع وتطيع أو}
\MsL{11}{AB03-F047v-L11}{تنظر وتستوي في القوّة. ونفس تربيع موضع الهيلاج}
\MsL{12}{AB03-F047v-L12}{الذي في الجهة التالية من البروج تقتل أيضًا وربّما}
\MsL{13}{AB03-F047v-L13}{فعل ذلك التسديس إذا كان نحسًا وكان في بروج كثيرة}
\MsL{14}{AB03-F047v-L14}{المطالع ويفعله \BF{B31v}التثليث إذا كان بحسًا وكان في بروج}
\MsL{15}{AB03-F047v-L15}{قصيرة المطالع. وإذا كان القمر الهيلاج قد يقتل موضع}
\MsL{16}{AB03-F047v-L16}{الشمس أيضًا وليس ينبغي أن يتوهّم أنّ عتده\PN{161} المواضع}
\MsL{17}{AB03-F047v-L17}{نقتل\PN{162} في كلّ وقت لا محالة بل إنّما يفعل ذلك إذا كانت}
\MsL{18}{AB03-F047v-L18}{منحوسة فقط وذلك أنّه يعوق عن ذلك أيضًا أن يكون}
\MsL{19}{AB03-F047v-L19}{الحدّ لأحد المسعدة بأن يلقى أحد المسعدة الشعاع إمّا}
\MsL{20}{AB03-F047v-L20}{على تربيع وإمّا على تثليث وإمّا على مقابلة إمّا إلى الدرجة}
\MsL{21}{AB03-F047v-L21}{المقابلة بعينها وإمّا إلى الدرج التي يتلوها\PN{163} بعد إلّا يتجاوّز\PN{164}}
\begin{PNotes}
\PNote{161}{AB03-n161}{B: \RL{\ArN عدّه}; T:\RL{\ArN هذه }}
\PNote{162}{AB03-n162}{B: \RL{\ArN تقتل}}
\PNote{163}{AB03-n163}{T: \RL{\ArN تتلوها}}
\PNote{164}{AB03-n164}{B: \RL{\ArN يتجاوّزها}}
\end{PNotes}
\Folio{048r}{48r}
\MsL{1}{AB03-F048r-L01}{في مطرح شعاع المشتري بأكثر من إثنتا عشرة درجة وفي مطرح}
\MsL{2}{AB03-F048r-L02}{شعاع الزهرة بأكثر من ثمان درجات. وكذلك أيضًا متى ألقى}
\MsL{3}{AB03-F048r-L03}{الهيلاج والكوكب الذي يلقاه التقاء جسمانيًّا وليكن\PN{165} عرضهما}
\MsL{4}{AB03-F048r-L04}{جميعًا عرضًا واحدًا بعينه. فإذا كانت الكواكب التي تخلّص}
\MsL{5}{AB03-F048r-L05}{وتعين أو التي يقبل كوكبين أو أكثر من ذلك في كلّ واحد}
\MsL{6}{AB03-F048r-L06}{من هذين الأمرين فينبغي أن ينظر إلى أغلب الصنفين}
\MsL{7}{AB03-F048r-L07}{بحسب الكثرة كثرة الأشياء التي تعينها وبحسب قوّتها أمّا}
\MsL{8}{AB03-F048r-L08}{بحسب الكثرة. فإذا كان أحد الصنفين أكثر من الآخر بشيء}
\MsL{9}{AB03-F048r-L09}{محسوس. وأمّا بحسب القوّة. فإذا كانت الكواكب التي تعين}
\MsL{10}{AB03-F048r-L10}{أو التي تقتل في مواضعها الملائمة لها ولم يكن الآخر}
\MsL{11}{AB03-F048r-L11}{في مثل هذه المواضع وبخاصّة إذا كانت كواكب أحد}
\MsL{12}{AB03-F048r-L12}{الصنفين شرقيّة وكواكب الصنف الآخر غربيّة وينبغي}
\MsL{13}{AB03-F048r-L13}{بالجملة إلّا يستعمل شيئًا ممّا كان منها تحت الشعاع لا في}
\MsL{14}{AB03-F048r-L14}{القتل ولا في المعونة. والخلاص إلّا أن يكون القمر هو}
\MsL{15}{AB03-F048r-L15}{الهيلاج فيكون موضع الشمس نفسه القاتل فيستعمل الكوكب}
\MsL{16}{AB03-F048r-L16}{النحس الذي معه إذا كان مفسدًا ولم يكن شيء من المسعدة}
\MsL{17}{AB03-F048r-L17}{مخلصًا له. وإذا كان السعد مكان ذلك النحس الذي}
\MsL{18}{AB03-F048r-L18}{قلنا مع الشمس كان مخلصًا له. وأمّا عدد السنين التي}
\MsL{19}{AB03-F048r-L19}{تكون وتحدث عن الأبعاد التي فيما بين موضع \BF{B32r}الهيلاج}
\MsL{20}{AB03-F048r-L20}{وبين الموضع القاتل فليس ينبغي أن نجعل أخذنا}
\MsL{21}{AB03-F048r-L21}{له على الإطلاق ولا كيف ما تفق\PN{166} على ما أخذ عن كثير}
\begin{PNotes}
\PNote{165}{AB03-n165}{T: \RL{\ArN لم يكن} (!)}
\PNote{166}{AB03-n166}{BT: \RL{\ArN اتّفق}}
\end{PNotes}
\Folio{048v}{48v}
\MsL{1}{AB03-F048v-L01}{من الناس أن زعموا أنّه يكون دائمًا عن أزمان [عن أرمان]}
\MsL{2}{AB03-F048v-L02}{مطالع درج ذلك البعد في الإقليم ولا أن نستعمل ذلك في شيء من}
\MsL{3}{AB03-F048v-L03}{الأوقات إلّا\PN{167} كان الطالع نفسه وهو موضع الهيلاج أو شيء}
\MsL{4}{AB03-F048v-L04}{من المواضع التي تصيّر إلى الطالع إذا كانت درجة الطالع}
\MsL{5}{AB03-F048v-L05}{الهيلاج. وأمّا إذا كان في وسط \{إلّا\} درجة وسط السماء فإنّا}
\MsL{6}{AB03-F048v-L06}{نستعمل في مقادير الأبعاد مطالع الفلك المستقيم وإذا}
\MsL{7}{AB03-F048v-L07}{كان موضع الهيلاج مائلًا عن الأوتاد إلى ناحية من النواحيّ}
\MsL{8}{AB03-F048v-L08}{استعملنا ما وصفناه في كتابنا الأوّل حيث ذكرنا مطالع}
\MsL{9}{AB03-F048v-L09}{البروج فيما بين أرباع الفلك حتّى نعلم من بعد كم زمان}
\MsL{10}{AB03-F048v-L10}{من أزمان معدّل النهار يصيّر الموضع القاتل التالي}
\MsL{11}{AB03-F048v-L11}{إلى موضع الهيلاج الذي هو الموضع المتقدّم حتّى يصل}
\MsL{12}{AB03-F048v-L12}{إلى الغارب فيكون مقادير الأبعاد بقدر معارب درجة}
\MsL{13}{AB03-F048v-L13}{الغارب إذا كان الهيلاج فيها. وذلك شيء قد شرحناه}
\MsL{14}{AB03-F048v-L14}{واستغنينا عن ذكره في هذا الموضع إذ كنّا قد فرغنا منه}
\MsL{15}{AB03-F048v-L15}{في الكتاب من أجل أنّه من الأشياء الهندسيّة التي كان}
\MsL{16}{AB03-F048v-L16}{يحبّ أن نشرح عملها فيه وقد بقي علينا تمييز كلّ واحد}
\MsL{17}{AB03-F048v-L17}{من أنواع الأشاء التي تتلقًّا الهيلاج وغروبه فما قدّمنا}
\MsL{18}{AB03-F048v-L18}{ذكره من الأشياء التي هي أقصر زمانًا ما كان منها قابلًا\PN{168}}
\MsL{19}{AB03-F048v-L19}{وما كان التي يقال لها \{طيمًا\} قطير\PN{169} فاس وهي الأوقات}
\MsL{20}{AB03-F048v-L20}{وسائر الأشياء التي يتلقّاه وتمييز ذلك يكون معرفة}
\MsL{21}{AB03-F048v-L21}{التلقّي الرديء الضارّ والجيّد التابع المعيّن الكائنين \LR{\Sg{〉}}على\LR{\Sg{〈}}}
\begin{PNotes}
\PNote{167}{AB03-n167}{B: add. \RL{\ArN إذا}}
\PNote{168}{AB03-n168}{B: \RL{\ArN قاتلًا}}
\PNote{169}{AB03-n169}{T: \RL{\ArN قليماقطر}}
\end{PNotes}
\Folio{049r}{49r}
\MsL{1}{AB03-F049r-L01}{الجهة التي ذكرنا آنفًا ومن قبل دخول السنين الآتيّة التي هي أوقات}
\MsL{2}{AB03-F049r-L02}{التحاويل فيما بعد التي تدلّ عليها ذلك الموضع الذي يتلقّا[ه] الهيلاج. وذلك أنّه}
\MsL{3}{AB03-F049r-L03}{إذا كانت مواضع التلقّي رديئة منحوسة وكانت مواضع الكواكب في دخول السنين التي}
\MsL{4}{AB03-F049r-L04}{تأتي من بعد \BF{B32v}تضرّ بالمواضع الرئيسيّة. فينبغي أن نظنّ أنّ الذي يكون عند ذلك موت}
\MsL{5}{AB03-F049r-L05}{بالحقيقة وإذا كان أحد هذين الأمرين صالحًا فينبغي أن نظنّ أنّ الذي يكون عند}
\MsL{6}{AB03-F049r-L06}{ذلك موت بالحقيقة وإذا كان أحد هذين الأمرين صالحًا فينبغي أن نظنّ أنّ}
\MsL{7}{AB03-F049r-L07}{الذي يكون هو الآفات الغطية الخطر. وإذا كان هذان الأمران جميعًا صالحين}
\MsL{8}{AB03-F049r-L08}{فينبغي أن نظنّ أنّ الذي يكون هو كسل وضعفًا ومضارّ ونفي بخروج أو مسير}
\MsL{9}{AB03-F049r-L09}{ويجوّز بدّل قوله يفي أن يقال انحطاط المراتب والغربة والأسفار وقد تعلم الأنواع}
\MsL{10}{AB03-F049r-L10}{الخاصّيّة التي تحدث من هذه الأشياء ومن مواضع الشيء التي تتلقّى لأمور}
\MsL{11}{AB03-F049r-L11}{المولّد وليس شيء يمنعنّا إذا شككنا في بعض الأوقات في المواضع التي يجب أن}
\MsL{12}{AB03-F049r-L12}{يكون لها الأمر في القتل أن تتفكّر وتنظرفي الأشياء التي تتلقّى الهيلاج في كلّ}
\MsL{13}{AB03-F049r-L13}{واحد من أمور المولود. فينبغي فيما يأتي فيها الأمر الموافق أمّا قد كان وقع وظهر}
\MsL{14}{AB03-F049r-L14}{من الأعراض وتتفقّد\PN{170} الأمر في جميعها إذا كان الأمر في تمام ما يحدث عنها متساويًّا}
\MsL{15}{AB03-F049r-L15}{في القوّة جميعها فنجعل فحصنا عن زيادتها ونقصانها على حسب ذلك. وينبغي أن}
\MsL{16}{AB03-F049r-L16}{نصحّح أقدار الاتّصالات وأبعاد ما بين الهيلاجات والكواكب بحسب}
\MsL{17}{AB03-F049r-L17}{عروض الكواكب على ما بيّنا في باب تحقيق الاتّصالات وأبعاد ما بين الكواكب}
\MsL{18}{AB03-F049r-L18}{قبل التسيير من الهيلاجات. فأمّا الذين ذهبوا المذهب الذي وصفنا فإنّهم قالوا}
\MsL{19}{AB03-F049r-L19}{في المعنى الذي يعطي السنّ وهو الكدخداه. أعني صاحب بيت وشرف ومثلّثة}
\MsL{20}{AB03-F049r-L20}{الهيلاج وقدّموا صاحب الحدّ أي هو لا نظر إلى الهيلاج فهو الكدخداه أو نظر}
\MsL{21}{AB03-F049r-L21}{إثنان جميعًا فأقر بهم بالدرجات مقدّمًا ومؤخرًا هو الكدخداه وهذا إذا}
\begin{PNotes}
\PNote{170}{AB03-n170}{B: \RL{\ArN نتفقّد}}
\end{PNotes}
\Folio{049v}{49v}
\MsL{1}{AB03-F049v-L01}{كان لكلّ واحد منهم خصّة في الطالع وسهم السعادة وبرج الاجتماع أو الامتلاء}
\MsL{2}{AB03-F049v-L02}{فهو إلى\PN{171} الوالي قرب نظره أو بعد بعد أن تنظر إلى الدرج فإن استوت شهادتهم}
\MsL{3}{AB03-F049v-L03}{فأقواهم \BF{B33r}موضعان. وإن كانوا في الموضع سواء فأقر بهم من الوتد أو ما يلي الوتد}
\MsL{4}{AB03-F049v-L04}{وإن كان بعضهم غربيًّا وبعضهم في حظوظ فالذي في حظوظه أوّلى والشرقيّ}
\MsL{5}{AB03-F049v-L05}{الآخر من الشمس أوّلًا وأعلم أنّ الشمس إذا كانت في الحمل والأسد تولت}
\MsL{6}{AB03-F049v-L06}{الهيلاجيّة والكدخداهيّة وكذلك القمر في الثور والسرطان وذلك أنّ}
\MsL{7}{AB03-F049v-L07}{لكلّ واحد منهما في برجيّة حظّين إذا لم يكن كدخداه ينظر. فإذا عرفت}
\MsL{8}{AB03-F049v-L08}{الكدخداه فإن كان في الطالع أو في وسط السماء مشرقًا في حظوظه أو}
\MsL{9}{AB03-F049v-L09}{في مشاكلته من التذكير والتأنيث والليليّ والنهاريّ بريئًا من \LR{\Sg{〉}}\LR{\Sg{〉}} الرجوع \LR{\Sg{〈}}\LR{\Sg{〈}}\PN{172}}
\MsL{10}{AB03-F049v-L10}{والاحتراق أو كان في الحادي عشر أو الخامس زائدًا على ما ذكرت بحلة}
\MsL{11}{AB03-F049v-L11}{واحدة فإنّه يدلّ على سنيه الوسطى. وإن كان في الثاني أو الثامن}
\MsL{12}{AB03-F049v-L12}{على ما وصفت في الطالع ووسط السماء وكان في السادس أو في الثاني}
\MsL{13}{AB03-F049v-L13}{عشر على ما وصفت في الحادي عشر والخامس فإنّه يدلّ على سنيه الصغرى}
\MsL{14}{AB03-F049v-L14}{وكلّما نقصت خلة انحطت مرتبة ممّا يدلّ عليها الكدخداه من أيّ موضع كان}
\MsL{15}{AB03-F049v-L15}{ومثال ذلك ما وصفت من كونه في الطالع ووسط السماء فإنّه إن لم يكن مشرقًا}
\MsL{16}{AB03-F049v-L16}{انحطت مرتبة فرجعت سنيه الكبرى إلى الوسطى وإن لم يكن مع ذلك في}
\MsL{17}{AB03-F049v-L17}{بعض حظوظه أو مشاكلته رجع إلى سنيه الصغرى وإن كان مع ذلك راجعًا}
\MsL{18}{AB03-F049v-L18}{رجع إلى سنيه صغرى شهورًا وإن كان مع ذلك راجعًا رجع إلى سنيه}
\MsL{19}{AB03-F049v-L19}{الصغرى شهورًا. وإن كان مع ذلك محترقًا رجع إلى سنيه الصغرى أيّامًا}
\MsL{20}{AB03-F049v-L20}{وقل في سائر هذه المواضع على هذا القياس وقالوا إنّ الرأس إذا كان}
\MsL{21}{AB03-F049v-L21}{مع الكدخداه بدون إثنتا عشرة درجة متقدّمًا أو متأخّرًا والكدخداه}
\begin{PNotes}
\PNote{171}{AB03-n171}{B: om.}
\PNote{172}{AB03-n172}{right margin}
\end{PNotes}
\Folio{050r}{50r}
\MsL{1}{AB03-F050r-L01}{سعد فإنّه يزيد مثل ربع سني الكدخداه التي عليها. وإن كان الذنب}
\MsL{2}{AB03-F050r-L02}{مكان الرأس فإنّه ينقصه مثل ذلك وإن كان الكدخداه نحسًا نقص من}
\MsL{3}{AB03-F050r-L03}{سنيه الرأس وزاد فيها الذنب على هذا القياس. وإن كانت الدرجات أقرب}
\MsL{4}{AB03-F050r-L04}{كان أوكد في كلّ الأمرين فإن كان الكدخداه في الحال التي لا \BF{B33v}تدلّ على شيء}
\MsL{5}{AB03-F050r-L05}{وكان المشتري والزهرة أو أحدهما في الطالع أو وسط السماء فإنّه يرجى للمولود}
\MsL{6}{AB03-F050r-L06}{مثل سنى أحدهما الصغرى أيهما اتّفق في الوتد الذي ذكرنا أن لا يكون \{منحرمًا\}}
\MsL{7}{AB03-F050r-L07}{من النحوس والقمر منحوس. وإذا عرفت سنى الكدخداه فانظر إلى الكوكب الناظر إلى}
\MsL{8}{AB03-F050r-L08}{كدخداه فإنّه أن نظر إليه سعد وكان يربا من النحوس والرجوع والاحتراق}
\MsL{9}{AB03-F050r-L09}{زاده مثل سنيه الصغرى. وإن نظر نحس من المقابلة أو التربيع نقصه مثل}
\MsL{10}{AB03-F050r-L10}{سنيه الصغرى. وإن كان للنحس في الميلاد حصص وهو مشرق في بعض}
\MsL{11}{AB03-F050r-L11}{حظوظه نقصه مثل سنيه الصغرى شهورًا وإن كان مع ذلك نظر تثليث}
\MsL{12}{AB03-F050r-L12}{أو تسديس نقصه أيّامًا فقط. وأمّا عطارد فإنّه إذا كان مع السعود و\{العطين\}}
\MsL{13}{AB03-F050r-L13}{مشرقًا أو مغربًا إعطأه مثل سنيه الصغرى. وإن كان مع النحوس والناقصين}
\MsL{14}{AB03-F050r-L14}{نقص مثل ذلك والشمس في المقابلة والمقارنة تنقص مثل النحوس}
\MsL{15}{AB03-F050r-L15}{وفي التربيع لا تريد الأشياء يسيّرًا والشمس أقطع ما تكون وأنحس على الكدخداه}
\MsL{16}{AB03-F050r-L16}{إذا كانت مع المرّيخ لاجتماعها\PN{173} حراريتهما. وإذا كان الكدخداه تحت الشعاع}
\MsL{17}{AB03-F050r-L17}{لم يقدر أن يقبل ما تريد السعود إلّا أن يكون \{صميمًا\} وإنّه حينئذ أسعد}
\MsL{18}{AB03-F050r-L18}{ما يكون نحسًا كان أو سعدًا. فإذا عرفت القمر من دلالة الكدخداه فسيّر الهيلاج}
\MsL{19}{AB03-F050r-L19}{بدرج المطالع بموضعه من الفلك على تلك الجهة إلى أجساد النحوس}
\MsL{20}{AB03-F050r-L20}{وشعاعاتها والذنب وتربيع القمر ومقابلته ومقارنته ودرجة الغارب}
\MsL{21}{AB03-F050r-L21}{إذا كان قد \{حاز\} ما تحت الأرض لكلّ درجة سنة فإن كانت تلك السنين}
\begin{PNotes}
\PNote{173}{AB03-n173}{B: \RL{\ArN لاجتماع}}
\end{PNotes}
\Folio{050v}{50v}
\MsL{1}{AB03-F050v-L01}{مشاكلة لسني الكدخداه أو قريبًا منها ولم يلق سعد شعاعه إلى ذلك}
\MsL{2}{AB03-F050v-L02}{الحدّ الذي انتهى فهه الهيلاج إلى شعاع النحس فإن المولود يهلك في تلك}
\MsL{3}{AB03-F050v-L03}{السنة لا سيّما إذا كان صاحب سنة المولود منحوسًا أو صاحب طالع منحوسًا}
\MsL{4}{AB03-F050v-L04}{من ذلك النحس وإن انتهى إلى نحس آخر \{يغيّر\} مشاكلة لسني الكدخداه كانت نكبة}
\MsL{5}{AB03-F050v-L05}{وكلّ ما انتهى \BF{B34r}الهيلاج إلى حسد أو شعاع نحس وألقى سعد شعاعه إلى}
\MsL{6}{AB03-F050v-L06}{تلك الدرجة أو لم يلق فهي سنة نكبة وأشد لذلك أن يكون النحس في}
\MsL{7}{AB03-F050v-L07}{مجرى درج الهيلاج يعني في العرض والقمر إذا كان هيلاجًا واتّصل بالطالع}
\MsL{8}{AB03-F050v-L08}{المنحوس أو بالشمس المنحوسة قبل\PN{174} وكذلك درجة.}
\Chap{I.25}{<25> في صورة بدن المولود وحليته ومزاجه}
\MsL{9}{AB03-F050v-L09}{\LR{\Sg{〉}}...\LR{\Sg{〈}}\PN{175} \Hd{\LR{\LatinF \Sg{〈}25\Sg{〉}}}\Hd{ في صورة بدن المولود وحليته ومزاجه} ولمّا كانت أحوال البدن متقدّمة}
\MsL{10}{AB03-F050v-L10}{لأحوال النفس بالطبيعة من أجل أن البدن أغلظ ينبغي لنا أن نتفقّد}
\MsL{11}{AB03-F050v-L11}{أمّا في أكثر الأمر فالأفق الشرقيّ والكوكب الذي يتلوه من الكواكب}
\MsL{12}{AB03-F050v-L12}{المتحيّرة أو المتولّيّة لتدبيره منها على الجهة التي قدّمنا ذكرها وإن}
\MsL{13}{AB03-F050v-L13}{كان هناك كوكب أيضًا. وأمّا على الأمر الجزئيّ\PN{176} أن تتفقّد من أمر القمر}
\MsL{14}{AB03-F050v-L14}{مثل ذلك أيضًا فإنّ صورة الأبدان وحالها إنّما يعلم من صورة هذين}
\MsL{15}{AB03-F050v-L15}{الموضعين جميعًا وصورة الكواكب المدبّرة لهما ومن امتزاج كلّ نوع منها ومن}
\MsL{16}{AB03-F050v-L16}{صورة أشكال الكواكب الثابتة أيضًا التي تطلع معها فأمّا بالذي يتقدّم}
\MsL{17}{AB03-F050v-L17}{قوّته منها قوّة غيره فهو الكواكب المتولّيّة للتدبير وبعدها الحال الذي}
\begin{PNotes}
\PNote{174}{AB03-n174}{B: \RL{\ArN قتل}}
\PNote{175}{AB03-n175}{here and in B: large blank space}
\PNote{176}{AB03-n176}{T: add. \RL{\ArN فينبغي}}
\end{PNotes}
\Folio{051r}{51r}
\MsL{1}{AB03-F051r-L01}{يخصّ المواضع أنفسها. وأمّا \BF{B34v}الشيء الجزئيّ الذي يمكن أن يخبر به أخبارًا}
\MsL{2}{AB03-F051r-L02}{بسيطًا مطلقًا من أمورها هو على ما أصف لك ما يقول في الكواكب أنّ زحل}
\MsL{3}{AB03-F051r-L03}{إذا وليّ التدبير وكان شرقيًّا كانت حلية المولود أنّه يكون العمل مخضّب\PN{177}}
\MsL{4}{AB03-F051r-L04}{البدن أسود الشعر جعد وجعد\PN{178} متكاثف متوسّط العينين معتدل العظم}
\MsL{5}{AB03-F051r-L05}{وأمّا مزاجه فيكون الغالب عليه البرد والرطوبة. وإذا كان غريبًا كان المولود}
\MsL{6}{AB03-F051r-L06}{آدم مهزولًا صغير الجسم بسط الشعر خفيفه حسن التأليف قياس بعضها}
\MsL{7}{AB03-F051r-L07}{إلى بعض على ما ينبغي أسود العينين وأمّا مزاجه فيكون الغالب عليه}
\MsL{8}{AB03-F051r-L08}{اليبس. وإذا كان المشتري هو المتولّيّ لهذه المواضع المذكّرة وكان شرقيًّا كان}
\MsL{9}{AB03-F051r-L09}{المولود أبيض اللون حسن الشعر متوسّط\PN{179} العينين حسن القامة والمقدار}
\MsL{10}{AB03-F051r-L10}{ذا وقار. وأمّا مزاجه فيكون الغالب عليه الحرارة والرطوبة وإذا كان غربيًّا}
\MsL{11}{AB03-F051r-L11}{كان المولود أبيض إلّا أنّه لا يكون حسن اللون مثل الذي ذكرنا ويكون سبط}
\MsL{12}{AB03-F051r-L12}{الشعر ويكون أنزع متوسّط العينين صغير الجسم وأمّا مزاجه فالغالب}
\MsL{13}{AB03-F051r-L13}{عليه الرطوبة. وإن\PN{180} كان المرّيخ أيضًا شرقيًّا ووليّ التدبير كانت الحلية فيما}
\MsL{14}{AB03-F051r-L14}{بين البياض والحمرة وكان حسن المقدار مخصب البدن أزرق العينين}
\MsL{15}{AB03-F051r-L15}{متكاثف الشعر متوسّطه. فأمّا المزاج فالغالب عليه الحرّ واليبس وإن}
\MsL{16}{AB03-F051r-L16}{كان غربيًّا \{كان\} المولود في حليته أحمر اللون حمرة مطلقة صغير البدن}
\MsL{17}{AB03-F051r-L17}{صغير العينين خفيف الشعر سبط أصهب وأمّا مزاجه فيكون الغالب}
\MsL{18}{AB03-F051r-L18}{عليه اليبس. وأمّا الزهرة فأفعالها شبيهة بأفعال المشتري إلّا أنّ}
\MsL{19}{AB03-F051r-L19}{الذي يحدث عنها يكون أجمل ويكون عليه قبول أكثر ويكون جماله بجمال}
\MsL{20}{AB03-F051r-L20}{النساء أشبه ويكون أحسن شكلًا وأحمد\PN{181} أخلاقًا \BF{B35r}وأنعم بدنًا وممّا يخصّها أن تجعل}
\MsL{21}{AB03-F051r-L21}{العينين شهلًا وبين حسنتين. وأمّا عطارد فإنّه إذا ولي التدبير وكان}
\begin{PNotes}
\PNote{177}{AB03-n177}{B: \RL{\ArN مخصب}}
\PNote{178}{AB03-n178}{B: add. \RL{\ArN ضوء}}
\PNote{179}{AB03-n179}{B: \Sg{〈}sic\Sg{〉} \RL{\ArN متوسط متوسّط}}
\PNote{180}{AB03-n180}{B: \RL{\ArN وإذا}}
\PNote{181}{AB03-n181}{T: \RL{\ArN وأجمل}}
\end{PNotes}
\Folio{051v}{51v}
\MsL{1}{AB03-F051v-L01}{شرقيًّا يكون المولود شبيهًا بلون العسل معتدل البدن حسن التأليف صغير}
\MsL{2}{AB03-F051v-L02}{العينين متوسّط الشعر. وأمّا مزاجه فيكون الغالب عليه الحرّ. وإذا كان غربيًّا كان}
\MsL{3}{AB03-F051v-L03}{المولود فيما بين الأدمة والصفرة مهزولًا دقيق الصوت غائر العينين حدقته}
\MsL{4}{AB03-F051v-L04}{شبيهة بحدق البقر\PN{182} مائلة إلى الحمرة. وأمّا مزاجه فيكون الغالب عليه اليبس}
\MsL{5}{AB03-F051v-L05}{وقد يعين كلّ واحد من هذه الكواكب الشمس والزهرة\PN{183} إذا شاركاها في الشكل}
\MsL{6}{AB03-F051v-L06}{أمّا الشمس فإنّها تعين على السنة\PN{184} والجمال وخصب البدن. وأمّا القمر فإنّه يعين على}
\MsL{7}{AB03-F051v-L07}{الاعتدال والعصف ورطوبة المزاج وبخاصّة إذا كان منصرفًا عنها. وأمّا على الأمر}
\MsL{8}{AB03-F051v-L08}{الجزئيّ فإنّ معونته\PN{185} تكون بخواصّ ما فيه من الضوء كما ذكرنا في أوّل هذا الكتاب}
\MsL{9}{AB03-F051v-L09}{وأيضًا فإنّ الكواكب بالجملة إذا كانت منسوبة إلى الغدوات وكانت ظهورها}
\MsL{10}{AB03-F051v-L10}{صيّرت الأبدان عظامًا. فإذا وقفت الوقوف الأوّل صيّرتها قويّة جلدة وإذا كانت}
\MsL{11}{AB03-F051v-L11}{\{راجعة\}\PN{186} صيّرتها معتدلة. وإذا وقفت الوقوف الثاني صيّرتها أضعف وإذا غابت}
\MsL{12}{AB03-F051v-L12}{صيّرتها خسيسة تنالها [[تنالها]] الآفات والشدائد وقد تعين أيضًا المواضع أنفسها}
\MsL{13}{AB03-F051v-L13}{التي يكون فيها على حلية البدن وصورته ومزاجه خاصّة كما قلنا آنفًا وأيضًا بالجملة}
\MsL{14}{AB03-F051v-L14}{فإنّ للربع الذي من الاعتدال الربيعيّ إلى الانقلاب الصيفيّ يصيّر المولود حسن اللون}
\MsL{15}{AB03-F051v-L15}{والمقدار مخصب البدن حسن العينين الغالب عليه الرطوبة والحرّ. وأمّا الربع}
\MsL{16}{AB03-F051v-L16}{الذي من الانقلاب الصيفيّ إلى الاعتدال الخريفيّ فإنّه يصيّر المولود متوسّط}
\MsL{17}{AB03-F051v-L17}{اللون معتدل في عظم \BF{B35v}البدن مخصب كبير العينين متكاثف الشعر جعده}
\MsL{18}{AB03-F051v-L18}{الغالب عليه الحرّ واليبس. وأمّا الربع الذي من الاعتدال الخريفيّ إلى الانقلاب}
\MsL{19}{AB03-F051v-L19}{الشتويّ فإنّه يصيّر لون المولود كلون العسل وتجعله قضيفًا دقيق الصوت}
\MsL{20}{AB03-F051v-L20}{عريض المنكبين متوسّط الشعر حسن العينين الغالب عليه اليبس والبرد. وأمّا}
\MsL{21}{AB03-F051v-L21}{الربع الذي من الانقلاب الشتويّ إلى الاعتدال الربيعيّ فإنّه يصيّر المولود أسود}
\begin{PNotes}
\PNote{182}{AB03-n182}{T: \RL{\ArN المغر}}
\PNote{183}{AB03-n183}{T: \RL{\ArN القمر}}
\PNote{184}{AB03-n184}{T: \RL{\ArN الحسنة}}
\PNote{185}{AB03-n185}{T: \RL{\ArN معرفته}}
\PNote{186}{AB03-n186}{B: \RL{\ArN \{قواجيد\}}; T: \RL{\ArN مقدّمة}}
\end{PNotes}
\Folio{052r}{52r}
\MsL{1}{AB03-F052r-L01}{اللون معتدل البدن سبط الشعر حفيفه حسن التأليف الغالب عليه البرد}
\MsL{2}{AB03-F052r-L02}{والرطوبة وأمّا على الأمر الجزئيّ فإنّ المواضع التي صورها شبيهة بصور}
\MsL{3}{AB03-F052r-L03}{الناس من الصور التي في فلك البروج والخارجة عنها تصيّر الأبدان حسنة}
\MsL{4}{AB03-F052r-L04}{التأليف معتدلة. وأمّا التي صورتها غير هذه الصورة فإنّها تغيّر الأبدان\PN{187} بحسب}
\MsL{5}{AB03-F052r-L05}{ما يلائم صورها وتشبّه الأعضاء بأعضائها مقدارًا ما من الشبه بأن تميلها إمّا}
\MsL{6}{AB03-F052r-L06}{[إمّا] إلى العظم والصغر وإمّا إلى القوّة والضعف وإمّا إلى جودة التأليف}
\MsL{7}{AB03-F052r-L07}{ورداءته يعني بذلك العقل الحسن. فالذي تميل فيها إلى العظم مثل الأسد}
\MsL{8}{AB03-F052r-L08}{والعذراء والرأس.\PN{188} والتي تميل إلى الصغر مثل السمكتين والسرطان والجدي منها}
\MsL{9}{AB03-F052r-L09}{وأيضًا فإنّ ما كان من مثل الحمل والثور والأسد يلي أعالمها\PN{189} ومقادمها فهي}
\MsL{10}{AB03-F052r-L10}{تميل الأعضاء إلى القوّة فتكون قوى البدن شديد أو ما كان منها يلي أسافلها}
\MsL{11}{AB03-F052r-L11}{ومؤاخرها يلي الأبدان إلى الضعف. وأمّا الذي على خلاف ذلك فمثل الرأس\PN{190}}
\MsL{12}{AB03-F052r-L12}{والدولو والعقرب فإن\PN{191} كان منها يلي مقادمها فإنّه يميل الأبدان إلى الضعيفة وما}
\MsL{13}{AB03-F052r-L13}{كان يلي مؤاخرها فإنّه يميل إلى خصب البدن وأيضًا فإنّ العذراء والميزان}
\MsL{14}{AB03-F052r-L14}{والرأس تميل الأبدان إلى الاعتدال \BF{B36r}وحسن التأليف والعقل. وأمّا العقرب}
\MsL{15}{AB03-F052r-L15}{[العقرب] والسمكتان والثور فإنّه يميلها إلى الخروج عن التأليف وعن الاعتدال}
\MsL{16}{AB03-F052r-L16}{الحمق وكذلك ينبغي أن تنظر في أمر سائر الصور وتقيس جميع ذلك بعضه إلى}
\MsL{17}{AB03-F052r-L17}{بعض فتحدث\PN{192} منه على الخواصّ التي تجتمع من امتزاج هذه الأشياء في حلية الأبدان}
\Chap{I.26}{<26> في العلل الأمراض العارضة في بدن المولود}
\MsL{18}{AB03-F052r-L18}{ومزاجاتها [إن شاء لله التعالى\PN{193}]. \Hd{\LR{\LatinF \Sg{〈}26\Sg{〉}}}\Hd{ في العلل الأمراض العارضة في بدن المولود}}
\MsL{19}{AB03-F052r-L19}{يجب علينا في هذا الموضع إذا أردنا علم جملة هذا الأمر أن تنظر\PN{194} في وتدي}
\MsL{20}{AB03-F052r-L20}{الأفق وهما الطالع والغارب وأن تنظر\PN{195} خاصّة في الغارب \LR{\Sg{〉}}\LR{\Sg{〉}} نفسه \LR{\Sg{〈}}\LR{\Sg{〈}}\PN{196} والموضع الذي قبل}
\MsL{21}{AB03-F052r-L21}{الغارب وهو بيت المرض وإن تتفقّد الكوكبين النحسين وتنظر كيف تجد حالهما}
\begin{PNotes}
\PNote{187}{AB03-n187}{T: \RL{\ArN اعتدال}}
\PNote{188}{AB03-n188}{T: \RL{\ArN الرامي}}
\PNote{189}{AB03-n189}{B: \RL{\ArN أعاليها}}
\PNote{190}{AB03-n190}{T: \RL{\ArN الرامي}}
\PNote{191}{AB03-n191}{B: add. \RL{\ArN ما}}
\PNote{192}{AB03-n192}{B: \RL{\ArN فتحدس}}
\PNote{193}{AB03-n193}{B: \RL{\ArN إن شاء الله زحل}}
\PNote{194}{AB03-n194}{B: \RL{\ArN ننظر}}
\PNote{195}{AB03-n195}{B: \RL{\ArN ننظر}}
\PNote{196}{AB03-n196}{left margin}
\end{PNotes}
\Folio{052v}{52v}
\MsL{1}{AB03-F052v-L01}{في قياسهما إلى هذه المواضع وذلك أنّهما متى كانا إذا قيسًا إلى الدرج التي}
\MsL{2}{AB03-F052v-L02}{ترتفع إلى المواضع التي ذكرنا ووجدًا جميعًا أو واحدًا منهما إمّا على الموضع نفسه}
\MsL{3}{AB03-F052v-L03}{وإمّا على تربيعه وإمّا على مقابلته. فينبغي أن تظنّ أن يصيب المولود آفات}
\MsL{4}{AB03-F052v-L04}{وعلل تدخل في بدنه لا سيّما إن كان النيّرين أو أحدهما كلاهما على الأوتاد على الجهة}
\MsL{5}{AB03-F052v-L05}{التي ذكرنا وهي إمّا على وتد واحد وإمّا على وتدين متقابلين فإنّه حينئذ \LR{\Sg{〉}}\LR{\Sg{〉}} ليس \LR{\Sg{〈}}\LR{\Sg{〈}}\PN{197} إنّما}
\MsL{6}{AB03-F052v-L06}{يقوى أحد الكواكب المنحسة إذا كان صاعدًا إليها فقط بعد النيّرين وكان}
\MsL{7}{AB03-F052v-L07}{النحس الآخر على الوتد دلّ على أحداث شتّى من الآفات والعلل التي ذكرنا مثل}
\MsL{8}{AB03-F052v-L08}{الآفات والعلل التي يدلّ عليها مواضع الأفق ومواضع البروج والكواكب المنحسة}
\MsL{9}{AB03-F052v-L09}{والمنحوسة والكواكب التي تشاركها في الشكل بل قد تقوى على ذلك أيضًا إذا كان}
\MsL{10}{AB03-F052v-L10}{صاعدًا قبل النيّرين وذلك أنّ أجزاء كلّ واحد من البروج التي تحيط بالجزئين}
\MsL{11}{AB03-F052v-L11}{\BF{B36v}التي\PN{198} تقبل الضرر من أجزاء الأفق تدلّ على الجزء والذي يحدث فيه ذلك الحادث}
\MsL{12}{AB03-F052v-L12}{من البدن. وهل ما يحدث في ذلك العضو آفة أو علّة أو كلاهما وأمّا طبائع}
\MsL{13}{AB03-F052v-L13}{الكواكب فإنّها تحدث أنواع الأعراض وأسبابها وذلك أنّ لزحل من الأجزاء الرئيسيّة}
\MsL{14}{AB03-F052v-L14}{من أجزاء الإنسان السمع الأيمن والطحال والمثانة والبلغم والمشتري اللمس والرئة}
\MsL{15}{AB03-F052v-L15}{والأضلاع والشرّيّانات\PN{199} والمني وللمرّيخ السمع الأيسر والكلى والعروق}
\MsL{16}{AB03-F052v-L16}{والمذاكرة وفي أكثر النسخ بدّل قوله المذاكير الأعضاء الصغار وللشمس البصر}
\MsL{17}{AB03-F052v-L17}{والدماغ والقلب والعصب وجميع أعضاء الجانب الأيمن. وللزهرة الشمّ}
\MsL{18}{AB03-F052v-L18}{والكبد والشحم ولعطارد الفكر والنظر\PN{200} واللسان والمرارة والمقعدة وللقمر والمذاقة}
\MsL{19}{AB03-F052v-L19}{والازدراد والمعدة والبطن والرحم وجميع أعضاء الجانب الأيسر وحدوث الآفة}
\MsL{20}{AB03-F052v-L20}{بالجملة يكون في أكثر الأمر إذا كانت الكواكب المنحسة التي هي السبب في ذلك}
\MsL{21}{AB03-F052v-L21}{شرقيّة وأمّا حدوث العلل فإنّه يكون إذا كانت هذه الكواكب غربيّة}
\begin{PNotes}
\PNote{197}{AB03-n197}{right margin}
\PNote{198}{AB03-n198}{B: \RL{\ArN الذي}}
\PNote{199}{AB03-n199}{B: \RL{\ArN شرّيّات}}
\PNote{200}{AB03-n200}{B:\RL{\ArN النطر والفكر }; T: \RL{\ArN النطق}}
\end{PNotes}
\Folio{053r}{53r}
\MsL{1}{AB03-F053r-L01}{وذلك لأنّ فيما بين كلّ واحد من هذين الأمرين فرق لأنّ الآفة إنّما تعرض}
\MsL{2}{AB03-F053r-L02}{مرّة واحدة ولا يمتد وجعها. وأمّا الأمراض\PN{201} فإنّها تعرض إمّا دائمًا أو\PN{202} بأدوار}
\MsL{3}{AB03-F053r-L03}{وأمّا الأشياء [[التي تدرك بها الأعراض]] التي تدرك بها الأعراض الجزئيّة فقد}
\MsL{4}{AB03-F053r-L04}{وجد لها لمّا تفقّدت ورصدت رصدًا خاصّيًّا وأحوالًا وأشكالًا تدلّ على الآفة}
\MsL{5}{AB03-F053r-L05}{وعلى العلّة وعلم ذلك من الأعراض التي تتبع وتلحق في أكثر الأمر بحسب وضع الكواكب}
\MsL{6}{AB03-F053r-L06}{المتشابهة الحال وذلك أنّ ذهاب البصر يحدث بإحدى العينين إذا كان}
\MsL{7}{AB03-F053r-L07}{القمر وحده على الأوتاد التي تقدّم ذكرها وكان الوقت مع ذلك وقت الاجتماع}
\MsL{8}{AB03-F053r-L08}{أو الاستقبال. وإذا كان على شكل آخر من أشكاله بقياسه إلى الشمس وكان}
\MsL{9}{AB03-F053r-L09}{متّصلًا بشيء من الكواكب اللطخيّة الشبيهة بالسحاب التي في \BF{B37r}فلك البروج كمثل}
\MsL{10}{AB03-F053r-L10}{الكوكب السحابيّ الذي في صدر السرطان وكمثل الثريّا من الثور وبرج الرأس}
\MsL{11}{AB03-F053r-L11}{وشولة العقرب وبما حول الصغيرة\PN{203} من أجزاء الأسد التي تسمّى الذوائب أو بالقليسين\PN{204}}
\MsL{12}{AB03-F053r-L12}{من الدلو يعني مصب الماء وفرع الدلو وهو المقبص. وإذا كان القمر على وتد من}
\MsL{13}{AB03-F053r-L13}{الأوتاد وكان غربيًّا وكان أو\PN{205} المرّيخ وزحل شرقيّين وكانا صائرين إليه من}
\MsL{14}{AB03-F053r-L14}{بعد وإذا كانت الشمس على بعض الأوتاد وكان هذان الكوكبان يطلعان}
\MsL{15}{AB03-F053r-L15}{ويصعدان قبلهما فإنّهما إن كانا مشاركين في الشكل للنيّرين جميعًا والنيّران}
\MsL{16}{AB03-F053r-L16}{جميعًا في برج واحد أو متقابلان على القطر كما قلنا وكانا في مشاركتهما\PN{206} منسوبين}
\MsL{17}{AB03-F053r-L17}{إلى الغدوات وفي مشاراكتهما للقمر منسوبين إلى العشيّات فإنّ الحادث}
\MsL{18}{AB03-F053r-L18}{يكون في العينين جميعًا. أمّا المرّيخ فإنّه يحدث عنه ما ذكرنا من ضربة تقع أو صدمة}
\MsL{19}{AB03-F053r-L19}{أو بالحديد أو بإحراق فإن شارك عطارد في الشكل فإنّ ذلك يحدث بصراع أو}
\MsL{20}{AB03-F053r-L20}{برياضة أو بما يفعله الشرار من الناس. وأمّا زحل فإنّ ذلك يحدث عنه}
\MsL{21}{AB03-F053r-L21}{بنزول الماء في العينين أو ببرد أو بزرقة أو ما أشبه ذلك وأيضًا فإنّ الزهرة}
\begin{PNotes}
\PNote{201}{AB03-n201}{T: \RL{\ArN الأعراض}}
\PNote{202}{AB03-n202}{B: \RL{\ArN وأمّا}}
\PNote{203}{AB03-n203}{T: \RL{\ArN الضفيرة}}
\PNote{204}{AB03-n204}{T: \RL{\ArN وبالقلبيس}}
\PNote{205}{AB03-n205}{B: om.}
\PNote{206}{AB03-n206}{B: add. \RL{\ArN الشكل؟}}
\end{PNotes}
\Folio{053v}{53v}
\MsL{1}{AB03-F053v-L01}{إن كانت على الأوتاد التي تقدّم ذكرها وبخاصّة إذا كانت على الوتد الغارب}
\MsL{2}{AB03-F053v-L02}{ثمّ كانت مع زحل أو مشاركته في الشكل أيضًا أو بدّلًا موضعيّهما في الطالع والغارب}
\MsL{3}{AB03-F053v-L03}{وكان المرّيخ عاليًا عليهما أو مقابلًا لهما على القطر فإنّه يعرض من ذلك للرجال}
\MsL{4}{AB03-F053v-L04}{أن\PN{207} يولد لهم وللنساء الإسقاط أو الولادة التي تكون قبل وقتها وقد يعرض لهنّ أيضًا}
\MsL{5}{AB03-F053v-L05}{تقطيع الأجنّة ويعرض ذلك خاصّة في السرطان والعذراء والجدي والقمر يتّصل}
\MsL{6}{AB03-F053v-L06}{بالمرّيخ من المشرق كأنّه يزيد في الطالع فإن شاركت الزهرة أيضًا عطارد دلّ}
\MsL{7}{AB03-F053v-L07}{على هذه الحال مع زحل وكان المرّيخ \BF{B37v}عاليًا عليهما أو كان مقابلًا لهما على القطر فإنّ}
\MsL{8}{AB03-F053v-L08}{المولود يكون خصيًّا أو خنثى أو مسدود المنفذ والمخرج. فإذا كانت هذه \{الأمر\}\PN{208}}
\MsL{9}{AB03-F053v-L09}{كذلك وكانت الشمس أيضًا مشاركة في الشكل وكان النيّران والزهرة مذكّرة وكان}
\MsL{10}{AB03-F053v-L10}{القمر غربيًّا وكانت النحوس في الدرج التي تتلوا فإنّ الذكر يكون مجنونًا أو تعرّض}
\MsL{11}{AB03-F053v-L11}{له آفة في مذاكيره وبخاصّة إذا كان في الحمل والأسد والعقرب والجدي والدلو.}
\MsL{12}{AB03-F053v-L12}{فأمّا الإناث فلا يولد لهنّ ويكنّ عواقم وربّما لم يسلم هؤلاء من آفة [[من آفة]] تعرض}
\MsL{13}{AB03-F053v-L13}{لهنّ في البصر والمولودون الذين تعتقل السنتهم ويكونون لثغًا أو غيرهم الذين}
\MsL{14}{AB03-F053v-L14}{يكون لهم زحل وعطارد مع الشمس على الأوتاد الذي ذكرنا وبخاصّة إذا كان}
\MsL{15}{AB03-F053v-L15}{عطارد أيضًا غربيًّا وكانا جميعًا مشاركين للقمر في الشكل فإذا اتّفق مع}
\MsL{16}{AB03-F053v-L16}{هذين الكوكبين المرّيخ كان من شأنه إن يحلّ\PN{209} اعتقال اللسان في أكثر الأمر منذ}
\MsL{17}{AB03-F053v-L17}{الوقت الذي يلقاه فيه القمر وأيضًا فإن النيّرين إن كانا جميعًا يصيّران إلى الكواكب}
\MsL{18}{AB03-F053v-L18}{المنحوسة والمنحسة على الأوتاد أو كانت المنحسة تقابل النيّرين على القطر وبخاصّة}
\MsL{19}{AB03-F053v-L19}{إن كان القمر في أحد العقدين أو في \{القسون\} وهو موضع الانعطاف}
\MsL{20}{AB03-F053v-L20}{من نهاية العرض فيما يرى لأن القسمون هو الشيء \{المنحرث\} أو في البروج}
\MsL{21}{AB03-F053v-L21}{المعلولة مثل الحمل والثور والسرطان والعقرب والجدي عرضت في البدن آفة من}
\begin{PNotes}
\PNote{207}{AB03-n207}{T: add. \RL{\ArN لا} (!)}
\PNote{208}{AB03-n208}{B: \RL{\ArN الأشياء}}
\PNote{209}{AB03-n209}{T: \RL{\ArN يجعل}}
\end{PNotes}
\Folio{054r}{54r}
\MsL{1}{AB03-F054r-L01}{من حدبة أو زمانة أو عرجة أو خلع أو استرخاء ومتى كانت المنحسة مع النيّرين}
\MsL{2}{AB03-F054r-L02}{عرضت هذه الآفات منذ وقت الولاد نفسه فإن كانت في وسط السماء أوكانت}
\MsL{3}{AB03-F054r-L03}{عالية على النيّرين وكان مسيرها في الطول معهما يعني عالية عليهما في العرض}
\MsL{4}{AB03-F054r-L04}{كانت هذه الآفات عن أعرض عظيمة ذوات خطر مثل السقوط من الأجراف}
\MsL{5}{AB03-F054r-L05}{ومن مواضع آخر ومن الوقوع في أيدي اللصوص أو من ذوات الأربع قوائم.}
\MsL{6}{AB03-F054r-L06}{فإن كان \BF{B38r}المرّيخ هو المستوليّ كانت هذه الآفات التي تعرض من حرق النار}
\MsL{7}{AB03-F054r-L07}{ومن الحرارة الرديّة\PN{210} أو من الوقوع في أيدي قطاع الطريق واللصوص. وإن كان}
\MsL{8}{AB03-F054r-L08}{المستوليّ وزحل كانت هذه الآفات من سقوط أو من غرق في الماء أو من تشنّج}
\MsL{9}{AB03-F054r-L09}{والآفات تعرض في أكثر الأمر إذا كان القمر على نقطتي الانقلابين والاعتدالين}
\MsL{10}{AB03-F054r-L10}{وإذا كان على نقطة الاعتدال الذي يبقى\PN{211} كانت الآفة التي تعرض إليهنّ خاصّة}
\MsL{11}{AB03-F054r-L11}{وإذا كان القمر على نقطة الانقلاب الصيفيّ كان ما يعرض \{القوابيّ\}. وإذا كان}
\MsL{12}{AB03-F054r-L12}{على نقطة الاعتدال الخريفيّ كان ما يعرض البرص. وإذا كان على نقطة الانقلاب}
\MsL{13}{AB03-F054r-L13}{الشتويّ كان ما يعرض في أكثر الآثار التي تسمّى العدس وما أشبهها. وأمّا}
\MsL{14}{AB03-F054r-L14}{الأعراض والعلل فمن شأنها أن تعرض إذا صارت النحوس على الحال التي ذكرنا}
\MsL{15}{AB03-F054r-L15}{في مشاركتها في الشكل للشمس أو للقمر إلّا أنّ ذلك على الخلاف أعني إذا شاركت}
\MsL{16}{AB03-F054r-L16}{النحوس القمر. فهي منسوبة إلى الغدوات والذي يحدث من الأمراض بالجملة هو ما أصف.}
\MsL{17}{AB03-F054r-L17}{أمّا زحل فإنّه يصيّر مزاج بطن المولود باردًا ويكون كثير البلغم صاحب موادّ تنصب}
\MsL{18}{AB03-F054r-L18}{إلى الأعضاء مهزولًا مسقامًا صاحب يرقان وقرحة في الأمعاء وسعال وقذف}
\MsL{19}{AB03-F054r-L19}{دم\PN{212} وقولنج وجذام. وأمّا النساء فيعرض لهنّ مع ذلك أوجاع الرحم. وأمّا المرّيخ}
\MsL{20}{AB03-F054r-L20}{فيصيّر المولود صاحب نفث الدم ووساوس سوداويّ وذات الرئة وجرب ويعرض}
\MsL{21}{AB03-F054r-L21}{له مع ذلك تأذّي بسبب بطّ أو كيّ يستعمل بسبب علل المواضع الخفيّة مثل}
\begin{PNotes}
\PNote{210}{AB03-n210}{T: \RL{\ArN الخراجات المؤذيّة}}
\PNote{211}{AB03-n211}{T: \RL{\ArN الربيعيّ}}
\PNote{212}{AB03-n212}{B: om.}
\end{PNotes}
\Folio{054v}{54v}
\MsL{1}{AB03-F054v-L01}{النواصير والبواسير وما ينبث وينتوا في البدن والقروح الحارّة}
\MsL{2}{AB03-F054v-L02}{الناريّة والقروح الساعيّة. وأمّا النساء فيعرض لهنّ مع ذلك إسقاط الجنين}
\MsL{3}{AB03-F054v-L03}{أو تقطعه أو تأكّله\PN{213} وقد يحدث للطبائع التي ذكرنا أنّها للكواكب المشتركة}
\MsL{4}{AB03-F054v-L04}{في الشكل وفي أعضاء البدن عللًا خاصّيّة والذي يعينها على تدبير الأشياء}
\MsL{5}{AB03-F054v-L05}{عطارد بخاصّة وذلك أنّه يعين زحل أويوافقه في التد بير وبخاصّة}
\MsL{6}{AB03-F054v-L06}{في أحداثه \BF{B38v}الانصباب المولد\PN{214} إلى الأعضاء والتأذّي بالرطوبات لا سيّما}
\MsL{7}{AB03-F054v-L07}{ما انحدر منها إلى الحنك والصدر والمعدة ويعيّن المرّيخ في أمر التجفيف}
\MsL{8}{AB03-F054v-L08}{وما يعرض منه مثل قشور القروح والخشكريشة والدبيلات والأورام}
\MsL{9}{AB03-F054v-L09}{التي تعرف بالحمرة والقوابي الخبيثة والماليخوليا\PN{215} والجنون والصرع وما أشبه}
\MsL{10}{AB03-F054v-L10}{ذلك وقد تكون أيضًا لعلل خواصّ بحسب اختلاف البروج التي تكون فيها}
\MsL{11}{AB03-F054v-L11}{الاشتراك في الشكل الذي تقدّم ذكره الذي يكون على الوتدين وذلك}
\MsL{12}{AB03-F054v-L12}{أنّ السرطان والجدي وبرج السمكتين وجميع البروج التي صورها شبيهة}
\MsL{13}{AB03-F054v-L13}{بصور الحيوان البريء والتي صورتها بصورة السمكة تحدث خاصّة العلل}
\MsL{14}{AB03-F054v-L14}{السعاعيّة والقوابي والنقرس والخنازير والنواصير والجذام أو ما أشبهها.}
\MsL{15}{AB03-F054v-L15}{وأمّا الرامي والتوأمان فيحدثان العلل التي يحدث فيها السقوط والصرع}
\MsL{16}{AB03-F054v-L16}{وإذا كانت الكواكب في الأجزاء الأواخر من البروج كان ما يحدث من العلل}
\MsL{17}{AB03-F054v-L17}{والآفات في أطراف البدن خاصّة ويعرض فيها ذلك بسبب فساد يصيبها}
\MsL{18}{AB03-F054v-L18}{وانصباب فضول يعرض منها الجذام ويعرض منها في أكثر الأمر النقرس في}
\MsL{19}{AB03-F054v-L19}{اليدين والرجلين وإذا كانت الأشياء على هذه الحال فإنّ الكواكب المسعدة إن}
\MsL{20}{AB03-F054v-L20}{لم يشاركها في الشكل شيء من الكواكب المنحسة التي هي سبب الفعل أو لم يشارك}
\MsL{21}{AB03-F054v-L21}{النيّرين إذا كانا على الأوتاد كانت الآفات والعلل العارضة مؤذيّة لا بدّ لها}
\begin{PNotes}
\PNote{213}{AB03-n213}{or: \RL{\ArN تآكل}}
\PNote{214}{AB03-n214}{B: \RL{\ArN الموادّ}}
\PNote{215}{AB03-n215}{B: \RL{\ArN الماولخوليا}; T: \RL{\ArN المالنخوليًّأ}}
\end{PNotes}
\Folio{055r}{55r}
\MsL{1}{AB03-F055r-L01}{وكذلك أيضًا متى شاركتها في الشكل وكانت المنحسة عالية على المسعدة}
\MsL{2}{AB03-F055r-L02}{أو قويّة فإن كانت المسعدة في أشكالها الشريفة\PN{216} وعلت على المنحسة التي هي}
\MsL{3}{AB03-F055r-L03}{سبب الفعل كانت الآفات حينئذ غير مستقبحة ولا ممّا فيه عار وكانت العلل}
\MsL{4}{AB03-F055r-L04}{عللًا خفّيّة سريعة السكون وربّما كانت سهلة التغيّر أيضًا. وذلك يكون}
\MsL{5}{AB03-F055r-L05}{إذا كانت المسعدة مشرقة وذلك أنّ من شأن المشتري أن يخفي الآفات}
\MsL{6}{AB03-F055r-L06}{وأنّ تسكّن العلل باستعانة ذوي العلل \BF{B39r}والآفات بالبأس\PN{217} عليها بسبب غنائهم}
\MsL{7}{AB03-F055r-L07}{أو بسبب مراتبهم فإذا كان مع عطارد كان تسكين للعلل\PN{218} باللأدويّة أو بعلاج من أطبّاء}
\MsL{8}{AB03-F055r-L08}{أخيار\PN{219}. وأمّا الزهرة فإنّها تحسّن الآفات بعض التحسين بأسباب من أسباب}
\MsL{9}{AB03-F055r-L09}{الله جلّ وعزّ أو من أسباب من وحي وتسكّن العلل بعلاجات تكون أسبابها}
\MsL{10}{AB03-F055r-L10}{من عند الله تعالى. فإن كان معها عطارد كان ذلك مع معونة وفائدة تنال}
\Chap{I.27}{<27> في حال النفس من المولود}
\MsL{11}{AB03-F055r-L11}{الذي به الآفة والعلل\PN{220} ممّا به من ذلك. \Hd{\LR{\LatinF \Sg{〈}27\Sg{〉}}}\Hd{ في حال النفس من المولود}}
\MsL{12}{AB03-F055r-L12}{أمّا أحوال النفس فإنّ ما كان منها في جزئها الناطق الفعليّ\PN{221} فإنّه يعلم أبدًا}
\MsL{13}{AB03-F055r-L13}{من الحال التي يرى عليها عطارد. وأمّا ما كان منها في جزئها الخلقي الذي}
\MsL{14}{AB03-F055r-L14}{ليس بناطق فإنّه يعلم من القمر ومن الكواكب التي تشاركه\PN{222} في الشكل باتّصالها}
\MsL{15}{AB03-F055r-L15}{وانصرافها عنه. وإذا كان الأمر في حركات النفس وأحوالها كثير الفنون فبالواجب}
\MsL{16}{AB03-F055r-L16}{ما صار فحصنا عنه لا يجري على الأمر المطلق البسيط ولا كيف ما اتّفق لكن يجري}
\MsL{17}{AB03-F055r-L17}{على تفقّد كثير وأرصاد مختلفة وذلك أنّ اختلاف البروج التي يكون فيها عطارد}
\MsL{18}{AB03-F055r-L18}{والقمر والكواكب المستوليّة عليها تعين على الحال في خواصّ النفس معونة}
\MsL{19}{AB03-F055r-L19}{كثيرة وكذلك أيضًا اشتراك أشكال الكواكب التي لها حصّة في الأمر الذي}
\MsL{20}{AB03-F055r-L20}{ذكرنا بقياسها إلى الشمس وإلى الأوتاد والحال أيضًا الخاصّيّة بطبيعة}
\MsL{21}{AB03-F055r-L21}{كلّ واحد من الكواكب تعين في حركات النفس. أمّا البروج بالجملة فإنّ المنقلبة}
\begin{PNotes}
\PNote{216}{AB03-n216}{T: \RL{\ArN الشرقيّة}}
\PNote{217}{AB03-n217}{T: \RL{\ArN بالناس}}
\PNote{218}{AB03-n218}{T: \RL{\ArN للهلل}}
\PNote{219}{AB03-n219}{T: \RL{\ArN الحبياد}}
\PNote{220}{AB03-n220}{T: \RL{\ArN العلّة}}
\PNote{221}{AB03-n221}{T: \RL{\ArN العقليّ}}
\PNote{222}{AB03-n222}{B: \RL{\ArN تشارك}}
\end{PNotes}
\Folio{055v}{55v}
\MsL{1}{AB03-F055v-L01}{منها تصيّر الأنفس محبّة لأمور الجموع والعامّة والمدن وتصيّرها محبّة للمرّيخ}
\MsL{2}{AB03-F055v-L02}{متشبّثة بأمور الله تعالى زكيّة محمودة الحركة بحّامة محسنة إلى غيرها حسنة}
\MsL{3}{AB03-F055v-L03}{الحدس والشبيه ذات علم بالقضاء \BF{B39v}من النجوم والعرافة وأمّا ذوات الجسدين}
\MsL{4}{AB03-F055v-L04}{من البروج فإنّها تصيّر الأنفس متفنّنة سهلة التغيّر يعسر الوقوف}
\MsL{5}{AB03-F055v-L05}{عليها طياشة سهلة التنقّل ذوات حالين عاشقة كثيرة الفنون محبّة للموسيقى}
\MsL{6}{AB03-F055v-L06}{متوانيّة فطنة ذوات ندامة. وأمّا البروج الثابتة فإنّها تصيّر الأنفس عادلة}
\MsL{7}{AB03-F055v-L07}{غير منخدعة ثابتة زكيّة رزينة فهمة صابرة محبّة للتعب قويّة ضابطة لذاتها}
\MsL{8}{AB03-F055v-L08}{حقوقة شديدة الماحلة محبّة للكرامة مشعبة عاصيّة غير مسعدة\PN{223}. وأمّا}
\MsL{9}{AB03-F055v-L09}{أحوال الكواكب وأشكالها فإنّ ما كان منها شرقيًّا طالعًا بخاصّيّة ما كان}
\MsL{10}{AB03-F055v-L10}{منها في مواجهة الخاصّيّة له فهو يصيّر النفس حرّة ساذجة منفردة بالرأي}
\MsL{11}{AB03-F055v-L11}{وتقال واقعة\PN{224} قويّة ذكية حادّة غير مؤذيّة عن شيء. وأمّا وقوف الكواكب المنسوبة}
\MsL{12}{AB03-F055v-L12}{إلى الغدوات وتوسّطها السماء فإنّه يصيّر النفس متفكّرة ثابتة ذاكرة رزينة}
\MsL{13}{AB03-F055v-L13}{فهيمة كثيرة الهمّة غير متغيّرة ولا منثنيّة\PN{225} ولا منخدعة مميّزة صاحبة عقل\PN{226}}
\MsL{14}{AB03-F055v-L14}{ذات فحص ومعرفة صحيحة. وأمّا طلوع الكواكب في أوّل الليل وغروبها فإنّه}
\MsL{15}{AB03-F055v-L15}{يصيّر النفس خفيفة طياشة ضعيفة لا تصبر على التعب قابلة للتأثير شقّيّة}
\MsL{16}{AB03-F055v-L16}{جبانة شكاكة تنقلب في الجرأة والجبن كالّة داهيّة عسيرة الحركة. وأمّا وقوف}
\MsL{17}{AB03-F055v-L17}{الكواكب المنسوبة إلى العشيّات وتوسّطها لمّا تحت الأرض من السماء وغروب}
\MsL{18}{AB03-F055v-L18}{الزهرة وعطارد المنسوبين إلى العشيّات إذا كان الوقت نهارًا وغروبهما}
\MsL{19}{AB03-F055v-L19}{المنسوب إلى الغدوات إذا كان الوقت ليلًا فإنّه يصيّر النفس زكيّة غافلة}
\MsL{20}{AB03-F055v-L20}{ليست بجيدّة\PN{227} الحفظ ولا كدودة ولا محبّة للتعب إلّا أنّها بحّاثة عن الأمور}
\MsL{21}{AB03-F055v-L21}{المستورة طالبة للأشياء الخفيّة مثل السحر والأسرار والعلوم\PN{228} العلويّة}
\begin{PNotes}
\PNote{223}{AB03-n223}{T: \RL{\ArN متنقّلة}}
\PNote{224}{AB03-n224}{B: \RL{\ArN ويقال وقاعه}}
\PNote{225}{AB03-n225}{T: \RL{\ArN متشبّثة}}
\PNote{226}{AB03-n226}{T: \RL{\ArN عمل}}
\PNote{227}{AB03-n227}{T: \RL{\ArN بحّبيدة}}
\PNote{228}{AB03-n228}{T: \RL{\ArN والأمور}}
\end{PNotes}
\Folio{056r}{56r}
\MsL{1}{AB03-F056r-L01}{وعلم الآلات والحيل وعلم العجائب والقضاء بالنجوم\PN{229} زاجرة للطير معبّرة للرؤيًا}
\MsL{2}{AB03-F056r-L02}{وما أشبه ذلك ومع ذلك أيضًا \BF{B40r}فإنّ الكواكب المتولّيّة \LR{\Sg{〉}}\LR{\Sg{〉}} لتدبير \LR{\Sg{〈}}\LR{\Sg{〈}}\PN{230} الأمور النفسانيّة إذا}
\MsL{3}{AB03-F056r-L03}{كانت في مواضعها وحيّزها الخاصّ لها والملائمين لها \LR{\Sg{〉}}\LR{\Sg{〉}} على ما بيّنًا \LR{\Sg{〈}}\LR{\Sg{〈}}\PN{231} فيما تقدّم صيّرت خواصّ}
\MsL{4}{AB03-F056r-L04}{النفس ظاهرة لا يمنعها من التمام مانع يعارضها منفردة كلّ واحدة بذاتها \{بحجة\}}
\MsL{5}{AB03-F056r-L05}{وبخاصّة إذا كانت كواكب أو احدة بأعيانها مستوليّة على الموضعين معًا. أعني}
\MsL{6}{AB03-F056r-L06}{أن تكون مشاركة في الشكل لموضع عطارد كيف ما كان ذلك أو تكون منصرفة}
\MsL{7}{AB03-F056r-L07}{عن القمر أو متّصلة به. فإذا لم يكن الأمر فيها على ما قلنا أو كانت في مواضع ليست}
\MsL{8}{AB03-F056r-L08}{خاصّة لها فإنّ طبائعها الخاصّيّة لمّا تصيّر أفعال النفس غير ظاهرة خفيّة غير}
\MsL{9}{AB03-F056r-L09}{تامّة ولا منجحة. وأمّا طبائع الكواكب التي تستولي أو تعلوا عليها فإنّها}
\MsL{10}{AB03-F056r-L10}{تصيّر أفعال النفس قويّة ضارة لمن تحت أيديهم كما أن الجوره والشرار الذين}
\MsL{11}{AB03-F056r-L11}{يحدّث ذلك فيهم من قبل مشاكلة الكواكب المنحسة إذا كانت هذه الكواكب}
\MsL{12}{AB03-F056r-L12}{هي المستوليّة صيّرت تركتهم\PN{232} إلى الأضرار بغيرهم حركة سهلة لا مانع لها ولا}
\MsL{13}{AB03-F056r-L13}{مشقّة فيها مشهورة. وإذا استولى على هذه الكواكب التي هي من حيّز المخالف لحيّزها}
\MsL{14}{AB03-F056r-L14}{صيّرتهم خاملي الذكر غير منجحين ملقين من العقوبات وكما أنّ الأخيار العدول}
\MsL{15}{AB03-F056r-L15}{الذين يعرض لهم ذلك بسبب مشاكلة الكواكب المسعدة إذا كانت}
\MsL{16}{AB03-F056r-L16}{غير مستعلٍ عليها فإنّهم يعرجون\PN{233} بالإحسان إلى غيرهم ويجدون ذلك من}
\MsL{17}{AB03-F056r-L17}{غير أن ينالهم الضرر من ذلك بل يكون عدلهم سبب منفعة لهم. وإذا كان استولى}
\MsL{18}{AB03-F056r-L18}{على هذه الكواكب ما يضادّها من الكواكب الحال على خلاف ذلك ما ذكرنا ومن}
\MsL{19}{AB03-F056r-L19}{أجل سكون هؤلاء ولينهم ومحبّتهم للناس ورحمتهم لهم يكونون ممّن يستخفّ به}
\MsL{20}{AB03-F056r-L20}{ويذمّ أو يظلمه كثير من الناس. فهذه هي الجهة الكلّيّة التي تفحص بها عن الجزء}
\MsL{21}{AB03-F056r-L21}{الخلقي. فأمّا الخواصّ والأحوال الجزئيّة التي تحدث عن طبيعة الكواكب}
\begin{PNotes}
\PNote{229}{AB03-n229}{B: add. \RL{\ArN وتكون متفلسفة}}
\PNote{230}{AB03-n230}{left margin}
\PNote{231}{AB03-n231}{left margin}
\PNote{232}{AB03-n232}{T: \RL{\ArN حركتهم}}
\PNote{233}{AB03-n233}{T: \RL{\ArN يفرحون}}
\end{PNotes}
\Folio{056v}{56v}
\MsL{1}{AB03-F056v-L01}{بحسب رئاستها فإنّا نذكرها ذكرًا \BF{B40v}مجملًا حتّى تنتهي إلى الامتزاج الذي}
\MsL{2}{AB03-F056v-L02}{يظهر على الأمر العامّيّ. أمّا إذا كان صاحب التدبير في أمر النفس زحل وحده}
\MsL{3}{AB03-F056v-L03}{وكان هو صاحب عطارد والقمر فإنّه إن كان في حاله في العالم وبقياسه إلى}
\MsL{4}{AB03-F056v-L04}{الأوتاد على أمر محمود فإنّه يصيّر المولود محبًّا للأديان\PN{234} قوى الرأي بعيد الغور\PN{235}}
\MsL{5}{AB03-F056v-L05}{منفردًا برأيه کدودًا متآمرًا متعدّيًا ذا يسار محبًّا للملوك غشومًا صاحب}
\MsL{6}{AB03-F056v-L06}{کنوز حسودًا وإذا\PN{236} كان على خلاف ذلك غير محمود الحال. فإنّه يصيّر المولود}
\MsL{7}{AB03-F056v-L07}{وسخًا حقيرًا صغير الهمّة غير مميّز منفردًا برأيه حسودًا فجبانًا معتزلًا عن الناس}
\MsL{8}{AB03-F056v-L08}{خبیث القول محبًّا للبوادي والنوح وقاحًا شقيّ البخت محبًّا للكدّ لا مودّة}
\MsL{9}{AB03-F056v-L09}{له غاشًّا لأوليائه لا سرور له مبغضًا للأديان\PN{237} . وإن كان مشاکلًا للمشتري}
\MsL{10}{AB03-F056v-L10}{فإنّ حاله إن كانت الحال المحمودة صيّر المولود خيرًا مكرمًا للمشائخ هاديًا ساكنًا}
\MsL{11}{AB03-F056v-L11}{مميّزًا محبًّا للملك كبير الهمّة مرزوقًا \{حسن\} الاختيار محبًّا لأوليائه قيّمًا حمولًا}
\MsL{12}{AB03-F056v-L12}{فیلسوفًا وإن كانت حاله على خلاف ذلك صيّر المولود غير معتاد للخير}
\MsL{13}{AB03-F056v-L13}{ذاهب العقل يعالج أمر الجنّ يأوي الهياكل مخبرًا بالغيب مبغضًا للأولاد}
\MsL{14}{AB03-F056v-L14}{لا صديق له يأوي إلى شراب غير مختلط بالناس وغير موثوق به جاهلًا رديئا}
\MsL{15}{AB03-F056v-L15}{آدم\PN{238} مكاشرًا ضعيفًا غير محبّ للكرامة متندّمًا كزًّا\PN{239} صعب اللقاء رديء الاختيار.}
\MsL{16}{AB03-F056v-L16}{وإن كان مشاکلًا للمرّيخ وهو على الحال الحمودة صغير المولود غير مميّز}
\MsL{17}{AB03-F056v-L17}{متعوبًا مذلًا \LR{\Sg{〉}}\LR{\Sg{〉}} متقلّبًا \LR{\Sg{〈}}\LR{\Sg{〈}}\PN{240} في الخيرة والجبن\PN{241} صعب العاملة لا يرحم متهاونًا خشنًا}
\MsL{18}{AB03-F056v-L18}{مقابلًا خطرًا بنفسه محبًّا للشعث غاشًّا مغتالًا رديء الحقد لا\PN{242} يؤثر في}
\MsL{19}{AB03-F056v-L19}{نفسه ممّا يعرضه منعمًا للناس منقلبًا عاصيًّا مبغضًا للرؤساء محبًّا للمراء}
\MsL{20}{AB03-F056v-L20}{والغلبة مفردًا\PN{243} رديء الأصل يعسر احتماله صعب المعاملة ثقيلًا فضًّا غليظًا}
\MsL{21}{AB03-F056v-L21}{مدّاحًا مسيئًا إلى الناس ظلومًا مهينًا مبغضًا للناس غير متغيّر ولا منتقل}
\begin{PNotes}
\PNote{234}{AB03-n234}{T: \RL{\ArN للأبدان }}
\PNote{235}{AB03-n235}{B: add. \{\RL{\ArN كثر}\}}
\PNote{236}{AB03-n236}{B: \RL{\ArN إن}}
\PNote{237}{AB03-n237}{T: \RL{\ArN للأبدان}}
\PNote{238}{AB03-n238}{T: \RL{\ArN أخضر}}
\PNote{239}{AB03-n239}{T: \RL{\ArN متبرّ ماكزًا }}
\PNote{240}{AB03-n240}{right margin}
\PNote{241}{AB03-n241}{T: \RL{\ArN الحجين}}
\PNote{242}{AB03-n242}{T: om.}
\PNote{243}{AB03-n243}{T: \RL{\ArN حقودًا}}
\end{PNotes}
\Folio{057r}{57r}
\MsL{1}{AB03-F057r-L01}{مدخلا لنفسه في أمور كثيرة إلّا أنّه سهل الرجوع \BF{B41r}صاحب}
\MsL{2}{AB03-F057r-L02}{عمل واجتهاد وهو بالجملة منجح وإن كان على خلاف تلك الحال}
\MsL{3}{AB03-F057r-L03}{صيّر المولود سلّا بًا قاطع طريق مرذولًا رديء الحال رديء}
\MsL{4}{AB03-F057r-L04}{الكسب لا يخاف الله ولا له مودّة شتّام مغتال خائن قاتل}
\MsL{5}{AB03-F057r-L05}{خبيث الفرج شديد قتّال للرجال ساحر يسلب الهياكل}
\MsL{6}{AB03-F057r-L06}{فاسق نبّاش وهو بالجملة رديء في جميع أحواله. وإن كان زحل}
\MsL{7}{AB03-F057r-L07}{مشاكل\PN{244} للزهرة فإنّه إن كان على الحال \LR{\Sg{〉}}\LR{\Sg{〉}} المحمودة \LR{\Sg{〈}}\LR{\Sg{〈}}\PN{245}}
\MsL{8}{AB03-F057r-L08}{صيّر المولود مبغضًا للنساء محبًّا للمشائخ والبراذين رديء}
\MsL{9}{AB03-F057r-L09}{اللقاء غير محبّ الذين محبًّا لمكتوماته كاهنًا متًّا لها محبًّا لأمور}
\MsL{10}{AB03-F057r-L10}{الله عزّ وجلّ متشبّثًا بها عفيفًا كثير الحياد مشاء محبًّا للحكمة}
\MsL{11}{AB03-F057r-L11}{ثقة أمينًا ضابطًا لنفسه مسخطًا غيورًا على النساء وإن}
\MsL{12}{AB03-F057r-L12}{كان على خلاف الحال المحمودة صيّر المولود شرها في}
\MsL{13}{AB03-F057r-L13}{الجماع نحسًا مرتكبًا للقبيح غير مميّز ولا طاهرًا في الجماع ولا}
\MsL{14}{AB03-F057r-L14}{تقيًّا مختدعًا للنساء وبخاصّة للقرابات منهنّ رکیکًا مذمومًا في}
\MsL{15}{AB03-F057r-L15}{جميع أحواله منتهكًا في الجماع مبغضًا للجميل متعوبًا رديء}
\MsL{16}{AB03-F057r-L16}{اللفظ متكبّرًا\PN{246} صاحب أمور \LR{\Sg{〉}}\LR{\Sg{〉}} خفيّة \LR{\Sg{〈}}\LR{\Sg{〈}}\PN{247} مرذولًا فاجرًا في أمر}
\MsL{17}{AB03-F057r-L17}{الجماع يركب القبيح ويرتكب منه الأمر الخارج عن المجرى}
\MsL{18}{AB03-F057r-L18}{الطبيعيّ محبًّا لفعل ذلك بالطاعن في السنّ وبالاترال \LR{\Sg{〉}}\LR{\LatinF sic}\LR{\Sg{〈}}\PN{248}}
\MsL{19}{AB03-F057r-L19}{من الناس على خلاف الشريعة وبالحيوان ولا يخاف الله}
\MsL{20}{AB03-F057r-L20}{عزّ وجلّ مستخفًّا بالأمور والهياكل مستهزئًا ساحرًا يدخل نفسه}
\MsL{21}{AB03-F057r-L21}{في كلّ شيء. فإن كان زحل مشاكلًا للكواكب عطارد فإنّه}
\begin{PNotes}
\PNote{244}{AB03-n244}{B: \RL{\ArN مشاكلًا}}
\PNote{245}{AB03-n245}{corrected interlinear}
\PNote{246}{AB03-n246}{T: \RL{\ArN سكيّرًا}}
\PNote{247}{AB03-n247}{interlinear}
\PNote{248}{AB03-n248}{T: \RL{\ArN بالأنذال}}
\end{PNotes}
\Folio{057v}{57v}
\MsL{1}{AB03-F057v-L01}{إن كان على الحال المحمودة صيّر المولود مبتّشًا بحّاثًا عن}
\MsL{2}{AB03-F057v-L02}{الأخبار طلّابًا للشرائع محبًّا للطبّ صاحب سرّ ينظر في الأمور}
\MsL{3}{AB03-F057v-L03}{الخفيّة المستورة صاحب عجائب مغالطًا ينظر ليومه خفيفًا صاحب}
\MsL{4}{AB03-F057v-L04}{تدبیر. فطنًا مستقصيًا مستيقظًا محبًّا للفهم والعمل منجحًا.}
\MsL{5}{AB03-F057v-L05}{وإن كان زحل على خلاف الحال المحمودة صيّر المولود حقودًا}
\MsL{6}{AB03-F057v-L06}{کدر النفس مكدورًا تعاويذ وخديعة خائنًا غير منجح. وأمّا المشتري}
\MsL{7}{AB03-F057v-L07}{فإنّه إن ولي التدبير لأمر النفس وحده فإنّه إن \BF{B41v}كان على}
\MsL{8}{AB03-F057v-L08}{الحال المحمودة صيّر أحوال النفس شبيهة بالأحوال التي ذكرنا}
\MsL{9}{AB03-F057v-L09}{إلّا أنّها تكون أضعف وأخمد عن غير تمييز مثل أنّه يكون بدّل}
\MsL{10}{AB03-F057v-L10}{كبير النفس مبذّرًا وبدّل النقيّ رديء المذهب وبدّل}
\MsL{11}{AB03-F057v-L11}{السخي جبانًا وبدّل العفيف معجبًا بنفسه وبدّل الحبّ للناس حسن}
\MsL{12}{AB03-F057v-L12}{الخلق وبدّل الجميل الأمر محبًّا للذّة وبدّل العظيم الهمّة تيّاها}
\MsL{13}{AB03-F057v-L13}{وبدّل الخير عن\PN{249} مميّز وما أشبه ذلك. فإن كان المشتري مشاکلًا}
\MsL{14}{AB03-F057v-L14}{للمرّيخ فإنّه إن كان في الحال المحمودة صيّر المولود خشنًا}
\MsL{15}{AB03-F057v-L15}{مخاصمًا صاحب حرب مدبّرًا صاحب حركة غير خاضع حادًّا}
\MsL{16}{AB03-F057v-L16}{صاحب عساكير ذااله\PN{250} صاحب عمل منكيًا منتقمًا محبًّا}
\MsL{17}{AB03-F057v-L17}{للغلبة مترأّسًا وقّاعًا على الأمور ودیعًا رجلًا باذلًا\PN{251} مظفّرًا}
\MsL{18}{AB03-F057v-L18}{كبير النفس محبًّا للكرامة غضوبًا مميّزًا صاحب أمر ونهي.}
\MsL{19}{AB03-F057v-L19}{فإن كان المشتري على خلاف ما ذكرنا كان المولود}
\MsL{20}{AB03-F057v-L20}{شتّامًا مخلّطًا فظًّا قليل الأعضاء شغابًا صاحب مراء}
\MsL{21}{AB03-F057v-L21}{منفردًا بجلده سعًا معجبًا عاصيًا\PN{252} سائلًا سريع الانتقال}
\begin{PNotes}
\PNote{249}{AB03-n249}{B: \RL{\ArN غير}}
\PNote{250}{AB03-n250}{T: \RL{\ArN دالّة}}
\PNote{251}{AB03-n251}{T: \RL{\ArN نازلًا}}
\PNote{252}{AB03-n252}{T: \RL{\ArN غاضبًا}}
\end{PNotes}
\Folio{058r}{58r}
\MsL{1}{AB03-F058r-L01}{خفيفًا صاحب ندامة لا يثبت على شيء واحد سخيفًا لا تمييز له}
\MsL{2}{AB03-F058r-L02}{ولا رأي ذاهب العقل باهيا ذامًا متلاقًا وهو بالجملة مختلف}
\MsL{3}{AB03-F058r-L03}{الأخلاق مضطرب الأحوال. فإن كان المشتري مشاكلًا للزهرة}
\MsL{4}{AB03-F058r-L04}{وهو على الحال المحمودة صيّر المولود تقیًّا متنعّمًا محبًّا للنظافة}
\MsL{5}{AB03-F058r-L05}{محبًّا للمنازعات\PN{253} محبًّا للنظر فيها محبًّا للغناء واللهو محبًّا للأكل}
\MsL{6}{AB03-F058r-L06}{حسن الأخلاق محسنًا رحيمًا سليم القلب محبًّا لله متعبّدًا}
\MsL{7}{AB03-F058r-L07}{محبًّا للجهاد عاقلًا وادًّا عفيفًا نقي النفس شكورًا مرزوقًأ محبًّا}
\MsL{8}{AB03-F058r-L08}{للنظر في الكتب مميّزًا يستعمل في الجماع الأمر المعتدل الحسن}
\MsL{9}{AB03-F058r-L09}{محبًّا للأقارب حسن الديانة محبًّا للعدل محبًّا للكرامة والنباهة}
\MsL{10}{AB03-F058r-L10}{وهو بالجملة خير فاضل. وإن كان المشتري على خلاف ما ذكرنا}
\MsL{11}{AB03-F058r-L11}{كان المولود مترفًا لذيذ العيش \BF{B42r}مؤنّث النفس رفّاضًا غضبه}
\MsL{12}{AB03-F058r-L12}{شبيه بغضب النساء منفاقًا مشغوفًا\PN{254} بالنساء صاحب عشق كثير}
\MsL{13}{AB03-F058r-L13}{الجماع منهمكًا فاجرًا محبًّا لزينة الشعر مقبحًا أرعن متلافًا محبًّا}
\MsL{14}{AB03-F058r-L14}{لعيب الناس سريع الكلام مؤنّثًا\PN{255} عقله كعقل المرأة\PN{256} الغالب عليه}
\MsL{15}{AB03-F058r-L15}{أمور الهياكل مقدّم في الأمور صاحب أسرار إلّا أنّه أمين لا شرّ}
\MsL{16}{AB03-F058r-L16}{عنده باشًّا قريبًا محمود المذاهب حرًّا\PN{257} في الأقوال\PN{258} التي تتصرّف}
\MsL{17}{AB03-F058r-L17}{فيها. فإذا شاكل المشتري\PN{259} عطارد وهو على الحال المحمودة}
\MsL{18}{AB03-F058r-L18}{صيّر المولود كثير النظر في الكتب محبًّا للقياس مهندسًا}
\MsL{19}{AB03-F058r-L19}{[[كثير النظر في الكتب محبًّا للقياس مهندسًا]] صاحب}
\MsL{20}{AB03-F058r-L20}{علوم تعليميّة شاعرًا خطيبًا ذكيًا عفيفًا محمود}
\MsL{21}{AB03-F058r-L21}{الرأي \LR{\LatinF ]}الرأي\LR{\LatinF [} حسن المشورة}
\begin{PNotes}
\PNote{253}{AB03-n253}{B: \RL{\ArN للمناغات}; T: \RL{\ArN للصناعات}}
\PNote{254}{AB03-n254}{T: \RL{\ArN معشوقًا}}
\PNote{255}{AB03-n255}{T: \RL{\ArN الألم متزینًا}}
\PNote{256}{AB03-n256}{T: \RL{\ArN النساء}}
\PNote{257}{AB03-n257}{BT: \RL{\ArN جدًّا}}
\PNote{258}{AB03-n258}{BT: \RL{\ArN الأحوال}}
\PNote{259}{AB03-n259}{T: \RL{\ArN هذا الكوكب}}
\end{PNotes}
\Folio{058v}{58v}
\MsL{1}{AB03-F058v-L01}{[سهو] \LR{\Sg{〉}}...\LR{\Sg{〈}}\PN{260}}
\begin{PNotes}
\PNote{260}{AB03-n260}{a blank page}
\end{PNotes}
\Folio{059r}{59r}
\MsL{1}{AB03-F059r-L01}{محسنًا قائمًا بالتدبير حسن الأخلاق صاحب جوائز محبًّا للجماعة}
\MsL{2}{AB03-F059r-L02}{جيّد الحدس سريع النجح صاحب سياسة حسن الديانة محبًّا}
\MsL{3}{AB03-F059r-L03}{لله طيّب الملك متودّدًا محبًّا للقرابة حسن الأدب محبًّا للحكمة}
\MsL{4}{AB03-F059r-L04}{ذا نباهة وقدر. فإذا كان المشتري على خلاف ما ذكرنا كان}
\MsL{5}{AB03-F059r-L05}{المولود خرقًا هكذا\PN{261} كثير الخطأ حقيرًا يعرض له الأشر شمس}
\MsL{6}{AB03-F059r-L06}{وهو شيء يعرض في القمر فيصيّر صاحبه إلى ذهاب العقل}
\MsL{7}{AB03-F059r-L07}{ويخبر عند ذلك بما يكون متشبّثًا بأمور الله عزّ وجلّ مهتمًا}
\MsL{8}{AB03-F059r-L08}{بأمور النفس يظنّ بنفسه أنه حكيم وهو عديم العقل معجبًا}
\MsL{9}{AB03-F059r-L09}{محتالًا مشعبذًا مضطرب الحركة كثير الحديث يكون صاحب}
\MsL{10}{AB03-F059r-L10}{تعليم كثير الشهوة. وأما إذا كان صاحب تدبير أمر النفس}
\MsL{11}{AB03-F059r-L11}{المرّيخ وحده كان على الحال المحمودة كان المولود قويًّا}
\MsL{12}{AB03-F059r-L12}{\{رئيسًا\} غضوبًا محبًّا للسلاح متفنّنًا بحدًّا صاحب وقائع}
\MsL{13}{AB03-F059r-L13}{مخاطرًا بنفسه غير خاضع متّكلًا على جلده حديدًا}
\MsL{14}{AB03-F059r-L14}{مقدامًا مستخفًّا بالأمور متغلّبًا\PN{262} ضابطًا متسخّطًا صاحب}
\MsL{15}{AB03-F059r-L15}{سياسة. فإذا \BF{B42v}كان على خلاف ما ذكرنا كان المولود فطنًا شتّامًا}
\MsL{16}{AB03-F059r-L16}{محبًّا لسفك الدماء محبًّا للشغب منفاقًا صيّاحًا محتالًا\PN{263}}
\MsL{17}{AB03-F059r-L17}{سخيفًا بتكيرًا سلّابًا لا رحمة له رديء الأفعال مضطربًا}
\MsL{18}{AB03-F059r-L18}{مجنونًا مبغضًا للقرابة لا يعرف الله عزّ وجلّ. فإن كان}
\MsL{19}{AB03-F059r-L19}{المرّيخ مشاکلًا للزهرة وهو على الحال المحمودة كان ثابتًا}
\MsL{20}{AB03-F059r-L20}{حسن المذهب محبًّا لأصحابه لذيذ العيش مسرورًا}
\MsL{21}{AB03-F059r-L21}{مزّاحًا ساذجًا حسن التأليف والشكل محبًّا للرفض صاحب}
\begin{PNotes}
\PNote{261}{AB03-n261}{B: \RL{\ArN هذاء}}
\PNote{262}{AB03-n262}{T: \RL{\ArN منقلبًا }}
\PNote{263}{AB03-n263}{B:\RL{\ArN محتاطًا} ; T: \RL{\ArN مختلطًا}}
\end{PNotes}
\Folio{059v}{59v}
\MsL{1}{AB03-F059v-L01}{عشق متنعّمًا صاحب دعوات صاحب رجلة سريع الميل}
\MsL{2}{AB03-F059v-L02}{إلى ارتكاب المحارم والآثام التي تكون من الجماع إلّا أنّه يكون}
\MsL{3}{AB03-F059v-L03}{منجحًا متحرّزًا غافلًا\PN{264} شديد التوبيخ صاحب تمييز وهو أيضًا}
\MsL{4}{AB03-F059v-L04}{محبّ للأحداث الذكورة والإناث منفاق سريع الغضب}
\MsL{5}{AB03-F059v-L05}{غیور وإن كان المرّيخ على خلاف ما ذكرنا صيّر المولود طباخًا\PN{265}}
\MsL{6}{AB03-F059v-L06}{كثير الجماع \{منهمكيًّا\}\PN{266} فيه مختلف الأحوال مستهزئًا بالناس}
\MsL{7}{AB03-F059v-L07}{فاجرًا كذّابًا غاشًّا ينخدع القريب والبعيد سريع إلى}
\MsL{8}{AB03-F059v-L08}{الشهوات ملولًا لها أيضًا مفسدًا للنساء المتزوّجات}
\MsL{9}{AB03-F059v-L09}{والأبكار صاحب حبّ حادًّا مضطربًا مغتالًا يخبث في}
\MsL{10}{AB03-F059v-L10}{إيمانه إلّا أنّه سريع الفضيحة فاسد العقل وربّما}
\MsL{11}{AB03-F059v-L11}{كان محبًّا للزينة متهوّرًا يرتكب منه القبيح نجسًا. فإذا}
\MsL{12}{AB03-F059v-L12}{شاكل المرّيخ عطارد وهو على الحال الحمودة كان المولود}
\MsL{13}{AB03-F059v-L13}{قائدًا \LR{\Sg{〉}}\LR{\Sg{〉}} صاحب \LR{\Sg{〈}}\LR{\Sg{〈}}\PN{267} دعاء \{ظابطًا\}\PN{268} سريع الحركة عزيزًا متنعّمًا محبًّا}
\MsL{14}{AB03-F059v-L14}{حكيمًا كدودًا صاحب أدب بريء اللسان خبیثًا غاشًّا لا}
\MsL{15}{AB03-F059v-L15}{ثبات له صاحب حیل رديء الأعمال سريع الفهم خدّاعًا}
\MsL{16}{AB03-F059v-L16}{مؤاتيًّا مغتالًا رديء المذهب كثير التفتيش محبًّا للشرّ}
\MsL{17}{AB03-F059v-L17}{إلّا أنّه منجح حسن المودّة\PN{269} لأشباهه حسن القبول لهم}
\MsL{18}{AB03-F059v-L18}{وهو بالجملة ضارّ لأعدائه محسن إلى أصدقائه وإذا كان}
\MsL{19}{AB03-F059v-L19}{\BF{B43r}المرّيخ على خلاف الحال المحمودة كان المولود منفاقًا عاصيًا}
\MsL{20}{AB03-F059v-L20}{صاحب ندامة متهوّنًا مضطرب الحركة كذّابًا لها لا}
\MsL{21}{AB03-F059v-L21}{يعرف الله عزّ وجلّ حانثًا في إيمانه خبيثًا صاحب شغب}
\begin{PNotes}
\PNote{264}{AB03-n264}{T: \RL{\ArN عاقل} \RL{\ArN متحرز}}
\PNote{265}{AB03-n265}{T: \RL{\ArN طقاحًا}}
\PNote{266}{AB03-n266}{i.e. \RL{\ArN منهمكًا}}
\PNote{267}{AB03-n267}{right margin}
\PNote{268}{AB03-n268}{B: \RL{\ArN ضابطًا} \RL{\ArN دعاء}; T: \RL{\ArN ضابطًا} \RL{\ArN دهاء}}
\PNote{269}{AB03-n269}{T: \RL{\ArN الولّدة }}
\end{PNotes}
\Folio{060r}{60r}
\MsL{1}{AB03-F060r-L01}{وإحراق مشهورًا بالشرّ شتّامًا قاطع طريق نقابًا متدنّسًا}
\MsL{2}{AB03-F060r-L02}{بالقتل صاحب تعاويذ ومکر ساحرًا خدّاعًا مشعبذًا قتّالًا}
\MsL{3}{AB03-F060r-L03}{للرجال. وإذا كان تدبير النفس للزهرة وحدها وحالها}
\MsL{4}{AB03-F060r-L04}{الحال الحمودة كان المولود ساكنًا خيّرًا متنعّمًا طاهرًا صاحب}
\MsL{5}{AB03-F060r-L05}{فكر محبًّا للرفض شديد الغيرة مبغضًا للشرّ محبًّا للصناعات}
\MsL{6}{AB03-F060r-L06}{محبًّا لله عزّ وجلّ جدًّا حسن الشكل حسن الأحوال جيّد}
\MsL{7}{AB03-F060r-L07}{الأحلام والرؤيا متودّدًا مجسّمًا\PN{270} رحيمًا حسن التصرف}
\MsL{8}{AB03-F060r-L08}{منجحًا وهو بالجملة مائل إلى أمور الجماع. فإن كانت الزهرة}
\MsL{9}{AB03-F060r-L09}{على خلاف ما ذكرنا كان المولود متوانيًّا صاحب عشق مؤنّثًا}
\MsL{10}{AB03-F060r-L10}{أحواله شبيهة بأحوال النساء لا جرأة له ولا تمييز خامرًا\PN{271}}
\MsL{11}{AB03-F060r-L11}{مذمومًا خامل الذكر بعيدًا\PN{272} كثيرًا. وإن شاكلت الزهرة عطارد}
\MsL{12}{AB03-F060r-L12}{وهي على الحال المحمودة كان المولود محبًّا للصناعات والحكمة}
\MsL{13}{AB03-F060r-L13}{ذا معرفة ذكيّة شاعرًا محبًّا للموسيقى محبًّا للجميل حسن الأخلاق}
\MsL{14}{AB03-F060r-L14}{متنعّمًا مترفًا مسرورًا محبًّا لأصدقائه حسن الدین محبًّا للحيل}
\MsL{15}{AB03-F060r-L15}{دهنًا جيّد الحدس مستقيم الطريقة محبًّا للتعليم يتعلّم من}
\MsL{16}{AB03-F060r-L16}{تلقاء نفسه يقتدي بأهل الفضل ويتشبّه بالخيار حسن}
\MsL{17}{AB03-F060r-L17}{اللفظ باشًّا في كلامه محبوبًا معتدل الأخلاق فاضلًا}
\MsL{18}{AB03-F060r-L18}{محبًّا للجهاد مميّزًا كبير الهمّة ممتنعًا عن مجامعة النساء}
\MsL{19}{AB03-F060r-L19}{مائلًا إلى الغلمان غيورًا وإذا كانت الزهرة على خلاف}
\MsL{20}{AB03-F060r-L20}{ذلك كان المولود خبیثًا كثير الجهل فاحش الكلام}
\MsL{21}{AB03-F060r-L21}{ذا وجهين ولسانين رديء الدين خدّاعًا مشغبًا كذّابًا}
\begin{PNotes}
\PNote{270}{AB03-n270}{T: \RL{\ArN محسنًا}}
\PNote{271}{AB03-n271}{T: \RL{\ArN جاحدًا}}
\PNote{272}{AB03-n272}{T: \RL{\ArN تغيّر}}
\end{PNotes}
\Folio{060v}{60v}
\MsL{1}{AB03-F060v-L01}{\BF{B43v}سعًا\PN{273} صاحب خبث في اليمين بعيد الغور في الشرّ غاشًّا لا}
\MsL{2}{AB03-F060v-L02}{ودّ له ولا ولي خدّاعًا للنساء مفسدًا للغلمان صاحب زينة مقبح}
\MsL{3}{AB03-F060v-L03}{صاحب فرّ\PN{274} وهجاء وتشنيع كثير يروم فعل جميع الأشياء}
\MsL{4}{AB03-F060v-L04}{وربّما كان يرائي بفعل إلّا بهذه الأسباب ليفسد بها وربّما}
\MsL{5}{AB03-F060v-L05}{جعلها بالحقيقة وهو صاحب أحوال قبيحة ويذم منه أحوال}
\MsL{6}{AB03-F060v-L06}{مختلفة. وإذا كان صاحب تدبير أمر النفس عطارد وحده}
\MsL{7}{AB03-F060v-L07}{وهو محمود الحال كان المولود قيّمًا فطنًا ذكيًا\PN{275} كثير الروايّة}
\MsL{8}{AB03-F060v-L08}{الأحاديث صاحب فوائد وتجارب ومنطق متكلّمًا في}
\MsL{9}{AB03-F060v-L09}{الأمور الطبيعيّة صاحب نظر ذكيًا متأدّیًا محسنًا صاحب}
\MsL{10}{AB03-F060v-L10}{فطنّ جيّد الحدس ذا علوم تعليميّة كتومًا للسرّ منجحًا. وإذا}
\MsL{11}{AB03-F060v-L11}{كان على خلاف الحال صيّر المولود مکّارًا سخيفًا مغالطًا}
\MsL{12}{AB03-F060v-L12}{راكبًا لرأيه خفيفًا سريع الانتقال سريع الرجوع أحمق}
\MsL{13}{AB03-F060v-L13}{جاهلًا كثير الخطأ كذّابًا مضطربًا في الأحوال لا ثبات}
\MsL{14}{AB03-F060v-L14}{له ولا أمانة عاصيًا ظلومًا وهو بالجملة كثير الغلط}
\MsL{15}{AB03-F060v-L15}{بين الخطأ مضطرب الأحوال. وإذا كان الأمر كذلك}
\MsL{16}{AB03-F060v-L16}{فينبغي أن يعلم أنّه يعيّن على الأشياء التي ذكرنا حال القمر}
\MsL{17}{AB03-F060v-L17}{أيضًا بعض المعونة وذلك أنّه إذا كان في مواضع الانعطاف}
\MsL{18}{AB03-F060v-L18}{من النهاية الجنوبيّة والشماليّة كان معيّنًا على تغيّر أخلاق}
\MsL{19}{AB03-F060v-L19}{النفس وكثرة حيلها وسرعة انتقالها. وإذا كان على}
\MsL{20}{AB03-F060v-L20}{نطاق البروج كان معيّنًا على حدّة الأخلاق ونفوذها}
\MsL{21}{AB03-F060v-L21}{في الأفعال وسرعة حركتها وأيضًا فإنّ النفس\PN{276} إذا كان}
\begin{PNotes}
\PNote{273}{AB03-n273}{T: \RL{\ArN سفها}}
\PNote{274}{AB03-n274}{T: \RL{\ArN ذمّ}}
\PNote{275}{AB03-n275}{T: \RL{\ArN ذهنا}}
\PNote{276}{AB03-n276}{T: \RL{\ArN القمر}}
\end{PNotes}
\Folio{061r}{61r}
\MsL{1}{AB03-F061r-L01}{طالعًا يعني إذا ظهر من الشعاع أوكان متزيّدًا في الضوء كان}
\MsL{2}{AB03-F061r-L02}{معيّنًا على أن يكون أكثر تزيّدًا وأشهر واضح واكد. وإذا كان}
\MsL{3}{AB03-F061r-L03}{ضوءه ناقصًا أو مستترًا تحت الشعاع كان معيّنًا على جمودها}
\MsL{4}{AB03-F061r-L04}{وکلالها وما يعرض فيها من التراخي والتخلّف والخمول}
\MsL{5}{AB03-F061r-L05}{وقد يعين على ما قلنا بعض المعونة الشمس أيضًا \BF{B44r}فإنّها}
\MsL{6}{AB03-F061r-L06}{إذا كانت كلت الكواكب المتولّيّ التدبير أمر النفس}
\MsL{7}{AB03-F061r-L07}{وذلك أنّه إذا شاكلته وخالط\PN{277} الحال المحمودة كانت}
\MsL{8}{AB03-F061r-L08}{أخلاق النفس أعدل وأبعد من الظلم وأكرم واحد\PN{278}}
\MsL{9}{AB03-F061r-L09}{وأحسن ديانة. وإذا كان الأمر على خلاف ذلك كانت}
\MsL{10}{AB03-F061r-L10}{أخلاق النفس أحسن وأنكد وأقسى وأكثر تفرد}
\MsL{11}{AB03-F061r-L11}{بالرأي وأذعر وأضعف مذهبًا وهي بالجملة عسرة العلاج\PN{279}}
\MsL{12}{AB03-F061r-L12}{وعلى هذا المذهب يكون تمييز أمر النفس بحسب امتزاج الكواكب}
\Chap{I.28}{<28> في الآفات النفسانيّة}
\MsL{13}{AB03-F061r-L13}{المتولّيّة والناظرة إليها. \Hd{\LR{\LatinF \Sg{〈}28\Sg{〉}}}\Hd{ في الآفات النفسانيّة} ينبغي لنا أن نتفقّد}
\MsL{14}{AB03-F061r-L14}{بالجملة في طلب علم الآفات النفسانيّة حال عطارد والقمر أحدهما}
\MsL{15}{AB03-F061r-L15}{عند الآخر وعند الأوتاد وعند الكواكب التي من شأنها أن تنحس. فإنّ}
\MsL{16}{AB03-F061r-L16}{عطارد والقمر إن كانا غير مرتبطين أحدهما بالآخر وكانوا يعلموا\PN{280}}
\MsL{17}{AB03-F061r-L17}{عليهما في الأفق الشرقيّ وكان يحيط بهما أو يقابلهما في الطول}
\MsL{18}{AB03-F061r-L18}{شيء من الكواكب التي حالها حال ضارة غير ملائمة فإنّه يعرض في}
\MsL{19}{AB03-F061r-L19}{الحال الخاصّيّة للنفس آفات \{فتون\} كثيرة وتميّز ذلك وتحصيله\PN{281} يكون}
\MsL{20}{AB03-F061r-L20}{من الأحوال الخاصّيّة التي ذكرناها من أحوال الكواكب التي تشاكل}
\MsL{21}{AB03-F061r-L21}{الواضع وقد ميّزنا أيضًا كثيرًا من الآفات النفسانيّة الضعيفة}
\begin{PNotes}
\PNote{277}{AB03-n277}{T: \RL{\ArN حاله}}
\PNote{278}{AB03-n278}{T: \RL{\ArN أجلّ}}
\PNote{279}{AB03-n279}{T: \RL{\ArN الصلاح}}
\PNote{280}{AB03-n280}{B: \RL{\ArN كان نعلوا}; T: \RL{\ArN كان يعلوا}}
\PNote{281}{AB03-n281}{T: \RL{\ArN تفصيله}}
\end{PNotes}
\Folio{061v}{61v}
\MsL{1}{AB03-F061v-L01}{بعض التمييز فيما تقدّم من قولنا في خواصّ النفس وقد يمكن أن يعرف}
\MsL{2}{AB03-F061v-L02}{قوّتها أو تزيّدها من إفراط أفعال الكواكب المنحسة وذلك أنّه}
\MsL{3}{AB03-F061v-L03}{لعّلى الإنسان أن \{سمّى\} الأخلاق التي لا تضبط آفات أن يكون}
\MsL{4}{AB03-F061v-L04}{مصيبًا إن كانت دون الحال المتوسّطة وإن كانت متجاوزة لها. وأمّا}
\MsL{5}{AB03-F061v-L05}{ما كان إفراطه منها شديدًا وكان بمنزلة مرض من الأمراض وكان}
\MsL{6}{AB03-F061v-L06}{حدوثه خارجًا \BF{B44v}عن الطبيعة كلّ الخروج وكان ذلك في الجزء العقليّ}
\MsL{7}{AB03-F061v-L07}{من النفس وفي الجزء القابل للتأثير فإن تعرّفه يكون على المثال}
\MsL{8}{AB03-F061v-L08}{الذي أصف. وذلك أن أصحاب الصرع في أكثر الأمر هم الذين}
\MsL{9}{AB03-F061v-L09}{هم لا يرتبط[[ون]] في مواليدهم القمر وعطارد كما قلنا أحدهما بالآخر}
\MsL{10}{AB03-F061v-L10}{ولا بالطالع ويكون مع ذلك في الوتد بالنهار زحل وبالليل \LR{\Sg{〉}}\LR{\Sg{〉}} بالمرّيخ \LR{\Sg{〈}}\LR{\Sg{〈}}\PN{282}}
\MsL{11}{AB03-F061v-L11}{هو صاحب هذه الحال. وأمّا المجانين فهم الذين يكون الأمر}
\MsL{12}{AB03-F061v-L12}{فيهم مع ذلك على ضدّ ما قلنا إن يكون صاحب الحال بالليل وبالنهار}
\MsL{13}{AB03-F061v-L13}{المرّيخ وخاصّة إذا كان في السرطان أو العذراء أو السمكتين وأمّا الذين}
\MsL{14}{AB03-F061v-L14}{يسمّون المضرّ وبين من الجزء والذين يقال لهم أوعر وقًا قالس\PN{283}}
\MsL{15}{AB03-F061v-L15}{وتفسير ذلك الذي يقلب على رؤوسهم الرطوبة فإنّهم يكونون}
\MsL{16}{AB03-F061v-L16}{إذا كان حال النحوس على ما وصفنا وكانت مستوليّة على القمر}
\MsL{17}{AB03-F061v-L17}{وهو الطالع من شعاع\PN{284} الشمس وكان المستوليّ عليه منها في وقت}
\MsL{18}{AB03-F061v-L18}{الاجتماع زحل وفي وقت الاستقبال المرّيخ وبخاصّة في الرأس\PN{285}}
\MsL{19}{AB03-F061v-L19}{والسمكتين فإذا كان الغالب على تلك الحال النحسان وحدهما}
\MsL{20}{AB03-F061v-L20}{على الجهة التي قدّمنا ذكرها صيّر الأمراض التي ذكرنا أنّها}
\MsL{21}{AB03-F061v-L21}{تنزل بجزء النفس الفعليّ\PN{286} لا برء لها إلّا أنّها لا تكون لا \LR{\Sg{〉}}\LR{\LatinF sic}\LR{\Sg{〈}} مستشنعة}
\begin{PNotes}
\PNote{282}{AB03-n282}{left margin; here and in B: \RL{\ArN القمر} corrected to \RL{\ArN المرّيخ}}
\PNote{283}{AB03-n283}{T: \RL{\ArN مافاليسا}}
\PNote{284}{AB03-n284}{B: \RL{\ArN الشعاع}}
\PNote{285}{AB03-n285}{T: \RL{\ArN الراعي}}
\PNote{286}{AB03-n286}{T: \RL{\ArN العقليّ}}
\end{PNotes}
\Folio{062r}{62r}
\MsL{1}{AB03-F062r-L01}{ولا مستهونة\PN{287} فإن شاكلها السعدان أعني المشتري والزهرة وكان}
\MsL{2}{AB03-F062r-L02}{النحسان على الوتد الغربيّ والسعدان على الوتد الشرقيّ صيّرت العلل}
\MsL{3}{AB03-F062r-L03}{القابلة للبرء إلّا أنّها تكون ميتة مستشنعة. فإن كان المشاكل لها}
\MsL{4}{AB03-F062r-L04}{المشتري كان البرء بعلاج طبّيّ إمّا بالتدبير وإمّا بالأدويّة وإن كان}
\MsL{5}{AB03-F062r-L05}{الذي يشاكلها الزهرة كان برء العلّة بوحي ومعونة من عند}
\MsL{6}{AB03-F062r-L06}{الله عزّ وجلّ. وإن كان\PN{288} النحوس على الأوتاد الشرقيّة وكانت}
\MsL{7}{AB03-F062r-L07}{النحوس\PN{289} غاربيّة صيّرت الأمراض ممّا لا بدّ لها وكانت مستشنعة}
\MsL{8}{AB03-F062r-L08}{مشهورة جدًّا. أمّا علل الصرع منها فيصيّرها دائمة غير مفارقة}
\MsL{9}{AB03-F062r-L09}{مع صيّاح وخطر فيحدث منه \BF{B45r}الموت. وأمّا علل الجنون والخروج}
\MsL{10}{AB03-F062r-L10}{عن العقل فيصيّر صاحبها لا يضبط نفسه ويعطي\PN{290} أقاربه \{ويلغوا\}}
\MsL{11}{AB03-F062r-L11}{ويفحش في القول وما أشبه ذلك. وأمّا كان منها من الأمراض}
\MsL{12}{AB03-F062r-L12}{التي تسمّى ضربة الجنّ وغلبة الرطوبة المؤذيّة على الرأس}
\MsL{13}{AB03-F062r-L13}{فيصيّرها من أسوسامس\PN{291} ويصيّرها مشهورة مع وثب على}
\MsL{14}{AB03-F062r-L14}{الناس وفدي وضرب وما أشبه ذلك من الفضائح وقد}
\MsL{15}{AB03-F062r-L15}{تعين بعض المواضع التي تكون فيها تلك الحال معونة}
\MsL{16}{AB03-F062r-L16}{خاصّيّة من ذلك أنّ موضع الشمس والمرّيخ تعين على حدوث}
\MsL{17}{AB03-F062r-L17}{الجنون خاصّة ومواضع المشتري وعطارد تعين على}
\MsL{18}{AB03-F062r-L18}{علّة الصرع ومواضع الزهرة تعين على الوحي وعلى الأخبار}
\MsL{19}{AB03-F062r-L19}{بالسرائر ومواضع زحل والقمر تعين على التأذّي بالرطوبات}
\MsL{20}{AB03-F062r-L20}{وعلى ضربة الجنّ وأكثر أصناف الأمراض الحادثة في جنّ}
\MsL{21}{AB03-F062r-L21}{النفس الفاعل في هذه الأنواع التي ذكرنا وعن مثل هذه}
\MsL{22}{AB03-F062r-L22}{الأحوال يكون حدوثها بحسب جملة طبائعها وأمّا}
\begin{PNotes}
\PNote{287}{AB03-n287}{T: \RL{\ArN مشهورة}}
\PNote{288}{AB03-n288}{B: \RL{\ArN كانت}}
\PNote{289}{AB03-n289}{BT: \RL{\ArN السعود}}
\PNote{290}{AB03-n290}{BT: \RL{\ArN يعصى}}
\PNote{291}{AB03-n291}{T: \RL{\ArN أنثوسيماسموس}}
\end{PNotes}
\Folio{062v}{62v}
\MsL{1}{AB03-F062v-L01}{الاختلاف الذي يحدث في الجنّ القابل للتدبير خاصّة فإنّه يظهر}
\MsL{2}{AB03-F062v-L02}{كثيرًا في زيادات الأمور الطبيعيّة التي تجنّس الذكر والأنثى}
\MsL{3}{AB03-F062v-L03}{ونقصانها ويدرك الفحص عنها بنحو شبيه بالنحو بعد أن نقيم}
\MsL{4}{AB03-F062v-L04}{الشمس مع القمر هاههنا مقام عطارد مع القمر هناك ويستعمل}
\MsL{5}{AB03-F062v-L05}{مشاكلة المرّيخ والزهرة لها. ونقول من بعد ذلك إنّه إن كان النيّران}
\MsL{6}{AB03-F062v-L06}{أو أحدهما في بروج مذكّرة فإن الرجال يفرطون في الأمر الطبيعيّ}
\MsL{7}{AB03-F062v-L07}{لهم وأمّا النساء فإنّهنّ يفرطن في الأمر الخارج عن الأمرالطبيعيّ}
\MsL{8}{AB03-F062v-L08}{لهنّ ويتجاوزن الطبيعة إلى رحلة النفس وقوّتها. وإن كان}
\MsL{9}{AB03-F062v-L09}{المرّيخ والزهرة أيضًا أو أحدهما أيّهما كان أو كلاهما جميعًا مذكّرين}
\MsL{10}{AB03-F062v-L10}{كان الرجال منهمكين في الجماع الذي يجري على الأمر الطبيعيّ \{خبرة\}}
\MsL{11}{AB03-F062v-L11}{مستهزئين عراعًا؟ أبدًا إلى الأمور القبيحة المخالفة للستة من}
\MsL{12}{AB03-F062v-L12}{أمور الجماع. وأمّا النساء فإنّهنّ يبصرن مشهيات\PN{292} في الجماع}
\MsL{13}{AB03-F062v-L13}{بالأمر الخارج عن الطبيعة \BF{B45v}طماخات\PN{293} كاللواتي يسمّين}
\MsL{14}{AB03-F062v-L14}{ساحقات وذلك أنّهنّ يفعلن مع النساء أفعال الرجال}
\MsL{15}{AB03-F062v-L15}{فإن كانت الزهرة وحدها هي المذكّرة كان ما يفعلن من ذلك}
\MsL{16}{AB03-F062v-L16}{سرًّا لا علانيّة. فإن كان المرّيخ أيضًا مذكّرًا كان ذلك علانيّة حتّى}
\MsL{17}{AB03-F062v-L17}{أنّهنّ في بعض الأوقات يظهرن أنّ النساء اللواتي معهنّ نساء}
\MsL{18}{AB03-F062v-L18}{لهنّ. فإن كان الأمر في الكواكب على خلاف ذلك أعني أن يكون}
\MsL{19}{AB03-F062v-L19}{النيّران في الحال التي ذكرنا في بروج مؤنّثة وحدها فإنّ}
\MsL{20}{AB03-F062v-L20}{النساء [[في]] يفرطن في الأمر الطبيعيّ أن تقبيح النفسيّ وتأنيثها}
\MsL{21}{AB03-F062v-L21}{فإن كانت الزهرة أيضًا مؤنّثة كان النساء أيضًا منهمكات في}
\MsL{22}{AB03-F062v-L22}{أمر الجماع فهو أحد مشتهرات مائلات إلى ما يفعل بهنّ من ذلك}
\begin{PNotes}
\PNote{292}{AB03-n292}{T: \RL{\ArN مستهزآت من}}
\PNote{293}{AB03-n293}{T: \RL{\ArN طقاحات}}
\end{PNotes}
\Folio{063r}{63r}
\MsL{1}{AB03-F063r-L01}{على الأمر الطبيعيّ في كلّ وقت ومع من نهايتها حتّى أنّهنّ لا يمتنعن}
\MsL{2}{AB03-F063r-L02}{في شيء من أمور الجماع فيه قبيحًا كان ذلك أو مخالفا للسنة وكان}
\MsL{3}{AB03-F063r-L03}{الرجال مخنّثين [[أو]] مسترخيين يميلون إلى الجماع الخارج}
\MsL{4}{AB03-F063r-L04}{عن الطبيعة لا يمتنعون أن يركب منهم القبيح إلّا أنّ ذلك}
\MsL{5}{AB03-F063r-L05}{يكون في سرّ \{وعياء\}. فإن كان المرّيخ أيضًا مؤنّثًا كان فعل كلّ واحد}
\MsL{6}{AB03-F063r-L06}{من هذين الصنفين هذه الأفعال أعني الزناء والنجاسة}
\MsL{7}{AB03-F063r-L07}{وارتكاب الأمور المذمومة وإتيان جميع الأشياء القبيحة}
\MsL{8}{AB03-F063r-L08}{علانيّة من غير استحياء حتّى أنّ الأمر يبلغ بهم في الشهوة}
\MsL{9}{AB03-F063r-L09}{إلى أن يشتهوا ويقرفوا وأحوال المرّيخ والزهرة الشرقيّة}
\MsL{10}{AB03-F063r-L10}{المنسوبة إلى الغدوات معيّنة على الرحلة والشهرة. وأمّا}
\MsL{11}{AB03-F063r-L11}{أحوالهما القريبة\PN{294} المنسوبة إلى العشیّات فمعيّنة على التأنيث}
\MsL{12}{AB03-F063r-L12}{والخمول. وكذلك أيضًا متى كان معها زحل كان معيّنًا على}
\MsL{13}{AB03-F063r-L13}{الرجاسة والنجاسة وعلى شدّة العار وذلك أنّ طبيعته معيّن}
\MsL{14}{AB03-F063r-L14}{على كلّ واحد من هذه الأشياء بأن كان معها المشتري كان}
\MsL{15}{AB03-F063r-L15}{معيّنًا على جمال الأمر والصيانة والاستحياء. فإن كان معها}
\MsL{16}{AB03-F063r-L16}{عطارد كان معيّنًا على شهرة الأمر وعلى سرعة حركة / /B46rالأحداث}
\Chap{I.29}{<29> في بخت المولود في المال}
\MsL{17}{AB03-F063r-L17}{وكثرة أنواعها وسوها\PN{295}. \Hd{\LR{\LatinF \Sg{〈}29\Sg{〉}}}\Hd{ في بخت المولود في المال}\PN{296} أمّا المال في أمور القنيّة كيف تكون فينبغي أن يؤخذ عليها}
\MsL{18}{AB03-F063r-L18}{من سهم البخت وحده الذي نعلمه أبدًا بأن ننظر إلى ما بين}
\MsL{19}{AB03-F063r-L19}{الشمس والقمر ونلقيه من الطالع كما بيّنا في باب العمر}
\MsL{20}{AB03-F063r-L20}{بالليل والنهار جميعًا لتلك الأسباب التي تقدّم ذكرها}
\begin{PNotes}
\PNote{294}{AB03-n294}{BT: \RL{\ArN الغربيّة}}
\PNote{295}{AB03-n295}{T: \RL{\ArN نموها}}
\PNote{296}{AB03-n296}{T: \RL{\ArN الملك} (wherever it occurs in this chapter)}
\end{PNotes}
\Folio{063v}{63v}
\MsL{1}{AB03-F063v-L01}{ووجه النظر في ذلك هو أن يأخذ الكواكب المتولّيّة لتدبير البرج}
\MsL{2}{AB03-F063v-L02}{الذي فيه هذا القسم وتنظر كيف قوّتهما وملازمتهما\PN{297} على الجهة التي حررنا\PN{298}}
\MsL{3}{AB03-F063v-L03}{فيما تقدّم وننظر أيضًا في الكواكب التي تشاركها في الشكل والتي}
\MsL{4}{AB03-F063v-L04}{تعلوا عليها إمّا من التي هي من ذلك الحيّز بعينه وإمّا من التي هي}
\MsL{5}{AB03-F063v-L05}{من الحيّز المخالف وذلك أنّه متى كانت المدبّرة للسهم قويّة كان}
\MsL{6}{AB03-F063v-L06}{المولود كثير المال وبخاصّة إن شهد لها النيّران شهادة ملائمة}
\MsL{7}{AB03-F063v-L07}{إلّا أنّ زحل يدلّ على أنّ ثروة المولود تكون من قبل البناء أو من}
\MsL{8}{AB03-F063v-L08}{العلامة\PN{299} أو من الملاحة وأمّا المشتري فمن الأمانة أو من القهرمة}
\MsL{9}{AB03-F063v-L09}{أو من مذاهب إلّا يمنة والعدالة. وأمّا المرّيخ فمن القيادة والسياسة}
\MsL{10}{AB03-F063v-L10}{وأمّا الزهرة فمن عطايا الأصدقاء أو من قبل النساء. وأمّا عطارد}
\MsL{11}{AB03-F063v-L11}{فيدلّ على أنّ الثروة تكون من العلم أو من التجارة. فإنّ شاكل}
\MsL{12}{AB03-F063v-L12}{زحل خاصّة سهم البخت ثمّ شاركه المشتري في الشكل دلّ على}
\MsL{13}{AB03-F063v-L13}{مواريث يرثها وبخاصّة متى كان ذلك في الأوتاد العلويّة}
\MsL{14}{AB03-F063v-L14}{وكان المشتري في برج ذي جسدين أو كان أيضًا متّصلًا بالقمر}
\MsL{15}{AB03-F063v-L15}{وذلك أنّه حينئذ يدلّ على أنّه سى\PN{300} ويرث موارث قوم غرباء}
\MsL{16}{AB03-F063v-L16}{والكواكب التي من حيّز الكواكب \BF{B46v}المدبّرة إن شهدت لها كان ما}
\MsL{17}{AB03-F063v-L17}{يملكه باقيًا محفوظًا. فإن غلبت عليه الكواكب التي من الحيّز المخالف}
\MsL{18}{AB03-F063v-L18}{على المواضع الرئيسيّة إذ كانت تابعة لها صائرة إليها دلّت على}
\MsL{19}{AB03-F063v-L19}{تلف المال والوقت الكلّيّ الذي فيه يكون ذلك يعلم من مجاز\PN{301}}
\MsL{20}{AB03-F063v-L20}{الكواكب التي هي سبب الفعل للأوتاد وللمواضع التي تصعدان}
\Chap{I.30}{<30> في بخت المولود في الحظوة}
\MsL{21}{AB03-F063v-L21}{الكواكب\LR{\LatinF .} \Hd{\LR{\LatinF \Sg{〈}30\Sg{〉}}}\Hd{ في بخت المولود في الحظوة} وأمّا أمور الحظوة}
\begin{PNotes}
\PNote{297}{AB03-n297}{T: \RL{\ArN ملائمتها}}
\PNote{298}{AB03-n298}{T: \RL{\ArN حدّدنا}}
\PNote{299}{AB03-n299}{T: \RL{\ArN الفلاحة}}
\PNote{300}{AB03-n300}{T: \RL{\ArN يتبنّى}}
\PNote{301}{AB03-n301}{T: \RL{\ArN محاذاة}}
\end{PNotes}
\Folio{064r}{64r}
\MsL{1}{AB03-F064r-L01}{والسعادة التي تليق بها فينبغي أن ننظر فيها من حال النيّرين ومن}
\MsL{2}{AB03-F064r-L02}{حال الكواكب المحتفّة بهما وننظر مع ذلك أيضًا في ملائماتها وذلك}
\MsL{3}{AB03-F064r-L03}{أنّ النيّرين إذا كانا جميعًا في بروج مذكّرة ثمّ كانا جميعًا على الأوتاد أو}
\MsL{4}{AB03-F064r-L04}{واحد منهما بخاصّة صاحب الحيّز منهما وكانت الكواكب المتحيّرة}
\MsL{5}{AB03-F064r-L05}{محتفّة بهما وكانت المحتفّة بالشمس منسوبة إلى الغدوات والمحتفّة}
\MsL{6}{AB03-F064r-L06}{بالقمر منسوبة إلى العشيّات وكان المولود ملكًا. فإن كنت الكواكب}
\MsL{7}{AB03-F064r-L07}{المختلفة\PN{302} بهما أيضًا على الأوتاد وكانت مشاركة في الشكل}
\MsL{8}{AB03-F064r-L08}{للوتد الذي فوق الأرض كان المولود عظيم الشأن قويًّا ملكًا على}
\MsL{9}{AB03-F064r-L09}{الدنيا وقد تزيد في سعادته متى كانت الكواكب المحتفّة متيامنه}
\MsL{10}{AB03-F064r-L10}{ونظنّ أنّه يوجد بالمتيامنه الشماليّة في العرض وكانت}
\MsL{11}{AB03-F064r-L11}{مشاركة في الشكل في الأوتاد التي فوق الأرض فإن كانت}
\MsL{12}{AB03-F064r-L12}{حال سائر الكواكب هذه الحال وكانت الشمس وحدها في}
\MsL{13}{AB03-F064r-L13}{برج مذكّر والقمر في برج مؤنّث وكان أحدهما على الوتد}
\MsL{14}{AB03-F064r-L14}{كان المولود صاحب سياسة مسلّطًا على أن يقتل وأن \{يسحن \PN{303}\}}
\MsL{15}{AB03-F064r-L15}{وإن لم تكن الكواكب \BF{B47r}محتفّة بهما على الأوتاد لكان المولود}
\MsL{16}{AB03-F064r-L16}{عظيم الشأن فقط وكانت مطنونة كحظوة من يملك على ناحية}
\MsL{17}{AB03-F064r-L17}{من النواحيّ أو كحظوة القهارمة أو كحظوة صاحب الجيش لا}
\MsL{18}{AB03-F064r-L18}{كخطوة القوّادين\PN{304} المدبّرين الرئاسة. فإن\PN{305} النيّران على الأوتاد}
\MsL{19}{AB03-F064r-L19}{وكان أكثر الكواكب المحتفّة بهما على الأوتاد أو كانت مشاركة}
\MsL{20}{AB03-F064r-L20}{للأوتاد في الشكل لم تكن حظوته نبيلة مشهورة بل يكون مقدّمًا}
\MsL{21}{AB03-F064r-L21}{في أمر تدبير البدن\PN{306} متوسّطًا في أمر تدبير المعاش فإن لم تكن}
\begin{PNotes}
\PNote{302}{AB03-n302}{T: \RL{\ArN المحتفّة}}
\PNote{303}{AB03-n303}{T: \RL{\ArN يستحيّ}}
\PNote{304}{AB03-n304}{BT: \RL{\ArN القوّاد}}
\PNote{305}{AB03-n305}{BT: add. \RL{\ArN لم يكون (!)}}
\PNote{306}{AB03-n306}{T: \RL{\ArN المدن}}
\end{PNotes}
\Folio{064v}{64v}
\MsL{1}{AB03-F064v-L01}{الكواكب المحتفّة بالنيّرين ملائمة للأوتاد كان حافظًا\PN{307} في الأفعال}
\MsL{2}{AB03-F064v-L02}{غير منجح وإذا لم تكن النيّران لا على الأوتاد ولا في بروج مذكّرة}
\MsL{3}{AB03-F064v-L03}{ولم تكن محتفّة بهما كان في غاية الذلّة وشقاء البخت بجملة}
\MsL{4}{AB03-F064v-L04}{المسلك الذي يسلك في الفحص عن الأمور هو على ما وصفت}
\MsL{5}{AB03-F064v-L05}{في زيادة الخطوط\PN{308} ونقصانها. فأمّا الأحوال التي هي فيما بين}
\MsL{6}{AB03-F064v-L06}{الزيادة النقصان التي هي كثيرة جدًّا فينبغي أن ننظر فيها}
\MsL{7}{AB03-F064v-L07}{من الأشياء التي توجد\PN{309} في ذلك النوع من التغيّر الجزئيّ الذي}
\MsL{8}{AB03-F064v-L08}{للنيّرين والكواكب التي تحفّهما التي تدبّر المحتفّة بهما}
\MsL{9}{AB03-F064v-L09}{وذلك أنّ المدبّرة إذا كانت هي صاحبة الحيّز أو هي السعود}
\MsL{10}{AB03-F064v-L10}{كانت الحظوة رأس\PN{310} وأثبت. وإذا كانت هي صاحبة الحيّز}
\MsL{11}{AB03-F064v-L11}{المخالف أو هي النحوس كانت الحظوظ\PN{311} مع حضوع وسهولة}
\MsL{12}{AB03-F064v-L12}{الزوال. فأمّا أمر الخطوة التي سيحدّث وقوعها فينبغي أن}
\MsL{13}{AB03-F064v-L13}{ننظر فيها من خاصّة حال الكواكب المحتفّة بالنيّرين}
\MsL{14}{AB03-F064v-L14}{وذلك أنّه\PN{312} كان المتولّيّ لهذا الأمر زحل كانت قوّة الحظوة}
\MsL{15}{AB03-F064v-L15}{في كثرة الملك وجميع الأموال. فإن كان المشتري أو الزهرة}
\MsL{16}{AB03-F064v-L16}{كانت قوّتها في سرور وعطايا وكرامات وكبر نفس فإن}
\MsL{17}{AB03-F064v-L17}{كان المرّيخ كانت في قيادة العساكر وفي خوف وفي غلبة}
\MsL{18}{AB03-F064v-L18}{ممّن هو تحته فإن كان عطارد كانت من فهم وأدب وعناية\LR{\LatinF .}}
\Chap{I.31}{<31> في حال صناعة المولود وعمله}
\MsL{19}{AB03-F064v-L19}{\BF{B47v}\Hd{\LR{\LatinF \Sg{〈}31\Sg{〉}}}\Hd{ في حال صناعة المولود وعمله} وأمّا حال صناعة}
\MsL{20}{AB03-F064v-L20}{المولود فإنّه يعلم من جهتين وهي من الشمس ومن البروج}
\MsL{21}{AB03-F064v-L21}{التي في وسط السماء وذلك أنّا نحتاج أن ننظر إلى الكواكب\PN{313}}
\begin{PNotes}
\PNote{307}{AB03-n307}{T: \RL{\ArN خاملًا}}
\PNote{308}{AB03-n308}{T: \RL{\ArN الحظوة}}
\PNote{309}{AB03-n309}{T: \RL{\ArN تؤخذ}}
\PNote{310}{AB03-n310}{BT: \RL{\ArN أرؤس}}
\PNote{311}{AB03-n311}{B: \RL{\ArN الحظوة}}
\PNote{312}{AB03-n312}{B: add. \RL{\ArN متى}}
\PNote{313}{AB03-n313}{B: \RL{\ArN الكو[[ا]]كب}}
\end{PNotes}
\Folio{065r}{65r}
\MsL{1}{AB03-F065r-L01}{الذي قد ظهر بالغدوات من تحت الشمس قريبًا وإلى الذي}
\MsL{2}{AB03-F065r-L02}{يكون في وسط السماء وبخاصّة متى اتّصل بالقمر. وإن وجدنا}
\MsL{3}{AB03-F065r-L03}{هاتين الحالتين لكوكب \LR{\Sg{〉}}\LR{\Sg{〉}} فينبغي لنا أن نستعمل هذا الكوكب وحده وكذلك أيضًا إذا كانت إحدى هاتين الحالتين لكوكب \LR{\Sg{〈}}\LR{\Sg{〈}} وإن\PN{314} لم تكن الأخرى لكوكب من}
\MsL{4}{AB03-F065r-L04}{الكواكب فإنّا نستعمل الكوكب الذي في إحدى الحالتين}
\MsL{5}{AB03-F065r-L05}{فقط وإن كان الذي قد ظهر بالغدوات من تحت الشمس قريبًا}
\MsL{6}{AB03-F065r-L06}{كوكب من الكواكب وكان الذي في وسط السماء يشاكل القمر كوكب}
\MsL{7}{AB03-F065r-L07}{غيره استعملناهما جميعًا وقدّمنا أوّلًا أن جحهما في كثرة عدد خطوط}
\MsL{8}{AB03-F065r-L08}{في التدبير على الجهة التي قدّمنا ذكرها وإن لم نجد كوكبًا قد ظهر}
\MsL{9}{AB03-F065r-L09}{من تحت الشعاع قريبًا بالغدوات ولا في وسط السماء فينبغي أن}
\MsL{10}{AB03-F065r-L10}{نأخذ صاحب وسط السماء إلّا أنّ ذلك في الأعمال الزمانيّة}
\MsL{11}{AB03-F065r-L11}{بأنّ من كانت الحال فيه هذه يكون في أكثر الأمر بطّالًا فهذا هو}
\MsL{12}{AB03-F065r-L12}{\{المسلك\} في تمييز أمر الكوًاكب التي تدبّر أمر الصناعة والعمل إليه.}
\MsL{13}{AB03-F065r-L13}{فأمّا حال الصناعة والعمل فإنّا نميّز من خاصّة حال ثلاثة كواكب}
\MsL{14}{AB03-F065r-L14}{أعني المرّيخ والزهرة وعطارد ومن حال البروج التي تتّفق منها}
\MsL{15}{AB03-F065r-L15}{وذلك أنّها إذا كان الفاعل للصناعة عطارد مثلًا كان المولود}
\MsL{16}{AB03-F065r-L16}{کاتبًا صاحب معاملات حاسبًا معلّمًا تاجرًا صيرفيًّا عرّافاً}
\MsL{17}{AB03-F065r-L17}{صاحب قضاء بالنجوم وصاحب ذبائح وهو بالجملة ممّن يعالج}
\MsL{18}{AB03-F065r-L18}{الكتاب والتفسير فالأحد فإن شهد له زحل كان مدبّر الأمور}
\MsL{19}{AB03-F065r-L19}{ومعبّرًا للأحلام \BF{B48r}أو مواظبًا على الهياكل بأسباب العرافة}
\MsL{20}{AB03-F065r-L20}{والابنو سماء سيموس\PN{315} وإن شهد له المشتري كان المولود}
\MsL{21}{AB03-F065r-L21}{واضع سرّ صاحب بلاغة وحكمة متغلّبًا\PN{316} مع قوم أشراف.}
\begin{PNotes}
\PNote{314}{AB03-n314}{B: om.}
\PNote{315}{AB03-n315}{T: \RL{\ArN أنثوسيما سموس}}
\PNote{316}{AB03-n316}{B: \RL{\ArN متقلّبًا}}
\end{PNotes}
\Folio{065v}{65v}
\MsL{1}{AB03-F065v-L01}{فإن كان صاحب الفعل الزهرة كان ممّن يعالج روائح الزهرة أو}
\MsL{2}{AB03-F065v-L02}{الطيب أو الشراب\PN{317} أو أصباغ أو آفاويّة أو أشياء مزيّنة مثل}
\MsL{3}{AB03-F065v-L03}{عطارد أو صاحب آكلة أو صاحب مواخير أو تاجر الشراب}
\MsL{4}{AB03-F065v-L04}{أو صيدلاني أو حائك أو بائع أفاويه أو مصوّرًا\PN{318} أو بزّازًا فإن شهد}
\MsL{5}{AB03-F065v-L05}{لها زحل كانت تجارته في الأشياء التي تستعمل \{للنتعم\} والزينة خدّاعًا}
\MsL{6}{AB03-F065v-L06}{ساحرًا وما أشبه ذلك. فإن شهد لهما المشتري كان المولود مصارعًا}
\MsL{7}{AB03-F065v-L07}{مكلّلًا مؤهلًا للكرامات وكان سبب رفعته من قبل النساء. فأمّا}
\MsL{8}{AB03-F065v-L08}{المرّيخ فإنّه إن شارك الشمس في الشكل كان ممّن صناعته}
\MsL{9}{AB03-F065v-L09}{من الصناعات التي تعالج بالحديد مثل تستعمل بالنار مثل}
\MsL{10}{AB03-F065v-L10}{الطبّاخ و\{سماك\}\PN{319} ووقّاد وصاحب نحاس أو صانع في المعادن}
\MsL{11}{AB03-F065v-L11}{وإن كان غير مشارك للشمس كانت صناعته من الصناعات}
\MsL{12}{AB03-F065v-L12}{التي تعالج بالحديد مثل تجارة السفن وغيرها من التجارة والفلاحة}
\MsL{13}{AB03-F065v-L13}{وقطع الحجارة ونحتها وقطع الحيوص؟\PN{320} وقصّ\PN{321} الخشب}
\MsL{14}{AB03-F065v-L14}{وما أشبه ذلك فإن شهد له زحل كان ملاحًا متاحًا \{كادحًا\} طباخًا}
\MsL{15}{AB03-F065v-L15}{مصوّرًا قيّمًا في الحمّام فإن شهد له المشتري كان جنديًّا شرطيًّا عشّارًا}
\MsL{16}{AB03-F065v-L16}{صاحب فندق \{مغريًا\} بناء ذباحًا وأيضًا فإنّه إن تولّى أمر الصناعة}
\MsL{17}{AB03-F065v-L17}{كوكبان ثمّ كان عطارد والزهرة كانت الصناعته الموسيقى وآلات}
\MsL{18}{AB03-F065v-L18}{الغناء واللهو والشعر وتأليف اللحون وبخاصّة متى بدّلًا موضعهما.}
\MsL{19}{AB03-F065v-L19}{\{فإن\}\PN{322} إن كان الأمر كذلك كان المولود ملهيًّا صاحب خیال أو نحاس}
\MsL{20}{AB03-F065v-L20}{الرقيق أو صانعًا لآلات اللهو فتّال الأوتار قاصًا\PN{323} مصوّرًا}
\MsL{21}{AB03-F065v-L21}{حائکًا واضع غناء فإن شهد لهما زحل أيضًا \BF{B48v}كانت صناعته مع هذه}
\begin{PNotes}
\PNote{317}{AB03-n317}{T: add. \RL{\ArN أو ألوان}}
\PNote{318}{AB03-n318}{T: add. \RL{\ArN أو صباغًا}}
\PNote{319}{AB03-n319}{T: \RL{\ArN سبّاك}}
\PNote{320}{AB03-n320}{B: \RL{\ArN الحصوص}; T: \RL{\ArN الفصوص}}
\PNote{321}{AB03-n321}{T: \RL{\ArN وشقّ}}
\PNote{322}{AB03-n322}{B: \RL{\ArN فإنّه}}
\PNote{323}{AB03-n323}{B: \RL{\ArN رقّاصًا}}
\end{PNotes}
\Folio{066r}{66r}
\MsL{1}{AB03-F066r-L01}{الأشياء المذكورة تاجرًا بائعا لحلي فإن شهد لهما المشتري كان صاحب}
\MsL{2}{AB03-F066r-L02}{خصومات متقلّدًا لأمور الحساب مواظبًا على مجامع الناس}
\MsL{3}{AB03-F066r-L03}{وأبواب السلطان معلّمًا للصبيان لها بأمور العامّة فإن كان}
\MsL{4}{AB03-F066r-L04}{المتولّيّ لأمور الصناعة عطارد والمرّيخ جميعًا كان صانعًا للتماثيل}
\MsL{5}{AB03-F066r-L05}{والسلاح نقّاشا للأشياء التي في الهياكل صانعًا لتماثيل صور الحيوان}
\MsL{6}{AB03-F066r-L06}{\{مضاهًا\}\PN{324} صاحب علاج باليد هما\PN{325} فاجرًا رديء الأفعال مزوّر \{الكتب\}}
\MsL{7}{AB03-F066r-L07}{فإن شهد لهما زحل كان المولود قتّالًا مختلسًا لثياب غيره سلّابًا قاطع}
\MsL{8}{AB03-F066r-L08}{طريق وصاحب غارات مكّارًا وإن شهد لهما المشتري كان محبًّا}
\MsL{9}{AB03-F066r-L09}{للسلاح أو محبًّا للزبارة\PN{326} ظابطًا\PN{327} خبيثًا للأعمال يتوكّل للغرباء ويكون}
\MsL{10}{AB03-F066r-L10}{كسبه من أمثال هذه الأشياء. فإن كان المتولّيّ لأمر الصناعة}
\MsL{11}{AB03-F066r-L11}{المرّيخ والزهرة جميعًا كان المولود صباغًا عطّارًا \{صانعًا\} للذهب}
\MsL{12}{AB03-F066r-L12}{والفضّة أو القلعيّ والأسرب فلّاحًا لاعبًا بالسلاح يعمل الأدويّة}
\MsL{13}{AB03-F066r-L13}{طبيبًا معالجًا بالأوديّة وإن شهد لهما زحل كان فيما على الحيوان}
\MsL{14}{AB03-F066r-L14}{الذي يذبح لله مكفّنًا للموتى نائحًا زامرًا على العيون صاحب}
\MsL{15}{AB03-F066r-L15}{\{أينوسنًا\PN{328} سيموس\} مقيمًا في مواضع فيها استدارة وسفك الدماء}
\MsL{16}{AB03-F066r-L16}{فإن شهد لهما المشتري كان مواظبًا على الهيكل صاحب زحر}
\MsL{17}{AB03-F066r-L17}{الطير قيّمًا يتزوّج من النساء ويتولّى خطب التزويج ويعيش من}
\MsL{18}{AB03-F066r-L18}{أشباه\PN{329} هذه الأشياء متنعّمًا رخيّ البال والصور التي تخصّ}
\MsL{19}{AB03-F066r-L19}{البروج التي تكون فيها الكواكب التي صورها صور الناس}
\MsL{20}{AB03-F066r-L20}{تعين معونة على جميع الصناعات التي تكون من المعرفة. فيحتاج}
\MsL{21}{AB03-F066r-L21}{إليها الناس والبروج التي صورها شبيهة بصور ذات الأربع}
\begin{PNotes}
\PNote{324}{AB03-n324}{B: \RL{\ArN \{مصارعًا\}}}
\PNote{325}{AB03-n325}{T: \RL{\ArN هجاء}}
\PNote{326}{AB03-n326}{T: \RL{\ArN للمبارزة}}
\PNote{327}{AB03-n327}{BT: \RL{\ArN ضابطا}}
\PNote{328}{AB03-n328}{T: \RL{\ArN أنثوسيمًا}}
\PNote{329}{AB03-n329}{T: \RL{\ArN أشياء}}
\end{PNotes}
\Folio{066v}{66v}
\MsL{1}{AB03-F066v-L01}{قوائم تعين على الصناعات التي \BF{B49r}تكون من الأشياء العونيّة والنجارات}
\MsL{2}{AB03-F066v-L02}{والبناء والتجارة وبروج الانقلابين والاعتدالين تعين على الصناعات}
\MsL{3}{AB03-F066v-L03}{التي تكون بالترجمة والنقل والهندسة والفلاحة وأمور الهياكل والبروج}
\MsL{4}{AB03-F066v-L04}{التي صورها شبيهة بصور الحيوان البريء وتعين على الصناعات}
\MsL{5}{AB03-F066v-L05}{التي تكون في الرطوبات وتستعمل فيها الرطوبات والحشائش وتجارة السفن}
\MsL{6}{AB03-F066v-L06}{والدفن واعمال الأشياء المملّحة أو المالحة. وإن كان صاحب موضع الصناعة}
\MsL{7}{AB03-F066v-L07}{القمر خاصّة وكان سيّره بعد الاجتماع مع عطارد ثمّ كان ذلك في الثور}
\MsL{8}{AB03-F066v-L08}{والجدي والسرطان كان المولود عرّافًا صاحب ذبائح ينظر في طساش\PN{330}}
\MsL{9}{AB03-F066v-L09}{الماء وإن كان في الرأس\PN{331} والسمكتين كان ممّن ينظر في الكهانة من}
\MsL{10}{AB03-F066v-L10}{الموتى ومن يحرّك الجنّ وإن كان في العذراء والعقرب كان ساحرًا}
\MsL{11}{AB03-F066v-L11}{منجّمًا مخبرًا بالغيب يتقدّم فيعرف الأشياء وإن كان في الميزان}
\MsL{12}{AB03-F066v-L12}{والحمل والأسد كان صاحب وحي معبّرًا للأحلام مغرمًا\PN{332}. فهذا هو}
\MsL{13}{AB03-F066v-L13}{المسلك الذي ينبغي أن يسلك في معرفة حال الصناعة\PN{333} ونوعها}
\MsL{14}{AB03-F066v-L14}{على جهة القياس. وأمّا مقدار قوّتها فينبغي أن ننظر فيه من قوّة}
\MsL{15}{AB03-F066v-L15}{الكواكب التي هي صاحبة التدبير وذلك أنّها إذا كانت}
\MsL{16}{AB03-F066v-L16}{شرقيّة أو على الأوتاد كان صناعاتهم صناعة رئاسيّة وقوّة}
\MsL{17}{AB03-F066v-L17}{وإن كانت غربيّة أو مائلة عن الأوتاد كانت صناعاتهم من}
\MsL{18}{AB03-F066v-L18}{تحت أيدي آخرین. وإذا علّت\PN{334} عليها السعود كانت نبيلة}
\MsL{19}{AB03-F066v-L19}{شريفة مريحة لا يخيب صاحبها ممدوحة. وإذا علّت\PN{335} عليها النحوس}
\MsL{20}{AB03-F066v-L20}{كانت نذلة خسيسة محسّرة يخيب صاحبها فإن كان زحل}
\MsL{21}{AB03-F066v-L21}{هو المضادّ عرض فيها برد وغلب عليها الإبطاء والتأخير}
\begin{PNotes}
\PNote{330}{AB03-n330}{BT: \RL{\ArN طساس}}
\PNote{331}{AB03-n331}{T: \RL{\ArN الرامي}}
\PNote{332}{AB03-n332}{T: \RL{\ArN معرفًا}}
\PNote{333}{AB03-n333}{B: \RL{\ArN الصناعات}}
\PNote{334}{AB03-n334}{T: \RL{\ArN غلبت}}
\PNote{335}{AB03-n335}{T: \RL{\ArN غلبت}}
\end{PNotes}
\Folio{067r}{67r}
\MsL{1}{AB03-F067r-L01}{وإن كان المرّيخ هو الذي يضادّ عرض فيها خطر وبشاعة. وإن كانا}
\MsL{2}{AB03-F067r-L02}{جميعًا متضادّين \BF{B49v}كانت هذه الأمور في غاية الرداءة في جملة العمر}
\MsL{3}{AB03-F067r-L03}{وتعرف زيادتها ونقصانها أيضًا من حال الكواكب التي هي بيت\PN{336} الحادث}
\Chap{I.32}{<32> في الازدواج}
\MsL{4}{AB03-F067r-L04}{بقياسها إلى الأوتاد المنسوبة إلى العشيّات. \Hd{\LR{\LatinF \Sg{〈}32\Sg{〉}}}\Hd{ في الازدواج} وإذا كان القول في الازدواج تابعًا لمّا وصفنا فينبغي أن يكون}
\MsL{5}{AB03-F067r-L05}{نظرنا في تزويج الرجال والنساء الجاري على السنة على ما أصف}
\MsL{6}{AB03-F067r-L06}{ينبغي أن ننظر في تزويج الرجال إلى حال القمر في وقت مولود وذلك}
\MsL{7}{AB03-F067r-L07}{أنّا ننظر أوّلًا إلى القمر فإن كان في الربعين الشرقيّين كان تزويج}
\MsL{8}{AB03-F067r-L08}{الرجال في حداثته أو يتزوّج بفتاة بعد أن قد طعن في السنّ. وإذا كان}
\MsL{9}{AB03-F067r-L09}{في الربعين الغربيّين كان بطيء التزويج أو يتزوّج بعجوز فإن كانت}
\MsL{10}{AB03-F067r-L10}{تحت الشعاع وشارك زحل في الشكل لم يتزوّج بتّة. ثمّ}
\MsL{11}{AB03-F067r-L11}{ننظر بعد ذلك على القمر في برج ذي صورة واحدة أو هو متّصل}
\MsL{12}{AB03-F067r-L12}{أيضًا بكوكب واحد فإنّه إن كان كذلك كان تزويجه بإمرأة واحدة}
\MsL{13}{AB03-F067r-L13}{وإن كان في برج ذي جسدين أو كان في برج كثير الصور وكان}
\MsL{14}{AB03-F067r-L14}{يتّصل بأكثر من كوكب واحد في برج واحد بعينه كان كثير التزويج.}
\MsL{15}{AB03-F067r-L15}{فإن كانت الكواكب التي يتّصل بها إمّا على جهة المقارنة وإمّا على جهة}
\MsL{16}{AB03-F067r-L16}{السعادة\PN{337} سعودًا كان النساء اللاتي يتزوّج بهنّ صالحات وإن كان}
\MsL{17}{AB03-F067r-L17}{نحوسًا كان النساء على ضدّ ذلك وذلك أنّه متى كان المتّصل به زحل}
\MsL{18}{AB03-F067r-L18}{كانت المرأة كدودة وعرّة الخلق وإن اتّصل به المشتري كانت}
\MsL{19}{AB03-F067r-L19}{عفيفة مدبّرة فإن اتّصل به المرّيخ كانت جريئة غير خاضعة}
\MsL{20}{AB03-F067r-L20}{وإن اتّصلت به الزهرة \BF{B50r}كانت جميلة بارّة\PN{338} سارّة. فإن اتّصل به}
\begin{PNotes}
\PNote{336}{AB03-n336}{T: \RL{\ArN سبب}}
\PNote{337}{AB03-n337}{T: \RL{\ArN الشهادة}}
\PNote{338}{AB03-n338}{T: \RL{\ArN باشّة}}
\end{PNotes}
\Folio{067v}{67v}
\MsL{1}{AB03-F067v-L01}{كان عطارد كانت المرأة فهمة بديعة\PN{339} وأيضًا فإنّ الزهرة إن كانت مع}
\MsL{2}{AB03-F067v-L02}{المشتري أو مع زحل أو مع عطارد كانت المرأة نافعة في أمر}
\MsL{3}{AB03-F067v-L03}{المعاش محبّة لزوجها ولأولادها فإن كانت مع المرّيخ كانت المرأة}
\MsL{4}{AB03-F067v-L04}{فاسدة لا ثبات لها ولا فهم ويجوّز بدّل قوله أنّ الزهرة إن}
\MsL{5}{AB03-F067v-L05}{كانت مع المشتري أن يقال إنّ القمر إذا اتّصل بالمشتري مع}
\MsL{6}{AB03-F067v-L06}{الزهرة. وأمّا في تزويج النساء فينبغي أن ننظر إلى حال الشمس في}
\MsL{7}{AB03-F067v-L07}{وقت مولدهنّ وذلك أنّ الشمس أيضًا متى كانت في الربعين}
\MsL{8}{AB03-F067v-L08}{الشرقيّين كان تزويج المرأة في أحداثها أو تتزوّج بحدث}
\MsL{9}{AB03-F067v-L09}{إذا طعنت في السنّ وإذا كان في الربعين الغربيّين كان}
\MsL{10}{AB03-F067v-L10}{تزويجها أو تتزوّج بشيخ وهي حدثة السنّ فإن كانت في}
\MsL{11}{AB03-F067v-L11}{برج ذي صورة واحدة ووجدت فيه كوكبًا واحدًا منسوبًا إلى}
\MsL{12}{AB03-F067v-L12}{الغدوات تزوّجت برجل واحد وإن كانت في بروج ذوات}
\MsL{13}{AB03-F067v-L13}{جسدين أو في بروج كثيرة الصور أو شاركه في الشكل أكثر من}
\MsL{14}{AB03-F067v-L14}{كوكب واحد من التي تنسب إلى الغدوات كانت كثيرة التزويج.}
\MsL{15}{AB03-F067v-L15}{وكذلك أيضًا إذا شارك زحل الشمس في الشكل كان الزوج}
\MsL{16}{AB03-F067v-L16}{صعبًا لا مودّة له غير خاضع وإن شاركتهما الزهرة كان ظاهرًا}
\MsL{17}{AB03-F067v-L17}{جميلًا وإن شاركهما عطارد كان نافعًا في أمر المعاش عمّالًا. فإن}
\MsL{18}{AB03-F067v-L18}{كانت الزهرة مع ذلك كان الزوج كليلًا متوقّيًّا للجماع وإن كانت}
\MsL{19}{AB03-F067v-L19}{الزهرة مع المرّيخ كان الزوج حادّ المزاج منهمكًا في الجماع فإن}
\MsL{20}{AB03-F067v-L20}{كانت مع عطارد كان مشفقًا على والده والذي ينبغي في هذا}
\MsL{21}{AB03-F067v-L21}{الموضع بالربعين الشرقيّين هو كما ذكرنا في صدر الكتاب.}
\begin{PNotes}
\PNote{339}{AB03-n339}{T: \RL{\ArN نزيهة}}
\end{PNotes}
\Folio{068r}{68r}
\MsL{1}{AB03-F068r-L01}{أمّا في الشمس فالربعان المتقدّمان للنقطة التي تطلع والنقطة}
\MsL{2}{AB03-F068r-L02}{التي تغرب من فلك البروج وفي القمر الربعان المتقدّمان فيما بين}
\MsL{3}{AB03-F068r-L03}{الاجتماع والامتلاء إلى أن ينتصف القمر في الضوء من جانبيّ الامتلاء.}
\MsL{4}{AB03-F068r-L04}{وأمّا الربعان الغربيان فالربعان المقابلان لهذين الربعين}
\MsL{5}{AB03-F068r-L05}{الذين تقدّم \BF{B50v}ذكرهما ويكون أمر التأمّل\PN{340} في أكثر الأمر ثابتًا باقيًا}
\MsL{6}{AB03-F068r-L06}{إذا اشترك النيّران في الشكل في المولدين جميعًا أعني مولد}
\MsL{7}{AB03-F068r-L07}{الرجل ومولد المرأة اشتراكًا متّفقًا أعني أن يكونا على تسديس}
\MsL{8}{AB03-F068r-L08}{أو تثليث وخاصّة متى تهيّأ أن يكون ذلك على التبديل لا سيّما إذا كان تهيّأ}
\MsL{9}{AB03-F068r-L09}{أن يكون القمر في مولد الرجل مشاركًا في الشكل للشمس وفي}
\MsL{10}{AB03-F068r-L10}{مولد المرأة مشاركة تابعًا ولم نحل أمر التأمّل بتّة يفترق}
\MsL{11}{AB03-F068r-L11}{من أدنى سبب إذا كان المواضع التي ذكرنا للنيّرين في}
\MsL{12}{AB03-F068r-L12}{بروج غير مرتبطة أو كانت متقابلة على القطر أو على}
\MsL{13}{AB03-F068r-L13}{تربيع. فإن كانت الكواكب المسعدة ترى النيّرين المتّفقين}
\MsL{14}{AB03-F068r-L14}{في الاشتراك في الشكل بقي أمر التزويج لذيذًا محبوبًا}
\MsL{15}{AB03-F068r-L15}{نافعًا فإن رأتهما المنحسة كان ثباته مع خصومة وكرامة}
\MsL{16}{AB03-F068r-L16}{وخسران وعلى هذا المثال أيضًا يكون الأمر في الأحوال}
\MsL{17}{AB03-F068r-L17}{التي ليست بمتّفقة فإن السعود متى شهدت للنيّرين ولم}
\MsL{18}{AB03-F068r-L18}{يفترق أمر التأمّل ولم ينحل بمرّة بل افتراق يكون فيه}
\MsL{19}{AB03-F068r-L19}{مراجعة وتذكر وتعطف وتدوم به المحبّة والألفة. فإن شهدت}
\MsL{20}{AB03-F068r-L20}{لها النحوس كان افتراقهما مع شرّ يجري فيما بينهما وقوى}
\MsL{21}{AB03-F068r-L21}{فإن صار معهما عطارد وكان الافتراق مع شهوة ولوم وعذل\PN{341}}
\begin{PNotes}
\PNote{340}{AB03-n340}{T: \RL{\ArN التأهّل} (wherever it occurs in this chapter)}
\PNote{341}{AB03-n341}{T: \RL{\ArN عدل}}
\end{PNotes}
\Folio{068v}{68v}
\MsL{1}{AB03-F068v-L01}{فإن كانت الزهرة معهما كان ذلك بحسب فجور أو سحر وما أشبههما.}
\MsL{2}{AB03-F068v-L02}{فأمّا الازدواج الذي يكون على جهة أخرى أيّ جهة كانت}
\MsL{3}{AB03-F068v-L03}{فينبغي أن نميّزه ونعلمه من نظرنا إلى الزهرة والمرّيخ وزحل. فإن}
\MsL{4}{AB03-F068v-L04}{كانت مع الكوكبين\PN{342} على الحال الملائمة كان أمر التزويج ملائمًا}
\MsL{5}{AB03-F068v-L05}{جاريًّا على السنة وذلك أنّ فيما بين الزهرة وفيما بين كلّ}
\MsL{6}{AB03-F068v-L06}{من هذين الكوكبين مناسبة ومناسبتها للمرّيخ تدلّ}
\MsL{7}{AB03-F068v-L07}{على أنّ الازدواج يكون مع فتاة والمناسبة صارت بينهما}
\MsL{8}{AB03-F068v-L08}{لأن شرف كلّ واحد منهما في برج من مثلّثه الآخر ومناسبتها}
\MsL{9}{AB03-F068v-L09}{لزحل تدلّ على أنّ التزوّج يكون مع عجوز والمناسبة}
\MsL{10}{AB03-F068v-L10}{صارت بينهما لأن شرف كلّ واحد منهما في برج من}
\MsL{11}{AB03-F068v-L11}{مثلّثه الآخر ومن قبل ذلك صارت الزهرة مع المرّيخ يحدث}
\MsL{12}{AB03-F068v-L12}{أحوال العشق مطلقًا. فإن كان معهما عطارد كان ذلك}
\MsL{13}{AB03-F068v-L13}{مشهورًا أيضًا فإن كانت \BF{B51r}في بروج مشتركة لهما مثل الجدي}
\MsL{14}{AB03-F068v-L14}{والسمكتين دلّت على أنّ المجامعة تكون في الإخوة والقرابات}
\MsL{15}{AB03-F068v-L15}{فإن حضر القمر في مولد الرجال كان الرجل ممّن يأتي أختين}
\MsL{16}{AB03-F068v-L16}{وإمرأتين بينهما قرابة وإن حضر في مولد النساء المشتري}
\MsL{17}{AB03-F068v-L17}{كانت المرأة ممّن ينكحها إخوان أو رجلان بينهما قرابة. فإن}
\MsL{18}{AB03-F068v-L18}{كانت أيضًا الزهرة مع زحل كان أمر التأهّل بالجملة لذيذًا}
\MsL{19}{AB03-F068v-L19}{حسن الثبات. وإن كان مع ذلك عطارد كان نافعًا وإن}
\MsL{20}{AB03-F068v-L20}{كان ذلك المرّيخ كان أمر التأهّل غير ثابت ضارًّا مع غيره}
\MsL{21}{AB03-F068v-L21}{فإن كانت حالها شبيهة بحالها كان الازدواج فيما بين}
\begin{PNotes}
\PNote{342}{AB03-n342}{T: \RL{\ArN النيّرين}}
\end{PNotes}
\Folio{069r}{69r}
\MsL{1}{AB03-F069r-L01}{المرأة وبين من سنّه مثل سنّها وإن كان أكثر تشريفًا\PN{343} منها تزوّجت المرأة}
\MsL{2}{AB03-F069r-L02}{منه و\PN{344} أسنّ منها والرجل بإمرأة أسنّ منه. فإن كانت الزهرة وزحل في}
\MsL{3}{AB03-F069r-L03}{بروج مشتركة لهما مثل الجدي والميزان وقع الازدواج بين القرابات}
\MsL{4}{AB03-F069r-L04}{وإن كان القمر على الحال التي ذكرنا وكانت هذه الحال في الطالع أو في}
\MsL{5}{AB03-F069r-L05}{وسط السماء نكح الذكر أمّه أو خالته أو إمرأة أبيه وجامعت المرأة ابنها\PN{345}}
\MsL{6}{AB03-F069r-L06}{أو ابن أختها أو زوج ابنتها. فإن كانت الشمس مع الحال التي}
\MsL{7}{AB03-F069r-L07}{ذكرنا أو كانت الكواكب غاربة نكح الذكر ابنته أو ابنة أخيه أو}
\MsL{8}{AB03-F069r-L08}{امرأة أبيه\PN{346} وجامعة المرأة والدها أوعمّها أو زوج أمّها فإن لم}
\MsL{9}{AB03-F069r-L09}{تكن هذه الحال التي قدّمنا ذكرها في بروج متحابّة والمتحابّة}
\MsL{10}{AB03-F069r-L10}{هي من جنس واحد [[توجد]] يعني المذكّرة والمؤنّثة لأن كلّ واحد}
\MsL{11}{AB03-F069r-L11}{منهما جنس والأشبه أن يكون الذي أرادها هناء المذكّرة.}
\MsL{12}{AB03-F069r-L12}{فإذا لم يكن في بروج المتحابّة التي ذكرها وكانت في مواضع}
\MsL{13}{AB03-F069r-L13}{\{مؤنّثة\} كان المولود منهمكًا في الجماع ليسهل عليه أن يرتكبه وأن}
\MsL{14}{AB03-F069r-L14}{يرتكب منه على كلّ جهة وبعض الصور إذا كانت فيها هذه الحال}
\MsL{15}{AB03-F069r-L15}{كان المولود مع ذلك نجسًا في أمر الجماع وهذه الصورة وهي}
\MsL{16}{AB03-F069r-L16}{معادلة\PN{347} الأسد والحمل والذي يسمّى باليونانيّة \BF{B51v}الشبيه بحرف\PN{348}}
\MsL{17}{AB03-F069r-L17}{اللام\PN{349} وهو هداعر وهو الذي آخر کواکب الدبران والذي يسمّى}
\MsL{18}{AB03-F069r-L18}{باليونانيّة فالفس وهو آخر الأسد ووجه الجدي. فإن كانت}
\MsL{19}{AB03-F069r-L19}{هذه الحال على الوتدين الأوّلين أعني الشرقيّ والجنوبيّ}
\MsL{20}{AB03-F069r-L20}{كان المولود مجاهدًا نافعًا له\PN{350} شاهدًا \{لها\} في مجامع الناس}
\begin{PNotes}
\PNote{343}{AB03-n343}{T: \RL{\ArN تشريقًا}}
\PNote{344}{AB03-n344}{B: \RL{\ArN من هو}}
\PNote{345}{AB03-n345}{T: \RL{\ArN أبيها}}
\PNote{346}{AB03-n346}{T: \RL{\ArN ابنه}}
\PNote{347}{AB03-n347}{T: \RL{\ArN المقاديم}}
\PNote{348}{AB03-n348}{B: \RL{\ArN بحذف}}
\PNote{349}{AB03-n349}{T:\RL{\ArN  الأمر}}
\PNote{350}{AB03-n350}{T: \RL{\ArN بأفعاله}}
\end{PNotes}
\Folio{069v}{69v}
\MsL{1}{AB03-F069v-L01}{وإن \{كلفت\}\PN{351} على الوتدين الآخرين أعني الغربيّ والشماليّ}
\MsL{2}{AB03-F069v-L02}{كان خصيًّا أو عاقرًا أو غير مثقوب فإن كان مع ذلك المرّيخ إن}
\MsL{3}{AB03-F069v-L03}{كان ذكرًا كان مقطوع الذكر والأنثيين إن كانت أنثى كانت}
\MsL{4}{AB03-F069v-L04}{سحّاقة وينبغي لنا أن نفحص عن حال المولود من الرجال في أمور}
\MsL{5}{AB03-F069v-L05}{الجماع من المرّيخ على هذه الجهة وذلك أنّه إذا فارق\PN{352} الزهرة وزحل}
\MsL{6}{AB03-F069v-L06}{وشهد له المشتري كان المولود ظاهرًا عفيفًا في أمور الجماع يقصد}
\MsL{7}{AB03-F069v-L07}{فيه إلى ما تحتاج إليه الطبيعة فقط. فإن اتّفق أن يكون مع}
\MsL{8}{AB03-F069v-L08}{زحل وحده كان متوقّيًّا للجماع عاجزًا عنه باردًا فيه. فإن شارك}
\MsL{9}{AB03-F069v-L09}{الزهرة والمشتري في الشكل كان سريع الحركة إلى ذلك شديد}
\MsL{10}{AB03-F069v-L10}{الشهوة له إلّا أنّه يضبط نفسه عنه ويصونها ويحذر\PN{353} الأشياء}
\MsL{11}{AB03-F069v-L11}{القبيحة. فإن كان مع الزهرة وحدها أو كان ثمّ المشتري}
\MsL{12}{AB03-F069v-L12}{من غير أن يكون معهما زحل كان شرّها كسلانا يحتال}
\MsL{13}{AB03-F069v-L13}{لنفسه اللذات من كلّ وجه. فإن كان أحد الكوكبين يعني المرّيخ}
\MsL{14}{AB03-F069v-L14}{والزهرة منسوبًا إلى العشيّات والآخر منسوبًا إلى الغدوات كان}
\MsL{15}{AB03-F069v-L15}{المولود مائلًا إلى الذكورة والإناث جميعًا بالسواء لا يميل إلى}
\MsL{16}{AB03-F069v-L16}{أحد الجنسين أكثر من الآخر. فإن كانا جميعًا منسوبين إلى}
\MsL{17}{AB03-F069v-L17}{العشيّات كان منهمكًا في مجامعة الإناث خاصّة. فإن كانت}
\MsL{18}{AB03-F069v-L18}{البروج مؤنّثة كان محتملًا أن يفعل به ذلك فإن كانا جميعًا}
\MsL{19}{AB03-F069v-L19}{منسوبين إلى الغدوات كان مائلًا إلى الذكورة فقط وإن كانت}
\MsL{20}{AB03-F069v-L20}{البروج مذكّرة كان مبتلى بجماع الذكور في أيّ قائمة\PN{354} كانوا. [[في]]}
\begin{PNotes}
\PNote{351}{AB03-n351}{T: \RL{\ArN كانت}}
\PNote{352}{AB03-n352}{T: \RL{\ArN قارن}}
\PNote{353}{AB03-n353}{T: \RL{\ArN ويحدر}}
\PNote{354}{AB03-n354}{T: \RL{\ArN قامة}}
\end{PNotes}
\Folio{070r}{70r}
\MsL{1}{AB03-F070r-L01}{\BF{B52r}وإن كانت الزهرة أكثر تغربًا كان ممّن يجامع السفل أو الإماء}
\MsL{2}{AB03-F070r-L02}{أو غيرهما ممّن هو غريب في الجنس مثل السودان الخزر والروم والترك}
\MsL{3}{AB03-F070r-L03}{وما أشبه ذلك أو من هو غير قرابة فإن كان المرّيخ أكثر تغربًا كان ممّن يجامع}
\MsL{4}{AB03-F070r-L04}{من هو أشرف منه مثل الذي يزني بمحصنة أو بمولاته وأمّا في أمر المواليد من النساء}
\MsL{5}{AB03-F070r-L05}{فينبغي أن ننظر فيه إلى الزهرة وذلك أنّها إذا شاركلت المشتري وعطارد في}
\MsL{6}{AB03-F070r-L06}{الشكل كانت المرأة عفيفة طاهرة في أمور الجماع فإن لم يكن هناك زحل وشاکلت}
\MsL{7}{AB03-F070r-L07}{عطارد كانت متحرّكة فيه إلى الجماع شبقة إلّا أنّها تكون متوقّيّة حذرة في أكثر}
\MsL{8}{AB03-F070r-L08}{الأمر متجنّبة\PN{355} للأمور القبيحة. فإن كانت الزهرة مع المرّيخ وحده أو شاركته في}
\MsL{9}{AB03-F070r-L09}{الشكل كانت المرأة شهمة\PN{356} منهمكة في الجماع حريصة عليه. فإن كان معهما المشتري}
\MsL{10}{AB03-F070r-L10}{وكان المرّيخ تحت الشعاع جامعت المرأة العبيد والسفلي والغرباء من قبيلتها}
\MsL{11}{AB03-F070r-L11}{فإن كانت الزهرة تحت الشعاع جامعت الأشراف أو مواليها على جهة الزناء}
\MsL{12}{AB03-F070r-L12}{والفجور. فإن كانت هذه الكواكب مؤنّثة في مواضعها أو في أشكالها كانت منهمكة}
\MsL{13}{AB03-F070r-L13}{في مجامعة النساء أيضًا فإن شاكل زحل هذه الأحوال الموصوفة ثمّ كان مؤنّثًا}
\MsL{14}{AB03-F070r-L14}{أيضًا كان سببًا للنجاسة والرجاسة وإن كان شرقيًّا مذكّرًا كانت المرأة متشرّفة}
\MsL{15}{AB03-F070r-L15}{من السطوح أو ممّن يعشّق المتشرّفين من السطوح والمشتري أبدًا يعيّن على تحسين}
\Chap{I.33}{<33> في الأولاد}
\MsL{16}{AB03-F070r-L16}{هذه الآفات وعطارد يعين على شهوتها\PN{357} وقبحها. \Hd{\LR{\LatinF \Sg{〈}33\Sg{〉}}}\Hd{ في الأولاد} ينبغي أن}
\MsL{17}{AB03-F070r-L17}{تنظر إلى الكواكب التي هي في المواضع\PN{358} الذي هو على سمت}
\MsL{18}{AB03-F070r-L18}{رؤوسنا \BF{B52v}أو في الموضع الذي يتلوه وهو موضع السعادة [[ أو في الموضع يتلوه ما هو موضع السعادة ]] أو في الموضع الذي يشاركه في الشكل}
\MsL{19}{AB03-F070r-L19}{فإن لم نجد ذلك فينبغي لنا أن نأخذ الكواكب التي تقابلها على القطر ونستعمل\PN{359}}
\MsL{20}{AB03-F070r-L20}{والمشتري والزهرة في إعطاء الأولاد والشمس والمرّيخ والزحل في عددهم أو في قلّتهم}
\begin{PNotes}
\PNote{355}{AB03-n355}{T: \RL{\ArN متجنّسة}}
\PNote{356}{AB03-n356}{T: \RL{\ArN شرهة}}
\PNote{357}{AB03-n357}{T: \RL{\ArN شهرتها}}
\PNote{358}{AB03-n358}{B: \RL{\ArN الموضع}}
\PNote{359}{AB03-n359}{B: \RL{\ArN [[مشتري القمر]]}; T: add. \RL{\ArN القمر}}
\end{PNotes}
\Folio{070v}{70v}
\MsL{1}{AB03-F070v-L01}{وتجعل عطارد مع أحدهما أعني مع الذي يشاركه في الشكل من الصنفين}
\MsL{2}{AB03-F070v-L02}{أيّهما كان وإذا كان شرقيًّا كان معطيًّا للأولاد وإذا كان غربيًّا كان عنه عدمهم.}
\MsL{3}{AB03-F070v-L03}{فإذا كانت الكواكب المعطيّة هي هذه بالجملة الكواكب وكانت وحدها كان المولود}
\MsL{4}{AB03-F070v-L04}{وحيدًا فإذا كانت في بروج ذوات جسدين مرتبة\PN{360} كان المولود دون\PN{361} إثنين.}
\MsL{5}{AB03-F070v-L05}{وكذلك أيضًا إذا كانت في بروج كثيرة الذرع مثل السمكتين والسرطان والعقرب}
\MsL{6}{AB03-F070v-L06}{كان ما يولد إثنين أو أكثر من إثنين وإن كانت مذكّرة في البروج وفي الأشكال}
\MsL{7}{AB03-F070v-L07}{التي لها بقياسها إلى الشمس كان ما يولد [[ صغورًا ]] ذكورًا فإن كانت مؤنّثة في}
\MsL{8}{AB03-F070v-L08}{البروج كان ما يولد إناثًا. فإن غلبت عليها النحوس أو كانت في بروج عواقر مثل الأسد}
\MsL{9}{AB03-F070v-L09}{والعذراء فإنّهم يولد لهم إلّا أنّهم يكون ذلك لخير ولا يكون لمّا يولد بقاء. فإن}
\MsL{10}{AB03-F070v-L10}{استولت الشمس والنحوس على هذه المواضع التي تقدّم ذكرها فإنّها إن كانت}
\MsL{11}{AB03-F070v-L11}{في بروج مذكّرة أو عواقر وغلبت عليها النحوس دلّت على عدم الأولاد البتّة. فأمّا}
\MsL{12}{AB03-F070v-L12}{إن كانت في بروج مؤنّثة أو كثيرة الزرع\PN{362} أو شهدت لها السعود دلّت على}
\MsL{13}{AB03-F070v-L13}{ولد به آفة أو ولد فصيّر العمر فإن كان للجزئين جميعًا حصّة في البروج}
\MsL{14}{AB03-F070v-L14}{المولوده كان عنها عدم الأولاد إمّا بعضهم وإمّا كلّهم ومقدار ذلك على حسب}
\MsL{15}{AB03-F070v-L15}{رجحان الكواكب التي شهدت لكلّ واحد من الجزئين على الآخر إن يوجد}
\MsL{16}{AB03-F070v-L16}{رجحانها بسبب كثرة الكواكب وإن اتّفق أن يكون ذلك بسبب قوّتها وإنّها أشدّ}
\MsL{17}{AB03-F070v-L17}{تشريفًا\PN{363} أو أشدّ تمكّنا من الوتد أو أشدّ استعلاء وفي \BF{B53r}أنّها صاعدة إلى الوتد. فإن}
\MsL{18}{AB03-F070v-L18}{كانت الكواكب هي أرباب البروج التي ذكرها المعطيّة للأولاد شرقيّة في مواضعها}
\MsL{19}{AB03-F070v-L19}{الخاصّيّة كان الولد شريفًا مشهورًا. فإن كانت غربيّة أو في المواضع التي هي}
\MsL{20}{AB03-F070v-L20}{للحيّز لمخالف لها كان المولود دليلًا غير مذكور. فإن كانت موافقة للطالع\PN{364}}
\MsL{21}{AB03-F070v-L21}{وسهم البخت كان المولود محبوبًا من والديه ويبقى حتّى يتزوّج ويرث مال}
\begin{PNotes}
\PNote{360}{AB03-n360}{T: \RL{\ArN مؤنّثة}}
\PNote{361}{AB03-n361}{B: \RL{\ArN المولودون}; T: om.}
\PNote{362}{AB03-n362}{T: \RL{\ArN الذرع}}
\PNote{363}{AB03-n363}{BT: \RL{\ArN تشريقًا}}
\PNote{364}{AB03-n364}{T: \RL{\ArN للساعة}}
\end{PNotes}
\Folio{071r}{71r}
\MsL{1}{AB03-F071r-L01}{والديه وإن كان غير مرتبطة أو مخالفة كان الولد مخاصمًا معاديًا لأبويه\PN{365} لا يرث}
\MsL{2}{AB03-F071r-L02}{ما لهما وأيضًا فإنّ الكواكب التي تعطي الأولاد إن اشتركت في الشكل اشتراك موافقة}
\MsL{3}{AB03-F071r-L03}{كان الإخوة متحابّين يكرم بعضهم بعضًا فإن كانت غير مرتبطة أو متقابلة على}
\MsL{4}{AB03-F071r-L04}{القطر كان المولود محبًّا لعداوة إخوته مغتالًا لهم. وأمّا أمور الأولاد الجزئيّة}
\MsL{5}{AB03-F071r-L05}{فإنّما يمكن للإنسان أن يعلمها إذا هو استعمل في كلّ واحد منهما الكو[[ا]]كب الذي هو}
\MsL{6}{AB03-F071r-L06}{سبيل عطاء\PN{366} الأولاد بأن يقيمه مقام الطالع فيعلم من سائر الأحوال التي توجد}
\Chap{I.34}{<34> في الأصدقاء والأعداء والعبيد}
\MsL{7}{AB03-F071r-L07}{أمور الأولاد الجزئيّة على الأمر المجمل كما نعلمها\PN{367} من المواليد. \Hd{\LR{\LatinF \Sg{〈}34\Sg{〉}}}\Hd{ في الأصدقاء والأعداء والعبيد}}
\MsL{8}{AB03-F071r-L08}{أمّا حالات المحبّة والبغض الذي يسمّى ما كان منها أعظم وأطول مدّة صداقة}
\MsL{9}{AB03-F071r-L09}{وعداوة ويسمّى ما كان منها أقلّ وأقصر مدّة المواصلة والهجران فينبغي أن يكون}
\MsL{10}{AB03-F071r-L10}{فحصنا عنها على هذه الجهة وهي أن تتفقّد فيما عظم منها المواضع التي هي أحقّ}
\MsL{11}{AB03-F071r-L11}{بالرئاسة في المولدين جميعًا يعني مولد المولود ومولد من تريد أن تعلم كيف}
\MsL{12}{AB03-F071r-L12}{يكون الحال بينه وبينه في الصداقة والعداوة إلى \BF{B53v}موضع الشمس وموضع القمر}
\MsL{13}{AB03-F071r-L13}{وموضع الطالع وموضع سهم السعادة وذلك أنّ هذه إذا كانت في برج واحد بعينه}
\MsL{14}{AB03-F071r-L14}{في المولودين أو كانت في مواضع متبادلة لها جميعًا وأمّا أكثرها وبخاصّة إذا كان}
\MsL{15}{AB03-F071r-L15}{من الطالعين هو نحو من سبع عشرة درجة فإنّه يكون عنها صداقة ثابتة لا تنحل}
\MsL{16}{AB03-F071r-L16}{ولا تتكدر. فإن كانت في مواضع غير مرتبطة أو مواضع متقابلة على القطر حدث}
\MsL{17}{AB03-F071r-L17}{عنها عداوة شديدة وتضادّ طول مدّته ويريد بقوله التبادل أن تكون شمس}
\MsL{18}{AB03-F071r-L18}{أحد المولودين في موضع قمر الآخر والقمر في موضع الشمس أو تكون تلك الحال في سهم}
\MsL{19}{AB03-F071r-L19}{السعادة وفي الطالع معهم. فإن لم يكن على أحد الوجهين الذي ذكرنا لكن كانت}
\MsL{20}{AB03-F071r-L20}{خاصّة في البروج التي تشترك في الأشكال فإنّها إن كانت في التي تشترك في التثليث}
\MsL{21}{AB03-F071r-L21}{أو التسديس كانت الصداقة أقلّ من الأوّل. فإن كانت في التي تشترك في التربيع}
\begin{PNotes}
\PNote{365}{AB03-n365}{T: \RL{\ArN لواليه}}
\PNote{366}{AB03-n366}{T: \RL{\ArN سبب إعطاء}}
\PNote{367}{AB03-n367}{T: \RL{\ArN نعملها}}
\end{PNotes}
\Folio{071v}{71v}
\MsL{1}{AB03-F071v-L01}{كانت العداوة أقلّ حتّى أنّه يعرض في بعض الأوقات من الصداقة قطع الكلام}
\MsL{2}{AB03-F071v-L02}{فيما بين الصديقين وقلّة المفاوضة وذلك يكون عند ممرّ النحوس بموضع الشكل}
\MsL{3}{AB03-F071v-L03}{ويعرض في العداوة صلح ومراجعة إذا صارت السعود إلى مواضع الأشكال ولأنّ}
\MsL{4}{AB03-F071v-L04}{أصناف المحبّة والبغضة ثلاثة وذلك أنّ أصحاب هاتين الحالتين إمّا أن يكونوا}
\MsL{5}{AB03-F071v-L05}{إنّما يصيّرون إليهما بمثل الهواء والرأي وإمّا بسبب منفعة أو مضرّة وإمّا}
\MsL{6}{AB03-F071v-L06}{بسبب لذّة أو غير ذلك\PN{368}. فإذا كانت جميع هذه المواضع التي تقدّم ذكرها أو بعضها}
\MsL{7}{AB03-F071v-L07}{ملائمة بعضها لبعض كان المجتمع في هذه الأنواع محبّة وإذا كانت \{لا يلادم\}\PN{369} بعضها}
\MsL{8}{AB03-F071v-L08}{بعضًا كان عنها عداوة فإن كانت الملائمة وحدها فيما بين مواضع النيّرين وحدهما}
\MsL{9}{AB03-F071v-L09}{كانت المحبّة باختيار ورأي وهذه هي المودّة الفاضلة الموثوق بها وكانت}
\MsL{10}{AB03-F071v-L10}{العداوة في غاية الرداءة ولم ينحل \BF{B54r}وكذلك أيضًا متى كانت الملائمة أو ضدّها}
\MsL{11}{AB03-F071v-L11}{فيما بين موضعيّ سهم السعادة كان ما يعرض من ذلك بسبب \PN{370}الّذة أو غمّ وقد}
\MsL{12}{AB03-F071v-L12}{ينبغي لنا أن نتفقّد أمر المواضع التي يكون فيها اشتراك الأشكال وننظر}
\MsL{13}{AB03-F071v-L13}{ما المستعلي عليها وما الذي يراها من الكواكب وذلك أنّ المواليد التي تكون}
\MsL{14}{AB03-F071v-L14}{فيها الكواكب المستعلية على مواضع الشكل المشترك في برج واحد بعينه بالتقريب}
\MsL{15}{AB03-F071v-L15}{أو في الذي يرتفع ويصيّر إليه فإنّه ينبغي أن تنسب أمر الصداقة وأمر العداوة}
\MsL{16}{AB03-F071v-L16}{في قوّتها وثباتها وصحّتها إلى تلك الكواكب. وأمّا المواليد التي تكون فيها رؤية\PN{371}}
\MsL{17}{AB03-F071v-L17}{الكواكب فاضلة نافعة في فعل الخير والقوّة فإنّما ينبغي أن تنسب منفعة الصداقة}
\MsL{18}{AB03-F071v-L18}{وصلاح أمر العداوة إلى تلك الكواكب التي ترى وأمّا أمر المواصلة والهجران}
\MsL{19}{AB03-F071v-L19}{الذين يكونان زيانا في بعض الأوقات ثمّ ينقصان. فينبغي أن تنظر فيهما إلى الكواكب}
\MsL{20}{AB03-F071v-L20}{في كلّ واحد من هذين المولودين أعني في أيّ الأزمنة يصيّر الهيلاج من كواكب أحد}
\MsL{21}{AB03-F071v-L21}{المولودين إلى مواضع كواكب الموضع الآخر فإنّه قد يكون بحسب ذلك محبّة}
\begin{PNotes}
\PNote{368}{AB03-n368}{T: \RL{\ArN لذّة أو غمّ}}
\PNote{369}{AB03-n369}{B: \RL{\ArN لا يلائم}}
\PNote{370}{AB03-n370}{seems like a line was skipped here cf. T}
\PNote{371}{AB03-n371}{T: \RL{\ArN رديئة}}
\end{PNotes}
\Folio{072r}{72r}
\MsL{1}{AB03-F072r-L01}{وبعض جزئيان وأقلّ ما يبقيان إلى وقت انقضاء ذلك وأكثر بقائهما إلى وقت}
\MsL{2}{AB03-F072r-L02}{لحوق كوكب آخر من الكواكب التي تعدّ\PN{372}. فزحل والمشتري إذا أمر أحدهما بموضع صاحبه}
\MsL{3}{AB03-F072r-L03}{حدث عنهما صداقة بسبب فلاحة أو ميراث. وأمّا زحل والمرّيخ إذا أمر أحدهما}
\MsL{4}{AB03-F072r-L04}{بمكان\PN{373} صاحبه فإنّه يحدث خصومة واضطرابًا يكون باختيار وأمّا زحل والزهرة}
\MsL{5}{AB03-F072r-L05}{إذا أمر أحدهما بالآخر فإنّه يحدث عنهما اتّصال \{وألفة\}\PN{374} بأسباب الأقارب إلّا أن}
\MsL{6}{AB03-F072r-L06}{ينقضي سريعًا. وأمّا زحل وعطارد فيحدث عنهما اشتراك في المعاش ومخالطة}
\MsL{7}{AB03-F072r-L07}{وأخذ وإعطاء وتجارة وأسرار وأمّا المشتري والمرّيخ فيحدث عنهما صحبة في}
\MsL{8}{AB03-F072r-L08}{أسباب الحظوة وبأسباب التدبير. وأمّا المشتري والزهرة فيحدث \BF{B54v}عنهما صداقة}
\MsL{9}{AB03-F072r-L09}{بأسباب النساء أو بأسباب العبادة والهياكل أو بأسباب الوحي وما أشبهها وأمّا}
\MsL{10}{AB03-F072r-L10}{المشتري وعطارد فيحدث عنهما معاشرة بأسباب المنطق والعلوم والفلسفة التي}
\MsL{11}{AB03-F072r-L11}{تكون باختيار. وأمّا المرّيخ والزهرة فيحدث عنهما اتّصال بأسباب العنق والفجور}
\MsL{12}{AB03-F072r-L12}{والزناء إلّا أنّ ذلك ليس ممّا يوثق به ولا يثبت مدّة طويلة وأمّا المرّيخ وعطارد}
\MsL{13}{AB03-F072r-L13}{فيحدّث عنهما عداوة وشهوة وخصومة بأسباب تعاويذ أو سحر. وأمّا الزهرة}
\MsL{14}{AB03-F072r-L14}{وعطارد فإنّه يحدث عنهما اشتراك في المعاش يكون بأسباب صناعة الموسيقي أو}
\MsL{15}{AB03-F072r-L15}{بأسباب كتب أو بأسباب النساء وأمّا الزيادة\PN{375} والنقصان في قوّة المواصلة والهجران.}
\MsL{16}{AB03-F072r-L16}{فينبغي أن يعلم من حال المواضع التي توجد\PN{376} إذا قيست بالأربعة المواضع الأوّل}
\MsL{17}{AB03-F072r-L17}{الرئيسيّة وذلك أنّها إذا كانت على الأوتاد ومع سهم السعادة أو مع النيّرين كان}
\MsL{18}{AB03-F072r-L18}{ما يدلّ عليه أظهر وأبين وإذا بعدت عنهما كان ذلك غير ظاهر والأمر في ضررها}
\MsL{19}{AB03-F072r-L19}{ونفعها يعلم من خاصّة حال الكواكب التي ترى في المواضع التي تقدّم ذكرها في الخير}
\MsL{20}{AB03-F072r-L20}{والشرّ. وأمّا القول في العبيد ومحبّة مواليهم لهم أو بعضهم فيعلم خاصّة من البروج}
\MsL{21}{AB03-F072r-L21}{التي تكون في موضع الشقاء وهو الثاني عشر من الطالع ومن الموافقة الطبيعيّة}
\begin{PNotes}
\PNote{372}{AB03-n372}{T: \RL{\ArN تتبع}}
\PNote{373}{AB03-n373}{B: \RL{\ArN بموضع}}
\PNote{374}{AB03-n374}{B: \RL{\ArN والفنه}; T: \RL{\ArN وإلف}}
\PNote{375}{AB03-n375}{T: \RL{\ArN الزهرة}}
\PNote{376}{AB03-n376}{T: \RL{\ArN تؤخذ}}
\end{PNotes}
\Folio{072v}{72v}
\MsL{1}{AB03-F072v-L01}{التي تكون للكواكب التي ترى هذا الموضع في وقت المولد نفسه وفي وقت يصيّر}
\MsL{2}{AB03-F072v-L02}{إليه أو عند المقابلة وبخاصّة إذا كانت الكواكب المستوليّة على البروج إمّا موافقة}
\MsL{3}{AB03-F072v-L03}{المواضع الرئيسيّة من مواليدهم في مشاركتها لها في الشكل أو مخالفتها له في ذلك.}
\Chap{I.35}{<35> في الغربة الأسفار}
\MsL{4}{AB03-F072v-L04}{\BF{B55r}\Hd{\LR{\LatinF \Sg{〈}35\Sg{〉}}}\Hd{ في الغربة الأسفار} وأمّا أمر الغربة فيعلم من حال النيّرين جميعًا عند الأوتاد}
\MsL{5}{AB03-F072v-L05}{لا سيّما من حال القمر وذلك أنّ القمر إذا كان غاربًا أومائلًا\PN{377} عن الأوتاد دلّ على غربة}
\MsL{6}{AB03-F072v-L06}{ونقلة من موضع إلى موضع وقد يفعل مثل ذلك المرّيخ في بعض الأوقات إذا كان}
\MsL{7}{AB03-F072v-L07}{في الغارب. وأمّا إذا مال عن موضع سمت الرؤوس يعني التاسع وكان بعده ذلك من النيّرين}
\MsL{8}{AB03-F072v-L08}{إمّا بعد القابلة وإمّا بعد التربيع فإن وقع سهم السعادة أيضًا في البرج الذي يحدث}
\MsL{9}{AB03-F072v-L09}{السفر كان جميع عمر الإنسان ونقلته وأعماله في غربة وإن نظرت السعود إلى المواضع}
\MsL{10}{AB03-F072v-L10}{التي ذكرنا أو كانت صائرة إليها تابعة لها كانت الأعمال التي تعالجها في الغربة}
\MsL{11}{AB03-F072v-L11}{نبيلة مرتجة وكان رجوعه من الغربة سريعًا لا يفوقه\PN{378} شيء وإن نظرت إليها النحوس}
\MsL{12}{AB03-F072v-L12}{كانت تلك الأعمال متعبة جالبة للضرّ\PN{379} ذوات خطر وكان رجوعه من الغربة بطيئًا.}
\MsL{13}{AB03-F072v-L13}{ويوجد\PN{380} امتزاج ذلك أبدًا بحسب ما عليه الكواكب التي تشارك هذه المواضع في الشكل}
\MsL{14}{AB03-F072v-L14}{كما بيّنّا فيما تقدّم في أكثر الأمر إذا وقع النيّران في ناحية الربعين الشرقيّين كان}
\MsL{15}{AB03-F072v-L15}{السفر في أكثر الأمر إلى المساكن التي في ناحية الشرق والجنوب وإن وقعا في الربعين}
\MsL{16}{AB03-F072v-L16}{الغربيين أو في الغارب نفسه كان السفر إلى ناحية\PN{381} الشمال والمغرب. فإن كانت البروج}
\MsL{17}{AB03-F072v-L17}{التي تحدث الغربة ذوات صورة واحدة إمّا هي في أنفسها وإمّا التي هي المتولّدة من}
\MsL{18}{AB03-F072v-L18}{الكواكب لتدبيرها كان السفر في الأحايين وفي بعض الأوقات فإن كانت ذوات}
\MsL{19}{AB03-F072v-L19}{جسدين أو ذوات صورتين كان السفر متّصلًا \{عند\}\PN{382} زمان طويل فإن كان صاحب}
\MsL{20}{AB03-F072v-L20}{المواضع التي تحدث السفر وموضعًا النيّرين الزهرة والمشتري لم يكن في السفر}
\MsL{21}{AB03-F072v-L21}{خطر وكان مع ذلك محبوبًا وذلك أنّه يكون عن توجيه ولاة البلدان إيّاهم}
\begin{PNotes}
\PNote{377}{AB03-n377}{T: \RL{\ArN أو مال}}
\PNote{378}{AB03-n378}{T: \RL{\ArN يعوقه}}
\PNote{379}{AB03-n379}{B: \RL{\ArN للضَرَر}}
\PNote{380}{AB03-n380}{BT: \RL{\ArN يؤخذ}}
\PNote{381}{AB03-n381}{T: \RL{\ArN نواحيّ}}
\PNote{382}{AB03-n382}{T: \RL{\ArN تمتدّ}}
\end{PNotes}
\Folio{073r}{73r}
\MsL{1}{AB03-F073r-L01}{\BF{B55v}بسبب حوائج أصدقائهم ويعينهم على ذلك جودة مزاج الهواء وغزارة الأشياء}
\MsL{2}{AB03-F073r-L02}{التي تصلح لهم فإن كان معها عطارد ينالهم كثيرًا من أجل السبب الذي ذكرنا منافع ونجح}
\MsL{3}{AB03-F073r-L03}{سعي وجوائز وكرامات فإن احتوى زحل والمرّيخ على النيّرين وبخاصّة متى كانا}
\MsL{4}{AB03-F073r-L04}{متقابلين في الطول في هذا الموضع وفي كلّ موضع فيه هذا القول متقابلين في الطول}
\MsL{5}{AB03-F073r-L05}{هي في مواضع واحد في الطول كان الذي ينالهم غير نافع ووقعوا في شدّة شديدة.}
\MsL{6}{AB03-F073r-L06}{فإن كانا في بروج رطبة كانت الشدّة بسبب ركوب الماء والغرق أو لسفر في طريق}
\MsL{7}{AB03-F073r-L07}{يعسر سلوكه ومواضع غير مؤهلة ولا مسكونة فإن كانا في بروج ثابتة كان ذلك بسبب}
\MsL{8}{AB03-F073r-L08}{سقوط من أجراف وهبوب رياح عاصفة وإن كانا في بروج الاعتدالين والانقلابين}
\MsL{9}{AB03-F073r-L09}{كان ذلك بسبب عوز الأشياء التي يحتاج إليها ونقصانها وحال ممرضة من أحوال الهواء}
\MsL{10}{AB03-F073r-L10}{وإن كانا في بروج صورها شبيهة بصور الناس كانت ذلك بسبب قطع طريق واغتيال}
\MsL{11}{AB03-F073r-L11}{وسلب فإن كانا في بروج صورها صور الحيوان البريء كان ذلك بسبب افتراس}
\MsL{12}{AB03-F073r-L12}{السباع أو زلازل. فإن كان هناك عطارد كان ذلك بسبب مزاج وقيعة فيها}
\MsL{13}{AB03-F073r-L13}{خطر وبسبب لسع الدوابّ والهوامّ ذوات السموم. وينبغي أن نجعل نظرنا في خواصّ}
\MsL{14}{AB03-F073r-L14}{أمر الأعراض نافعة كانت أو ضارّة من المواضع الرئيسيّة التي هي السبب للأعمال أو}
\Chap{I.36}{<36> في منيّة المولود}
\MsL{15}{AB03-F073r-L15}{للمال ولأمر البدن أو للحظوة بسبب الحال الذي تقدّم ذكرها. \Hd{\LR{\LatinF \Sg{〈}36\Sg{〉}}}\Hd{ في منيّة المولود}}
\MsL{16}{AB03-F073r-L16}{ولمّا كان الباقي بعد الذي ذكرنا من أمور المواليد الفحص عن حال المنيّة فينبغي}
\MsL{17}{AB03-F073r-L17}{أن نعلم أوّلًا الأشياء التي قلناها في العمر هل يكون القتل عن التلقّي وإلقاء الشعاع}
\MsL{18}{AB03-F073r-L18}{أوعن نزول المتولّيّ لأمر العمر \BF{B56r}إلى الوتد الغارب وذلك أنّه إن كان القتل عن القاء}
\MsL{19}{AB03-F073r-L19}{الشعاع وعن تلقّي ما يتلقّاه فيجب أن تتفقّد موضع التقاء الهيلاج والشعاع أو}
\MsL{20}{AB03-F073r-L20}{النحس حيث ما التقياه. وإن كان من قبل مصيّره إلى درجة الغارب فينبغي أن تتفقّد}
\MsL{21}{AB03-F073r-L21}{نفس موضع الغارب وذلك أنّه على حسب الكواكب التي تتلوا المواضع التي ذكرنا ن}
\Folio{073v}{73v}
\MsL{1}{AB03-F073v-L01}{كان يتلوها كوكب أو على حسب الكواكب التي تصيّر إليها قبل غيرها تكون الحال في}
\MsL{2}{AB03-F073v-L02}{المنيّة يعني إذا كان القتل عن القاء الشعاع وعن تلقّي ما يتلقّى من النحوس الكواكب التي}
\MsL{3}{AB03-F073v-L03}{تكون تالية لمواضع التلقّي وهي التي تطلع بعد تلك المواضع وتغيب من بعد أيضًا. فإن}
\MsL{4}{AB03-F073v-L04}{كان عند مصيّر الهيلاج نزوله إلى درجة الغارب الكواكب التي تغرب قبل غروب}
\MsL{5}{AB03-F073v-L05}{درجة الغارب وهي الكواكب المتقدّمة لدرجة الغارب إذا أخذ مع طبائعها في معرفة}
\MsL{6}{AB03-F073v-L06}{الأعراض واختلافها خواصّ الكواكب التي تشبه الشكل إذا كان القتل عن تلقّيها وإن}
\MsL{7}{AB03-F073v-L07}{كان عن القاء شعاعها فخاصّة حال المواضع القاتلة التي هي لها في موضع البروج من}
\MsL{8}{AB03-F073v-L08}{جهة طبائع الحدود فزحل إذا كان صاحب المنيّة كان حدوثها عن أمراض مزمنة}
\MsL{9}{AB03-F073v-L09}{وسلّ ونزلات وذوبان والحمّى التي يكون معها نافض وأوجاع الطحال والاستسقاء}
\MsL{10}{AB03-F073v-L10}{والدرف أو عن علل الأرحام وجميع العلل التي تكون عن غلبة البرد وكثرته وإذا كان}
\MsL{11}{AB03-F073v-L11}{صاحب المنيّة المشتري كان الموت عن وذبحة وذات الرئة وسكتة وتشنّج وصداع}
\MsL{12}{AB03-F073v-L12}{وعلل القلب وجميع العلل التي تكون عن ريح مفرطة أو رائحة منتنة. وإذا كان صاحبها}
\MsL{13}{AB03-F073v-L13}{المرّيخ كان ذلك عن حمّيات دائمة والحمّيات التي تسمّى \{انمرطرلطا\}\PN{383}. وهي شطر}
\MsL{14}{AB03-F073v-L14}{الغبّ وعن الضربة التي تكون بغتة وأوجاع الكلّي وعلل قذف الدم وانفجار}
\MsL{15}{AB03-F073v-L15}{وإسقاط الأجنّة والولادة والحمرة والعلل المهلكة وجميع الأمراض التي تجلب\PN{384} حمّى}
\MsL{16}{AB03-F073v-L16}{ملتهبة \BF{B56v}وإفراط المرارة وإن كان صحابها الزهرة كان ذلك عن علّة المعدة والكبد}
\MsL{17}{AB03-F073v-L17}{والفؤاد واختلاف الدم وعن الأورام الساعيّة والنواصيّة وشرب الأدويّة}
\MsL{18}{AB03-F073v-L18}{وجميع الأشياء التي تعرض فيها الأعراض عن كثرة الرطوبة أو فسادها. وإذا}
\MsL{19}{AB03-F073v-L19}{كان صاحب المنيّة عطارد كان الموت عن جنون وذهاب العقل والوسواس}
\MsL{20}{AB03-F073v-L20}{السوداويّ وسقوط وصرع وعلل السعال والقذف وجميع العلل التي تكون عمن كثرة}
\MsL{21}{AB03-F073v-L21}{اليبس أو فساده والذين يموتون موته\PN{385} أنفسهم تكون مفارقتهم للحيوة على هذه الجهات}
\begin{PNotes}
\PNote{383}{AB03-n383}{B: \RL{\ArN امطر طاوس}; T: \RL{\ArN طرناوس}}
\PNote{384}{AB03-n384}{B: add. \RL{\ArN بسبب}}
\PNote{385}{AB03-n385}{T: \RL{\ArN ميتة}}
\end{PNotes}
\Folio{074r}{74r}
\MsL{1}{AB03-F074r-L01}{المذكورة إمّا عن جميعها وإمّا عن بعضها إذا كانت الكواكب التي أمر تريد الموت إليها}
\MsL{2}{AB03-F074r-L02}{على حال الطبيعة الخاصّيّة والملائمة \PN{386}ولم يعل عليها شيء من الكواكب التي تكون المنيَّة منيّة\PN{387}}
\MsL{3}{AB03-F074r-L03}{متوء صعبة وتكون المنيّة منيّة عنف واضطراب وهي المنيّة التي تكون بالقتل أو [[يعرض]] يغيّر ذلك}
\MsL{4}{AB03-F074r-L04}{من عوارض البلاء إذا كان المستوليّ على موضع القتل المنحسان جميعًا إمّا بأن يكونا في تلك}
\MsL{5}{AB03-F074r-L05}{المواضع وإمّا بأن يكونا على تربيعها أو في مقابلتهما على القطر أو يكون أحدهما أيّهما اتّفق}
\MsL{6}{AB03-F074r-L06}{مستوليًّا على الشمس أو على القمر أو عليهما جميعًا ووجدنا في نسخة أخرى إذا كان ناظرًا}
\MsL{7}{AB03-F074r-L07}{إلى النيّرين أو أحدهما أيّهما أو المنيّة السوء تكون عن اجتماع النحسين. وأمّا عظم المنيّة فيكون}
\MsL{8}{AB03-F074r-L08}{عن شهادة النيّرين وحال المنيّة أيضًا تكون عن رؤية سائر الكواكب وعن البروج التي}
\MsL{9}{AB03-F074r-L09}{تكون فيها النحسان وذلك أنّ زحل إذا كان على تربيع الشمس أو مقابلتها وهو على}
\MsL{10}{AB03-F074r-L10}{خلاف حيّز. أمّا إذا كان في البروج الثابتة فيحدث عنه الموت بسبب رجم}
\MsL{11}{AB03-F074r-L11}{العوامّ أو عن حنق أو عن عذاب وكذلك أيضًا يكون متى كان غاريبًا\PN{388} والقمر تابعًا له وذلك}
\MsL{12}{AB03-F074r-L12}{إن كان زحل في السابع والقمر بعده في الثامن أو يغيب بعده. أمّا إذا كان في البروج}
\MsL{13}{AB03-F074r-L13}{التي صورها شبيهة بصورالحيوان فإنّ الموت يكون عن افتراس السباع فإن شهد}
\MsL{14}{AB03-F074r-L14}{المشتري يعني أن نظر وكان منحوسًا \BF{B57r}ألقى ذلك للسباع في مواضع اجتماع الناس أو في}
\MsL{15}{AB03-F074r-L15}{أيّام شهورة فإن استقبل أحد النيّرين أيّهما اتّفق ثمّ كان مقابل الطالع أو في الطالع}
\MsL{16}{AB03-F074r-L16}{كان الموت في السجون. فإن شارك عطارد في الشكل وخاصّة في موضع صورة الحياة}
\MsL{17}{AB03-F074r-L17}{التي في الكرة أو في البروج التي صورها صورة الحيوان البريء كان الموت عن}
\MsL{18}{AB03-F074r-L18}{لسع ذوات السموم ونهشها. فإن كانت معها الزهرة كان ذلك عن شرب السموم القاتلة}
\MsL{19}{AB03-F074r-L19}{ومن اغتيال النساء وإن كان الموضع القتل في برج العذراء أو في السمكتين أو في البروج الرطبة}
\MsL{20}{AB03-F074r-L20}{وشارك القمر في الشكل كان ذلك عن غرق في الماء وإن كانت في الموضع الذي يسمّى السفينة}
\MsL{21}{AB03-F074r-L21}{كان الموت عن أهوال البحر وإن كان مع الشمس في برج الانقلابين أو في البروج التي}
\begin{PNotes}
\PNote{386}{AB03-n386}{B: \RL{\ArN أو الملائمة}}
\PNote{387}{AB03-n387}{B: \RL{\ArN المنيَّة عنها منيّة}}
\PNote{388}{AB03-n388}{B: \RL{\ArN غربيًّا}}
\end{PNotes}
\Folio{074v}{74v}
\MsL{1}{AB03-F074v-L01}{صورها صور ذوات الأربع قوائم أو قابل المرّيخ مكان الشمس كان بسبب يقع عليه}
\MsL{2}{AB03-F074v-L02}{وفي نسخة أخرى أو قابل المرّيخ مكان الشمس أو القمر أو كان على تربيعهما في غير حيّزهما فإن}
\MsL{3}{AB03-F074v-L03}{كان في وسط السماء أو في وتد الأرض كان الموت عن سقطة من موضع عال مثال موضع}
\MsL{4}{AB03-F074v-L04}{مرتفع من الأرض أو مثل النيّر و\{الجب\} وما يشبهه إذا كانت في وتد الأرض. وأمّا المرّيخ فإنّه}
\MsL{5}{AB03-F074v-L05}{إن كان على تربيع أو مقابلة الشمس والقمر\PN{389} على خلاف أمر الحيّز أمّا إذا كان في البروج التي صورها}
\MsL{6}{AB03-F074v-L06}{على صور الناس كان ذلك عن ذبح يعرض وفي وشغب يقع بين قوم وأهل قبيلتهم أو عن حروب}
\MsL{7}{AB03-F074v-L07}{أو عن قتل للإنسان نفسه بأسباب النساء وعن قتل النساء وذلك يكون إذا شهدت الزهرة}
\MsL{8}{AB03-F074v-L08}{فإن كان عطارد مشاركًا له في الشكل كان الهلاك على أيدي اللصوص وقطّاع الطريق في البحر والدعار\LR{\Sg{〉}}ة\LR{\Sg{〈}}.}
\MsL{9}{AB03-F074v-L09}{فإن كان في البروج المقطوعة\PN{390} الأعضاء التي ليست بتامّة الصور أو كانت مع رأس الغول الذي يسمّى}
\MsL{10}{AB03-F074v-L10}{برسياس يعني الفارس الممسك رأس الغول وفي نسخة أخرى في البروج المقطوعة الأعضاء أو كان}
\MsL{11}{AB03-F074v-L11}{طلوعها مع رأس الغول كان الموت بضرب العنق أو بتقطيع الأعضاء. فإن كان في العقرب}
\MsL{12}{AB03-F074v-L12}{والثور كان الموت بالحريق أو بالقطع \BF{B57v}أو \{ببطء\}\PN{391} الأطبّاء أو تشنّج فإن كان في وسط السماء أو في}
\MsL{13}{AB03-F074v-L13}{ما يقابل وسط السماء يعني وتد الأرض كان الموت بالصلب وبخاصّة إذا كان مع \{فيغاوس\}\PN{392}}
\MsL{14}{AB03-F074v-L14}{وهو الملتهب واندرهميس فإن كان في البرج الغارب أو الطالع أحرق الإنسان حيّا}
\MsL{15}{AB03-F074v-L15}{فإن كان في بروج ذوات أربع قوائم كان ذلك عن سقوط وتكسّر وسحب فإن شهد له المشتري}
\MsL{16}{AB03-F074v-L16}{وكان منحوسًا كان الهلاك بوجوب حكم مشهور وغضب السلاطين والملوك. فإن اجتمع}
\MsL{17}{AB03-F074v-L17}{النحسان جميعًا أو قابل أحدهما الآخر على حال من الأحوال التي ذكرنا أنّها سبب الموت كانا}
\MsL{18}{AB03-F074v-L18}{معيّنين خاصّة على المايّة\PN{393} السوء وصاحب الأمر في الحال\PN{394} الموت وهو صاحب الأمر القاتل وتكون}
\MsL{19}{AB03-F074v-L19}{أعراض المنيّة مضعّفة مثنى\PN{395} فإمّا في الكيفيّة وإمّا في الكمّية إذا كان للنحسين حصص في}
\MsL{20}{AB03-F074v-L20}{المواضع القاتلة ومن كانت الكواكب فيه هذه الحال لم يدفن ولم يوارى. فأمّا السباع ذوات المخالب}
\MsL{21}{AB03-F074v-L21}{من الطير فإنّها تأكّله\PN{396} إذا كان النحسان في بروج صورها صور هذه الأشياء التي ذكرنا ولم}
\MsL{22}{AB03-F074v-L22}{يشهد للمواضع القاتلة شيء من السعود وهو تحت الأرض يعني السعد ويكون الموت في الغربة}
\MsL{23}{AB03-F074v-L23}{إذا وقعت الكواكب المستوليّة على المواضع القاتلة في المواضع القاتلة وخاصّة إذا كان القمر في المواضع التي ذكرنا وكان على تربيعهما أو مقابلًا لها.}
\begin{PNotes}
\PNote{389}{AB03-n389}{B: \RL{\ArN أو القمر}}
\PNote{390}{AB03-n390}{T: \RL{\ArN منقطعة }}
\PNote{391}{AB03-n391}{T: \RL{\ArN بنظر}}
\PNote{392}{AB03-n392}{T: \RL{\ArN قيفوس}}
\PNote{393}{AB03-n393}{B: \RL{\ArN المنيّة}}
\PNote{394}{AB03-n394}{B: \RL{\ArN حال}}
\PNote{395}{AB03-n395}{T: \RL{\ArN مثنّاة}}
\PNote{396}{AB03-n396}{or: \RL{\ArN تآكل}}
\end{PNotes}
\Folio{075r}{75r}
\Chap{I.37}{<37> في قسمة الأزمان التي للمولود}
\MsL{1}{AB03-F075r-L01}{\Hd{\LR{\LatinF \Sg{〈}37\Sg{〉}}}\Hd{ في قسمة الأزمان التي للمولود} أمّا إذ قد انتهينا إلى هذا الموضع على طريق الجملة في كلّ واحد من هذه الأنواع فقد يجب أن}
\MsL{2}{AB03-F075r-L02}{نضيف إلى تلك الحال بعينها جميع الأشياء التي يجب النظر فيها من قسمة أزمان العمر على المذهب}
\MsL{3}{AB03-F075r-L03}{الطبيعيّ وما يليق بالمذاهب التي ذكرناها في الأمور الجزئيّة وكما أنّه قد تقدّم\PN{397} قضاء تامّ وحتّم}
\MsL{4}{AB03-F075r-L04}{\BF{B58r}على الأمور الجزئيّة مأخوذ من المواضع التي قد ذكرت في القول في أمر المواليد على أنّ أعظم منها}
\MsL{5}{AB03-F075r-L05}{والحتم الذي يأخذ من البلدان أنفسها على ما بيّنا في الكتاب وحال الهويّ والأسباب الكلّيّة}
\MsL{6}{AB03-F075r-L06}{وهي قضايًّا تدخل فيها الأشياء الجزائيّة التي ينظر فيها من أمر المواليد على جهة الأمر المجمل مثل}
\MsL{7}{AB03-F075r-L07}{صورة الأبدان وحليتها وخواصّ حالات الأنفس واختلاف العادات والسنين. وكان}
\MsL{8}{AB03-F075r-L08}{واجبًا على من نظر في مثل هذه الأشياء من المذهب الطبيعيّ أوّلًا بسبب الأوّل الذي هو}
\MsL{9}{AB03-F075r-L09}{أوّل ولحقّ من غيره ليلًا يغلط أحد بسبب اشتباه المواليد. فنقول مثالًا أنّ المولود في}
\MsL{10}{AB03-F075r-L10}{بلاد الحبشة يكون أبيض اللون سبط\PN{398} الشعر وإنّ المولود في بلاد جرمانيّا أو في بلاد علاطبًا}
\MsL{11}{AB03-F075r-L11}{يكون آدم اللون جعد الشعر وإن هؤلاء إنث\PN{399} أخلاقًا أو محبّين للقياس والنظر في الأمور}
\MsL{12}{AB03-F075r-L12}{وإنّ الذين في بلاد الآفاس\PN{400} أنفسهم وحشية وهم غير متأدبين أو يغلط في ذلك أيضًا}
\MsL{13}{AB03-F075r-L13}{بسبب السنين والعادات التي تخصّ بعض الأمور فنقول في أمر التأهّل والتزويج}
\MsL{14}{AB03-F075r-L14}{نقول مثلًا أنّ الرجل من أهل أنطاليّا يتزوّج بأخته وكان الواجب أن يزيد في قوله}
\MsL{15}{AB03-F075r-L15}{أنّه يكون ذلك إذا كان المولود مصريًّا أو يقول في المصريّ إنّه يتزوّج بأمّه وكان الواجب}
\MsL{16}{AB03-F075r-L16}{أن يزيد في قوله إذا كان المولود فارسيًّا وبالجملة. فينبغي أن نأخذ أوّلًا للحالات}
\MsL{17}{AB03-F075r-L17}{الكلّيّة في الحتم والقضيّة التي تقضيها ثمّ نصيّر إلى تلك الحالات الجزئيّة التي}
\MsL{18}{AB03-F075r-L18}{نقضي بها ليتبيّن من ذلك الأمر في الزيادات والنقصانات التي تعرض في ذلك}
\begin{PNotes}
\PNote{397}{AB03-n397}{B: \RL{\ArN تقدّمه}}
\PNote{398}{AB03-n398}{T: \RL{\ArN الوسط}}
\PNote{399}{AB03-n399}{B: \RL{\ArN انس}}
\PNote{400}{AB03-n400}{B: \RL{\ArN الاس}}
\end{PNotes}
\Folio{075v}{75v}
\MsL{1}{AB03-F075v-L01}{وكذلك أيضًا يجب صورة في قسمة أزمان المولود وأوقاته أن تعلم\PN{401} أوّلًا من اختلاف}
\MsL{2}{AB03-F075v-L02}{الأسنان الزمانيّة وموافقها لكلّ واحد من الأحوال\PN{402} التي تحدث وأن تتفقّد أمرها}
\MsL{3}{AB03-F075v-L03}{ليلًا يغلط في وقت من الأوقات في الأعراض العامّيّة والبسيطة التي ينظر فيها في المواليد}
\MsL{4}{AB03-F075v-L04}{فنقول إنّ الطفل يباشر الأعمال أو يتزوّج أو يفعل شيئًا \BF{B58v}من الأشياء التي إنّما يفعلها من هو}
\MsL{5}{AB03-F075v-L05}{أسنّ منه وإنّ الشيخ الغاني ولد كذا أو يفعل شيئًا من الأفعال التي إنّما يفعلها الأحداث لكن}
\MsL{6}{AB03-F075v-L06}{ينبغي بالجملة أن تضيف الأشياء التي تظهر لنا من النظر في أمر الأزمان التي تشاكلها ويمكن}
\MsL{7}{AB03-F075v-L07}{أن يحدث فيها من الأسنان الموافقة لذلك وذلك أنّ ها هنا مسلكًا ما ومذهبًا لجميع}
\MsL{8}{AB03-F075v-L08}{الاختلافات الزمانيّة ومأخذًا واحدًا بعينه لأنّ ما للطبيعة الكلّيّة التي تعمّ جميع الناس}
\MsL{9}{AB03-F075v-L09}{مشبهًا ونظير لترتيب الكواكب السبعة المتحيّرة وهذا المسلك والمأخذ تبتدئ}
\MsL{10}{AB03-F075v-L10}{فيها من السنّ الأوّل ومن الكرة التي هي أقرب إلينا أعني أو كرة القمر وتنتهي إلى آخر الأسنان}
\MsL{11}{AB03-F075v-L11}{وإلى آخر الكواكب\PN{403} التي تسمّى كرة زحل. ويعرض بالحقيقة في كلّ واحد الأسنان الأمر الملائم}
\MsL{12}{AB03-F075v-L12}{لطبيعة الكواكب المتحيّرة الذي صار نظيرًا لتلك السنّ وذلك شيء ينبغي أن نفتقّده}
\MsL{13}{AB03-F075v-L13}{لنعلم منه أمر كلّ واحد من هذه الأزمان بكلّيّته. وأمّا اختلافاتها الجزئيّة فإنّها تعلم}
\MsL{14}{AB03-F075v-L14}{من الخواصّ التي توجد وذلك أنّ القمر لمّا كان صاحب تدبير سنّ الطفوليّة وهو الأربع}
\MsL{15}{AB03-F075v-L15}{سنون الأولى في أكثر الأمر التي هي العدد الذي يسمّى ذلك الأربع سنين صار بدن الطفل}
\MsL{16}{AB03-F075v-L16}{فيها رطبًا غير مستحكم سريع النمو وصار غذاؤه في أكثر الأمر مائيًا وصارت سجيته سهلة}
\MsL{17}{AB03-F075v-L17}{التغيّر وصارت نفسه غير مستحكمة ولا قويّة على ما يليق بالأعراض التي تعرض من قبل}
\MsL{18}{AB03-F075v-L18}{الجزء والفاعل من \{نفس\} الطفل. ولمّا كان صاحب تدبير العشر سنين التي بعد هذه}
\MsL{19}{AB03-F075v-L19}{التي من الصبًّا وهي السنّ الثانية للكوكب الثاني وهو عطارد إذا كان هو صاحب نصف}
\MsL{20}{AB03-F075v-L20}{عدد العشرين سنة صار المولود مبتدئًا في هذه السنين بأن يقوى فيها فهم النفس}
\MsL{21}{AB03-F075v-L21}{ونطقها وكأنّما يحمل ويغرس فيه غروس التعاليم وأصول الأخلاق ويتبيّن فيها}
\begin{PNotes}
\PNote{401}{AB03-n401}{B: \RL{\ArN نعلم}}
\PNote{402}{AB03-n402}{T: \RL{\ArN الأحداث}}
\PNote{403}{AB03-n403}{B: add. \RL{\ArN المتحيّرة}}
\end{PNotes}
\Folio{076r}{76r}
\MsL{1}{AB03-F076r-L01}{من خواصّ الأعمال التي يعملها \BF{B59r}التعلم والأدب عن أصناف الرياضة}
\MsL{2}{AB03-F076r-L02}{الأولى. ولمّا كان صاحب تدبير اليقين\PN{404} الثالثة وهي من\PN{405} الفتيان}
\MsL{3}{AB03-F076r-L03}{في الثمان السنين التي تأتي بعد السنين التي ذكرناها للزهرة}
\MsL{4}{AB03-F076r-L04}{على مثال العدد المساويّ يسمّى دورها الخاصّ لمّا ابتدأت بسبب}
\MsL{5}{AB03-F076r-L05}{ذلك حركة مجاريّ المنيّ بامتلائها وتحرّكت إلى أمور الجماع وهذا}
\MsL{6}{AB03-F076r-L06}{هو الوقت الذي يكون فيه خاصّة في الأنفس تشبيهه بالكلب}
\MsL{7}{AB03-F076r-L07}{وقوّة\PN{406} الضبط والعشق ومحبّة لكلّ ما تهيّأ من أمور الجماع}
\MsL{8}{AB03-F076r-L08}{وهيجان وشبق والانخداع والعمّيّ\PN{407} فيما نسخ له منها. ولمّا كان}
\MsL{9}{AB03-F076r-L09}{صاحب السنّ الرابعة وهي السنّ الوسطى في المرتبة التي هي}
\MsL{10}{AB03-F076r-L10}{سنّ الأحداث صاحب الكرة الوسطى أعني الشمس سنين عددها}
\MsL{11}{AB03-F076r-L11}{مثل عدد التسع عشرة سنة التي هي دورها الخاصّ في الاجتماعات}
\MsL{12}{AB03-F076r-L12}{صيّر النفس مستوليّة على الأعمال قادرة عليها تابعة إلى تدبير}
\MsL{13}{AB03-F076r-L13}{المعاش والنيل\PN{408} والاستقامة والانتقال من أحوال الهزل}
\MsL{14}{AB03-F076r-L14}{واللعب وكثيرة الخطأ إلى صيانة النفس والوقار ومحبّة الكرامة}
\MsL{15}{AB03-F076r-L15}{ثمّ من بعد الشمس يتولّي المرّيخ السنّ الخامسة وهي السنّ التي}
\MsL{16}{AB03-F076r-L16}{تكون فيها الرحلة خمسة عشر سنة على عدد سنين دورها الخاصّ}
\MsL{17}{AB03-F076r-L17}{له. فيحدث صعوبة المعاش وبلايا وآفات ويحدث\PN{409} في النفس}
\MsL{18}{AB03-F076r-L18}{والبدن غمومًا وكأنّه يحسن ويفهم أمر انحطاط البدن ويدعوا}
\MsL{19}{AB03-F076r-L19}{الإنسان إلى أن يعمل شيئًا من الأعمال التي لها قدر يعتز به\PN{410} ويناله}
\MsL{20}{AB03-F076r-L20}{منها التعب ويحرص على أن يفعل ذلك من قبل أن يقرب من الموت}
\begin{PNotes}
\PNote{404}{AB03-n404}{B: \RL{\ArN اليقين} corrected to \RL{\ArN السنّ}; T: \RL{\ArN السنّ}}
\PNote{405}{AB03-n405}{T: \RL{\ArN سنّ}}
\PNote{406}{AB03-n406}{T: \RL{\ArN قلّة}}
\PNote{407}{AB03-n407}{T: \RL{\ArN والعناء}}
\PNote{408}{AB03-n408}{T: \RL{\ArN الميل}}
\PNote{409}{AB03-n409}{T: \RL{\ArN ويسنّ}}
\PNote{410}{AB03-n410}{T: \RL{\ArN يعتدّ به}}
\end{PNotes}
\Folio{076v}{76v}
\MsL{1}{AB03-F076v-L01}{ثمّ السادس من الكواكب وهو المشتري له سنّ الكهول وهي إثنتا عشرة سنة بعدد}
\MsL{2}{AB03-F076v-L02}{سني دوره الخاصّ وحينئذ ينصرف الإنسان عن مباشرة الأعمال بنفسه}
\MsL{3}{AB03-F076v-L03}{والكدّ والاضطراب وركوب الخطر ويلزم مكان ذلك حسن المذهب وتقدّم\PN{411} العنايّة}
\MsL{4}{AB03-F076v-L04}{بالأمور ينزه نفسه عن القبيح ويتفكّر في الأمور ويعظ ويعزى ويصيّر في هذا الوقت خاصّة}
\MsL{5}{AB03-F076v-L05}{معيّنًا باكتساب الكرامة والحرّيّة مع الحياء والأمر الجميل ويكون \BF{B59v}بآخره لزحل السنّ الأخيرة وهي}
\MsL{6}{AB03-F076v-L06}{سنّ الشيخوخة إلى آخر العمر وهي السنّ التي يعرض فيها البرد وعسر حركات البدن والنفس}
\MsL{7}{AB03-F076v-L07}{والنهوض والشهوات والتنعّم وسرعة الانحطاط الذي يعرض بالطبيعة وهي الحال التي}
\MsL{8}{AB03-F076v-L08}{ينتهي فيها عمر الإنسان ويقع عليه الكتابة\PN{412} والضعف والآفات وعسر الاحتمال الجميع الأمور}
\MsL{9}{AB03-F076v-L09}{على حسب فتور جميع حركاته. فهذه صفة خواصّ أحوال الأزمان التي ترى على حسب}
\MsL{10}{AB03-F076v-L10}{الأمر الكلّيّ العامّ من أمر الطبيعة. ويدي فيها بالقمر فإعطأه أقلّ الأجزاء وهو العشرون}
\MsL{11}{AB03-F076v-L11}{دورة الأصغر ودور الشمس الأصغر الذي عرف لهما باجتماع في العودات. ثمّ يجاوز ذلك}
\MsL{12}{AB03-F076v-L12}{إلى أعظم الأجزاء الذي هو النصف فأخذ من دور عطارد النصف ثمّ انتهى إلى الأجزاء التامّة}
\MsL{13}{AB03-F076v-L13}{فاستعمل دور كلّ كوكب كما هو فصارت جميع هذه الأسنان ثمان وتسعون سنة فإن}
\MsL{14}{AB03-F076v-L14}{تجاوزها المولود عاش إلى سنّ القمر. فإن تجاوز عاد إلى سنّ عطارد فإن تجاوز عاد إلى}
\MsL{15}{AB03-F076v-L15}{سنّ الزهرة وهو تمام مائة وعشرون سنة التي هي سنّ الشمس الكبرى ولذلك ما \{جعلتم\}\PN{413}}
\MsL{16}{AB03-F076v-L16}{سنيها على هذا المقدار فأمّا صفة الخواصّ الجزئيّة التي للأزمان فإنّها لمّا كان ينبغي أن}
\MsL{17}{AB03-F076v-L17}{تحسب خواصّ المواليد جعلنا أخذنا لمّا كان منها على الجهة الكلّيّة التي ذكرنا من الهيلاجات}
\MsL{18}{AB03-F076v-L18}{الرئيسيّة إلّا أنّا نجعل أخذنا لذلك من جميعها لأمر واحد منها كما جعلناه\PN{414} في مدّة القمر\LR{\LatinF .}}
\MsL{19}{AB03-F076v-L19}{فأمّا الهيلاج الذي يأخذ من الطالع فإنّا نستعمله في الأعراض التي تعرض في البدن\PN{415}}
\MsL{20}{AB03-F076v-L20}{والغربة. فأمّا الذي يؤخذ من سهم البخت ففي أمور المال والقنيّة وأمّا الذي}
\begin{PNotes}
\PNote{411}{AB03-n411}{B: \RL{\ArN يقدّم}}
\PNote{412}{AB03-n412}{B: \RL{\ArN الكأبة}}
\PNote{413}{AB03-n413}{B: \RL{\ArN جعلت}}
\PNote{414}{AB03-n414}{T: \RL{\ArN فعلنا}}
\PNote{415}{AB03-n415}{T: \RL{\ArN للبلدن}}
\end{PNotes}
\Folio{077r}{77r}
\MsL{1}{AB03-F077r-L01}{يؤخذ من القمر ففي أحوال النفس والتأهّل. وأمّا الذي يؤخذ من الشمس}
\MsL{2}{AB03-F077r-L02}{ففي أمور الحظوة والشرف وأمّا الذي يؤخذ من وسط السماء ففي سائر حالات}
\MsL{3}{AB03-F077r-L03}{التدبير الجزئيّة مثل الأعمال والمصادقة وأمر الأولاد\LR{\LatinF .} فإنّا إذا فعلنا\PN{416} كذلك}
\MsL{4}{AB03-F077r-L04}{لم يكن صاحب الأمر كوكبًا واحدًا إمّا سعدًا وإمّا نحسًا وهذا هو الصواب}
\MsL{5}{AB03-F077r-L05}{إذ كانت \BF{B60r}قد تعرض كثيرًا في زمان واحد بعينه أعراض متضادّة كما أنّه}
\MsL{6}{AB03-F077r-L06}{قد يصاب الإنسان ببعض قراباته ويرث منه مالًا أو يصيبه مرض ويناله}
\MsL{7}{AB03-F077r-L07}{في تلك الحال من مرضه خطوة وفائدة أو شرف ويكون متعطّلًا من العمل}
\MsL{8}{AB03-F077r-L08}{وتولد\PN{417} له أولاد وما أشبه ذلك من الأعراض التي تعرض كثيرًا\LR{\LatinF . }وذلك أنّه}
\MsL{9}{AB03-F077r-L09}{ليس الحال في البدن والنفس والمال والحظوة والأعراض الجيّدة والرديئة}
\MsL{10}{AB03-F077r-L10}{حالًا واحدة حتّى يجب ضرورة أن ينال الإنسان في جميع هذه الأشياء}
\MsL{11}{AB03-F077r-L11}{النجح أو الخيبة لكنّ لعلّه أن يكون ممكنًا أن يعرض له شيء يشبه ذلك في}
\MsL{12}{AB03-F077r-L12}{الأوقات التي هي غاية السعادة والشقاوة إذا اتّفق أن يجتمع في جميع هذه}
\MsL{13}{AB03-F077r-L13}{الهيلاجات أو في أكثرها أن يلقاه\PN{418} السعود كلّها أو النحوس وكلّما\PN{419} يعرض}
\MsL{14}{AB03-F077r-L14}{ذلك من قبل تعصير [[ال]]طبيعة الإنسان عن بلوغ كلّ واحدة من هاتين}
\MsL{15}{AB03-F077r-L15}{الغايتين أو من قبل سرعة ميلها وخروجها عن الاعتدال بسبب تغيّر الأشياء}
\MsL{16}{AB03-F077r-L16}{الصالحة والرديئة فعلى هذه الجهة التي ذكرنا نميّز ونفصّل أمر مواضع}
\MsL{17}{AB03-F077r-L17}{الهيلاجات. وأمّا المواضع التي تتلقّى الهيلاجات فليس إنّما ينبغي أن نستعمل منها}
\MsL{18}{AB03-F077r-L18}{غير المواضع التي تدلّ على القتل فقط كما يفعل\PN{420} في مدّة العمر بل ينبغي أن}
\MsL{19}{AB03-F077r-L19}{نستعملها كلّها جملة ولا ينبغي أيضًا أن نقتصر منها على ما كان نلقّيه تلقاء جسمانيًّا}
\MsL{20}{AB03-F077r-L20}{أو على مقابلة أو على تربيع فقط لكنّا نستعمل أيضًا ما كان يلقّيه من التسديس}
\MsL{21}{AB03-F077r-L21}{والتثليث والمناظرة المستويّة في القوّة والعالية الإمرة والمنحفظة}
\begin{PNotes}
\PNote{416}{AB03-n416}{T: \RL{\ArN عملنا}}
\PNote{417}{AB03-n417}{B: \RL{\ArN يولد}}
\PNote{418}{AB03-n418}{B: \RL{\ArN يلقاها}}
\PNote{419}{AB03-n419}{B: \RL{\ArN كلّها وقلّ ما}}
\PNote{420}{AB03-n420}{B: \RL{\ArN نفعل}}
\end{PNotes}
\Folio{077v}{77v}
\MsL{1}{AB03-F077v-L01}{ونعطي الأزمان أوّلًا في كلّ واحد من الهيلاجات الكوكب الذي هو في درجة الهيلاج نفسها}
\MsL{2}{AB03-F077v-L02}{أو مشارك لها في الشكل. فإن لم يكن في تلك الدرجة كوكب ولا شاركها في الشكل فينبغي}
\MsL{3}{AB03-F077v-L03}{أن نعطي بدله أقرب ما يتقدّم تلك الدرجة من الكواكب أو من الشعاع الذي يلقّيه}
\MsL{4}{AB03-F077v-L04}{أحد الكواكب من بعض الأشكال إلى أن \{ينتهي\} إلى الكوكب الذي يرى الدرجة التي تتلوا}
\MsL{5}{AB03-F077v-L05}{بعد تلك الدرجة على الجهة التي تتلوا من البروج ثمّ \{يقتر\}\PN{421} الزمان الذي بعده الكوكب}
\MsL{6}{AB03-F077v-L06}{والذي ينتهي إليه وإلى\PN{422} شعاعه وكذلك تفعل في سائر الهيلاجات وتجعل\PN{423} بعض ما تستعمله}
\MsL{7}{AB03-F077v-L07}{من الكواكب ونقيمه مقام \BF{B60v}المدبّر مع الكواكب\PN{424} الذي له قسمة الأزمان أزمان الكواكب التي}
\MsL{8}{AB03-F077v-L08}{هي أصحاب حدود الأجزاء إلى أن ينتهي\PN{425} إليها في السنين وينبغي أن نعطي درج الأبعاد}
\MsL{9}{AB03-F077v-L09}{درجات من السنين. أمّا ما في الهيلاج الذي يأخذ من الطالع بعدد من عدد السنين مساويًّا}
\MsL{10}{AB03-F077v-L10}{لعدد أزمان مطالع الأقاليم التي تتّفق بين درجة الطالع والكوكب الذي يتلوا درجة}
\MsL{11}{AB03-F077v-L11}{الطالع أو شعاعه أعني الإقليم والبلد الذي ولد فيه المولود. وأمّا الهيلاج الذي}
\MsL{12}{AB03-F077v-L12}{يؤخذ من وسط السماء بسنين عددها مساوٍ لعدد أزمان مطالع البروج في}
\MsL{13}{AB03-F077v-L13}{الفلك المستقيم والأمر في الهيلاجات الباقيّة التي تؤخذ من المواضع الباقيّة تجري}
\MsL{14}{AB03-F077v-L14}{على هذا القياس وعلى حسب قربها من الأوتاد من المشارق أو المغارب وفي}
\MsL{15}{AB03-F077v-L15}{ممرّها وسط السماء على ما بيّنا ورسمنا في معرفة مطالع االبروج فيما بين أرباع}
\MsL{16}{AB03-F077v-L16}{الفلك في صدر هذا الكتاب وعلى هذه الأشياء التي ذكرنا تعلم الأشياء المتولّيّة}
\MsL{17}{AB03-F077v-L17}{لأمر الزمان الكلّيّ وذلك أن بطليموس إنّما جعل صاحب الأمر صاحب التلقّي وأقام}
\MsL{18}{AB03-F077v-L18}{صاحب الحدّ مقام المدبّر وهو دون صاحب التلقّي. وهذا هو الأصواب إذا كانت}
\MsL{19}{AB03-F077v-L19}{شعاعات الكواكب أقوى من الحدود وكانت مع ذلك أطول زمانًا من أزمان}
\MsL{20}{AB03-F077v-L20}{الحدّ. فكما أنّه قد يعرض مثلًا أن يقع السنون في شعاع كوكب سعد من أحد الهياليج}
\MsL{21}{AB03-F077v-L21}{حدود كثيرة إلى أن يتلقّاه شعاع آخر فتكون حاله بتلك المدّة إمّا محمودة}
\begin{PNotes}
\PNote{421}{AB03-n421}{B: \RL{\ArN ىفطر}; T: \RL{\ArN نعطي}}
\PNote{422}{AB03-n422}{B: \RL{\ArN أو إلى}}
\PNote{423}{AB03-n423}{B: \RL{\ArN نجعل}}
\PNote{424}{AB03-n424}{B: \RL{\ArN الكوكب}}
\PNote{425}{AB03-n425}{B: \RL{\ArN تنتهي}}
\end{PNotes}
\Folio{078r}{78r}
\MsL{1}{AB03-F078r-L01}{في ذلك الفن دائمًا وإمّا مذمومة بالجملة\PN{426} وأمّا المتولّيّة لأمر السنين فإنّا نعلمها}
\MsL{2}{AB03-F078r-L02}{بأن نأخذ عدد السنين التامّة التي مضت من وقت المولد ونلقيها من كلّ واحد من}
\MsL{3}{AB03-F078r-L03}{المواضع الهيلاجات من البروج إلى الجهة التي تتلوا من البروج لكلّ برج منه سنة}
\MsL{4}{AB03-F078r-L04}{شمسيّة ونأخذ مدبّر البرج الذي ينتهي إليه العدد في السنة التي تتحوّل للمولود}
\MsL{5}{AB03-F078r-L05}{وتنفعل في الشهور أيضًا مثل هذا الفعل بعينه فتأخذ عدد الشهور من شهر المولود\PN{427}}
\MsL{6}{AB03-F078r-L06}{الذي ينتهي إليه وكذلك نفعل في الأيّام أيضًا. فنأخذ عدد الأيّام من يوم المولد}
\MsL{7}{AB03-F078r-L07}{والتحويل من الشهر والساعة من اليوم فنلقيه من موضع الشهور لكلّ برج يومين}
\MsL{8}{AB03-F078r-L08}{\BF{B61r}وسدس بالتقريب وإنّما نفعل ذلك لتكون السنة ثلاثة عشر شهرًا فيعود الشهر الثالث}
\MsL{9}{AB03-F078r-L09}{إلى البرج الذي كان إليه تدبير السنة وكذلك يصيّر الشهر الواحد من هذه الشهر}
\MsL{10}{AB03-F078r-L10}{ثلاثة عشر برجًا فينتهي آخر الشهر إلى البرج المتولّيّ لتدبير الشهر على ما يليق وشاكل}
\MsL{11}{AB03-F078r-L11}{هذا المذهب وسنضع جدولًا مفردًا لعدد أيّام الشهور التحويليّة والأيّام الخاصّيّة}
\MsL{12}{AB03-F078r-L12}{منها وحصّة كلّ بروج من\PN{428} البروج بعد فراغنا من هذا القول. وينبغي أن ننظر أوّلًا}
\MsL{13}{AB03-F078r-L13}{في حلول الكواكب في مواضع الأزمنة إذ كان ذلك معيّنًا على معرفة الأحداث}
\MsL{14}{AB03-F078r-L14}{التي تحدث في الأزمنة معونة ليست باليسيرة فتنظر في حلول المشتري وزحل والبروج}
\MsL{15}{AB03-F078r-L15}{المدبّرة للسنين وفي حلول المرّيخ والزهرة والشمس وعطارد في البروج المدبّرة والشهور}
\MsL{16}{AB03-F078r-L16}{وننظر في حلول القمر بمواضع الأيّام من البروج ونستعمل أصحاب الأزمان}
\MsL{17}{AB03-F078r-L17}{الكلّيّة من الكواكب في كون الحادث وتمامه وهي أصحاب الأسنان والتلقّي}
\MsL{18}{AB03-F078r-L18}{ونستعمل أصحاب الأزمان الجزئيّة التي تعدد تلك في المعونة على ذلك}
\MsL{19}{AB03-F078r-L19}{أو خلاف المعونة على حسب ملائمة طبائع الأشياء التي نجد في المواضع}
\MsL{20}{AB03-F078r-L20}{فتكون سببًا للأعراض القويّة والضعيفة يعني التي تكون سببًا}
\MsL{21}{AB03-F078r-L21}{لقوّة الأعراض وضعفها أو بعدها عن الملائمة وذلك أن الذي يدلّ على}
\begin{PNotes}
\PNote{426}{AB03-n426}{T: \RL{\ArN بالحمل}}
\PNote{427}{AB03-n427}{B: add. \RL{\ArN ونلقيه من المواضع التي تدبير السنة إليها من ذلك البرج فنأخذ مدبّر البرج}}
\PNote{428}{AB03-n428}{B: om.}
\end{PNotes}
\Folio{078v}{78v}
\MsL{1}{AB03-F078v-L01}{خاصّة الحال الكلّيّة وعلى مدّة الزمان وهو موضع الهيلاج وصاحب}
\MsL{2}{AB03-F078v-L02}{الأزمان الكلّيّة مع أصحاب الحدود لأنّ كلّ واحد من هذه الكواكب قد}
\MsL{3}{AB03-F078v-L03}{شاكل في نفس الموانع مواضع الكواكب\PN{429} الذي كان أمر النيّرين إليهما}
\MsL{4}{AB03-F078v-L04}{منذ أوّل الأمر. فأمّا العرض هل هو صالح أم ضدّ ذلك فيعلم من خاصّة}
\MsL{5}{AB03-F078v-L05}{حال طبيعيّة أصحاب الزمان ومزاجها في فعلها للخير أو الشرّ وما}
\MsL{6}{AB03-F078v-L06}{شاكلتها منذ أوّل الأمر للشيء أو مخالفتها لها.\PN{430} وأمّا في أيّ الأزمان}
\MsL{7}{AB03-F078v-L07}{خاصّة \BF{B61v}ثلثين ويظهر الحادث فيعلم من حال البروج المتولّيّة لأمر السنة}
\MsL{8}{AB03-F078v-L08}{والشهر إذا قيست بالواضع التي هي أسباب الحادث ومن الحال التي تكون}
\MsL{9}{AB03-F078v-L09}{عند حلول الكواكب والشمس في البروج المتولّيّة لأمر السنين}
\MsL{10}{AB03-F078v-L10}{والشهور وذلك أن الكواكب التي توافق المواضع التي ذكرتا منذ}
\MsL{11}{AB03-F078v-L11}{وقت المولود ويشاركها في الشكل عند الحلول مشاركة توافقه فهي}
\MsL{12}{AB03-F078v-L12}{فاعلة للخير في الأمر المقصود كما إنّها إذا خالفتها كانت فاعلة}
\MsL{13}{AB03-F078v-L13}{للشرّ والمشاركة الموافقة هي من التثليث أو التسديس وأمّا}
\MsL{14}{AB03-F078v-L14}{التي تشاركها مشاركة غير موافقة بل\PN{431} نقابلها على القطر وتكون\PN{432} على}
\MsL{15}{AB03-F078v-L15}{تربيعها وتكون على خلاف أمر الخير فهي سبب الشرّ فأنّا في سائر}
\MsL{16}{AB03-F078v-L16}{الأحوال فليس تحدث شيئًا فإن استولت على الأزمان وعلى الحلول}
\MsL{17}{AB03-F078v-L17}{كواكب واحدة بأعيانها أفرطت طبيعة الحادث واشتدت إلى الخير كان}
\MsL{18}{AB03-F078v-L18}{ميلها أو إلى السرّ وأخرى أن يكون ذلك متى لم يكن استيلاؤها}
\MsL{19}{AB03-F078v-L19}{على السبب من جهة ما هي أصحاب الأزمان فقط لكن من جهة}
\MsL{20}{AB03-F078v-L20}{أنّها أيضًا كانت المدبّرة في أمر المولود وقد يكون المولود منجحًا أو}
\MsL{21}{AB03-F078v-L21}{خائبًا في جميع الأشياء معًا. أمّا إذا كان جميع أنواع الهيلاجات أو أكثرها}
\begin{PNotes}
\PNote{429}{AB03-n429}{B: \RL{\ArN الكوكب}}
\PNote{430}{AB03-n430}{B: \RL{\ArN له}}
\PNote{431}{AB03-n431}{B: \RL{\ArN بأن}}
\PNote{432}{AB03-n432}{B: \RL{\ArN أو تكون}}
\end{PNotes}
\Folio{079r}{79r}
\MsL{1}{AB03-F079r-L01}{في موضع واحد بعينه وأمّا إذا كانت متفرّقة وكان الزمان الواحد بعينه فإن اتّفقت فيه جميع}
\MsL{2}{AB03-F079r-L02}{الأشياء التي تتلقّا سعده أو جميعها منحسة أو كان أكثرها كذلك فهذا هو المذهب الذي}
\MsL{3}{AB03-F079r-L03}{يفحص به عن أمور الأزمان على النحو الذي يليق بالأمر الطبيعيّ. فأمّا المذهب الجزئيّة التي}
\MsL{4}{AB03-F079r-L04}{يفحص بها عن أحوال\PN{433} الأحداث الزمانيّة فإنّه لمّا كان إحصاء ما يعرض منها كثير الفنون}
\MsL{5}{AB03-F079r-L05}{عسر الوصف لا سيّما في هذا الموضع وجب أن يلغى ذكرها من أجل التعرّض الذي}
\MsL{6}{AB03-F079r-L06}{قصدناه في هذا الكتاب من شرح فعل جملة طبيعيّة الكواكب وملائمة \BF{B62r}الأشياء الجزئيّة على تلك}
\Chap{I.38}{<38> وإذا أردت}
\MsL{7}{AB03-F079r-L07}{المذهب\LR{\LatinF . }\Hd{\LR{\LatinF \Sg{〈}38\Sg{〉}}}\Hd{ وإذا أردت} إذا أردت أن تعرف أوّل الشهر الذي تريد من شهور التحويل}
\MsL{8}{AB03-F079r-L08}{فاطلب في جدول الشهور مثل الذي معك من عدد الشهور الثلاثة وخذ ما تحته من عدد الأيّام والساعات}
\MsL{9}{AB03-F079r-L09}{والكسور فما حصل فزده على الأيّام والساعات الماضيّة من وقت المولد إلى وقت طالع التحويل واطلب}
\MsL{10}{AB03-F079r-L10}{في جدول الشهور مثل الذي معك في عدد الشهور التامّة وخذ ما تحته من عدد الأيّام والساعات والكسور}
\MsL{11}{AB03-F079r-L11}{فما حصل فزده على الأيّام والساعات الماضيّة من شهر المولد في وقت المولد أو في وقت التحويل}
\MsL{12}{AB03-F079r-L12}{ممّا تريد أن تعمل عليه فما بلغ فألقه من ذلك الشهر وخذ لكلّ شهر عدد أيّامه لشهور الروم التي}
\MsL{13}{AB03-F079r-L13}{عليها تعمل في تحاويل السنين فما بقي من الأيّام والساعات المعتدلة الماضيّة من الشهر الرومي}
\MsL{14}{AB03-F079r-L14}{الذي ينتهي إليه الذي تلك الأيّام منه فهو أوّل ذلك الشهر الذي أردت من شهور التحويل}
\MsL{15}{AB03-F079r-L15}{تبعد\PN{434} من البروج الذي انتهت اليه تلك السنة من موضع كلّ هيلاج من الهياليج التي}
\MsL{16}{AB03-F079r-L16}{هي الطالع والشمس والقمر وسهم السعادة ووسط السماء على توالي البروج لكلّ شهر}
\MsL{17}{AB03-F079r-L17}{أو برج فالبرج الذي ينتهي إليه هو المدبّر لذلك الشكل\PN{435} وموضع كلّ واحد من الهياليج في أصل}
\Chap{I.39}{<39> وإذا أردت}
\MsL{18}{AB03-F079r-L18}{المولد هو موضع ذلك الشكل\PN{436} من ذلك البرج. \Hd{\LR{\LatinF \Sg{〈}39\Sg{〉}}}\Hd{ وإذا أردت} أن تعرف موضع اليوم الذي تريد}
\MsL{19}{AB03-F079r-L19}{من الشهر الذي أنت فيه فأدخل بعدد الأيّام التامّة الماضيّة من شهر المولد أو من شهر}
\MsL{20}{AB03-F079r-L20}{التحويل إلى جدول الزيادات في الأيّام وخذ ما تحته من الأجزاء والدقائق \BF{B62v}فما بلغ فزده على}
\MsL{21}{AB03-F079r-L21}{الهيلاج من برج الشهر فما بلغ فألقه من برج الشهر لكلّ برج ثلاثين درجة فحيث انتهت}
\begin{PNotes}
\PNote{433}{AB03-n433}{B: \RL{\ArN حال}}
\PNote{434}{AB03-n434}{B: \RL{\ArN بعد}}
\PNote{435}{AB03-n435}{B: \RL{\ArN الشهر} (!)}
\PNote{436}{AB03-n436}{B: \RL{\ArN الشهر} (!)}
\end{PNotes}
\Folio{079v}{79v}
\MsL{1}{AB03-F079v-L01}{من البروج فهو موضع أوّل اليوم الذي أردت من البرج الذي يقع ذلك العدد فيه لكلّ واحد}
\MsL{2}{AB03-F079v-L02}{من الهياليج الخمسة وتنظر إلى الكوكب في ذلك الموضع من الكواكب المتحيّرة والذي\PN{437}}
\Chap{I.40}{<40> وقد قسم}
\MsL{3}{AB03-F079v-L03}{يشاركه في الشكل والكوكب الذي يدبّر ذلك البرج. \Hd{\LR{\LatinF \Sg{〈}40\Sg{〉}}}\Hd{ وقد قسم} آخرون السنة على إثنا}
\MsL{4}{AB03-F079v-L04}{عشر شهرًا على مصيّر الشمس. فجعلوا ابتداء كلّ شهر من مسير الشمس من البرج الذي}
\MsL{5}{AB03-F079v-L05}{ينتهي إليه إلى مثل البروج التي كانت فيها في الأصل من البرج الذي كانت فيه وقالوا}
\MsL{6}{AB03-F079v-L06}{تنظر إلى مقدار ما فارقت الشمس موضعها في أصل المولد من الدرج فتزيّده على طالع}
\MsL{7}{AB03-F079v-L07}{المولد وغيره من الهياليج وتنظر ما يوافق من أجزاء البروج فهو موضع ذلك اليوم.}
\MsL{8}{AB03-F079v-L08}{ويجب أن نفعل مثل ذلك ببروج الأشياء في السنين أيضًا. وقال آخرون تأخذ}
\MsL{9}{AB03-F079v-L09}{من الشمس إلى القمر في كلّ يوم فتلقّيه من طالع أصل المولود \{تقميّة\} للمولود في}
\MsL{10}{AB03-F079v-L10}{ذلك اليوم مقام سهم السعادة. تمّ كتاب الأربعاء شرح البتّانيّ محمّد الله وسنة}
\MsL{11}{AB03-F079v-L11}{في ثالث عشر شهر شوّال المبارك سنة تسع وثلاث\LR{\Sg{〉}}ين\LR{\Sg{〈}} وتسعمائة}
\MsL{12}{AB03-F079v-L12}{والحمد \{لله وسلّم\}. \PN{438}}
\begin{PNotes}
\PNote{437}{AB03-n437}{B: \RL{\ArN أو الذي}}
\PNote{438}{AB03-n438}{B: \RL{\ArN تمّ كتاب الأربعاء شرح البتّانيّ محمّد الله ومنه علقه أحمد بن تميم أحمد حامدًا الله على نعمه ومصلّيًّا على نبيّه محمّد وآله ومسلّمًا وهو حسبنا ونعم الوكيل}}
\end{PNotes}
\Folio{080r}{80r}
\par\medskip\hypertarget{AB03-F080r-T01}{}\begin{center}\footnotesize\setlength{\tabcolsep}{2.5pt}\renewcommand{\arraystretch}{1.05}\resizebox{\ifdim\width>\linewidth\linewidth\else\width\fi}{!}{\begin{tabular}{|c|c|c|c|c|c|c|c|}
\hline
\shortstack{\RL{\ArT كسور}\\\RL{\ArT الساعات}} & \shortstack{\RL{\ArT ساعات}\\\RL{\ArT معتدلات}} & \shortstack{\RL{\ArT أيّام}\\\RL{\ArT الشهور}} & \RL{\ArT شهور التحويل} & \multirow{32}{*}{} & \multicolumn{3}{|c|}{\shortstack{\RL{\ArT جدول أيّام الشهور وما يزاد على}\\\RL{\ArT درجات سهم التحويل\PN{439}}}} \\
\cline{1-4}\cline{6-8}
\RL{\ArT كسور} & \RL{\ArT ساعات} & \RL{\ArT أيّام} & \RL{\ArT شهور} &  & \RL{\ArT دقائق} & \RL{\ArT درج} & \RL{\ArT أيّام الشهور} \\
\cline{1-4}\cline{6-8}
\multirow{2}{*}{\RL{\ArT يز}} & \multirow{2}{*}{\RL{\ArT ب}} & \multirow{2}{*}{\RL{\ArT كح}} & \multirow{2}{*}{\RL{\ArT ا}} &  & \RL{\ArT نج} & \RL{\ArT يج} & \RL{\ArT ا} \\
\cline{6-8}
 &  &  &  &  & \RL{\ArT مو} & \RL{\ArT ك} & \RL{\ArT ب} \\
\cline{1-4}\cline{6-8}
\multirow{2}{*}{\RL{\ArT له}} & \multirow{2}{*}{\RL{\ArT د}} & \multirow{2}{*}{\RL{\ArT نو}} & \multirow{2}{*}{\RL{\ArT ب}} &  & \RL{\ArT لط} & \RL{\ArT ما} & \RL{\ArT ج} \\
\cline{6-8}
 &  &  &  &  & \RL{\ArT لب} & \RL{\ArT نه} & \RL{\ArT د} \\
\cline{1-4}\cline{6-8}
\multirow{2}{*}{\RL{\ArT نب}} & \multirow{2}{*}{\RL{\ArT و}} & \multirow{2}{*}{\RL{\ArT فد}} & \multirow{2}{*}{\RL{\ArT ج}} &  & \RL{\ArT كه} & \RL{\ArT --} & \RL{\ArT ه} \\
\cline{6-8}
 &  &  &  &  & \RL{\ArT يز} & \RL{\ArT --} & \RL{\ArT و} \\
\cline{1-4}\cline{6-8}
\multirow{2}{*}{\RL{\ArT ي}} & \multirow{2}{*}{\RL{\ArT ط}} & \multirow{2}{*}{\RL{\ArT قيب}} & \multirow{2}{*}{\RL{\ArT د}} &  & \RL{\ArT \{ك\}} & \RL{\ArT ضا} & \RL{\ArT ز} \\
\cline{6-8}
 &  &  &  &  & \RL{\ArT ج} & \RL{\ArT \{قح\}} & \RL{\ArT ح} \\
\cline{1-4}\cline{6-8}
\multirow{2}{*}{\RL{\ArT كز}} & \multirow{2}{*}{\RL{\ArT يا}} & \multirow{2}{*}{\RL{\ArT قم}} & \multirow{2}{*}{\RL{\ArT ه}} &  & \RL{\ArT نو} & \RL{\ArT قكا} & \RL{\ArT ط} \\
\cline{6-8}
 &  &  &  &  & \RL{\ArT مط} & \RL{\ArT قلح} & \RL{\ArT ي} \\
\cline{1-4}\cline{6-8}
\multirow{2}{*}{\RL{\ArT مه}} & \multirow{2}{*}{\RL{\ArT يج}} & \multirow{2}{*}{\RL{\ArT قصح}} & \multirow{2}{*}{\RL{\ArT و}} &  & \RL{\ArT مب} & \RL{\ArT قلب} & \RL{\ArT يا} \\
\cline{6-8}
 &  &  &  &  & \RL{\ArT له} & \RL{\ArT قصو} & \RL{\ArT يب} \\
\cline{1-4}\cline{6-8}
\multirow{2}{*}{\RL{\ArT ب}} & \multirow{2}{*}{\RL{\ArT يو}} & \multirow{2}{*}{\RL{\ArT قضو}} & \multirow{2}{*}{\RL{\ArT ز}} &  & \RL{\ArT كح} & \RL{\ArT قف} & \RL{\ArT يج} \\
\cline{6-8}
 &  &  &  &  & \RL{\ArT كا} & \RL{\ArT قضد} & \RL{\ArT يد} \\
\cline{1-4}\cline{6-8}
\multirow{2}{*}{\RL{\ArT كد}} & \multirow{2}{*}{\RL{\ArT يح}} & \multirow{2}{*}{\RL{\ArT ركد}} & \multirow{2}{*}{\RL{\ArT ح}} &  & \RL{\ArT يد} & \RL{\ArT ب ح} & \RL{\ArT يه} \\
\cline{6-8}
 &  &  &  &  & \RL{\ArT ز} & \RL{\ArT ركب} & \RL{\ArT يو} \\
\cline{1-4}\cline{6-8}
\multirow{2}{*}{\RL{\ArT لز}} & \multirow{2}{*}{\RL{\ArT ك}} & \multirow{2}{*}{\RL{\ArT ريب}} & \multirow{2}{*}{\RL{\ArT ط}} &  & \RL{\ArT ۰} & \RL{\ArT رلو} & \RL{\ArT يز} \\
\cline{6-8}
 &  &  &  &  & \RL{\ArT نب} & \RL{\ArT رمصا} & \RL{\ArT يح} \\
\cline{1-4}\cline{6-8}
\multirow{2}{*}{\RL{\ArT نه}} & \multirow{2}{*}{\RL{\ArT كب}} & \multirow{2}{*}{\RL{\ArT رف}} & \multirow{2}{*}{\RL{\ArT ي}} &  & \RL{\ArT مه} & \RL{\ArT رصج} & \RL{\ArT يط} \\
\cline{6-8}
 &  &  &  &  & \RL{\ArT لح} & \RL{\ArT رعز} & \RL{\ArT ك} \\
\cline{1-4}\cline{6-8}
\multirow{2}{*}{\RL{\ArT يب}} & \multirow{2}{*}{\RL{\ArT ا}} & \multirow{2}{*}{\RL{\ArT سط}} & \multirow{2}{*}{\RL{\ArT يا}} &  & \RL{\ArT لا} & \RL{\ArT رضا} & \RL{\ArT كا} \\
\cline{6-8}
 &  &  &  &  & \RL{\ArT كد} & \RL{\ArT \{سيه\}} & \RL{\ArT كب} \\
\cline{1-4}\cline{6-8}
\multirow{2}{*}{\RL{\ArT ك}} & \multirow{2}{*}{\RL{\ArT ج}} & \multirow{2}{*}{\RL{\ArT سلز}} & \multirow{2}{*}{\RL{\ArT يب}} &  & \RL{\ArT يز} & \RL{\ArT سيط} & \RL{\ArT كج} \\
\cline{6-8}
 &  &  &  &  & \RL{\ArT ي} & \RL{\ArT سلج} & \RL{\ArT كد} \\
\cline{1-4}\cline{6-8}
\multirow{2}{*}{\RL{\ArT مز}} & \multirow{2}{*}{\RL{\ArT ه}} & \multirow{2}{*}{\RL{\ArT سصه}} & \multirow{2}{*}{\RL{\ArT يج}} &  & \RL{\ArT ج} & \RL{\ArT سمز} & \RL{\ArT كه} \\
\cline{6-8}
 &  &  &  &  & \RL{\ArT نو} & \RL{\ArT ۰} & \RL{\ArT كو} \\
\cline{1-4}\cline{6-8}
\multicolumn{4}{|c|}{\multirow{4}{*}{}} &  & \RL{\ArT مح} & \RL{\ArT يد} & \RL{\ArT كز} \\
\cline{6-8}
\multicolumn{4}{|c|}{} &  & \RL{\ArT ما} & \RL{\ArT كح} & \RL{\ArT كح} \\
\cline{6-8}
\multicolumn{4}{|c|}{} &  & \RL{\ArT ؟} & \RL{\ArT ل} & \RL{\ArT \{كا مز كح\}} \\
\cline{6-8}
\multicolumn{4}{|c|}{} &  & \RL{\ArT ۰} & \RL{\ArT \{۰\}} & \RL{\ArT ؟} \\
\hline
\end{tabular}}\end{center}
\begin{PNotes}
\PNote{439}{AB03-n439}{This table presents multiples of ca. 13;52,54,30° per day. Scribal mistakes have not been corrected. For a copy with different errors, see Chabás, J., \& Goldstein, B. R. (2012). A Survey of European Astronomical Tables in the Late Middle Ages. Leiden, Niederlande: Brill. doi: https://doi.org/10.1163/9789004230590, pp. 100, 221.}
\end{PNotes}
\Folio{080v}{80v}
\par\medskip\hypertarget{AB03-F080v-T01}{}\begin{center}\footnotesize\setlength{\tabcolsep}{2.5pt}\renewcommand{\arraystretch}{1.05}\resizebox{\ifdim\width>\linewidth\linewidth\else\width\fi}{!}{\begin{tabular}{|c|c|c|c|c|c|c|c|c|c|c|c|c|c|}
\hline
\multicolumn{3}{|c|}{\shortstack{\RL{\ArT أرباب المثلّثات بالنهار}\\\RL{\ArT على مذهب بطلميوس}}} & \multirow{25}{*}{} & \multicolumn{3}{|c|}{\RL{\ArT الوجوه على مذهب بطلميوس}} & \multirow{25}{*}{} & \multicolumn{5}{|c|}{\RL{\ArT الحدود على مذهب سيّد الحكمة بطلميوس}} & \RL{\ArT أسماء البروج} \\
\cline{1-3}\cline{5-7}\cline{9-14}
\multirow{2}{*}{\RL{\ArT خ}} & \multirow{2}{*}{\RL{\ArT \{ي\}}} & \multirow{2}{*}{\RL{\ArT س}} &  & \multirow{2}{*}{\RL{\ArT ي}} & \multirow{2}{*}{\RL{\ArT س}} & \multirow{2}{*}{\RL{\ArT خ}} &  & \RL{\ArT ل} & \RL{\ArT خ} & \RL{\ArT د} & \RL{\ArT ة} & \RL{\ArT ي} & \multirow{2}{*}{\RL{\ArT الحمل}} \\
\cline{9-13}
 &  &  &  &  &  &  &  & \RL{\ArT د} & \RL{\ArT ه} & \RL{\ArT ز} & \RL{\ArT ح} & \RL{\ArT و} &  \\
\cline{1-3}\cline{5-7}\cline{9-14}
\multirow{2}{*}{\RL{\ArT ل}} & \multirow{2}{*}{\RL{\ArT ر}} & \multirow{2}{*}{\RL{\ArT ة}} &  & \multirow{2}{*}{\RL{\ArT ل}} & \multirow{2}{*}{\RL{\ArT د}} & \multirow{2}{*}{\RL{\ArT ة}} &  & \RL{\ArT خ} & \RL{\ArT ل} & \RL{\ArT ي} & \RL{\ArT د} & \RL{\ArT ة} & \multirow{2}{*}{\RL{\ArT الثور}} \\
\cline{9-13}
 &  &  &  &  &  &  &  & \RL{\ArT و} & \RL{\ArT ب} & \RL{\ArT ز} & \RL{\ArT ز} & \RL{\ArT ح} &  \\
\cline{1-3}\cline{5-7}\cline{9-14}
\multirow{2}{*}{\RL{\ArT ي}} & \multirow{2}{*}{\RL{\ArT د}} & \multirow{2}{*}{\RL{\ArT ل}} &  & \multirow{2}{*}{\RL{\ArT ل}} & \multirow{2}{*}{\RL{\ArT ة}} & \multirow{2}{*}{\RL{\ArT د}} &  & \RL{\ArT ل} & \RL{\ArT خ} & \RL{\ArT ة} & \RL{\ArT ي} & \RL{\ArT د} & \multirow{2}{*}{\RL{\ArT الجوزاء}} \\
\cline{9-13}
 &  &  &  &  &  &  &  & \RL{\ArT د} & \RL{\ArT و} & \RL{\ArT ز} & \RL{\ArT و} & \RL{\ArT ز} &  \\
\cline{1-3}\cline{5-7}\cline{9-14}
\multirow{2}{*}{\shortstack{\RL{\ArT ة}\\\RL{\ArT ر}}} & \multirow{2}{*}{\RL{\ArT خ}} & \multirow{2}{*}{\RL{\ArT خ}} &  & \multirow{2}{*}{\RL{\ArT ي}} & \multirow{2}{*}{\RL{\ArT خ}} & \multirow{2}{*}{\RL{\ArT ر}} &  & \RL{\ArT ل} & \RL{\ArT ة} & \RL{\ArT د} & \RL{\ArT ي} & \RL{\ArT خ} & \multirow{2}{*}{\RL{\ArT السرطان}} \\
\cline{9-13}
 &  &  &  &  &  &  &  & \RL{\ArT ج} & \RL{\ArT ز} & \RL{\ArT ز} & \RL{\ArT ز} & \RL{\ArT و} &  \\
\cline{1-3}\cline{5-7}\cline{9-14}
\multirow{2}{*}{\RL{\ArT خ}} & \multirow{2}{*}{\RL{\ArT ي}} & \multirow{2}{*}{\RL{\ArT س}} &  & \multirow{2}{*}{\RL{\ArT خ}} & \multirow{2}{*}{\RL{\ArT ي}} & \multirow{2}{*}{\RL{\ArT س}} &  & \RL{\ArT ل} & \RL{\ArT ة} & \RL{\ArT خ} & \RL{\ArT د} & \RL{\ArT ي} & \multirow{2}{*}{\RL{\ArT الأسد}} \\
\cline{9-13}
 &  &  &  &  &  &  &  & \RL{\ArT و} & \RL{\ArT و} & \RL{\ArT ه} & \RL{\ArT ز} & \RL{\ArT و} &  \\
\cline{1-3}\cline{5-7}\cline{9-14}
\multirow{2}{*}{\RL{\ArT ل}} & \multirow{2}{*}{\RL{\ArT ر}} & \multirow{2}{*}{\RL{\ArT ة}} &  & \multirow{2}{*}{\RL{\ArT ة}} & \multirow{2}{*}{\RL{\ArT ل}} & \multirow{2}{*}{\RL{\ArT د}} &  & \RL{\ArT خ} & \RL{\ArT ل} & \RL{\ArT ي} & \RL{\ArT ة} & \RL{\ArT د} & \multirow{2}{*}{\RL{\ArT السنبلة}} \\
\cline{9-13}
 &  &  &  &  &  &  &  & \RL{\ArT و} & \RL{\ArT و} & \RL{\ArT ه} & \RL{\ArT و} & \RL{\ArT ز} &  \\
\cline{1-3}\cline{5-7}\cline{9-14}
\multirow{2}{*}{\RL{\ArT ي}} & \multirow{2}{*}{\RL{\ArT د}} & \multirow{2}{*}{\RL{\ArT ل}} &  & \multirow{2}{*}{\RL{\ArT د}} & \multirow{2}{*}{\RL{\ArT ل}} & \multirow{2}{*}{\RL{\ArT ة}} &  & \RL{\ArT خ} & \RL{\ArT ي} & \RL{\ArT د} & \RL{\ArT ة} & \RL{\ArT ل} & \multirow{2}{*}{\RL{\ArT الميزان}} \\
\cline{9-13}
 &  &  &  &  &  &  &  & \RL{\ArT و} & \RL{\ArT ح} & \RL{\ArT ه} & \RL{\ArT ه} & \RL{\ArT و} &  \\
\cline{1-3}\cline{5-7}\cline{9-14}
\multirow{2}{*}{\shortstack{\RL{\ArT ة}\\\RL{\ArT ر}}} & \multirow{2}{*}{\RL{\ArT خ}} & \multirow{2}{*}{\RL{\ArT خ}} &  & \multirow{2}{*}{\RL{\ArT ر}} & \multirow{2}{*}{\RL{\ArT ي}} & \multirow{2}{*}{\RL{\ArT خ}} &  & \RL{\ArT ل} & \RL{\ArT د} & \RL{\ArT ي} & \RL{\ArT ة} & \RL{\ArT خ} & \multirow{2}{*}{\RL{\ArT العقرب}} \\
\cline{9-13}
 &  &  &  &  &  &  &  & \RL{\ArT ج} & \RL{\ArT و} & \RL{\ArT ح} & \RL{\ArT ز} & \RL{\ArT و} &  \\
\cline{1-3}\cline{5-7}\cline{9-14}
\multirow{2}{*}{\RL{\ArT خ}} & \multirow{2}{*}{\RL{\ArT ي}} & \multirow{2}{*}{\RL{\ArT س}} &  & \multirow{2}{*}{\RL{\ArT س}} & \multirow{2}{*}{\RL{\ArT خ}} & \multirow{2}{*}{\RL{\ArT ي}} &  & \RL{\ArT خ} & \RL{\ArT ل} & \RL{\ArT د} & \RL{\ArT ة} & \RL{\ArT ي} & \multirow{2}{*}{\RL{\ArT القوس}} \\
\cline{9-13}
 &  &  &  &  &  &  &  & \RL{\ArT ه} & \RL{\ArT و} & \RL{\ArT ه} & \RL{\ArT و} & \RL{\ArT ح} &  \\
\cline{1-3}\cline{5-7}\cline{9-14}
\multirow{2}{*}{\RL{\ArT ل}} & \multirow{2}{*}{\RL{\ArT د}} & \multirow{2}{*}{\RL{\ArT ة}} &  & \multirow{2}{*}{\RL{\ArT د}} & \multirow{2}{*}{\RL{\ArT ة}} & \multirow{2}{*}{\RL{\ArT ل}} &  & \RL{\ArT خ} & \RL{\ArT ل} & \RL{\ArT ي} & \RL{\ArT د} & \RL{\ArT ة} & \multirow{2}{*}{\RL{\ArT الجدي}} \\
\cline{9-13}
 &  &  &  &  &  &  &  & \RL{\ArT ه} & \RL{\ArT و} & \RL{\ArT ز} & \RL{\ArT و} & \RL{\ArT و} &  \\
\cline{1-3}\cline{5-7}\cline{9-14}
\multirow{2}{*}{\RL{\ArT ي}} & \multirow{2}{*}{\RL{\ArT د}} & \multirow{2}{*}{\RL{\ArT ل}} &  & \multirow{2}{*}{\RL{\ArT ة}} & \multirow{2}{*}{\RL{\ArT د}} & \multirow{2}{*}{\RL{\ArT ل}} &  & \RL{\ArT خ} & \RL{\ArT ي} & \RL{\ArT ة} & \RL{\ArT د} & \RL{\ArT ل} & \multirow{2}{*}{\RL{\ArT الدلو}} \\
\cline{9-13}
 &  &  &  &  &  &  &  & \RL{\ArT ه} & \RL{\ArT ه} & \RL{\ArT ح} & \RL{\ArT و} & \RL{\ArT و} &  \\
\cline{1-3}\cline{5-7}\cline{9-14}
\multirow{2}{*}{\shortstack{\RL{\ArT ة}\\\RL{\ArT ر}}} & \multirow{2}{*}{\RL{\ArT خ}} & \multirow{2}{*}{\RL{\ArT خ}} &  & \multirow{2}{*}{\RL{\ArT خ}} & \multirow{2}{*}{\RL{\ArT ر}} & \multirow{2}{*}{\RL{\ArT ي}} &  & \RL{\ArT ل} & \RL{\ArT خ} & \RL{\ArT د} & \RL{\ArT ي} & \RL{\ArT ة} & \multirow{2}{*}{\RL{\ArT الحوت}} \\
\cline{9-13}
 &  &  &  &  &  &  &  & \RL{\ArT ه} & \RL{\ArT ه} & \RL{\ArT و} & \RL{\ArT و} & \RL{\ArT ح} &  \\
\hline
\end{tabular}}\end{center}
\end{document}
